\section{逐题推导与实验工作坊}

这一节把前面的题号索引继续展开。每道题都从自己的题面模型出发，先说明慢在哪里，再说明状态怎样保存、哪些候选被安全排除、哪些边界会打破错误实现。题型共性只作为背景，真正要读的是每道题自己的瓶颈和反例。

\subsection{数组与原地维护工作坊}

先写一个不会漏答案的枚举或辅助数组版本，再观察哪些位置已经被前缀、后缀、慢指针或边界确定。优化时把数组划成语义清楚的区域，每次循环只扩大已确认区域。

\subsubsection{LeetCode 5 最长回文子串}

先从位置关系入手。本题的单点模型是回文围绕中心对称，中心可以是字符也可以是两个字符之间。朴素版可以为每个位置重新寻找最终状态，或者用辅助数组把答案先做出来；它虽然慢，却能说明哪些位置必须被覆盖。本题真正浪费的部分是中心扩展避免枚举子串后再逐个检查；每次扩展只比较新增左右边界，后面的优化只能消掉这部分重复，不能改变输出语义。

优化时把数组切成几段：已经确认的前缀、未知区、尚未处理的后缀，或者左侧贡献和右侧贡献。本题落到代码上就是中心扩展避免枚举子串后再逐个检查；每次扩展只比较新增左右边界。循环中每改一个下标，都要能说出哪一段扩大了，哪一段缩小了。放到“数组与原地维护”这条主线里看，关键原则是：优化要把数组切成若干语义明确的区间：已经处理好的前缀、未知区、已经确定的后缀，或者左侧贡献和右侧贡献。双指针、前后缀数组、原地交换和稳定写入，本质都是让每个元素只进入少数几次状态转移。

正确性可以按区间不变量证明：已处理部分满足题意，未知部分还没承诺，循环每次只把一个元素移入正确区域。对本题来说，回文围绕中心对称，中心可以是字符也可以是两个字符之间给出了答案范围，中心扩展避免枚举子串后再逐个检查；每次扩展只比较新增左右边界给出了推进方式。这道题的边界不能留到最后补丁式处理，实验时把偶数长度回文、单字符、全相同字符、多个同长答案列成表，再打印关键下标和有效区间。若某个元素被写入后又被未来步骤破坏，说明区间不变量没有成立。

C++ 实现上，C++ 中优先使用 \code{std::vector} 和清楚的整型下标；涉及乘积、面积、和或距离时提前考虑 \code{long long}。如果题目要求原地修改，返回长度、容器容量和有效区间要分开讲。如果追问变成“若要求回文子序列，应转成区间 DP，不能用中心扩展”，先判断原来的状态是否仍然覆盖新答案集合，再决定是微调边界、改 value 语义，还是换成另一类结构。

\subsubsection{LeetCode 26 删除有序数组重复项}

先从位置关系入手。本题的单点模型是有序数组让重复元素相邻，慢指针表示下一个唯一值写入位置。朴素版可以为每个位置重新寻找最终状态，或者用辅助数组把答案先做出来；它虽然慢，却能说明哪些位置必须被覆盖。本题真正浪费的部分是快指针扫描，只有发现新值才写入慢指针并推进，后面的优化只能消掉这部分重复，不能改变输出语义。

优化时把数组切成几段：已经确认的前缀、未知区、尚未处理的后缀，或者左侧贡献和右侧贡献。本题落到代码上就是快指针扫描，只有发现新值才写入慢指针并推进。循环中每改一个下标，都要能说出哪一段扩大了，哪一段缩小了。放到“数组与原地维护”这条主线里看，关键原则是：优化要把数组切成若干语义明确的区间：已经处理好的前缀、未知区、已经确定的后缀，或者左侧贡献和右侧贡献。双指针、前后缀数组、原地交换和稳定写入，本质都是让每个元素只进入少数几次状态转移。

正确性可以按区间不变量证明：已处理部分满足题意，未知部分还没承诺，循环每次只把一个元素移入正确区域。对本题来说，有序数组让重复元素相邻，慢指针表示下一个唯一值写入位置给出了答案范围，快指针扫描，只有发现新值才写入慢指针并推进给出了推进方式。这道题的边界不能留到最后补丁式处理，实验时把空数组、全重复、无重复、返回长度不等于物理容量列成表，再打印关键下标和有效区间。若某个元素被写入后又被未来步骤破坏，说明区间不变量没有成立。

C++ 实现上，C++ 中优先使用 \code{std::vector} 和清楚的整型下标；涉及乘积、面积、和或距离时提前考虑 \code{long long}。如果题目要求原地修改，返回长度、容器容量和有效区间要分开讲。如果追问变成“若允许每个元素最多保留两次，只需改变写入条件为和前两个比较”，先判断原来的状态是否仍然覆盖新答案集合，再决定是微调边界、改 value 语义，还是换成另一类结构。

\subsubsection{LeetCode 27 移除元素}

先从位置关系入手。本题的单点模型是原地过滤，慢指针左侧始终是不等于目标值的稳定前缀。朴素版可以为每个位置重新寻找最终状态，或者用辅助数组把答案先做出来；它虽然慢，却能说明哪些位置必须被覆盖。本题真正浪费的部分是保持顺序用快慢指针；不保持顺序可用左右交换减少写入，后面的优化只能消掉这部分重复，不能改变输出语义。

优化时把数组切成几段：已经确认的前缀、未知区、尚未处理的后缀，或者左侧贡献和右侧贡献。本题落到代码上就是保持顺序用快慢指针；不保持顺序可用左右交换减少写入。循环中每改一个下标，都要能说出哪一段扩大了，哪一段缩小了。放到“数组与原地维护”这条主线里看，关键原则是：优化要把数组切成若干语义明确的区间：已经处理好的前缀、未知区、已经确定的后缀，或者左侧贡献和右侧贡献。双指针、前后缀数组、原地交换和稳定写入，本质都是让每个元素只进入少数几次状态转移。

正确性可以按区间不变量证明：已处理部分满足题意，未知部分还没承诺，循环每次只把一个元素移入正确区域。对本题来说，原地过滤，慢指针左侧始终是不等于目标值的稳定前缀给出了答案范围，保持顺序用快慢指针；不保持顺序可用左右交换减少写入给出了推进方式。这道题的边界不能留到最后补丁式处理，实验时把所有元素都被移除、没有元素被移除、目标值在尾部密集列成表，再打印关键下标和有效区间。若某个元素被写入后又被未来步骤破坏，说明区间不变量没有成立。

C++ 实现上，C++ 中优先使用 \code{std::vector} 和清楚的整型下标；涉及乘积、面积、和或距离时提前考虑 \code{long long}。如果题目要求原地修改，返回长度、容器容量和有效区间要分开讲。如果追问变成“工程里要先确认是否要求稳定顺序，不同要求对应不同写法”，先判断原来的状态是否仍然覆盖新答案集合，再决定是微调边界、改 value 语义，还是换成另一类结构。

\subsubsection{LeetCode 31 下一个排列}

先从位置关系入手。本题的单点模型是字典序结构由最长非递增后缀决定，枢轴是后缀前一个位置。朴素版可以为每个位置重新寻找最终状态，或者用辅助数组把答案先做出来；它虽然慢，却能说明哪些位置必须被覆盖。本题真正浪费的部分是从右找第一个上升点，再从右找刚好更大的元素交换，最后反转后缀，后面的优化只能消掉这部分重复，不能改变输出语义。

优化时把数组切成几段：已经确认的前缀、未知区、尚未处理的后缀，或者左侧贡献和右侧贡献。本题落到代码上就是从右找第一个上升点，再从右找刚好更大的元素交换，最后反转后缀。循环中每改一个下标，都要能说出哪一段扩大了，哪一段缩小了。放到“数组与原地维护”这条主线里看，关键原则是：优化要把数组切成若干语义明确的区间：已经处理好的前缀、未知区、已经确定的后缀，或者左侧贡献和右侧贡献。双指针、前后缀数组、原地交换和稳定写入，本质都是让每个元素只进入少数几次状态转移。

正确性可以按区间不变量证明：已处理部分满足题意，未知部分还没承诺，循环每次只把一个元素移入正确区域。对本题来说，字典序结构由最长非递增后缀决定，枢轴是后缀前一个位置给出了答案范围，从右找第一个上升点，再从右找刚好更大的元素交换，最后反转后缀给出了推进方式。这道题的边界不能留到最后补丁式处理，实验时把完全降序、重复数字、只有一个元素、交换后后缀必须最小化列成表，再打印关键下标和有效区间。若某个元素被写入后又被未来步骤破坏，说明区间不变量没有成立。

C++ 实现上，C++ 中优先使用 \code{std::vector} 和清楚的整型下标；涉及乘积、面积、和或距离时提前考虑 \code{long long}。如果题目要求原地修改，返回长度、容器容量和有效区间要分开讲。如果追问变成“若要求上一个排列，比较方向和后缀处理对称变化”，先判断原来的状态是否仍然覆盖新答案集合，再决定是微调边界、改 value 语义，还是换成另一类结构。

\subsubsection{LeetCode 54 螺旋矩阵}

先从位置关系入手。本题的单点模型是四个边界收缩模拟一圈一圈访问。朴素版可以为每个位置重新寻找最终状态，或者用辅助数组把答案先做出来；它虽然慢，却能说明哪些位置必须被覆盖。本题真正浪费的部分是每走完一条边就收缩对应边界，并在走反方向前检查是否仍合法，后面的优化只能消掉这部分重复，不能改变输出语义。

优化时把数组切成几段：已经确认的前缀、未知区、尚未处理的后缀，或者左侧贡献和右侧贡献。本题落到代码上就是每走完一条边就收缩对应边界，并在走反方向前检查是否仍合法。循环中每改一个下标，都要能说出哪一段扩大了，哪一段缩小了。放到“数组与原地维护”这条主线里看，关键原则是：优化要把数组切成若干语义明确的区间：已经处理好的前缀、未知区、已经确定的后缀，或者左侧贡献和右侧贡献。双指针、前后缀数组、原地交换和稳定写入，本质都是让每个元素只进入少数几次状态转移。

正确性可以按区间不变量证明：已处理部分满足题意，未知部分还没承诺，循环每次只把一个元素移入正确区域。对本题来说，四个边界收缩模拟一圈一圈访问给出了答案范围，每走完一条边就收缩对应边界，并在走反方向前检查是否仍合法给出了推进方式。这道题的边界不能留到最后补丁式处理，实验时把单行、单列、空矩阵、最后一圈只剩一个元素列成表，再打印关键下标和有效区间。若某个元素被写入后又被未来步骤破坏，说明区间不变量没有成立。

C++ 实现上，C++ 中优先使用 \code{std::vector} 和清楚的整型下标；涉及乘积、面积、和或距离时提前考虑 \code{long long}。如果题目要求原地修改，返回长度、容器容量和有效区间要分开讲。如果追问变成“若改成生成矩阵，边界逻辑相同，只是读变成写”，先判断原来的状态是否仍然覆盖新答案集合，再决定是微调边界、改 value 语义，还是换成另一类结构。

\subsubsection{LeetCode 75 颜色分类}

先从位置关系入手。本题的单点模型是三指针维护零区、一号区、未知区和二区。朴素版可以为每个位置重新寻找最终状态，或者用辅助数组把答案先做出来；它虽然慢，却能说明哪些位置必须被覆盖。本题真正浪费的部分是遇到二时只移动右边界，不移动扫描指针，因为换回来的元素未知，后面的优化只能消掉这部分重复，不能改变输出语义。

优化时把数组切成几段：已经确认的前缀、未知区、尚未处理的后缀，或者左侧贡献和右侧贡献。本题落到代码上就是遇到二时只移动右边界，不移动扫描指针，因为换回来的元素未知。循环中每改一个下标，都要能说出哪一段扩大了，哪一段缩小了。放到“数组与原地维护”这条主线里看，关键原则是：优化要把数组切成若干语义明确的区间：已经处理好的前缀、未知区、已经确定的后缀，或者左侧贡献和右侧贡献。双指针、前后缀数组、原地交换和稳定写入，本质都是让每个元素只进入少数几次状态转移。

正确性可以按区间不变量证明：已处理部分满足题意，未知部分还没承诺，循环每次只把一个元素移入正确区域。对本题来说，三指针维护零区、一号区、未知区和二区给出了答案范围，遇到二时只移动右边界，不移动扫描指针，因为换回来的元素未知给出了推进方式。这道题的边界不能留到最后补丁式处理，实验时把全零、全二、已经有序、逆序、交换后漏检查列成表，再打印关键下标和有效区间。若某个元素被写入后又被未来步骤破坏，说明区间不变量没有成立。

C++ 实现上，C++ 中优先使用 \code{std::vector} 和清楚的整型下标；涉及乘积、面积、和或距离时提前考虑 \code{long long}。如果题目要求原地修改，返回长度、容器容量和有效区间要分开讲。如果追问变成“这是荷兰国旗问题，能推广到三类分区的原地维护”，先判断原来的状态是否仍然覆盖新答案集合，再决定是微调边界、改 value 语义，还是换成另一类结构。

\subsubsection{LeetCode 88 合并两个有序数组}

先从位置关系入手。本题的单点模型是第一个数组尾部有空位，适合从后往前写入最终最大值。朴素版可以为每个位置重新寻找最终状态，或者用辅助数组把答案先做出来；它虽然慢，却能说明哪些位置必须被覆盖。本题真正浪费的部分是逆向填充避免覆盖尚未读取的 nums1 元素，后面的优化只能消掉这部分重复，不能改变输出语义。

优化时把数组切成几段：已经确认的前缀、未知区、尚未处理的后缀，或者左侧贡献和右侧贡献。本题落到代码上就是逆向填充避免覆盖尚未读取的 nums1 元素。循环中每改一个下标，都要能说出哪一段扩大了，哪一段缩小了。放到“数组与原地维护”这条主线里看，关键原则是：优化要把数组切成若干语义明确的区间：已经处理好的前缀、未知区、已经确定的后缀，或者左侧贡献和右侧贡献。双指针、前后缀数组、原地交换和稳定写入，本质都是让每个元素只进入少数几次状态转移。

正确性可以按区间不变量证明：已处理部分满足题意，未知部分还没承诺，循环每次只把一个元素移入正确区域。对本题来说，第一个数组尾部有空位，适合从后往前写入最终最大值给出了答案范围，逆向填充避免覆盖尚未读取的 nums1 元素给出了推进方式。这道题的边界不能留到最后补丁式处理，实验时把第二个数组为空、第一个有效元素为空、重复值、尾部剩余列成表，再打印关键下标和有效区间。若某个元素被写入后又被未来步骤破坏，说明区间不变量没有成立。

C++ 实现上，C++ 中优先使用 \code{std::vector} 和清楚的整型下标；涉及乘积、面积、和或距离时提前考虑 \code{long long}。如果题目要求原地修改，返回长度、容器容量和有效区间要分开讲。如果追问变成“若没有额外空位，只能新开数组或使用更复杂原地归并”，先判断原来的状态是否仍然覆盖新答案集合，再决定是微调边界、改 value 语义，还是换成另一类结构。

\subsubsection{LeetCode 189 轮转数组}

先从位置关系入手。本题的单点模型是右移 k 位等价于把数组切成两段后换序。朴素版可以为每个位置重新寻找最终状态，或者用辅助数组把答案先做出来；它虽然慢，却能说明哪些位置必须被覆盖。本题真正浪费的部分是三次反转能原地完成：整体、前 k、后 n 减 k，后面的优化只能消掉这部分重复，不能改变输出语义。

优化时把数组切成几段：已经确认的前缀、未知区、尚未处理的后缀，或者左侧贡献和右侧贡献。本题落到代码上就是三次反转能原地完成：整体、前 k、后 n 减 k。循环中每改一个下标，都要能说出哪一段扩大了，哪一段缩小了。放到“数组与原地维护”这条主线里看，关键原则是：优化要把数组切成若干语义明确的区间：已经处理好的前缀、未知区、已经确定的后缀，或者左侧贡献和右侧贡献。双指针、前后缀数组、原地交换和稳定写入，本质都是让每个元素只进入少数几次状态转移。

正确性可以按区间不变量证明：已处理部分满足题意，未知部分还没承诺，循环每次只把一个元素移入正确区域。对本题来说，右移 k 位等价于把数组切成两段后换序给出了答案范围，三次反转能原地完成：整体、前 k、后 n 减 k给出了推进方式。这道题的边界不能留到最后补丁式处理，实验时把k 为零、k 大于 n、空数组、单元素列成表，再打印关键下标和有效区间。若某个元素被写入后又被未来步骤破坏，说明区间不变量没有成立。

C++ 实现上，C++ 中优先使用 \code{std::vector} 和清楚的整型下标；涉及乘积、面积、和或距离时提前考虑 \code{long long}。如果题目要求原地修改，返回长度、容器容量和有效区间要分开讲。如果追问变成“环状替换也可做，但要处理最大公约数个循环”，先判断原来的状态是否仍然覆盖新答案集合，再决定是微调边界、改 value 语义，还是换成另一类结构。

\subsubsection{LeetCode 238 除自身以外数组的乘积}

先从位置关系入手。本题的单点模型是不能用除法时，答案由左侧乘积和右侧乘积组成。朴素版可以为每个位置重新寻找最终状态，或者用辅助数组把答案先做出来；它虽然慢，却能说明哪些位置必须被覆盖。本题真正浪费的部分是先写入左侧前缀积，再从右向左乘右侧后缀积，后面的优化只能消掉这部分重复，不能改变输出语义。

优化时把数组切成几段：已经确认的前缀、未知区、尚未处理的后缀，或者左侧贡献和右侧贡献。本题落到代码上就是先写入左侧前缀积，再从右向左乘右侧后缀积。循环中每改一个下标，都要能说出哪一段扩大了，哪一段缩小了。放到“数组与原地维护”这条主线里看，关键原则是：优化要把数组切成若干语义明确的区间：已经处理好的前缀、未知区、已经确定的后缀，或者左侧贡献和右侧贡献。双指针、前后缀数组、原地交换和稳定写入，本质都是让每个元素只进入少数几次状态转移。

正确性可以按区间不变量证明：已处理部分满足题意，未知部分还没承诺，循环每次只把一个元素移入正确区域。对本题来说，不能用除法时，答案由左侧乘积和右侧乘积组成给出了答案范围，先写入左侧前缀积，再从右向左乘右侧后缀积给出了推进方式。这道题的边界不能留到最后补丁式处理，实验时把一个零、多个零、负数、乘积溢出列成表，再打印关键下标和有效区间。若某个元素被写入后又被未来步骤破坏，说明区间不变量没有成立。

C++ 实现上，C++ 中优先使用 \code{std::vector} 和清楚的整型下标；涉及乘积、面积、和或距离时提前考虑 \code{long long}。如果题目要求原地修改，返回长度、容器容量和有效区间要分开讲。如果追问变成“这个题展示输出数组可作为状态存储的一部分”，先判断原来的状态是否仍然覆盖新答案集合，再决定是微调边界、改 value 语义，还是换成另一类结构。

\subsubsection{LeetCode 283 移动零}

先从位置关系入手。本题的单点模型是稳定保留非零元素顺序，把零集中到尾部。朴素版可以为每个位置重新寻找最终状态，或者用辅助数组把答案先做出来；它虽然慢，却能说明哪些位置必须被覆盖。本题真正浪费的部分是慢指针写入下一个非零位置，最后补零或用交换，后面的优化只能消掉这部分重复，不能改变输出语义。

优化时把数组切成几段：已经确认的前缀、未知区、尚未处理的后缀，或者左侧贡献和右侧贡献。本题落到代码上就是慢指针写入下一个非零位置，最后补零或用交换。循环中每改一个下标，都要能说出哪一段扩大了，哪一段缩小了。放到“数组与原地维护”这条主线里看，关键原则是：优化要把数组切成若干语义明确的区间：已经处理好的前缀、未知区、已经确定的后缀，或者左侧贡献和右侧贡献。双指针、前后缀数组、原地交换和稳定写入，本质都是让每个元素只进入少数几次状态转移。

正确性可以按区间不变量证明：已处理部分满足题意，未知部分还没承诺，循环每次只把一个元素移入正确区域。对本题来说，稳定保留非零元素顺序，把零集中到尾部给出了答案范围，慢指针写入下一个非零位置，最后补零或用交换给出了推进方式。这道题的边界不能留到最后补丁式处理，实验时把全零、无零、零在开头结尾、稳定顺序要求列成表，再打印关键下标和有效区间。若某个元素被写入后又被未来步骤破坏，说明区间不变量没有成立。

C++ 实现上，C++ 中优先使用 \code{std::vector} 和清楚的整型下标；涉及乘积、面积、和或距离时提前考虑 \code{long long}。如果题目要求原地修改，返回长度、容器容量和有效区间要分开讲。如果追问变成“若目标是移动指定值，模型与移除元素相同”，先判断原来的状态是否仍然覆盖新答案集合，再决定是微调边界、改 value 语义，还是换成另一类结构。

\subsubsection{LeetCode 581 最短无序连续子数组}

先从位置关系入手。本题的单点模型是要找破坏全局有序性的最左和最右位置。朴素版可以为每个位置重新寻找最终状态，或者用辅助数组把答案先做出来；它虽然慢，却能说明哪些位置必须被覆盖。本题真正浪费的部分是从左维护历史最大值确定右边界，从右维护历史最小值确定左边界，后面的优化只能消掉这部分重复，不能改变输出语义。

优化时把数组切成几段：已经确认的前缀、未知区、尚未处理的后缀，或者左侧贡献和右侧贡献。本题落到代码上就是从左维护历史最大值确定右边界，从右维护历史最小值确定左边界。循环中每改一个下标，都要能说出哪一段扩大了，哪一段缩小了。放到“数组与原地维护”这条主线里看，关键原则是：优化要把数组切成若干语义明确的区间：已经处理好的前缀、未知区、已经确定的后缀，或者左侧贡献和右侧贡献。双指针、前后缀数组、原地交换和稳定写入，本质都是让每个元素只进入少数几次状态转移。

正确性可以按区间不变量证明：已处理部分满足题意，未知部分还没承诺，循环每次只把一个元素移入正确区域。对本题来说，要找破坏全局有序性的最左和最右位置给出了答案范围，从左维护历史最大值确定右边界，从右维护历史最小值确定左边界给出了推进方式。这道题的边界不能留到最后补丁式处理，实验时把已经有序、逆序、重复值、局部乱序列成表，再打印关键下标和有效区间。若某个元素被写入后又被未来步骤破坏，说明区间不变量没有成立。

C++ 实现上，C++ 中优先使用 \code{std::vector} 和清楚的整型下标；涉及乘积、面积、和或距离时提前考虑 \code{long long}。如果题目要求原地修改，返回长度、容器容量和有效区间要分开讲。如果追问变成“排序比较法更直观但复杂度更高”，先判断原来的状态是否仍然覆盖新答案集合，再决定是微调边界、改 value 语义，还是换成另一类结构。

\subsubsection{LeetCode 647 回文子串}

先从位置关系入手。本题的单点模型是每个回文都有唯一中心或中心间隙。朴素版可以为每个位置重新寻找最终状态，或者用辅助数组把答案先做出来；它虽然慢，却能说明哪些位置必须被覆盖。本题真正浪费的部分是对每个中心向两边扩展，扩展成功一次就计数一次，后面的优化只能消掉这部分重复，不能改变输出语义。

优化时把数组切成几段：已经确认的前缀、未知区、尚未处理的后缀，或者左侧贡献和右侧贡献。本题落到代码上就是对每个中心向两边扩展，扩展成功一次就计数一次。循环中每改一个下标，都要能说出哪一段扩大了，哪一段缩小了。放到“数组与原地维护”这条主线里看，关键原则是：优化要把数组切成若干语义明确的区间：已经处理好的前缀、未知区、已经确定的后缀，或者左侧贡献和右侧贡献。双指针、前后缀数组、原地交换和稳定写入，本质都是让每个元素只进入少数几次状态转移。

正确性可以按区间不变量证明：已处理部分满足题意，未知部分还没承诺，循环每次只把一个元素移入正确区域。对本题来说，每个回文都有唯一中心或中心间隙给出了答案范围，对每个中心向两边扩展，扩展成功一次就计数一次给出了推进方式。这道题的边界不能留到最后补丁式处理，实验时把空串、全相同字符、偶数回文、奇数回文列成表，再打印关键下标和有效区间。若某个元素被写入后又被未来步骤破坏，说明区间不变量没有成立。

C++ 实现上，C++ 中优先使用 \code{std::vector} 和清楚的整型下标；涉及乘积、面积、和或距离时提前考虑 \code{long long}。如果题目要求原地修改，返回长度、容器容量和有效区间要分开讲。如果追问变成“Manacher 算法可线性解决，但中心扩展更适合面试基础”，先判断原来的状态是否仍然覆盖新答案集合，再决定是微调边界、改 value 语义，还是换成另一类结构。

\subsection{滑动窗口与双指针工作坊}

先枚举所有窗口，再确认右端扩展和左端收缩是否具有单调性。只有窗口状态可以增量维护，且移动左端不会漏掉未来更优答案时，滑动窗口才成立。

\subsubsection{LeetCode 3 无重复字符的最长子串}

先枚举所有连续窗口，明确左端、右端和窗口内容怎样决定答案。本题的模型是窗口保存当前无重复子串，右端每次加入一个字符；暴力版会反复统计相邻窗口中相同的内容。慢点具体落在用上次出现位置直接跳过无效左端，左边界只能向右不能回退，所以后面要证明左端或右端移动时不会漏掉可能答案。

窗口题只在条件有单调性或窗口长度固定时成立。本题用用上次出现位置直接跳过无效左端，左边界只能向右不能回退维护窗口，含义是右端只带来一个新字符或新元素，左端只移走一个旧元素；如果输入条件改变导致状态不再单调，就必须换方法。放到“滑动窗口与双指针”这条主线里看，关键原则是：滑动窗口成立必须有单调性或固定长度边界。右端扩展让某个条件单调变好或变坏，左端收缩可以恢复合法性或缩短答案。若数组含负数、状态会来回震荡，就不能硬套窗口。

证明重点是排除左端。若某个左端被移过，它对应的窗口已经被完整检查，或继续保留只会让窗口更长、更不合法。本题的窗口保存当前无重复子串，右端每次加入一个字符和用上次出现位置直接跳过无效左端，左边界只能向右不能回退必须一起说明这一点。这道题的边界不能留到最后补丁式处理，实验时重点跑空串、全相同字符、重复字符出现在窗口左侧之外、扩展字符集，并打印左右端、窗口计数和答案。若左端发生回退，或者窗口失去合法性后仍更新答案，就能很快定位错误。

C++ 实现上，窗口计数若字符集固定，可以用 \code{std::array<int, 128>} 或 26 长度数组；字符集不固定再用哈希表。左闭右闭和左闭右开选一种坚持到底，避免窗口长度和删除位置错位。如果追问变成“若要求恰好包含 k 种字符，窗口条件从无重复变成种类计数”，先判断原来的状态是否仍然覆盖新答案集合，再决定是微调边界、改 value 语义，还是换成另一类结构。

\subsubsection{LeetCode 11 盛最多水的容器}

先枚举所有连续窗口，明确左端、右端和窗口内容怎样决定答案。本题的模型是左右指针代表当前选取的两条边，面积由短板和宽度共同决定；暴力版会反复统计相邻窗口中相同的内容。慢点具体落在每次移动短板，因为移动长板只会缩小宽度且不会提高短板，所以后面要证明左端或右端移动时不会漏掉可能答案。

窗口题只在条件有单调性或窗口长度固定时成立。本题用每次移动短板，因为移动长板只会缩小宽度且不会提高短板维护窗口，含义是右端只带来一个新字符或新元素，左端只移走一个旧元素；如果输入条件改变导致状态不再单调，就必须换方法。放到“滑动窗口与双指针”这条主线里看，关键原则是：滑动窗口成立必须有单调性或固定长度边界。右端扩展让某个条件单调变好或变坏，左端收缩可以恢复合法性或缩短答案。若数组含负数、状态会来回震荡，就不能硬套窗口。

证明重点是排除左端。若某个左端被移过，它对应的窗口已经被完整检查，或继续保留只会让窗口更长、更不合法。本题的左右指针代表当前选取的两条边，面积由短板和宽度共同决定和每次移动短板，因为移动长板只会缩小宽度且不会提高短板必须一起说明这一点。这道题的边界不能留到最后补丁式处理，实验时重点跑两端高度相等、数组长度为二、面积乘法可能溢出，并打印左右端、窗口计数和答案。若左端发生回退，或者窗口失去合法性后仍更新答案，就能很快定位错误。

C++ 实现上，窗口计数若字符集固定，可以用 \code{std::array<int, 128>} 或 26 长度数组；字符集不固定再用哈希表。左闭右闭和左闭右开选一种坚持到底，避免窗口长度和删除位置错位。如果追问变成“接雨水计算每个位置贡献，不能把本题双指针证明直接搬过去”，先判断原来的状态是否仍然覆盖新答案集合，再决定是微调边界、改 value 语义，还是换成另一类结构。

\subsubsection{LeetCode 15 三数之和}

先枚举所有连续窗口，明确左端、右端和窗口内容怎样决定答案。本题的模型是排序后固定第一个数，剩余两数转成有序双指针；暴力版会反复统计相邻窗口中相同的内容。慢点具体落在和太小移动左指针，和太大移动右指针；固定值和左右值都要跳过重复，所以后面要证明左端或右端移动时不会漏掉可能答案。

窗口题只在条件有单调性或窗口长度固定时成立。本题用和太小移动左指针，和太大移动右指针；固定值和左右值都要跳过重复维护窗口，含义是右端只带来一个新字符或新元素，左端只移走一个旧元素；如果输入条件改变导致状态不再单调，就必须换方法。放到“滑动窗口与双指针”这条主线里看，关键原则是：滑动窗口成立必须有单调性或固定长度边界。右端扩展让某个条件单调变好或变坏，左端收缩可以恢复合法性或缩短答案。若数组含负数、状态会来回震荡，就不能硬套窗口。

证明重点是排除左端。若某个左端被移过，它对应的窗口已经被完整检查，或继续保留只会让窗口更长、更不合法。本题的排序后固定第一个数，剩余两数转成有序双指针和和太小移动左指针，和太大移动右指针；固定值和左右值都要跳过重复必须一起说明这一点。这道题的边界不能留到最后补丁式处理，实验时重点跑全零、重复元素很多、没有答案、结果顺序不要求但组合不能重复，并打印左右端、窗口计数和答案。若左端发生回退，或者窗口失去合法性后仍更新答案，就能很快定位错误。

C++ 实现上，窗口计数若字符集固定，可以用 \code{std::array<int, 128>} 或 26 长度数组；字符集不固定再用哈希表。左闭右闭和左闭右开选一种坚持到底，避免窗口长度和删除位置错位。如果追问变成“四数之和再外层固定一个数，并用 long long 防止目标和溢出”，先判断原来的状态是否仍然覆盖新答案集合，再决定是微调边界、改 value 语义，还是换成另一类结构。

\subsubsection{LeetCode 16 最接近的三数之和}

先枚举所有连续窗口，明确左端、右端和窗口内容怎样决定答案。本题的模型是排序后固定一数，双指针寻找最接近目标的两数和；暴力版会反复统计相邻窗口中相同的内容。慢点具体落在每次根据当前和与目标的大小移动指针，同时更新最小绝对差，所以后面要证明左端或右端移动时不会漏掉可能答案。

窗口题只在条件有单调性或窗口长度固定时成立。本题用每次根据当前和与目标的大小移动指针，同时更新最小绝对差维护窗口，含义是右端只带来一个新字符或新元素，左端只移走一个旧元素；如果输入条件改变导致状态不再单调，就必须换方法。放到“滑动窗口与双指针”这条主线里看，关键原则是：滑动窗口成立必须有单调性或固定长度边界。右端扩展让某个条件单调变好或变坏，左端收缩可以恢复合法性或缩短答案。若数组含负数、状态会来回震荡，就不能硬套窗口。

证明重点是排除左端。若某个左端被移过，它对应的窗口已经被完整检查，或继续保留只会让窗口更长、更不合法。本题的排序后固定一数，双指针寻找最接近目标的两数和和每次根据当前和与目标的大小移动指针，同时更新最小绝对差必须一起说明这一点。这道题的边界不能留到最后补丁式处理，实验时重点跑差值相等、负数、目标很大、三数和溢出，并打印左右端、窗口计数和答案。若左端发生回退，或者窗口失去合法性后仍更新答案，就能很快定位错误。

C++ 实现上，窗口计数若字符集固定，可以用 \code{std::array<int, 128>} 或 26 长度数组；字符集不固定再用哈希表。左闭右闭和左闭右开选一种坚持到底，避免窗口长度和删除位置错位。如果追问变成“若要求返回所有最接近组合，需要记录同差值并去重”，先判断原来的状态是否仍然覆盖新答案集合，再决定是微调边界、改 value 语义，还是换成另一类结构。

\subsubsection{LeetCode 18 四数之和}

先枚举所有连续窗口，明确左端、右端和窗口内容怎样决定答案。本题的模型是两层固定加一层双指针，关键是剪枝和去重；暴力版会反复统计相邻窗口中相同的内容。慢点具体落在排序提供单调性，外层可用最小可能和和最大可能和提前跳过，所以后面要证明左端或右端移动时不会漏掉可能答案。

窗口题只在条件有单调性或窗口长度固定时成立。本题用排序提供单调性，外层可用最小可能和和最大可能和提前跳过维护窗口，含义是右端只带来一个新字符或新元素，左端只移走一个旧元素；如果输入条件改变导致状态不再单调，就必须换方法。放到“滑动窗口与双指针”这条主线里看，关键原则是：滑动窗口成立必须有单调性或固定长度边界。右端扩展让某个条件单调变好或变坏，左端收缩可以恢复合法性或缩短答案。若数组含负数、状态会来回震荡，就不能硬套窗口。

证明重点是排除左端。若某个左端被移过，它对应的窗口已经被完整检查，或继续保留只会让窗口更长、更不合法。本题的两层固定加一层双指针，关键是剪枝和去重和排序提供单调性，外层可用最小可能和和最大可能和提前跳过必须一起说明这一点。这道题的边界不能留到最后补丁式处理，实验时重点跑重复元素、目标超出 int、答案很多导致输出空间主导复杂度，并打印左右端、窗口计数和答案。若左端发生回退，或者窗口失去合法性后仍更新答案，就能很快定位错误。

C++ 实现上，窗口计数若字符集固定，可以用 \code{std::array<int, 128>} 或 26 长度数组；字符集不固定再用哈希表。左闭右闭和左闭右开选一种坚持到底，避免窗口长度和删除位置错位。如果追问变成“k 数之和可以递归降维，但工程实现要控制重复和边界”，先判断原来的状态是否仍然覆盖新答案集合，再决定是微调边界、改 value 语义，还是换成另一类结构。

\subsubsection{LeetCode 76 最小覆盖子串}

先枚举所有连续窗口，明确左端、右端和窗口内容怎样决定答案。本题的模型是窗口内字符计数必须覆盖目标计数，满足后尽量收缩左端；暴力版会反复统计相邻窗口中相同的内容。慢点具体落在用 formed 记录满足需求的字符种类，避免每次全量比较，所以后面要证明左端或右端移动时不会漏掉可能答案。

窗口题只在条件有单调性或窗口长度固定时成立。本题用用 formed 记录满足需求的字符种类，避免每次全量比较维护窗口，含义是右端只带来一个新字符或新元素，左端只移走一个旧元素；如果输入条件改变导致状态不再单调，就必须换方法。放到“滑动窗口与双指针”这条主线里看，关键原则是：滑动窗口成立必须有单调性或固定长度边界。右端扩展让某个条件单调变好或变坏，左端收缩可以恢复合法性或缩短答案。若数组含负数、状态会来回震荡，就不能硬套窗口。

证明重点是排除左端。若某个左端被移过，它对应的窗口已经被完整检查，或继续保留只会让窗口更长、更不合法。本题的窗口内字符计数必须覆盖目标计数，满足后尽量收缩左端和用 formed 记录满足需求的字符种类，避免每次全量比较必须一起说明这一点。这道题的边界不能留到最后补丁式处理，实验时重点跑目标有重复字符、源串缺字符、最小窗口在开头或结尾、大小写，并打印左右端、窗口计数和答案。若左端发生回退，或者窗口失去合法性后仍更新答案，就能很快定位错误。

C++ 实现上，窗口计数若字符集固定，可以用 \code{std::array<int, 128>} 或 26 长度数组；字符集不固定再用哈希表。左闭右闭和左闭右开选一种坚持到底，避免窗口长度和删除位置错位。如果追问变成“若字符集固定，数组计数更快；若是 Unicode，哈希表更通用”，先判断原来的状态是否仍然覆盖新答案集合，再决定是微调边界、改 value 语义，还是换成另一类结构。

\subsubsection{LeetCode 167 两数之和 II 输入有序数组}

先枚举所有连续窗口，明确左端、右端和窗口内容怎样决定答案。本题的模型是有序数组中左右指针的和随指针移动单调变化；暴力版会反复统计相邻窗口中相同的内容。慢点具体落在当前和太小左指针右移，太大右指针左移，等于目标返回，所以后面要证明左端或右端移动时不会漏掉可能答案。

窗口题只在条件有单调性或窗口长度固定时成立。本题用当前和太小左指针右移，太大右指针左移，等于目标返回维护窗口，含义是右端只带来一个新字符或新元素，左端只移走一个旧元素；如果输入条件改变导致状态不再单调，就必须换方法。放到“滑动窗口与双指针”这条主线里看，关键原则是：滑动窗口成立必须有单调性或固定长度边界。右端扩展让某个条件单调变好或变坏，左端收缩可以恢复合法性或缩短答案。若数组含负数、状态会来回震荡，就不能硬套窗口。

证明重点是排除左端。若某个左端被移过，它对应的窗口已经被完整检查，或继续保留只会让窗口更长、更不合法。本题的有序数组中左右指针的和随指针移动单调变化和当前和太小左指针右移，太大右指针左移，等于目标返回必须一起说明这一点。这道题的边界不能留到最后补丁式处理，实验时重点跑两个数在边界、重复值、题目下标从一开始、保证有解，并打印左右端、窗口计数和答案。若左端发生回退，或者窗口失去合法性后仍更新答案，就能很快定位错误。

C++ 实现上，窗口计数若字符集固定，可以用 \code{std::array<int, 128>} 或 26 长度数组；字符集不固定再用哈希表。左闭右闭和左闭右开选一种坚持到底，避免窗口长度和删除位置错位。如果追问变成“无序版本通常用哈希表，不能直接套双指针”，先判断原来的状态是否仍然覆盖新答案集合，再决定是微调边界、改 value 语义，还是换成另一类结构。

\subsubsection{LeetCode 209 长度最小的子数组}

先枚举所有连续窗口，明确左端、右端和窗口内容怎样决定答案。本题的模型是正数数组让窗口和随右端增加、随左端减少；暴力版会反复统计相邻窗口中相同的内容。慢点具体落在右端扩展到满足目标后，不断左移缩短并更新答案，所以后面要证明左端或右端移动时不会漏掉可能答案。

窗口题只在条件有单调性或窗口长度固定时成立。本题用右端扩展到满足目标后，不断左移缩短并更新答案维护窗口，含义是右端只带来一个新字符或新元素，左端只移走一个旧元素；如果输入条件改变导致状态不再单调，就必须换方法。放到“滑动窗口与双指针”这条主线里看，关键原则是：滑动窗口成立必须有单调性或固定长度边界。右端扩展让某个条件单调变好或变坏，左端收缩可以恢复合法性或缩短答案。若数组含负数、状态会来回震荡，就不能硬套窗口。

证明重点是排除左端。若某个左端被移过，它对应的窗口已经被完整检查，或继续保留只会让窗口更长、更不合法。本题的正数数组让窗口和随右端增加、随左端减少和右端扩展到满足目标后，不断左移缩短并更新答案必须一起说明这一点。这道题的边界不能留到最后补丁式处理，实验时重点跑不存在答案、单元素达到目标、全正是必要条件，并打印左右端、窗口计数和答案。若左端发生回退，或者窗口失去合法性后仍更新答案，就能很快定位错误。

C++ 实现上，窗口计数若字符集固定，可以用 \code{std::array<int, 128>} 或 26 长度数组；字符集不固定再用哈希表。左闭右闭和左闭右开选一种坚持到底，避免窗口长度和删除位置错位。如果追问变成“若含负数，需要前缀和加单调队列”，先判断原来的状态是否仍然覆盖新答案集合，再决定是微调边界、改 value 语义，还是换成另一类结构。

\subsubsection{LeetCode 438 找到字符串中所有字母异位词}

先枚举所有连续窗口，明确左端、右端和窗口内容怎样决定答案。本题的模型是固定长度窗口内字符频次等于目标频次即为答案；暴力版会反复统计相邻窗口中相同的内容。慢点具体落在右进左出维护计数，窗口长度达到目标长度后检查差异，所以后面要证明左端或右端移动时不会漏掉可能答案。

窗口题只在条件有单调性或窗口长度固定时成立。本题用右进左出维护计数，窗口长度达到目标长度后检查差异维护窗口，含义是右端只带来一个新字符或新元素，左端只移走一个旧元素；如果输入条件改变导致状态不再单调，就必须换方法。放到“滑动窗口与双指针”这条主线里看，关键原则是：滑动窗口成立必须有单调性或固定长度边界。右端扩展让某个条件单调变好或变坏，左端收缩可以恢复合法性或缩短答案。若数组含负数、状态会来回震荡，就不能硬套窗口。

证明重点是排除左端。若某个左端被移过，它对应的窗口已经被完整检查，或继续保留只会让窗口更长、更不合法。本题的固定长度窗口内字符频次等于目标频次即为答案和右进左出维护计数，窗口长度达到目标长度后检查差异必须一起说明这一点。这道题的边界不能留到最后补丁式处理，实验时重点跑目标比源串长、重复字符、字符集固定、答案重叠，并打印左右端、窗口计数和答案。若左端发生回退，或者窗口失去合法性后仍更新答案，就能很快定位错误。

C++ 实现上，窗口计数若字符集固定，可以用 \code{std::array<int, 128>} 或 26 长度数组；字符集不固定再用哈希表。左闭右闭和左闭右开选一种坚持到底，避免窗口长度和删除位置错位。如果追问变成“维护差异计数可以避免每次比较整个数组”，先判断原来的状态是否仍然覆盖新答案集合，再决定是微调边界、改 value 语义，还是换成另一类结构。

\subsubsection{LeetCode 567 字符串的排列}

先枚举所有连续窗口，明确左端、右端和窗口内容怎样决定答案。本题的模型是目标串任意排列出现在源串中，等价于固定长度窗口频次相同；暴力版会反复统计相邻窗口中相同的内容。慢点具体落在窗口长度固定为目标长度，右进左出维护字符计数差异，所以后面要证明左端或右端移动时不会漏掉可能答案。

窗口题只在条件有单调性或窗口长度固定时成立。本题用窗口长度固定为目标长度，右进左出维护字符计数差异维护窗口，含义是右端只带来一个新字符或新元素，左端只移走一个旧元素；如果输入条件改变导致状态不再单调，就必须换方法。放到“滑动窗口与双指针”这条主线里看，关键原则是：滑动窗口成立必须有单调性或固定长度边界。右端扩展让某个条件单调变好或变坏，左端收缩可以恢复合法性或缩短答案。若数组含负数、状态会来回震荡，就不能硬套窗口。

证明重点是排除左端。若某个左端被移过，它对应的窗口已经被完整检查，或继续保留只会让窗口更长、更不合法。本题的目标串任意排列出现在源串中，等价于固定长度窗口频次相同和窗口长度固定为目标长度，右进左出维护字符计数差异必须一起说明这一点。这道题的边界不能留到最后补丁式处理，实验时重点跑目标比源串长、重复字符、字符集固定、窗口重叠，并打印左右端、窗口计数和答案。若左端发生回退，或者窗口失去合法性后仍更新答案，就能很快定位错误。

C++ 实现上，窗口计数若字符集固定，可以用 \code{std::array<int, 128>} 或 26 长度数组；字符集不固定再用哈希表。左闭右闭和左闭右开选一种坚持到底，避免窗口长度和删除位置错位。如果追问变成“与找到所有异位词同型，只是返回布尔值”，先判断原来的状态是否仍然覆盖新答案集合，再决定是微调边界、改 value 语义，还是换成另一类结构。

\subsubsection{LeetCode 862 和至少为 K 的最短子数组}

先枚举所有连续窗口，明确左端、右端和窗口内容怎样决定答案。本题的模型是数组允许负数，普通滑动窗口的和不再随左端移动单调变化；暴力版会反复统计相邻窗口中相同的内容。慢点具体落在在前缀和数组上用单调队列维护可能成为最优左端的下标，所以后面要证明左端或右端移动时不会漏掉可能答案。

窗口题只在条件有单调性或窗口长度固定时成立。本题用在前缀和数组上用单调队列维护可能成为最优左端的下标维护窗口，含义是右端只带来一个新字符或新元素，左端只移走一个旧元素；如果输入条件改变导致状态不再单调，就必须换方法。放到“滑动窗口与双指针”这条主线里看，关键原则是：滑动窗口成立必须有单调性或固定长度边界。右端扩展让某个条件单调变好或变坏，左端收缩可以恢复合法性或缩短答案。若数组含负数、状态会来回震荡，就不能硬套窗口。

证明重点是排除左端。若某个左端被移过，它对应的窗口已经被完整检查，或继续保留只会让窗口更长、更不合法。本题的数组允许负数，普通滑动窗口的和不再随左端移动单调变化和在前缀和数组上用单调队列维护可能成为最优左端的下标必须一起说明这一点。这道题的边界不能留到最后补丁式处理，实验时重点跑负数、答案不存在、K 为负或零、队尾支配关系，并打印左右端、窗口计数和答案。若左端发生回退，或者窗口失去合法性后仍更新答案，就能很快定位错误。

C++ 实现上，窗口计数若字符集固定，可以用 \code{std::array<int, 128>} 或 26 长度数组；字符集不固定再用哈希表。左闭右闭和左闭右开选一种坚持到底，避免窗口长度和删除位置错位。如果追问变成“这是从窗口题升级到前缀和加单调队列的代表”，先判断原来的状态是否仍然覆盖新答案集合，再决定是微调边界、改 value 语义，还是换成另一类结构。

\subsubsection{LeetCode 1438 绝对差不超过限制的最长连续子数组}

先枚举所有连续窗口，明确左端、右端和窗口内容怎样决定答案。本题的模型是窗口合法性取决于窗口内最大值和最小值之差；暴力版会反复统计相邻窗口中相同的内容。慢点具体落在两个单调队列分别维护最大和最小，超过限制时移动左端并清理过期下标，所以后面要证明左端或右端移动时不会漏掉可能答案。

窗口题只在条件有单调性或窗口长度固定时成立。本题用两个单调队列分别维护最大和最小，超过限制时移动左端并清理过期下标维护窗口，含义是右端只带来一个新字符或新元素，左端只移走一个旧元素；如果输入条件改变导致状态不再单调，就必须换方法。放到“滑动窗口与双指针”这条主线里看，关键原则是：滑动窗口成立必须有单调性或固定长度边界。右端扩展让某个条件单调变好或变坏，左端收缩可以恢复合法性或缩短答案。若数组含负数、状态会来回震荡，就不能硬套窗口。

证明重点是排除左端。若某个左端被移过，它对应的窗口已经被完整检查，或继续保留只会让窗口更长、更不合法。本题的窗口合法性取决于窗口内最大值和最小值之差和两个单调队列分别维护最大和最小，超过限制时移动左端并清理过期下标必须一起说明这一点。这道题的边界不能留到最后补丁式处理，实验时重点跑limit 为零、重复元素、最大最小同时过期、窗口收缩多次，并打印左右端、窗口计数和答案。若左端发生回退，或者窗口失去合法性后仍更新答案，就能很快定位错误。

C++ 实现上，窗口计数若字符集固定，可以用 \code{std::array<int, 128>} 或 26 长度数组；字符集不固定再用哈希表。左闭右闭和左闭右开选一种坚持到底，避免窗口长度和删除位置错位。如果追问变成“它说明窗口状态可以是动态极值，不只是频次或总和”，先判断原来的状态是否仍然覆盖新答案集合，再决定是微调边界、改 value 语义，还是换成另一类结构。

\subsection{前缀和与哈希表工作坊}

先用双层枚举区间，再把区间条件改写成两个前缀状态的关系。哈希表保存历史前缀的次数、最早位置或存在性，具体保存什么由题目目标决定。

\subsubsection{LeetCode 1 两数之和}

先用双层循环枚举区间或候选对，再把条件写成历史状态和当前状态的关系。本题的模型是当前元素只需要寻找唯一补数，历史值到下标的映射就是最小状态；暴力版的问题不是不会找，而是每个当前位置都回头扫太多历史。优化目标就是把先查再插入，避免同一个下标被用两次；若数组有序才考虑双指针变成一次查表或一次状态读取。

哈希表保存的是当前下标之前的历史，不是随手放进去的缓存。本题需要先查再插入，避免同一个下标被用两次；若数组有序才考虑双指针，所以要先决定 value 是次数、最早位置、最近位置还是布尔值。这个决定直接改变答案类型。放到“前缀和与哈希表”这条主线里看，关键原则是：前缀和把区间问题改写成两个前缀状态的差；哈希表把寻找匹配前缀改成查表。频次表、最早位置表、最近位置表对应不同目标：计数、最长区间、最短区间不能混用。

证明回到代数关系：任意合法答案都能对应到一个当前状态和一个历史状态；算法查到的历史状态也一定构成合法答案。本题的当前元素只需要寻找唯一补数，历史值到下标的映射就是最小状态说明要找什么，先查再插入，避免同一个下标被用两次；若数组有序才考虑双指针说明历史表怎样保存它。这道题的边界不能留到最后补丁式处理，实验时把重复数字、目标为偶数且补数等于当前值、没有答案时返回空结果加入对拍集合，用平方暴力和优化版比较。失败时先看初始历史状态、查询更新顺序和哈希表 value 类型。

C++ 实现上，哈希表查询优先用 \code{find}，只在确实要插入或累加时使用 \code{operator[]}。前缀和、区间和、计数乘法常需要 \code{long long}；模运算还要处理负余数归一化。如果追问变成“若改成返回所有不重复数对，要先排序或用频次表处理重复”，先判断原来的状态是否仍然覆盖新答案集合，再决定是微调边界、改 value 语义，还是换成另一类结构。

\subsubsection{LeetCode 49 字母异位词分组}

先用双层循环枚举区间或候选对，再把条件写成历史状态和当前状态的关系。本题的模型是异位词共享同一个规范表示，可以是排序字符串或字符计数；暴力版的问题不是不会找，而是每个当前位置都回头扫太多历史。优化目标就是把哈希表按规范键分桶，避免两两比较字符串变成一次查表或一次状态读取。

哈希表保存的是当前下标之前的历史，不是随手放进去的缓存。本题需要哈希表按规范键分桶，避免两两比较字符串，所以要先决定 value 是次数、最早位置、最近位置还是布尔值。这个决定直接改变答案类型。放到“前缀和与哈希表”这条主线里看，关键原则是：前缀和把区间问题改写成两个前缀状态的差；哈希表把寻找匹配前缀改成查表。频次表、最早位置表、最近位置表对应不同目标：计数、最长区间、最短区间不能混用。

证明回到代数关系：任意合法答案都能对应到一个当前状态和一个历史状态；算法查到的历史状态也一定构成合法答案。本题的异位词共享同一个规范表示，可以是排序字符串或字符计数说明要找什么，哈希表按规范键分桶，避免两两比较字符串说明历史表怎样保存它。这道题的边界不能留到最后补丁式处理，实验时把空字符串、重复单词、字符集不止小写字母、计数键必须无歧义加入对拍集合，用平方暴力和优化版比较。失败时先看初始历史状态、查询更新顺序和哈希表 value 类型。

C++ 实现上，哈希表查询优先用 \code{find}，只在确实要插入或累加时使用 \code{operator[]}。前缀和、区间和、计数乘法常需要 \code{long long}；模运算还要处理负余数归一化。如果追问变成“若字符串很长且字符集固定，计数键通常比排序键更稳定”，先判断原来的状态是否仍然覆盖新答案集合，再决定是微调边界、改 value 语义，还是换成另一类结构。

\subsubsection{LeetCode 128 最长连续序列}

先用双层循环枚举区间或候选对，再把条件写成历史状态和当前状态的关系。本题的模型是连续段只需要从没有前驱的数开始扩展；暴力版的问题不是不会找，而是每个当前位置都回头扫太多历史。优化目标就是把哈希集合判断 x 减一 是否存在，跳过非起点避免重复扫描变成一次查表或一次状态读取。

哈希表保存的是当前下标之前的历史，不是随手放进去的缓存。本题需要哈希集合判断 x 减一 是否存在，跳过非起点避免重复扫描，所以要先决定 value 是次数、最早位置、最近位置还是布尔值。这个决定直接改变答案类型。放到“前缀和与哈希表”这条主线里看，关键原则是：前缀和把区间问题改写成两个前缀状态的差；哈希表把寻找匹配前缀改成查表。频次表、最早位置表、最近位置表对应不同目标：计数、最长区间、最短区间不能混用。

证明回到代数关系：任意合法答案都能对应到一个当前状态和一个历史状态；算法查到的历史状态也一定构成合法答案。本题的连续段只需要从没有前驱的数开始扩展说明要找什么，哈希集合判断 x 减一 是否存在，跳过非起点避免重复扫描说明历史表怎样保存它。这道题的边界不能留到最后补丁式处理，实验时把重复元素、负数、空数组、整数边界加入对拍集合，用平方暴力和优化版比较。失败时先看初始历史状态、查询更新顺序和哈希表 value 类型。

C++ 实现上，哈希表查询优先用 \code{find}，只在确实要插入或累加时使用 \code{operator[]}。前缀和、区间和、计数乘法常需要 \code{long long}；模运算还要处理负余数归一化。如果追问变成“若数据流动态插入，可以用并查集或区间合并维护长度”，先判断原来的状态是否仍然覆盖新答案集合，再决定是微调边界、改 value 语义，还是换成另一类结构。

\subsubsection{LeetCode 202 快乐数}

先用双层循环枚举区间或候选对，再把条件写成历史状态和当前状态的关系。本题的模型是反复变换数字会进入一或循环，状态空间最终有限；暴力版的问题不是不会找，而是每个当前位置都回头扫太多历史。优化目标就是把用集合检测重复状态，或用快慢指针看成函数图变成一次查表或一次状态读取。

哈希表保存的是当前下标之前的历史，不是随手放进去的缓存。本题需要用集合检测重复状态，或用快慢指针看成函数图，所以要先决定 value 是次数、最早位置、最近位置还是布尔值。这个决定直接改变答案类型。放到“前缀和与哈希表”这条主线里看，关键原则是：前缀和把区间问题改写成两个前缀状态的差；哈希表把寻找匹配前缀改成查表。频次表、最早位置表、最近位置表对应不同目标：计数、最长区间、最短区间不能混用。

证明回到代数关系：任意合法答案都能对应到一个当前状态和一个历史状态；算法查到的历史状态也一定构成合法答案。本题的反复变换数字会进入一或循环，状态空间最终有限说明要找什么，用集合检测重复状态，或用快慢指针看成函数图说明历史表怎样保存它。这道题的边界不能留到最后补丁式处理，实验时把输入为一、进入循环、位平方和函数、负数不在题意内加入对拍集合，用平方暴力和优化版比较。失败时先看初始历史状态、查询更新顺序和哈希表 value 类型。

C++ 实现上，哈希表查询优先用 \code{find}，只在确实要插入或累加时使用 \code{operator[]}。前缀和、区间和、计数乘法常需要 \code{long long}；模运算还要处理负余数归一化。如果追问变成“快慢指针版本能训练把数值迭代看成链表”，先判断原来的状态是否仍然覆盖新答案集合，再决定是微调边界、改 value 语义，还是换成另一类结构。

\subsubsection{LeetCode 217 存在重复元素}

先用双层循环枚举区间或候选对，再把条件写成历史状态和当前状态的关系。本题的模型是判断是否出现过是哈希集合最直接的职责；暴力版的问题不是不会找，而是每个当前位置都回头扫太多历史。优化目标就是把扫描时若插入前已经存在则立即返回真变成一次查表或一次状态读取。

哈希表保存的是当前下标之前的历史，不是随手放进去的缓存。本题需要扫描时若插入前已经存在则立即返回真，所以要先决定 value 是次数、最早位置、最近位置还是布尔值。这个决定直接改变答案类型。放到“前缀和与哈希表”这条主线里看，关键原则是：前缀和把区间问题改写成两个前缀状态的差；哈希表把寻找匹配前缀改成查表。频次表、最早位置表、最近位置表对应不同目标：计数、最长区间、最短区间不能混用。

证明回到代数关系：任意合法答案都能对应到一个当前状态和一个历史状态；算法查到的历史状态也一定构成合法答案。本题的判断是否出现过是哈希集合最直接的职责说明要找什么，扫描时若插入前已经存在则立即返回真说明历史表怎样保存它。这道题的边界不能留到最后补丁式处理，实验时把空数组、全唯一、重复在末尾、值域很小加入对拍集合，用平方暴力和优化版比较。失败时先看初始历史状态、查询更新顺序和哈希表 value 类型。

C++ 实现上，哈希表查询优先用 \code{find}，只在确实要插入或累加时使用 \code{operator[]}。前缀和、区间和、计数乘法常需要 \code{long long}；模运算还要处理负余数归一化。如果追问变成“若允许修改输入，排序后看相邻可以降低额外空间”，先判断原来的状态是否仍然覆盖新答案集合，再决定是微调边界、改 value 语义，还是换成另一类结构。

\subsubsection{LeetCode 242 有效的字母异位词}

先用双层循环枚举区间或候选对，再把条件写成历史状态和当前状态的关系。本题的模型是两个字符串是异位词等价于每个字符出现次数相同；暴力版的问题不是不会找，而是每个当前位置都回头扫太多历史。优化目标就是把固定小写字母集时用定长计数数组，字符集未知时用哈希表变成一次查表或一次状态读取。

哈希表保存的是当前下标之前的历史，不是随手放进去的缓存。本题需要固定小写字母集时用定长计数数组，字符集未知时用哈希表，所以要先决定 value 是次数、最早位置、最近位置还是布尔值。这个决定直接改变答案类型。放到“前缀和与哈希表”这条主线里看，关键原则是：前缀和把区间问题改写成两个前缀状态的差；哈希表把寻找匹配前缀改成查表。频次表、最早位置表、最近位置表对应不同目标：计数、最长区间、最短区间不能混用。

证明回到代数关系：任意合法答案都能对应到一个当前状态和一个历史状态；算法查到的历史状态也一定构成合法答案。本题的两个字符串是异位词等价于每个字符出现次数相同说明要找什么，固定小写字母集时用定长计数数组，字符集未知时用哈希表说明历史表怎样保存它。这道题的边界不能留到最后补丁式处理，实验时把长度不同、空串、Unicode 字符、计数减到负数加入对拍集合，用平方暴力和优化版比较。失败时先看初始历史状态、查询更新顺序和哈希表 value 类型。

C++ 实现上，哈希表查询优先用 \code{find}，只在确实要插入或累加时使用 \code{operator[]}。前缀和、区间和、计数乘法常需要 \code{long long}；模运算还要处理负余数归一化。如果追问变成“异位词分组把这个二元比较扩展成规范 key 分桶”，先判断原来的状态是否仍然覆盖新答案集合，再决定是微调边界、改 value 语义，还是换成另一类结构。

\subsubsection{LeetCode 325 和等于 k 的最长子数组长度}

先用双层循环枚举区间或候选对，再把条件写成历史状态和当前状态的关系。本题的模型是最长区间要求在同一前缀差关系下尽量拉大左右端距离；暴力版的问题不是不会找，而是每个当前位置都回头扫太多历史。优化目标就是把哈希表保存每个前缀和最早出现下标，当前前缀查询 prefix 减 k变成一次查表或一次状态读取。

哈希表保存的是当前下标之前的历史，不是随手放进去的缓存。本题需要哈希表保存每个前缀和最早出现下标，当前前缀查询 prefix 减 k，所以要先决定 value 是次数、最早位置、最近位置还是布尔值。这个决定直接改变答案类型。放到“前缀和与哈希表”这条主线里看，关键原则是：前缀和把区间问题改写成两个前缀状态的差；哈希表把寻找匹配前缀改成查表。频次表、最早位置表、最近位置表对应不同目标：计数、最长区间、最短区间不能混用。

证明回到代数关系：任意合法答案都能对应到一个当前状态和一个历史状态；算法查到的历史状态也一定构成合法答案。本题的最长区间要求在同一前缀差关系下尽量拉大左右端距离说明要找什么，哈希表保存每个前缀和最早出现下标，当前前缀查询 prefix 减 k说明历史表怎样保存它。这道题的边界不能留到最后补丁式处理，实验时把从下标零开始、负数、前缀重复时不能覆盖最早位置、答案不存在加入对拍集合，用平方暴力和优化版比较。失败时先看初始历史状态、查询更新顺序和哈希表 value 类型。

C++ 实现上，哈希表查询优先用 \code{find}，只在确实要插入或累加时使用 \code{operator[]}。前缀和、区间和、计数乘法常需要 \code{long long}；模运算还要处理负余数归一化。如果追问变成“计数版本保存频次，最长版本保存最早位置，这是核心差别”，先判断原来的状态是否仍然覆盖新答案集合，再决定是微调边界、改 value 语义，还是换成另一类结构。

\subsubsection{LeetCode 454 四数相加 II}

先用双层循环枚举区间或候选对，再把条件写成历史状态和当前状态的关系。本题的模型是四个数组的和为零可以拆成两个二数组和互为相反数；暴力版的问题不是不会找，而是每个当前位置都回头扫太多历史。优化目标就是把先统计 A 和 B 的所有两数和频次，再枚举 C 和 D 查询相反数变成一次查表或一次状态读取。

哈希表保存的是当前下标之前的历史，不是随手放进去的缓存。本题需要先统计 A 和 B 的所有两数和频次，再枚举 C 和 D 查询相反数，所以要先决定 value 是次数、最早位置、最近位置还是布尔值。这个决定直接改变答案类型。放到“前缀和与哈希表”这条主线里看，关键原则是：前缀和把区间问题改写成两个前缀状态的差；哈希表把寻找匹配前缀改成查表。频次表、最早位置表、最近位置表对应不同目标：计数、最长区间、最短区间不能混用。

证明回到代数关系：任意合法答案都能对应到一个当前状态和一个历史状态；算法查到的历史状态也一定构成合法答案。本题的四个数组的和为零可以拆成两个二数组和互为相反数说明要找什么，先统计 A 和 B 的所有两数和频次，再枚举 C 和 D 查询相反数说明历史表怎样保存它。这道题的边界不能留到最后补丁式处理，实验时把重复值、负数、答案数量很大、两数和范围加入对拍集合，用平方暴力和优化版比较。失败时先看初始历史状态、查询更新顺序和哈希表 value 类型。

C++ 实现上，哈希表查询优先用 \code{find}，只在确实要插入或累加时使用 \code{operator[]}。前缀和、区间和、计数乘法常需要 \code{long long}；模运算还要处理负余数归一化。如果追问变成“这是 meet-in-the-middle 思想，比四层枚举少两个维度”，先判断原来的状态是否仍然覆盖新答案集合，再决定是微调边界、改 value 语义，还是换成另一类结构。

\subsubsection{LeetCode 523 连续的子数组和}

先用双层循环枚举区间或候选对，再把条件写成历史状态和当前状态的关系。本题的模型是长度至少二的子数组和为 k 的倍数，等价于两个前缀和同余；暴力版的问题不是不会找，而是每个当前位置都回头扫太多历史。优化目标就是把哈希表保存某个余数最早出现下标，当前余数若重复且距离至少二则成立变成一次查表或一次状态读取。

哈希表保存的是当前下标之前的历史，不是随手放进去的缓存。本题需要哈希表保存某个余数最早出现下标，当前余数若重复且距离至少二则成立，所以要先决定 value 是次数、最早位置、最近位置还是布尔值。这个决定直接改变答案类型。放到“前缀和与哈希表”这条主线里看，关键原则是：前缀和把区间问题改写成两个前缀状态的差；哈希表把寻找匹配前缀改成查表。频次表、最早位置表、最近位置表对应不同目标：计数、最长区间、最短区间不能混用。

证明回到代数关系：任意合法答案都能对应到一个当前状态和一个历史状态；算法查到的历史状态也一定构成合法答案。本题的长度至少二的子数组和为 k 的倍数，等价于两个前缀和同余说明要找什么，哈希表保存某个余数最早出现下标，当前余数若重复且距离至少二则成立说明历史表怎样保存它。这道题的边界不能留到最后补丁式处理，实验时把k 为零、负数余数、长度限制、从开头开始的区间加入对拍集合，用平方暴力和优化版比较。失败时先看初始历史状态、查询更新顺序和哈希表 value 类型。

C++ 实现上，哈希表查询优先用 \code{find}，只在确实要插入或累加时使用 \code{operator[]}。前缀和、区间和、计数乘法常需要 \code{long long}；模运算还要处理负余数归一化。如果追问变成“保存最早下标是为了满足长度条件并保留最长距离”，先判断原来的状态是否仍然覆盖新答案集合，再决定是微调边界、改 value 语义，还是换成另一类结构。

\subsubsection{LeetCode 525 连续数组}

先用双层循环枚举区间或候选对，再把条件写成历史状态和当前状态的关系。本题的模型是把零看成负一后，零和一数量相等等价于前缀和相同；暴力版的问题不是不会找，而是每个当前位置都回头扫太多历史。优化目标就是把哈希表保存每个前缀和最早下标，重复出现时更新最长长度变成一次查表或一次状态读取。

哈希表保存的是当前下标之前的历史，不是随手放进去的缓存。本题需要哈希表保存每个前缀和最早下标，重复出现时更新最长长度，所以要先决定 value 是次数、最早位置、最近位置还是布尔值。这个决定直接改变答案类型。放到“前缀和与哈希表”这条主线里看，关键原则是：前缀和把区间问题改写成两个前缀状态的差；哈希表把寻找匹配前缀改成查表。频次表、最早位置表、最近位置表对应不同目标：计数、最长区间、最短区间不能混用。

证明回到代数关系：任意合法答案都能对应到一个当前状态和一个历史状态；算法查到的历史状态也一定构成合法答案。本题的把零看成负一后，零和一数量相等等价于前缀和相同说明要找什么，哈希表保存每个前缀和最早下标，重复出现时更新最长长度说明历史表怎样保存它。这道题的边界不能留到最后补丁式处理，实验时把全零、全一、从下标零开始、不要覆盖最早位置加入对拍集合，用平方暴力和优化版比较。失败时先看初始历史状态、查询更新顺序和哈希表 value 类型。

C++ 实现上，哈希表查询优先用 \code{find}，只在确实要插入或累加时使用 \code{operator[]}。前缀和、区间和、计数乘法常需要 \code{long long}；模运算还要处理负余数归一化。如果追问变成“这是前缀和最早位置语义，不是频次语义”，先判断原来的状态是否仍然覆盖新答案集合，再决定是微调边界、改 value 语义，还是换成另一类结构。

\subsubsection{LeetCode 560 和为 K 的子数组}

先用双层循环枚举区间或候选对，再把条件写成历史状态和当前状态的关系。本题的模型是当前前缀和减去目标值，就是需要匹配的历史前缀；暴力版的问题不是不会找，而是每个当前位置都回头扫太多历史。优化目标就是把哈希表保存前缀和出现次数，先查询再更新变成一次查表或一次状态读取。

哈希表保存的是当前下标之前的历史，不是随手放进去的缓存。本题需要哈希表保存前缀和出现次数，先查询再更新，所以要先决定 value 是次数、最早位置、最近位置还是布尔值。这个决定直接改变答案类型。放到“前缀和与哈希表”这条主线里看，关键原则是：前缀和把区间问题改写成两个前缀状态的差；哈希表把寻找匹配前缀改成查表。频次表、最早位置表、最近位置表对应不同目标：计数、最长区间、最短区间不能混用。

证明回到代数关系：任意合法答案都能对应到一个当前状态和一个历史状态；算法查到的历史状态也一定构成合法答案。本题的当前前缀和减去目标值，就是需要匹配的历史前缀说明要找什么，哈希表保存前缀和出现次数，先查询再更新说明历史表怎样保存它。这道题的边界不能留到最后补丁式处理，实验时把负数、目标为零、从下标零开始的区间、重复前缀加入对拍集合，用平方暴力和优化版比较。失败时先看初始历史状态、查询更新顺序和哈希表 value 类型。

C++ 实现上，哈希表查询优先用 \code{find}，只在确实要插入或累加时使用 \code{operator[]}。前缀和、区间和、计数乘法常需要 \code{long long}；模运算还要处理负余数归一化。如果追问变成“若要求最长区间，哈希表应保存最早位置而不是次数”，先判断原来的状态是否仍然覆盖新答案集合，再决定是微调边界、改 value 语义，还是换成另一类结构。

\subsubsection{LeetCode 974 和可被 K 整除的子数组}

先用双层循环枚举区间或候选对，再把条件写成历史状态和当前状态的关系。本题的模型是区间和能被 K 整除等价于两个前缀和余数相同；暴力版的问题不是不会找，而是每个当前位置都回头扫太多历史。优化目标就是把统计每个余数出现次数，当前余数能和所有历史同余前缀组成答案变成一次查表或一次状态读取。

哈希表保存的是当前下标之前的历史，不是随手放进去的缓存。本题需要统计每个余数出现次数，当前余数能和所有历史同余前缀组成答案，所以要先决定 value 是次数、最早位置、最近位置还是布尔值。这个决定直接改变答案类型。放到“前缀和与哈希表”这条主线里看，关键原则是：前缀和把区间问题改写成两个前缀状态的差；哈希表把寻找匹配前缀改成查表。频次表、最早位置表、最近位置表对应不同目标：计数、最长区间、最短区间不能混用。

证明回到代数关系：任意合法答案都能对应到一个当前状态和一个历史状态；算法查到的历史状态也一定构成合法答案。本题的区间和能被 K 整除等价于两个前缀和余数相同说明要找什么，统计每个余数出现次数，当前余数能和所有历史同余前缀组成答案说明历史表怎样保存它。这道题的边界不能留到最后补丁式处理，实验时把负数余数归一化、K 为一、从开头开始、重复余数很多加入对拍集合，用平方暴力和优化版比较。失败时先看初始历史状态、查询更新顺序和哈希表 value 类型。

C++ 实现上，哈希表查询优先用 \code{find}，只在确实要插入或累加时使用 \code{operator[]}。前缀和、区间和、计数乘法常需要 \code{long long}；模运算还要处理负余数归一化。如果追问变成“如果问最长长度，保存最早位置而不是次数”，先判断原来的状态是否仍然覆盖新答案集合，再决定是微调边界、改 value 语义，还是换成另一类结构。

\subsubsection{LeetCode 1074 元素和为目标值的子矩阵数量}

先用双层循环枚举区间或候选对，再把条件写成历史状态和当前状态的关系。本题的模型是二维矩阵子区域和可以固定上下边界后转成一维子数组和；暴力版的问题不是不会找，而是每个当前位置都回头扫太多历史。优化目标就是把枚举行边界，压缩每列区间和，再用前缀和哈希统计目标和子数组变成一次查表或一次状态读取。

哈希表保存的是当前下标之前的历史，不是随手放进去的缓存。本题需要枚举行边界，压缩每列区间和，再用前缀和哈希统计目标和子数组，所以要先决定 value 是次数、最早位置、最近位置还是布尔值。这个决定直接改变答案类型。放到“前缀和与哈希表”这条主线里看，关键原则是：前缀和把区间问题改写成两个前缀状态的差；哈希表把寻找匹配前缀改成查表。频次表、最早位置表、最近位置表对应不同目标：计数、最长区间、最短区间不能混用。

证明回到代数关系：任意合法答案都能对应到一个当前状态和一个历史状态；算法查到的历史状态也一定构成合法答案。本题的二维矩阵子区域和可以固定上下边界后转成一维子数组和说明要找什么，枚举行边界，压缩每列区间和，再用前缀和哈希统计目标和子数组说明历史表怎样保存它。这道题的边界不能留到最后补丁式处理，实验时把单行、单列、负数、目标为零、矩阵较大加入对拍集合，用平方暴力和优化版比较。失败时先看初始历史状态、查询更新顺序和哈希表 value 类型。

C++ 实现上，哈希表查询优先用 \code{find}，只在确实要插入或累加时使用 \code{operator[]}。前缀和、区间和、计数乘法常需要 \code{long long}；模运算还要处理负余数归一化。如果追问变成“这是二维前缀思想和一维哈希计数的组合”，先判断原来的状态是否仍然覆盖新答案集合，再决定是微调边界、改 value 语义，还是换成另一类结构。

\subsubsection{LeetCode 1248 统计优美子数组}

先用双层循环枚举区间或候选对，再把条件写成历史状态和当前状态的关系。本题的模型是恰好 K 个奇数可转成前缀奇数计数差为 K；暴力版的问题不是不会找，而是每个当前位置都回头扫太多历史。优化目标就是把哈希表保存某个奇数计数出现次数，当前计数查询 count - K变成一次查表或一次状态读取。

哈希表保存的是当前下标之前的历史，不是随手放进去的缓存。本题需要哈希表保存某个奇数计数出现次数，当前计数查询 count - K，所以要先决定 value 是次数、最早位置、最近位置还是布尔值。这个决定直接改变答案类型。放到“前缀和与哈希表”这条主线里看，关键原则是：前缀和把区间问题改写成两个前缀状态的差；哈希表把寻找匹配前缀改成查表。频次表、最早位置表、最近位置表对应不同目标：计数、最长区间、最短区间不能混用。

证明回到代数关系：任意合法答案都能对应到一个当前状态和一个历史状态；算法查到的历史状态也一定构成合法答案。本题的恰好 K 个奇数可转成前缀奇数计数差为 K说明要找什么，哈希表保存某个奇数计数出现次数，当前计数查询 count - K说明历史表怎样保存它。这道题的边界不能留到最后补丁式处理，实验时把K 为零、全偶数、全奇数、从开头开始的区间加入对拍集合，用平方暴力和优化版比较。失败时先看初始历史状态、查询更新顺序和哈希表 value 类型。

C++ 实现上，哈希表查询优先用 \code{find}，只在确实要插入或累加时使用 \code{operator[]}。前缀和、区间和、计数乘法常需要 \code{long long}；模运算还要处理负余数归一化。如果追问变成“也可用至多 K 减至多 K-1 的滑动窗口”，先判断原来的状态是否仍然覆盖新答案集合，再决定是微调边界、改 value 语义，还是换成另一类结构。

\subsection{二分查找与答案空间工作坊}

先线性扫描或线性试答案，再找出能把候选一分为二的单调谓词。二分的核心不是 mid 写法，而是被排除的一半确实不可能包含答案。

\subsubsection{LeetCode 33 搜索旋转排序数组}

先用线性扫描或线性试答案保证正确，再找边界。本题的模型是旋转数组每次二分至少有一半保持有序；二分不是为了套循环，而是因为判断哪一半有序，再看目标是否落在有序半边范围内给出了一条能排除候选的边界线。若这条边界线说不清，宁愿先保留线性版本。

二分的状态不是数组本身，而是候选区间。本题用判断哪一半有序，再看目标是否落在有序半边范围内缩小区间；每次更新左右边界时，都要能说清被丢掉的一半为什么没有答案。边界不变量比 mid 公式更重要。放到“二分查找与答案空间”这条主线里看，关键原则是：二分前必须先定义谓词：某个位置是否已经满足边界条件，或者某个答案是否可行。数组二分找的是边界，答案二分找的是最小可行值或最大可行值。

证明二分时，不要只说复杂度是对数。要证明当判断哪一半有序，再看目标是否落在有序半边范围内成立或不成立时，被排除的一侧确实没有目标。本题的旋转数组每次二分至少有一半保持有序提供候选空间，循环只是在这个空间里保留答案。这道题的边界不能留到最后补丁式处理，实验时覆盖未旋转、长度一、目标不存在、边界等号、存在重复元素时退化，逐轮打印 left、mid、right 和谓词值。只要有一次被排除区间仍含答案，二分不变量就错了。

C++ 实现上，中点写成 \code{left + (right - left) / 2}；平方、乘法和总量比较避免溢出。答案二分最好把检查函数单独写出，让单调谓词可读。如果追问变成“若重复元素很多，需要在无法判断有序半边时收缩边界”，先判断原来的状态是否仍然覆盖新答案集合，再决定是微调边界、改 value 语义，还是换成另一类结构。

\subsubsection{LeetCode 34 排序数组中查找首尾位置}

先用线性扫描或线性试答案保证正确，再找边界。本题的模型是首尾位置可以拆成两个边界：第一个大于等于和第一个大于；二分不是为了套循环，而是因为不要找到一个目标后向两边线性扩展，否则重复元素很多时退化给出了一条能排除候选的边界线。若这条边界线说不清，宁愿先保留线性版本。

二分的状态不是数组本身，而是候选区间。本题用不要找到一个目标后向两边线性扩展，否则重复元素很多时退化缩小区间；每次更新左右边界时，都要能说清被丢掉的一半为什么没有答案。边界不变量比 mid 公式更重要。放到“二分查找与答案空间”这条主线里看，关键原则是：二分前必须先定义谓词：某个位置是否已经满足边界条件，或者某个答案是否可行。数组二分找的是边界，答案二分找的是最小可行值或最大可行值。

证明二分时，不要只说复杂度是对数。要证明当不要找到一个目标后向两边线性扩展，否则重复元素很多时退化成立或不成立时，被排除的一侧确实没有目标。本题的首尾位置可以拆成两个边界：第一个大于等于和第一个大于提供候选空间，循环只是在这个空间里保留答案。这道题的边界不能留到最后补丁式处理，实验时覆盖目标不存在、目标在开头或结尾、数组为空、全是目标，逐轮打印 left、mid、right 和谓词值。只要有一次被排除区间仍含答案，二分不变量就错了。

C++ 实现上，中点写成 \code{left + (right - left) / 2}；平方、乘法和总量比较避免溢出。答案二分最好把检查函数单独写出，让单调谓词可读。如果追问变成“这个模型能迁移到计数某值出现次数和插入区间边界”，先判断原来的状态是否仍然覆盖新答案集合，再决定是微调边界、改 value 语义，还是换成另一类结构。

\subsubsection{LeetCode 35 搜索插入位置}

先用线性扫描或线性试答案保证正确，再找边界。本题的模型是题目等价于寻找第一个大于等于目标的位置；二分不是为了套循环，而是因为统一用 lower bound 语义，不在相等时提前返回也能保持边界一致给出了一条能排除候选的边界线。若这条边界线说不清，宁愿先保留线性版本。

二分的状态不是数组本身，而是候选区间。本题用统一用 lower bound 语义，不在相等时提前返回也能保持边界一致缩小区间；每次更新左右边界时，都要能说清被丢掉的一半为什么没有答案。边界不变量比 mid 公式更重要。放到“二分查找与答案空间”这条主线里看，关键原则是：二分前必须先定义谓词：某个位置是否已经满足边界条件，或者某个答案是否可行。数组二分找的是边界，答案二分找的是最小可行值或最大可行值。

证明二分时，不要只说复杂度是对数。要证明当统一用 lower bound 语义，不在相等时提前返回也能保持边界一致成立或不成立时，被排除的一侧确实没有目标。本题的题目等价于寻找第一个大于等于目标的位置提供候选空间，循环只是在这个空间里保留答案。这道题的边界不能留到最后补丁式处理，实验时覆盖目标小于所有元素、目标大于所有元素、重复值、空数组，逐轮打印 left、mid、right 和谓词值。只要有一次被排除区间仍含答案，二分不变量就错了。

C++ 实现上，中点写成 \code{left + (right - left) / 2}；平方、乘法和总量比较避免溢出。答案二分最好把检查函数单独写出，让单调谓词可读。如果追问变成“手写二分时用左闭右开可直接返回 left”，先判断原来的状态是否仍然覆盖新答案集合，再决定是微调边界、改 value 语义，还是换成另一类结构。

\subsubsection{LeetCode 69 x 的平方根}

先用线性扫描或线性试答案保证正确，再找边界。本题的模型是答案是最大的整数 r，使得 r 的平方不超过 x；二分不是为了套循环，而是因为在整数范围上二分右边界，比较时用除法或 long long 避免平方溢出给出了一条能排除候选的边界线。若这条边界线说不清，宁愿先保留线性版本。

二分的状态不是数组本身，而是候选区间。本题用在整数范围上二分右边界，比较时用除法或 long long 避免平方溢出缩小区间；每次更新左右边界时，都要能说清被丢掉的一半为什么没有答案。边界不变量比 mid 公式更重要。放到“二分查找与答案空间”这条主线里看，关键原则是：二分前必须先定义谓词：某个位置是否已经满足边界条件，或者某个答案是否可行。数组二分找的是边界，答案二分找的是最小可行值或最大可行值。

证明二分时，不要只说复杂度是对数。要证明当在整数范围上二分右边界，比较时用除法或 long long 避免平方溢出成立或不成立时，被排除的一侧确实没有目标。本题的答案是最大的整数 r，使得 r 的平方不超过 x提供候选空间，循环只是在这个空间里保留答案。这道题的边界不能留到最后补丁式处理，实验时覆盖x 为零、一、非完全平方、接近 int 最大值，逐轮打印 left、mid、right 和谓词值。只要有一次被排除区间仍含答案，二分不变量就错了。

C++ 实现上，中点写成 \code{left + (right - left) / 2}；平方、乘法和总量比较避免溢出。答案二分最好把检查函数单独写出，让单调谓词可读。如果追问变成“若要求浮点精度，可以二分实数区间或用牛顿迭代”，先判断原来的状态是否仍然覆盖新答案集合，再决定是微调边界、改 value 语义，还是换成另一类结构。

\subsubsection{LeetCode 74 搜索二维矩阵}

先用线性扫描或线性试答案保证正确，再找边界。本题的模型是矩阵每行有序且下一行首元素大于上一行末元素，可视为一维有序数组；二分不是为了套循环，而是因为把下标 mid 映射到 row 和 col，再做标准二分给出了一条能排除候选的边界线。若这条边界线说不清，宁愿先保留线性版本。

二分的状态不是数组本身，而是候选区间。本题用把下标 mid 映射到 row 和 col，再做标准二分缩小区间；每次更新左右边界时，都要能说清被丢掉的一半为什么没有答案。边界不变量比 mid 公式更重要。放到“二分查找与答案空间”这条主线里看，关键原则是：二分前必须先定义谓词：某个位置是否已经满足边界条件，或者某个答案是否可行。数组二分找的是边界，答案二分找的是最小可行值或最大可行值。

证明二分时，不要只说复杂度是对数。要证明当把下标 mid 映射到 row 和 col，再做标准二分成立或不成立时，被排除的一侧确实没有目标。本题的矩阵每行有序且下一行首元素大于上一行末元素，可视为一维有序数组提供候选空间，循环只是在这个空间里保留答案。这道题的边界不能留到最后补丁式处理，实验时覆盖空矩阵、单行、单列、目标小于全部或大于全部，逐轮打印 left、mid、right 和谓词值。只要有一次被排除区间仍含答案，二分不变量就错了。

C++ 实现上，中点写成 \code{left + (right - left) / 2}；平方、乘法和总量比较避免溢出。答案二分最好把检查函数单独写出，让单调谓词可读。如果追问变成“若只是行列分别有序但行之间不连续，应改用右上角排除法”，先判断原来的状态是否仍然覆盖新答案集合，再决定是微调边界、改 value 语义，还是换成另一类结构。

\subsubsection{LeetCode 81 搜索旋转排序数组 II}

先用线性扫描或线性试答案保证正确，再找边界。本题的模型是旋转数组存在重复值时，二分可能无法判断哪一半有序；二分不是为了套循环，而是因为当左右和中点值相同导致信息不足时，收缩边界；否则仍按有序半边排除给出了一条能排除候选的边界线。若这条边界线说不清，宁愿先保留线性版本。

二分的状态不是数组本身，而是候选区间。本题用当左右和中点值相同导致信息不足时，收缩边界；否则仍按有序半边排除缩小区间；每次更新左右边界时，都要能说清被丢掉的一半为什么没有答案。边界不变量比 mid 公式更重要。放到“二分查找与答案空间”这条主线里看，关键原则是：二分前必须先定义谓词：某个位置是否已经满足边界条件，或者某个答案是否可行。数组二分找的是边界，答案二分找的是最小可行值或最大可行值。

证明二分时，不要只说复杂度是对数。要证明当当左右和中点值相同导致信息不足时，收缩边界；否则仍按有序半边排除成立或不成立时，被排除的一侧确实没有目标。本题的旋转数组存在重复值时，二分可能无法判断哪一半有序提供候选空间，循环只是在这个空间里保留答案。这道题的边界不能留到最后补丁式处理，实验时覆盖全相同元素、目标不存在、重复值跨过旋转点、退化到线性，逐轮打印 left、mid、right 和谓词值。只要有一次被排除区间仍含答案，二分不变量就错了。

C++ 实现上，中点写成 \code{left + (right - left) / 2}；平方、乘法和总量比较避免溢出。答案二分最好把检查函数单独写出，让单调谓词可读。如果追问变成“这题说明二分依赖信息量，重复值会削弱每次排除一半的能力”，先判断原来的状态是否仍然覆盖新答案集合，再决定是微调边界、改 value 语义，还是换成另一类结构。

\subsubsection{LeetCode 153 寻找旋转排序数组中的最小值}

先用线性扫描或线性试答案保证正确，再找边界。本题的模型是最小值是旋转断点，右端点可帮助判断中点在哪一段；二分不是为了套循环，而是因为若 mid 大于 right，最小值在右侧；否则在左侧含 mid给出了一条能排除候选的边界线。若这条边界线说不清，宁愿先保留线性版本。

二分的状态不是数组本身，而是候选区间。本题用若 mid 大于 right，最小值在右侧；否则在左侧含 mid缩小区间；每次更新左右边界时，都要能说清被丢掉的一半为什么没有答案。边界不变量比 mid 公式更重要。放到“二分查找与答案空间”这条主线里看，关键原则是：二分前必须先定义谓词：某个位置是否已经满足边界条件，或者某个答案是否可行。数组二分找的是边界，答案二分找的是最小可行值或最大可行值。

证明二分时，不要只说复杂度是对数。要证明当若 mid 大于 right，最小值在右侧；否则在左侧含 mid成立或不成立时，被排除的一侧确实没有目标。本题的最小值是旋转断点，右端点可帮助判断中点在哪一段提供候选空间，循环只是在这个空间里保留答案。这道题的边界不能留到最后补丁式处理，实验时覆盖未旋转、长度一、旋转一位、无重复条件，逐轮打印 left、mid、right 和谓词值。只要有一次被排除区间仍含答案，二分不变量就错了。

C++ 实现上，中点写成 \code{left + (right - left) / 2}；平方、乘法和总量比较避免溢出。答案二分最好把检查函数单独写出，让单调谓词可读。如果追问变成“有重复时遇到相等只能收缩 right，最坏退化”，先判断原来的状态是否仍然覆盖新答案集合，再决定是微调边界、改 value 语义，还是换成另一类结构。

\subsubsection{LeetCode 154 寻找旋转排序数组中的最小值 II}

先用线性扫描或线性试答案保证正确，再找边界。本题的模型是重复元素会让 mid 和 right 相等时无法判断最小值方向；二分不是为了套循环，而是因为相等时只能安全地丢掉一个 right，因为它不是唯一必要候选给出了一条能排除候选的边界线。若这条边界线说不清，宁愿先保留线性版本。

二分的状态不是数组本身，而是候选区间。本题用相等时只能安全地丢掉一个 right，因为它不是唯一必要候选缩小区间；每次更新左右边界时，都要能说清被丢掉的一半为什么没有答案。边界不变量比 mid 公式更重要。放到“二分查找与答案空间”这条主线里看，关键原则是：二分前必须先定义谓词：某个位置是否已经满足边界条件，或者某个答案是否可行。数组二分找的是边界，答案二分找的是最小可行值或最大可行值。

证明二分时，不要只说复杂度是对数。要证明当相等时只能安全地丢掉一个 right，因为它不是唯一必要候选成立或不成立时，被排除的一侧确实没有目标。本题的重复元素会让 mid 和 right 相等时无法判断最小值方向提供候选空间，循环只是在这个空间里保留答案。这道题的边界不能留到最后补丁式处理，实验时覆盖全相等、重复值跨断点、最小值出现多次，逐轮打印 left、mid、right 和谓词值。只要有一次被排除区间仍含答案，二分不变量就错了。

C++ 实现上，中点写成 \code{left + (right - left) / 2}；平方、乘法和总量比较避免溢出。答案二分最好把检查函数单独写出，让单调谓词可读。如果追问变成“这题说明二分依赖的信息不足时会退化，不是所有旋转题都严格对数”，先判断原来的状态是否仍然覆盖新答案集合，再决定是微调边界、改 value 语义，还是换成另一类结构。

\subsubsection{LeetCode 162 寻找峰值}

先用线性扫描或线性试答案保证正确，再找边界。本题的模型是比较 mid 和 mid 加一 的斜率能判断某侧一定存在峰值；二分不是为了套循环，而是因为上升则右侧有峰，下降则左侧含 mid 有峰给出了一条能排除候选的边界线。若这条边界线说不清，宁愿先保留线性版本。

二分的状态不是数组本身，而是候选区间。本题用上升则右侧有峰，下降则左侧含 mid 有峰缩小区间；每次更新左右边界时，都要能说清被丢掉的一半为什么没有答案。边界不变量比 mid 公式更重要。放到“二分查找与答案空间”这条主线里看，关键原则是：二分前必须先定义谓词：某个位置是否已经满足边界条件，或者某个答案是否可行。数组二分找的是边界，答案二分找的是最小可行值或最大可行值。

证明二分时，不要只说复杂度是对数。要证明当上升则右侧有峰，下降则左侧含 mid 有峰成立或不成立时，被排除的一侧确实没有目标。本题的比较 mid 和 mid 加一 的斜率能判断某侧一定存在峰值提供候选空间，循环只是在这个空间里保留答案。这道题的边界不能留到最后补丁式处理，实验时覆盖长度一、峰在边界、多个峰、相邻不相等假设，逐轮打印 left、mid、right 和谓词值。只要有一次被排除区间仍含答案，二分不变量就错了。

C++ 实现上，中点写成 \code{left + (right - left) / 2}；平方、乘法和总量比较避免溢出。答案二分最好把检查函数单独写出，让单调谓词可读。如果追问变成“二维峰值需要按行列最大值进行更复杂二分”，先判断原来的状态是否仍然覆盖新答案集合，再决定是微调边界、改 value 语义，还是换成另一类结构。

\subsubsection{LeetCode 240 搜索二维矩阵 II}

先用线性扫描或线性试答案保证正确，再找边界。本题的模型是右上角同时具有向左变小、向下变大的排除能力；二分不是为了套循环，而是因为当前值大于目标就左移，小于目标就下移给出了一条能排除候选的边界线。若这条边界线说不清，宁愿先保留线性版本。

二分的状态不是数组本身，而是候选区间。本题用当前值大于目标就左移，小于目标就下移缩小区间；每次更新左右边界时，都要能说清被丢掉的一半为什么没有答案。边界不变量比 mid 公式更重要。放到“二分查找与答案空间”这条主线里看，关键原则是：二分前必须先定义谓词：某个位置是否已经满足边界条件，或者某个答案是否可行。数组二分找的是边界，答案二分找的是最小可行值或最大可行值。

证明二分时，不要只说复杂度是对数。要证明当当前值大于目标就左移，小于目标就下移成立或不成立时，被排除的一侧确实没有目标。本题的右上角同时具有向左变小、向下变大的排除能力提供候选空间，循环只是在这个空间里保留答案。这道题的边界不能留到最后补丁式处理，实验时覆盖空矩阵、单行、单列、目标在角落、重复值，逐轮打印 left、mid、right 和谓词值。只要有一次被排除区间仍含答案，二分不变量就错了。

C++ 实现上，中点写成 \code{left + (right - left) / 2}；平方、乘法和总量比较避免溢出。答案二分最好把检查函数单独写出，让单调谓词可读。如果追问变成“从左上角无法一次排除一行或一列”，先判断原来的状态是否仍然覆盖新答案集合，再决定是微调边界、改 value 语义，还是换成另一类结构。

\subsubsection{LeetCode 410 分割数组的最大值}

先用线性扫描或线性试答案保证正确，再找边界。本题的模型是连续分成 m 段后最大段和越大，越容易完成分割；二分不是为了套循环，而是因为二分最大段和上限，检查函数贪心累计，超过上限就新开一段给出了一条能排除候选的边界线。若这条边界线说不清，宁愿先保留线性版本。

二分的状态不是数组本身，而是候选区间。本题用二分最大段和上限，检查函数贪心累计，超过上限就新开一段缩小区间；每次更新左右边界时，都要能说清被丢掉的一半为什么没有答案。边界不变量比 mid 公式更重要。放到“二分查找与答案空间”这条主线里看，关键原则是：二分前必须先定义谓词：某个位置是否已经满足边界条件，或者某个答案是否可行。数组二分找的是边界，答案二分找的是最小可行值或最大可行值。

证明二分时，不要只说复杂度是对数。要证明当二分最大段和上限，检查函数贪心累计，超过上限就新开一段成立或不成立时，被排除的一侧确实没有目标。本题的连续分成 m 段后最大段和越大，越容易完成分割提供候选空间，循环只是在这个空间里保留答案。这道题的边界不能留到最后补丁式处理，实验时覆盖m 为一、m 等于数组长度、单个元素超过猜测上限、总和溢出，逐轮打印 left、mid、right 和谓词值。只要有一次被排除区间仍含答案，二分不变量就错了。

C++ 实现上，中点写成 \code{left + (right - left) / 2}；平方、乘法和总量比较避免溢出。答案二分最好把检查函数单独写出，让单调谓词可读。如果追问变成“也可用 DP 求精确最优，但答案二分更适合大输入”，先判断原来的状态是否仍然覆盖新答案集合，再决定是微调边界、改 value 语义，还是换成另一类结构。

\subsubsection{LeetCode 540 有序数组中的单一元素}

先用线性扫描或线性试答案保证正确，再找边界。本题的模型是除一个元素外其余元素都成对出现，单一元素会改变配对下标的奇偶规律；二分不是为了套循环，而是因为把 mid 调整到偶数位，若 nums[mid] 等于 nums[mid+1]，单一元素在右侧，否则在左侧给出了一条能排除候选的边界线。若这条边界线说不清，宁愿先保留线性版本。

二分的状态不是数组本身，而是候选区间。本题用把 mid 调整到偶数位，若 nums[mid] 等于 nums[mid+1]，单一元素在右侧，否则在左侧缩小区间；每次更新左右边界时，都要能说清被丢掉的一半为什么没有答案。边界不变量比 mid 公式更重要。放到“二分查找与答案空间”这条主线里看，关键原则是：二分前必须先定义谓词：某个位置是否已经满足边界条件，或者某个答案是否可行。数组二分找的是边界，答案二分找的是最小可行值或最大可行值。

证明二分时，不要只说复杂度是对数。要证明当把 mid 调整到偶数位，若 nums[mid] 等于 nums[mid+1]，单一元素在右侧，否则在左侧成立或不成立时，被排除的一侧确实没有目标。本题的除一个元素外其余元素都成对出现，单一元素会改变配对下标的奇偶规律提供候选空间，循环只是在这个空间里保留答案。这道题的边界不能留到最后补丁式处理，实验时覆盖单一元素在开头或结尾、数组长度一、mid 加一越界，逐轮打印 left、mid、right 和谓词值。只要有一次被排除区间仍含答案，二分不变量就错了。

C++ 实现上，中点写成 \code{left + (right - left) / 2}；平方、乘法和总量比较避免溢出。答案二分最好把检查函数单独写出，让单调谓词可读。如果追问变成“异或也能求值，但二分利用有序性达到对数时间”，先判断原来的状态是否仍然覆盖新答案集合，再决定是微调边界、改 value 语义，还是换成另一类结构。

\subsubsection{LeetCode 704 二分查找}

先用线性扫描或线性试答案保证正确，再找边界。本题的模型是基础题的重点是固定区间语义；二分不是为了套循环，而是因为左闭右开写法能统一 lower bound 和存在性检查给出了一条能排除候选的边界线。若这条边界线说不清，宁愿先保留线性版本。

二分的状态不是数组本身，而是候选区间。本题用左闭右开写法能统一 lower bound 和存在性检查缩小区间；每次更新左右边界时，都要能说清被丢掉的一半为什么没有答案。边界不变量比 mid 公式更重要。放到“二分查找与答案空间”这条主线里看，关键原则是：二分前必须先定义谓词：某个位置是否已经满足边界条件，或者某个答案是否可行。数组二分找的是边界，答案二分找的是最小可行值或最大可行值。

证明二分时，不要只说复杂度是对数。要证明当左闭右开写法能统一 lower bound 和存在性检查成立或不成立时，被排除的一侧确实没有目标。本题的基础题的重点是固定区间语义提供候选空间，循环只是在这个空间里保留答案。这道题的边界不能留到最后补丁式处理，实验时覆盖空数组、长度一、目标在边界、目标不存在，逐轮打印 left、mid、right 和谓词值。只要有一次被排除区间仍含答案，二分不变量就错了。

C++ 实现上，中点写成 \code{left + (right - left) / 2}；平方、乘法和总量比较避免溢出。答案二分最好把检查函数单独写出，让单调谓词可读。如果追问变成“所有复杂二分都应该能退回到这道题的边界不变量”，先判断原来的状态是否仍然覆盖新答案集合，再决定是微调边界、改 value 语义，还是换成另一类结构。

\subsubsection{LeetCode 875 爱吃香蕉的珂珂}

先用线性扫描或线性试答案保证正确，再找边界。本题的模型是速度越大，总耗时越小，答案具有单调可行性；二分不是为了套循环，而是因为二分速度，检查函数用向上取整累加小时数给出了一条能排除候选的边界线。若这条边界线说不清，宁愿先保留线性版本。

二分的状态不是数组本身，而是候选区间。本题用二分速度，检查函数用向上取整累加小时数缩小区间；每次更新左右边界时，都要能说清被丢掉的一半为什么没有答案。边界不变量比 mid 公式更重要。放到“二分查找与答案空间”这条主线里看，关键原则是：二分前必须先定义谓词：某个位置是否已经满足边界条件，或者某个答案是否可行。数组二分找的是边界，答案二分找的是最小可行值或最大可行值。

证明二分时，不要只说复杂度是对数。要证明当二分速度，检查函数用向上取整累加小时数成立或不成立时，被排除的一侧确实没有目标。本题的速度越大，总耗时越小，答案具有单调可行性提供候选空间，循环只是在这个空间里保留答案。这道题的边界不能留到最后补丁式处理，实验时覆盖速度下界、堆很大、h 等于堆数、乘法加法溢出，逐轮打印 left、mid、right 和谓词值。只要有一次被排除区间仍含答案，二分不变量就错了。

C++ 实现上，中点写成 \code{left + (right - left) / 2}；平方、乘法和总量比较避免溢出。答案二分最好把检查函数单独写出，让单调谓词可读。如果追问变成“船运包裹和分割数组都是同类最小可行上限”，先判断原来的状态是否仍然覆盖新答案集合，再决定是微调边界、改 value 语义，还是换成另一类结构。

\subsubsection{LeetCode 1011 在 D 天内送达包裹}

先用线性扫描或线性试答案保证正确，再找边界。本题的模型是容量越大，需要天数越少，存在最小可行容量；二分不是为了套循环，而是因为下界是最大包裹，上界是总重量，检查函数按顺序贪心装船给出了一条能排除候选的边界线。若这条边界线说不清，宁愿先保留线性版本。

二分的状态不是数组本身，而是候选区间。本题用下界是最大包裹，上界是总重量，检查函数按顺序贪心装船缩小区间；每次更新左右边界时，都要能说清被丢掉的一半为什么没有答案。边界不变量比 mid 公式更重要。放到“二分查找与答案空间”这条主线里看，关键原则是：二分前必须先定义谓词：某个位置是否已经满足边界条件，或者某个答案是否可行。数组二分找的是边界，答案二分找的是最小可行值或最大可行值。

证明二分时，不要只说复杂度是对数。要证明当下界是最大包裹，上界是总重量，检查函数按顺序贪心装船成立或不成立时，被排除的一侧确实没有目标。本题的容量越大，需要天数越少，存在最小可行容量提供候选空间，循环只是在这个空间里保留答案。这道题的边界不能留到最后补丁式处理，实验时覆盖D 为一、D 等于包裹数、单个包裹超过猜测容量，逐轮打印 left、mid、right 和谓词值。只要有一次被排除区间仍含答案，二分不变量就错了。

C++ 实现上，中点写成 \code{left + (right - left) / 2}；平方、乘法和总量比较避免溢出。答案二分最好把检查函数单独写出，让单调谓词可读。如果追问变成“它和分割数组最大值都是连续分段最大和模型”，先判断原来的状态是否仍然覆盖新答案集合，再决定是微调边界、改 value 语义，还是换成另一类结构。

\subsection{栈、单调栈与表达式工作坊}

先对每个位置向左或向右扫描，再识别还没有被解决的候选。栈保存等待未来元素解决的对象，新元素到来时一次性结算被它支配的一批旧候选。

\subsubsection{LeetCode 20 有效的括号}

先让每个位置向左或向右暴力寻找答案，观察哪些候选一直等到未来才被解决。本题的模型是括号匹配要求最近打开的左括号先被关闭；慢点集中在栈保存尚未匹配的左括号，遇到右括号必须和栈顶类型一致。栈的作用不是保存所有历史，而是保存仍未被当前输入解决的那一小段历史。

栈里应保存还没确定答案的对象。本题用栈保存尚未匹配的左括号，遇到右括号必须和栈顶类型一致触发结算；弹出时要说清当前元素是右边界、左边界还是运算触发点。栈里存值还是下标，要由距离、宽度和过期判断决定。放到“栈、单调栈与表达式”这条主线里看，关键原则是：当新元素到来时，它会一次性解决栈顶一串被它支配的旧元素。被弹出的元素以后再也不会参与比较，因此总体线性。

正确性来自弹出时机。一旦栈保存尚未匹配的左括号，遇到右括号必须和栈顶类型一致触发，栈顶对象的答案已经由当前边界和新的栈顶共同确定；未来元素不会再改变它。本题的括号匹配要求最近打开的左括号先被关闭说明栈中对象等待的是什么。这道题的边界不能留到最后补丁式处理，实验时准备空串、以右括号开头、交叉匹配、结束后栈非空，打印每次入栈、弹栈和结算对象。重点观察相等元素、空栈、哨兵和最后一次结算。

C++ 实现上，栈题多数时候用 \code{std::vector<int>} 保存下标，比适配器更方便查看栈顶和计算宽度。表达式题要注意左右操作数顺序，柱状图题要注意哨兵和宽度。如果追问变成“若加入通配符星号，需要额外维护可行区间或双栈”，先判断原来的状态是否仍然覆盖新答案集合，再决定是微调边界、改 value 语义，还是换成另一类结构。

\subsubsection{LeetCode 32 最长有效括号}

先让每个位置向左或向右暴力寻找答案，观察哪些候选一直等到未来才被解决。本题的模型是有效括号子串需要知道每段非法边界之后的最早可连接位置；慢点集中在栈保存下标，先放入哨兵负一；匹配后用当前下标减栈顶得到长度。栈的作用不是保存所有历史，而是保存仍未被当前输入解决的那一小段历史。

栈里应保存还没确定答案的对象。本题用栈保存下标，先放入哨兵负一；匹配后用当前下标减栈顶得到长度触发结算；弹出时要说清当前元素是右边界、左边界还是运算触发点。栈里存值还是下标，要由距离、宽度和过期判断决定。放到“栈、单调栈与表达式”这条主线里看，关键原则是：当新元素到来时，它会一次性解决栈顶一串被它支配的旧元素。被弹出的元素以后再也不会参与比较，因此总体线性。

正确性来自弹出时机。一旦栈保存下标，先放入哨兵负一；匹配后用当前下标减栈顶得到长度触发，栈顶对象的答案已经由当前边界和新的栈顶共同确定；未来元素不会再改变它。本题的有效括号子串需要知道每段非法边界之后的最早可连接位置说明栈中对象等待的是什么。这道题的边界不能留到最后补丁式处理，实验时准备空串、全左括号、全右括号、多个断开的合法段，打印每次入栈、弹栈和结算对象。重点观察相等元素、空栈、哨兵和最后一次结算。

C++ 实现上，栈题多数时候用 \code{std::vector<int>} 保存下标，比适配器更方便查看栈顶和计算宽度。表达式题要注意左右操作数顺序，柱状图题要注意哨兵和宽度。如果追问变成“DP 写法也可做，但栈写法更直接表达最近非法边界”，先判断原来的状态是否仍然覆盖新答案集合，再决定是微调边界、改 value 语义，还是换成另一类结构。

\subsubsection{LeetCode 42 接雨水}

先让每个位置向左或向右暴力寻找答案，观察哪些候选一直等到未来才被解决。本题的模型是每个位置的水由左侧最高和右侧最高的较小者决定；慢点集中在前后缀、双指针、单调栈都是不同状态维护方式；要能说明各自结算对象。栈的作用不是保存所有历史，而是保存仍未被当前输入解决的那一小段历史。

栈里应保存还没确定答案的对象。本题用前后缀、双指针、单调栈都是不同状态维护方式；要能说明各自结算对象触发结算；弹出时要说清当前元素是右边界、左边界还是运算触发点。栈里存值还是下标，要由距离、宽度和过期判断决定。放到“栈、单调栈与表达式”这条主线里看，关键原则是：当新元素到来时，它会一次性解决栈顶一串被它支配的旧元素。被弹出的元素以后再也不会参与比较，因此总体线性。

正确性来自弹出时机。一旦前后缀、双指针、单调栈都是不同状态维护方式；要能说明各自结算对象触发，栈顶对象的答案已经由当前边界和新的栈顶共同确定；未来元素不会再改变它。本题的每个位置的水由左侧最高和右侧最高的较小者决定说明栈中对象等待的是什么。这道题的边界不能留到最后补丁式处理，实验时准备单调递增、单调递减、平台高度、多个凹槽、数组长度不足三，打印每次入栈、弹栈和结算对象。重点观察相等元素、空栈、哨兵和最后一次结算。

C++ 实现上，栈题多数时候用 \code{std::vector<int>} 保存下标，比适配器更方便查看栈顶和计算宽度。表达式题要注意左右操作数顺序，柱状图题要注意哨兵和宽度。如果追问变成“若柱子变成二维网格，需要最小堆从边界向内扩展”，先判断原来的状态是否仍然覆盖新答案集合，再决定是微调边界、改 value 语义，还是换成另一类结构。

\subsubsection{LeetCode 84 柱状图中最大的矩形}

先让每个位置向左或向右暴力寻找答案，观察哪些候选一直等到未来才被解决。本题的模型是每根柱子作为最矮高度时，需要知道左右第一个更矮位置；慢点集中在单调递增栈在遇到更矮柱时结算被弹出柱子的最大宽度。栈的作用不是保存所有历史，而是保存仍未被当前输入解决的那一小段历史。

栈里应保存还没确定答案的对象。本题用单调递增栈在遇到更矮柱时结算被弹出柱子的最大宽度触发结算；弹出时要说清当前元素是右边界、左边界还是运算触发点。栈里存值还是下标，要由距离、宽度和过期判断决定。放到“栈、单调栈与表达式”这条主线里看，关键原则是：当新元素到来时，它会一次性解决栈顶一串被它支配的旧元素。被弹出的元素以后再也不会参与比较，因此总体线性。

正确性来自弹出时机。一旦单调递增栈在遇到更矮柱时结算被弹出柱子的最大宽度触发，栈顶对象的答案已经由当前边界和新的栈顶共同确定；未来元素不会再改变它。本题的每根柱子作为最矮高度时，需要知道左右第一个更矮位置说明栈中对象等待的是什么。这道题的边界不能留到最后补丁式处理，实验时准备递增、递减、相等高度、哨兵零、宽度计算，打印每次入栈、弹栈和结算对象。重点观察相等元素、空栈、哨兵和最后一次结算。

C++ 实现上，栈题多数时候用 \code{std::vector<int>} 保存下标，比适配器更方便查看栈顶和计算宽度。表达式题要注意左右操作数顺序，柱状图题要注意哨兵和宽度。如果追问变成“最大矩形可以作为二维矩阵最大矩形的行压缩子问题”，先判断原来的状态是否仍然覆盖新答案集合，再决定是微调边界、改 value 语义，还是换成另一类结构。

\subsubsection{LeetCode 85 最大矩形}

先让每个位置向左或向右暴力寻找答案，观察哪些候选一直等到未来才被解决。本题的模型是每一行都可以把连续的一当成柱状图高度；慢点集中在逐行更新每列高度，然后把当前行转成柱状图最大矩形问题。栈的作用不是保存所有历史，而是保存仍未被当前输入解决的那一小段历史。

栈里应保存还没确定答案的对象。本题用逐行更新每列高度，然后把当前行转成柱状图最大矩形问题触发结算；弹出时要说清当前元素是右边界、左边界还是运算触发点。栈里存值还是下标，要由距离、宽度和过期判断决定。放到“栈、单调栈与表达式”这条主线里看，关键原则是：当新元素到来时，它会一次性解决栈顶一串被它支配的旧元素。被弹出的元素以后再也不会参与比较，因此总体线性。

正确性来自弹出时机。一旦逐行更新每列高度，然后把当前行转成柱状图最大矩形问题触发，栈顶对象的答案已经由当前边界和新的栈顶共同确定；未来元素不会再改变它。本题的每一行都可以把连续的一当成柱状图高度说明栈中对象等待的是什么。这道题的边界不能留到最后补丁式处理，实验时准备全零、全一、单行、单列、宽度和高度结算，打印每次入栈、弹栈和结算对象。重点观察相等元素、空栈、哨兵和最后一次结算。

C++ 实现上，栈题多数时候用 \code{std::vector<int>} 保存下标，比适配器更方便查看栈顶和计算宽度。表达式题要注意左右操作数顺序，柱状图题要注意哨兵和宽度。如果追问变成“这是二维问题压成一维单调栈的典型做法”，先判断原来的状态是否仍然覆盖新答案集合，再决定是微调边界、改 value 语义，还是换成另一类结构。

\subsubsection{LeetCode 150 逆波兰表达式求值}

先让每个位置向左或向右暴力寻找答案，观察哪些候选一直等到未来才被解决。本题的模型是后缀表达式把优先级和括号都编码进 token 顺序；慢点集中在遇到数字入栈，遇到运算符弹出右操作数和左操作数计算。栈的作用不是保存所有历史，而是保存仍未被当前输入解决的那一小段历史。

栈里应保存还没确定答案的对象。本题用遇到数字入栈，遇到运算符弹出右操作数和左操作数计算触发结算；弹出时要说清当前元素是右边界、左边界还是运算触发点。栈里存值还是下标，要由距离、宽度和过期判断决定。放到“栈、单调栈与表达式”这条主线里看，关键原则是：当新元素到来时，它会一次性解决栈顶一串被它支配的旧元素。被弹出的元素以后再也不会参与比较，因此总体线性。

正确性来自弹出时机。一旦遇到数字入栈，遇到运算符弹出右操作数和左操作数计算触发，栈顶对象的答案已经由当前边界和新的栈顶共同确定；未来元素不会再改变它。本题的后缀表达式把优先级和括号都编码进 token 顺序说明栈中对象等待的是什么。这道题的边界不能留到最后补丁式处理，实验时准备负数 token、多位数、除法向零截断、减法顺序，打印每次入栈、弹栈和结算对象。重点观察相等元素、空栈、哨兵和最后一次结算。

C++ 实现上，栈题多数时候用 \code{std::vector<int>} 保存下标，比适配器更方便查看栈顶和计算宽度。表达式题要注意左右操作数顺序，柱状图题要注意哨兵和宽度。如果追问变成“中缀表达式求值需要另一个运算符栈处理优先级”，先判断原来的状态是否仍然覆盖新答案集合，再决定是微调边界、改 value 语义，还是换成另一类结构。

\subsubsection{LeetCode 155 最小栈}

先让每个位置向左或向右暴力寻找答案，观察哪些候选一直等到未来才被解决。本题的模型是普通栈只能看到栈顶，不能常数时间知道历史最小值；慢点集中在辅助栈同步保存每个深度对应的最小值。栈的作用不是保存所有历史，而是保存仍未被当前输入解决的那一小段历史。

栈里应保存还没确定答案的对象。本题用辅助栈同步保存每个深度对应的最小值触发结算；弹出时要说清当前元素是右边界、左边界还是运算触发点。栈里存值还是下标，要由距离、宽度和过期判断决定。放到“栈、单调栈与表达式”这条主线里看，关键原则是：当新元素到来时，它会一次性解决栈顶一串被它支配的旧元素。被弹出的元素以后再也不会参与比较，因此总体线性。

正确性来自弹出时机。一旦辅助栈同步保存每个深度对应的最小值触发，栈顶对象的答案已经由当前边界和新的栈顶共同确定；未来元素不会再改变它。本题的普通栈只能看到栈顶，不能常数时间知道历史最小值说明栈中对象等待的是什么。这道题的边界不能留到最后补丁式处理，实验时准备重复最小值、弹出最小值、空栈操作约定，打印每次入栈、弹栈和结算对象。重点观察相等元素、空栈、哨兵和最后一次结算。

C++ 实现上，栈题多数时候用 \code{std::vector<int>} 保存下标，比适配器更方便查看栈顶和计算宽度。表达式题要注意左右操作数顺序，柱状图题要注意哨兵和宽度。如果追问变成“也可保存差值压缩空间，但可读性和溢出处理更复杂”，先判断原来的状态是否仍然覆盖新答案集合，再决定是微调边界、改 value 语义，还是换成另一类结构。

\subsubsection{LeetCode 224 基本计算器}

先让每个位置向左或向右暴力寻找答案，观察哪些候选一直等到未来才被解决。本题的模型是括号会开启新的表达式上下文，括号结束时要回到上一层；慢点集中在栈保存进入括号前的结果和符号，当前层算完后乘以上层符号并合并。栈的作用不是保存所有历史，而是保存仍未被当前输入解决的那一小段历史。

栈里应保存还没确定答案的对象。本题用栈保存进入括号前的结果和符号，当前层算完后乘以上层符号并合并触发结算；弹出时要说清当前元素是右边界、左边界还是运算触发点。栈里存值还是下标，要由距离、宽度和过期判断决定。放到“栈、单调栈与表达式”这条主线里看，关键原则是：当新元素到来时，它会一次性解决栈顶一串被它支配的旧元素。被弹出的元素以后再也不会参与比较，因此总体线性。

正确性来自弹出时机。一旦栈保存进入括号前的结果和符号，当前层算完后乘以上层符号并合并触发，栈顶对象的答案已经由当前边界和新的栈顶共同确定；未来元素不会再改变它。本题的括号会开启新的表达式上下文，括号结束时要回到上一层说明栈中对象等待的是什么。这道题的边界不能留到最后补丁式处理，实验时准备多位数、前导空格、一元负号、嵌套括号，打印每次入栈、弹栈和结算对象。重点观察相等元素、空栈、哨兵和最后一次结算。

C++ 实现上，栈题多数时候用 \code{std::vector<int>} 保存下标，比适配器更方便查看栈顶和计算宽度。表达式题要注意左右操作数顺序，柱状图题要注意哨兵和宽度。如果追问变成“基本计算器 II 没有括号但有乘除优先级，需要不同的栈语义”，先判断原来的状态是否仍然覆盖新答案集合，再决定是微调边界、改 value 语义，还是换成另一类结构。

\subsubsection{LeetCode 227 基本计算器 II}

先让每个位置向左或向右暴力寻找答案，观察哪些候选一直等到未来才被解决。本题的模型是乘除优先级高于加减，可以在读到下一个运算符时结算上一项；慢点集中在栈保存已经确定符号的项，乘除直接修改栈顶，加减压入正负项。栈的作用不是保存所有历史，而是保存仍未被当前输入解决的那一小段历史。

栈里应保存还没确定答案的对象。本题用栈保存已经确定符号的项，乘除直接修改栈顶，加减压入正负项触发结算；弹出时要说清当前元素是右边界、左边界还是运算触发点。栈里存值还是下标，要由距离、宽度和过期判断决定。放到“栈、单调栈与表达式”这条主线里看，关键原则是：当新元素到来时，它会一次性解决栈顶一串被它支配的旧元素。被弹出的元素以后再也不会参与比较，因此总体线性。

正确性来自弹出时机。一旦栈保存已经确定符号的项，乘除直接修改栈顶，加减压入正负项触发，栈顶对象的答案已经由当前边界和新的栈顶共同确定；未来元素不会再改变它。本题的乘除优先级高于加减，可以在读到下一个运算符时结算上一项说明栈中对象等待的是什么。这道题的边界不能留到最后补丁式处理，实验时准备多位数、空格、除法向零截断、最后一个数字需要结算，打印每次入栈、弹栈和结算对象。重点观察相等元素、空栈、哨兵和最后一次结算。

C++ 实现上，栈题多数时候用 \code{std::vector<int>} 保存下标，比适配器更方便查看栈顶和计算宽度。表达式题要注意左右操作数顺序，柱状图题要注意哨兵和宽度。如果追问变成“若加入括号，需要再引入上下文栈或递归下降解析”，先判断原来的状态是否仍然覆盖新答案集合，再决定是微调边界、改 value 语义，还是换成另一类结构。

\subsubsection{LeetCode 232 用栈实现队列}

先让每个位置向左或向右暴力寻找答案，观察哪些候选一直等到未来才被解决。本题的模型是两个栈分别负责输入顺序和输出顺序；慢点集中在只有输出栈为空时，才把输入栈整体倒过去，保证均摊常数。栈的作用不是保存所有历史，而是保存仍未被当前输入解决的那一小段历史。

栈里应保存还没确定答案的对象。本题用只有输出栈为空时，才把输入栈整体倒过去，保证均摊常数触发结算；弹出时要说清当前元素是右边界、左边界还是运算触发点。栈里存值还是下标，要由距离、宽度和过期判断决定。放到“栈、单调栈与表达式”这条主线里看，关键原则是：当新元素到来时，它会一次性解决栈顶一串被它支配的旧元素。被弹出的元素以后再也不会参与比较，因此总体线性。

正确性来自弹出时机。一旦只有输出栈为空时，才把输入栈整体倒过去，保证均摊常数触发，栈顶对象的答案已经由当前边界和新的栈顶共同确定；未来元素不会再改变它。本题的两个栈分别负责输入顺序和输出顺序说明栈中对象等待的是什么。这道题的边界不能留到最后补丁式处理，实验时准备连续 peek、空队列、交替 push pop，打印每次入栈、弹栈和结算对象。重点观察相等元素、空栈、哨兵和最后一次结算。

C++ 实现上，栈题多数时候用 \code{std::vector<int>} 保存下标，比适配器更方便查看栈顶和计算宽度。表达式题要注意左右操作数顺序，柱状图题要注意哨兵和宽度。如果追问变成“用队列实现栈则要通过轮转保持最近元素在队首”，先判断原来的状态是否仍然覆盖新答案集合，再决定是微调边界、改 value 语义，还是换成另一类结构。

\subsubsection{LeetCode 239 滑动窗口最大值}

先让每个位置向左或向右暴力寻找答案，观察哪些候选一直等到未来才被解决。本题的模型是每个窗口最大值来自仍未过期且没有被新元素支配的候选；慢点集中在单调队列保存下标，队尾小于等于新值的候选被弹出。栈的作用不是保存所有历史，而是保存仍未被当前输入解决的那一小段历史。

栈里应保存还没确定答案的对象。本题用单调队列保存下标，队尾小于等于新值的候选被弹出触发结算；弹出时要说清当前元素是右边界、左边界还是运算触发点。栈里存值还是下标，要由距离、宽度和过期判断决定。放到“栈、单调栈与表达式”这条主线里看，关键原则是：当新元素到来时，它会一次性解决栈顶一串被它支配的旧元素。被弹出的元素以后再也不会参与比较，因此总体线性。

正确性来自弹出时机。一旦单调队列保存下标，队尾小于等于新值的候选被弹出触发，栈顶对象的答案已经由当前边界和新的栈顶共同确定；未来元素不会再改变它。本题的每个窗口最大值来自仍未过期且没有被新元素支配的候选说明栈中对象等待的是什么。这道题的边界不能留到最后补丁式处理，实验时准备k 为一、k 等于数组长度、重复最大值、队首过期，打印每次入栈、弹栈和结算对象。重点观察相等元素、空栈、哨兵和最后一次结算。

C++ 实现上，栈题多数时候用 \code{std::vector<int>} 保存下标，比适配器更方便查看栈顶和计算宽度。表达式题要注意左右操作数顺序，柱状图题要注意哨兵和宽度。如果追问变成“受限子序列和使用同样队列维护最近 k 个 DP 最大值”，先判断原来的状态是否仍然覆盖新答案集合，再决定是微调边界、改 value 语义，还是换成另一类结构。

\subsubsection{LeetCode 394 字符串解码}

先让每个位置向左或向右暴力寻找答案，观察哪些候选一直等到未来才被解决。本题的模型是嵌套结构需要保存上一层字符串和重复次数；慢点集中在遇到左括号入栈当前状态，右括号弹出并拼接展开。栈的作用不是保存所有历史，而是保存仍未被当前输入解决的那一小段历史。

栈里应保存还没确定答案的对象。本题用遇到左括号入栈当前状态，右括号弹出并拼接展开触发结算；弹出时要说清当前元素是右边界、左边界还是运算触发点。栈里存值还是下标，要由距离、宽度和过期判断决定。放到“栈、单调栈与表达式”这条主线里看，关键原则是：当新元素到来时，它会一次性解决栈顶一串被它支配的旧元素。被弹出的元素以后再也不会参与比较，因此总体线性。

正确性来自弹出时机。一旦遇到左括号入栈当前状态，右括号弹出并拼接展开触发，栈顶对象的答案已经由当前边界和新的栈顶共同确定；未来元素不会再改变它。本题的嵌套结构需要保存上一层字符串和重复次数说明栈中对象等待的是什么。这道题的边界不能留到最后补丁式处理，实验时准备多位数字、嵌套、空片段、数字后必须跟括号，打印每次入栈、弹栈和结算对象。重点观察相等元素、空栈、哨兵和最后一次结算。

C++ 实现上，栈题多数时候用 \code{std::vector<int>} 保存下标，比适配器更方便查看栈顶和计算宽度。表达式题要注意左右操作数顺序，柱状图题要注意哨兵和宽度。如果追问变成“递归下降解析更通用，但迭代栈更符合工程栈控制”，先判断原来的状态是否仍然覆盖新答案集合，再决定是微调边界、改 value 语义，还是换成另一类结构。

\subsubsection{LeetCode 402 移掉 K 位数字}

先让每个位置向左或向右暴力寻找答案，观察哪些候选一直等到未来才被解决。本题的模型是要删除 k 位使剩余数字最小，本质是让高位尽量小；慢点集中在单调递增栈中遇到更小数字时，弹出前面较大的数字并消耗删除次数。栈的作用不是保存所有历史，而是保存仍未被当前输入解决的那一小段历史。

栈里应保存还没确定答案的对象。本题用单调递增栈中遇到更小数字时，弹出前面较大的数字并消耗删除次数触发结算；弹出时要说清当前元素是右边界、左边界还是运算触发点。栈里存值还是下标，要由距离、宽度和过期判断决定。放到“栈、单调栈与表达式”这条主线里看，关键原则是：当新元素到来时，它会一次性解决栈顶一串被它支配的旧元素。被弹出的元素以后再也不会参与比较，因此总体线性。

正确性来自弹出时机。一旦单调递增栈中遇到更小数字时，弹出前面较大的数字并消耗删除次数触发，栈顶对象的答案已经由当前边界和新的栈顶共同确定；未来元素不会再改变它。本题的要删除 k 位使剩余数字最小，本质是让高位尽量小说明栈中对象等待的是什么。这道题的边界不能留到最后补丁式处理，实验时准备前导零、k 未用完、k 等于长度、数字已经递增，打印每次入栈、弹栈和结算对象。重点观察相等元素、空栈、哨兵和最后一次结算。

C++ 实现上，栈题多数时候用 \code{std::vector<int>} 保存下标，比适配器更方便查看栈顶和计算宽度。表达式题要注意左右操作数顺序，柱状图题要注意哨兵和宽度。如果追问变成“去除重复字母也用单调栈，但还要保证每个字符保留一次”，先判断原来的状态是否仍然覆盖新答案集合，再决定是微调边界、改 value 语义，还是换成另一类结构。

\subsubsection{LeetCode 496 下一个更大元素 I}

先让每个位置向左或向右暴力寻找答案，观察哪些候选一直等到未来才被解决。本题的模型是第二个数组提供每个值右侧第一个更大值的全局关系；慢点集中在单调栈预处理值到下一个更大值映射，再回答第一个数组查询。栈的作用不是保存所有历史，而是保存仍未被当前输入解决的那一小段历史。

栈里应保存还没确定答案的对象。本题用单调栈预处理值到下一个更大值映射，再回答第一个数组查询触发结算；弹出时要说清当前元素是右边界、左边界还是运算触发点。栈里存值还是下标，要由距离、宽度和过期判断决定。放到“栈、单调栈与表达式”这条主线里看，关键原则是：当新元素到来时，它会一次性解决栈顶一串被它支配的旧元素。被弹出的元素以后再也不会参与比较，因此总体线性。

正确性来自弹出时机。一旦单调栈预处理值到下一个更大值映射，再回答第一个数组查询触发，栈顶对象的答案已经由当前边界和新的栈顶共同确定；未来元素不会再改变它。本题的第二个数组提供每个值右侧第一个更大值的全局关系说明栈中对象等待的是什么。这道题的边界不能留到最后补丁式处理，实验时准备不存在更大值、递减数组、值唯一假设，打印每次入栈、弹栈和结算对象。重点观察相等元素、空栈、哨兵和最后一次结算。

C++ 实现上，栈题多数时候用 \code{std::vector<int>} 保存下标，比适配器更方便查看栈顶和计算宽度。表达式题要注意左右操作数顺序，柱状图题要注意哨兵和宽度。如果追问变成“若值不唯一，应按下标而不是值建映射”，先判断原来的状态是否仍然覆盖新答案集合，再决定是微调边界、改 value 语义，还是换成另一类结构。

\subsubsection{LeetCode 503 下一个更大元素 II}

先让每个位置向左或向右暴力寻找答案，观察哪些候选一直等到未来才被解决。本题的模型是环形数组可以通过遍历两倍下标模拟绕回；慢点集中在第一遍入栈，第二遍只帮助未解决下标找到答案。栈的作用不是保存所有历史，而是保存仍未被当前输入解决的那一小段历史。

栈里应保存还没确定答案的对象。本题用第一遍入栈，第二遍只帮助未解决下标找到答案触发结算；弹出时要说清当前元素是右边界、左边界还是运算触发点。栈里存值还是下标，要由距离、宽度和过期判断决定。放到“栈、单调栈与表达式”这条主线里看，关键原则是：当新元素到来时，它会一次性解决栈顶一串被它支配的旧元素。被弹出的元素以后再也不会参与比较，因此总体线性。

正确性来自弹出时机。一旦第一遍入栈，第二遍只帮助未解决下标找到答案触发，栈顶对象的答案已经由当前边界和新的栈顶共同确定；未来元素不会再改变它。本题的环形数组可以通过遍历两倍下标模拟绕回说明栈中对象等待的是什么。这道题的边界不能留到最后补丁式处理，实验时准备全递减、全相等、长度一、取模下标，打印每次入栈、弹栈和结算对象。重点观察相等元素、空栈、哨兵和最后一次结算。

C++ 实现上，栈题多数时候用 \code{std::vector<int>} 保存下标，比适配器更方便查看栈顶和计算宽度。表达式题要注意左右操作数顺序，柱状图题要注意哨兵和宽度。如果追问变成“不要在第二遍重复入栈，否则可能死循环或重复计算”，先判断原来的状态是否仍然覆盖新答案集合，再决定是微调边界、改 value 语义，还是换成另一类结构。

\subsubsection{LeetCode 739 每日温度}

先让每个位置向左或向右暴力寻找答案，观察哪些候选一直等到未来才被解决。本题的模型是每一天等待右侧第一个更高温度；慢点集中在单调递减栈保存尚未找到答案的下标。栈的作用不是保存所有历史，而是保存仍未被当前输入解决的那一小段历史。

栈里应保存还没确定答案的对象。本题用单调递减栈保存尚未找到答案的下标触发结算；弹出时要说清当前元素是右边界、左边界还是运算触发点。栈里存值还是下标，要由距离、宽度和过期判断决定。放到“栈、单调栈与表达式”这条主线里看，关键原则是：当新元素到来时，它会一次性解决栈顶一串被它支配的旧元素。被弹出的元素以后再也不会参与比较，因此总体线性。

正确性来自弹出时机。一旦单调递减栈保存尚未找到答案的下标触发，栈顶对象的答案已经由当前边界和新的栈顶共同确定；未来元素不会再改变它。本题的每一天等待右侧第一个更高温度说明栈中对象等待的是什么。这道题的边界不能留到最后补丁式处理，实验时准备没有更高温、相等温度、递增、递减，打印每次入栈、弹栈和结算对象。重点观察相等元素、空栈、哨兵和最后一次结算。

C++ 实现上，栈题多数时候用 \code{std::vector<int>} 保存下标，比适配器更方便查看栈顶和计算宽度。表达式题要注意左右操作数顺序，柱状图题要注意哨兵和宽度。如果追问变成“下一个更大元素与它同型，只是输出值或距离不同”，先判断原来的状态是否仍然覆盖新答案集合，再决定是微调边界、改 value 语义，还是换成另一类结构。

\subsubsection{LeetCode 895 最大频率栈}

先让每个位置向左或向右暴力寻找答案，观察哪些候选一直等到未来才被解决。本题的模型是弹出时既要频率最高，又要在同频中最近入栈；慢点集中在哈希表记录值频率，频率到栈记录达到该频率的入栈顺序，维护当前最大频率。栈的作用不是保存所有历史，而是保存仍未被当前输入解决的那一小段历史。

栈里应保存还没确定答案的对象。本题用哈希表记录值频率，频率到栈记录达到该频率的入栈顺序，维护当前最大频率触发结算；弹出时要说清当前元素是右边界、左边界还是运算触发点。栈里存值还是下标，要由距离、宽度和过期判断决定。放到“栈、单调栈与表达式”这条主线里看，关键原则是：当新元素到来时，它会一次性解决栈顶一串被它支配的旧元素。被弹出的元素以后再也不会参与比较，因此总体线性。

正确性来自弹出时机。一旦哈希表记录值频率，频率到栈记录达到该频率的入栈顺序，维护当前最大频率触发，栈顶对象的答案已经由当前边界和新的栈顶共同确定；未来元素不会再改变它。本题的弹出时既要频率最高，又要在同频中最近入栈说明栈中对象等待的是什么。这道题的边界不能留到最后补丁式处理，实验时准备重复 push、pop 后最大频率下降、同频最近性、空栈不在题意内，打印每次入栈、弹栈和结算对象。重点观察相等元素、空栈、哨兵和最后一次结算。

C++ 实现上，栈题多数时候用 \code{std::vector<int>} 保存下标，比适配器更方便查看栈顶和计算宽度。表达式题要注意左右操作数顺序，柱状图题要注意哨兵和宽度。如果追问变成“它是栈和哈希表按频率维度组合的设计题”，先判断原来的状态是否仍然覆盖新答案集合，再决定是微调边界、改 value 语义，还是换成另一类结构。

\subsection{堆与动态极值工作坊}

先每轮扫描所有候选找极值，再把动态极值查询交给堆。堆只保证堆顶，不保证全局遍历有序；有过期元素时要在使用堆顶前清理。

\subsubsection{LeetCode 23 合并 K 个升序链表}

先每轮扫描所有候选找极值，确认答案选择的对象。本题的模型是每条链表当前头节点都是可能的下一个最小元素；暴力版慢在动态候选集合不断变化，却每次重新找最值。本题的优化点是小顶堆保存每条非空链的当前头，弹出最小节点后把它的后继入堆，因此堆中只能放足以比较下一步选择的轻量状态。

堆只保证堆顶，不保证内部全局有序。本题的小顶堆保存每条非空链的当前头，弹出最小节点后把它的后继入堆要转成堆顶可回答的问题；如果元素会过期，还要在真正使用堆顶前清理或跳过旧状态。放到“堆与动态极值”这条主线里看，关键原则是：堆把插入和弹出极值控制在对数复杂度。Top K 用固定大小堆维护边界，合并多路序列用堆保存每一路当前头部，数据流中位数用双堆维护两半。

证明堆题时，要说明堆覆盖了所有可能的下一步候选，并且堆顶过期时不会被使用。本题的小顶堆保存每条非空链的当前头，弹出最小节点后把它的后继入堆是选择顺序，每条链表当前头节点都是可能的下一个最小元素是候选集合来源。这道题的边界不能留到最后补丁式处理，实验时覆盖空链表数组、某条链为空、重复值、堆里保存节点指针还是值，记录堆顶是否过期、比较器是否让正确候选在堆顶、K 或窗口大小为边界值时是否稳定。

C++ 实现上，\code{std::priority_queue} 默认大顶堆，小顶堆需要比较器。堆中保存大对象时尽量保存下标或轻量记录；若不能删除任意元素，就用懒删除或过期检查。如果追问变成“也可以分治两两合并，堆更适合 K 较大且链长不均的场景”，先判断原来的状态是否仍然覆盖新答案集合，再决定是微调边界、改 value 语义，还是换成另一类结构。

\subsubsection{LeetCode 215 数组中的第 K 个最大元素}

先每轮扫描所有候选找极值，确认答案选择的对象。本题的模型是只关心第 K 大，不需要完整排序；暴力版慢在动态候选集合不断变化，却每次重新找最值。本题的优化点是大小为 K 的小顶堆保存当前前 K 大边界；快速选择可做到平均线性，因此堆中只能放足以比较下一步选择的轻量状态。

堆只保证堆顶，不保证内部全局有序。本题的大小为 K 的小顶堆保存当前前 K 大边界；快速选择可做到平均线性要转成堆顶可回答的问题；如果元素会过期，还要在真正使用堆顶前清理或跳过旧状态。放到“堆与动态极值”这条主线里看，关键原则是：堆把插入和弹出极值控制在对数复杂度。Top K 用固定大小堆维护边界，合并多路序列用堆保存每一路当前头部，数据流中位数用双堆维护两半。

证明堆题时，要说明堆覆盖了所有可能的下一步候选，并且堆顶过期时不会被使用。本题的大小为 K 的小顶堆保存当前前 K 大边界；快速选择可做到平均线性是选择顺序，只关心第 K 大，不需要完整排序是候选集合来源。这道题的边界不能留到最后补丁式处理，实验时覆盖K 为一、K 为 n、重复值、随机 pivot，记录堆顶是否过期、比较器是否让正确候选在堆顶、K 或窗口大小为边界值时是否稳定。

C++ 实现上，\code{std::priority_queue} 默认大顶堆，小顶堆需要比较器。堆中保存大对象时尽量保存下标或轻量记录；若不能删除任意元素，就用懒删除或过期检查。如果追问变成“若数据流持续到来，堆方案比一次性快速选择更自然”，先判断原来的状态是否仍然覆盖新答案集合，再决定是微调边界、改 value 语义，还是换成另一类结构。

\subsubsection{LeetCode 295 数据流的中位数}

先每轮扫描所有候选找极值，确认答案选择的对象。本题的模型是数据流需要动态维护较小一半和较大一半；暴力版慢在动态候选集合不断变化，却每次重新找最值。本题的优化点是大顶堆保存较小一半，小顶堆保存较大一半，大小差不超过一，因此堆中只能放足以比较下一步选择的轻量状态。

堆只保证堆顶，不保证内部全局有序。本题的大顶堆保存较小一半，小顶堆保存较大一半，大小差不超过一要转成堆顶可回答的问题；如果元素会过期，还要在真正使用堆顶前清理或跳过旧状态。放到“堆与动态极值”这条主线里看，关键原则是：堆把插入和弹出极值控制在对数复杂度。Top K 用固定大小堆维护边界，合并多路序列用堆保存每一路当前头部，数据流中位数用双堆维护两半。

证明堆题时，要说明堆覆盖了所有可能的下一步候选，并且堆顶过期时不会被使用。本题的大顶堆保存较小一半，小顶堆保存较大一半，大小差不超过一是选择顺序，数据流需要动态维护较小一半和较大一半是候选集合来源。这道题的边界不能留到最后补丁式处理，实验时覆盖偶数个元素、负数、重复值、平均值溢出，记录堆顶是否过期、比较器是否让正确候选在堆顶、K 或窗口大小为边界值时是否稳定。

C++ 实现上，\code{std::priority_queue} 默认大顶堆，小顶堆需要比较器。堆中保存大对象时尽量保存下标或轻量记录；若不能删除任意元素，就用懒删除或过期检查。如果追问变成“若还要删除任意元素，需要延迟删除哈希表”，先判断原来的状态是否仍然覆盖新答案集合，再决定是微调边界、改 value 语义，还是换成另一类结构。

\subsubsection{LeetCode 347 前 K 个高频元素}

先每轮扫描所有候选找极值，确认答案选择的对象。本题的模型是先统计频次，再只维护前 K 个候选；暴力版慢在动态候选集合不断变化，却每次重新找最值。本题的优化点是小顶堆大小超过 K 就弹出最低频；桶排序利用频次上界，因此堆中只能放足以比较下一步选择的轻量状态。

堆只保证堆顶，不保证内部全局有序。本题的小顶堆大小超过 K 就弹出最低频；桶排序利用频次上界要转成堆顶可回答的问题；如果元素会过期，还要在真正使用堆顶前清理或跳过旧状态。放到“堆与动态极值”这条主线里看，关键原则是：堆把插入和弹出极值控制在对数复杂度。Top K 用固定大小堆维护边界，合并多路序列用堆保存每一路当前头部，数据流中位数用双堆维护两半。

证明堆题时，要说明堆覆盖了所有可能的下一步候选，并且堆顶过期时不会被使用。本题的小顶堆大小超过 K 就弹出最低频；桶排序利用频次上界是选择顺序，先统计频次，再只维护前 K 个候选是候选集合来源。这道题的边界不能留到最后补丁式处理，实验时覆盖K 等于不同元素数、频次相同、结果顺序不要求，记录堆顶是否过期、比较器是否让正确候选在堆顶、K 或窗口大小为边界值时是否稳定。

C++ 实现上，\code{std::priority_queue} 默认大顶堆，小顶堆需要比较器。堆中保存大对象时尽量保存下标或轻量记录；若不能删除任意元素，就用懒删除或过期检查。如果追问变成“若要按频次和字典序排序，比较器要同时处理两个键”，先判断原来的状态是否仍然覆盖新答案集合，再决定是微调边界、改 value 语义，还是换成另一类结构。

\subsubsection{LeetCode 480 滑动窗口中位数}

先每轮扫描所有候选找极值，确认答案选择的对象。本题的模型是窗口移动时既要加入新元素，也要删除滑出的旧元素并查询中位数；暴力版慢在动态候选集合不断变化，却每次重新找最值。本题的优化点是双堆维护左右两半，延迟删除表记录已离窗但尚未到堆顶的元素，因此堆中只能放足以比较下一步选择的轻量状态。

堆只保证堆顶，不保证内部全局有序。本题的双堆维护左右两半，延迟删除表记录已离窗但尚未到堆顶的元素要转成堆顶可回答的问题；如果元素会过期，还要在真正使用堆顶前清理或跳过旧状态。放到“堆与动态极值”这条主线里看，关键原则是：堆把插入和弹出极值控制在对数复杂度。Top K 用固定大小堆维护边界，合并多路序列用堆保存每一路当前头部，数据流中位数用双堆维护两半。

证明堆题时，要说明堆覆盖了所有可能的下一步候选，并且堆顶过期时不会被使用。本题的双堆维护左右两半，延迟删除表记录已离窗但尚未到堆顶的元素是选择顺序，窗口移动时既要加入新元素，也要删除滑出的旧元素并查询中位数是候选集合来源。这道题的边界不能留到最后补丁式处理，实验时覆盖重复值、偶数窗口、平均值溢出、堆顶过期，记录堆顶是否过期、比较器是否让正确候选在堆顶、K 或窗口大小为边界值时是否稳定。

C++ 实现上，\code{std::priority_queue} 默认大顶堆，小顶堆需要比较器。堆中保存大对象时尽量保存下标或轻量记录；若不能删除任意元素，就用懒删除或过期检查。如果追问变成“普通 priority\_queue 不能任意删除，所以必须解释懒删除”，先判断原来的状态是否仍然覆盖新答案集合，再决定是微调边界、改 value 语义，还是换成另一类结构。

\subsubsection{LeetCode 621 任务调度器}

先每轮扫描所有候选找极值，确认答案选择的对象。本题的模型是相同任务之间有冷却间隔，核心受最高频任务骨架限制；暴力版慢在动态候选集合不断变化，却每次重新找最值。本题的优化点是可用大顶堆模拟，也可用最高频公式计算最短时间，因此堆中只能放足以比较下一步选择的轻量状态。

堆只保证堆顶，不保证内部全局有序。本题的可用大顶堆模拟，也可用最高频公式计算最短时间要转成堆顶可回答的问题；如果元素会过期，还要在真正使用堆顶前清理或跳过旧状态。放到“堆与动态极值”这条主线里看，关键原则是：堆把插入和弹出极值控制在对数复杂度。Top K 用固定大小堆维护边界，合并多路序列用堆保存每一路当前头部，数据流中位数用双堆维护两半。

证明堆题时，要说明堆覆盖了所有可能的下一步候选，并且堆顶过期时不会被使用。本题的可用大顶堆模拟，也可用最高频公式计算最短时间是选择顺序，相同任务之间有冷却间隔，核心受最高频任务骨架限制是候选集合来源。这道题的边界不能留到最后补丁式处理，实验时覆盖冷却为零、多个最高频、空闲被其他任务填满，记录堆顶是否过期、比较器是否让正确候选在堆顶、K 或窗口大小为边界值时是否稳定。

C++ 实现上，\code{std::priority_queue} 默认大顶堆，小顶堆需要比较器。堆中保存大对象时尽量保存下标或轻量记录；若不能删除任意元素，就用懒删除或过期检查。如果追问变成“若任务有不同释放时间或执行时长，就变成更一般的调度问题”，先判断原来的状态是否仍然覆盖新答案集合，再决定是微调边界、改 value 语义，还是换成另一类结构。

\subsubsection{LeetCode 632 最小区间覆盖 K 个列表}

先每轮扫描所有候选找极值，确认答案选择的对象。本题的模型是每个列表至少取一个元素时，当前区间由最小当前值和最大当前值决定；暴力版慢在动态候选集合不断变化，却每次重新找最值。本题的优化点是小顶堆保存各列表当前元素，同时维护当前最大值；弹出最小所在列表并推进，因此堆中只能放足以比较下一步选择的轻量状态。

堆只保证堆顶，不保证内部全局有序。本题的小顶堆保存各列表当前元素，同时维护当前最大值；弹出最小所在列表并推进要转成堆顶可回答的问题；如果元素会过期，还要在真正使用堆顶前清理或跳过旧状态。放到“堆与动态极值”这条主线里看，关键原则是：堆把插入和弹出极值控制在对数复杂度。Top K 用固定大小堆维护边界，合并多路序列用堆保存每一路当前头部，数据流中位数用双堆维护两半。

证明堆题时，要说明堆覆盖了所有可能的下一步候选，并且堆顶过期时不会被使用。本题的小顶堆保存各列表当前元素，同时维护当前最大值；弹出最小所在列表并推进是选择顺序，每个列表至少取一个元素时，当前区间由最小当前值和最大当前值决定是候选集合来源。这道题的边界不能留到最后补丁式处理，实验时覆盖某个列表为空、重复值、区间长度相同的 tie-breaker、列表耗尽，记录堆顶是否过期、比较器是否让正确候选在堆顶、K 或窗口大小为边界值时是否稳定。

C++ 实现上，\code{std::priority_queue} 默认大顶堆，小顶堆需要比较器。堆中保存大对象时尽量保存下标或轻量记录；若不能删除任意元素，就用懒删除或过期检查。如果追问变成“这是多路归并和覆盖窗口的结合”，先判断原来的状态是否仍然覆盖新答案集合，再决定是微调边界、改 value 语义，还是换成另一类结构。

\subsubsection{LeetCode 703 数据流中的第 K 大元素}

先每轮扫描所有候选找极值，确认答案选择的对象。本题的模型是数据不断到达，只关心当前第 K 大而非完整排序；暴力版慢在动态候选集合不断变化，却每次重新找最值。本题的优化点是维护大小为 K 的小顶堆，堆顶就是当前第 K 大元素，因此堆中只能放足以比较下一步选择的轻量状态。

堆只保证堆顶，不保证内部全局有序。本题的维护大小为 K 的小顶堆，堆顶就是当前第 K 大元素要转成堆顶可回答的问题；如果元素会过期，还要在真正使用堆顶前清理或跳过旧状态。放到“堆与动态极值”这条主线里看，关键原则是：堆把插入和弹出极值控制在对数复杂度。Top K 用固定大小堆维护边界，合并多路序列用堆保存每一路当前头部，数据流中位数用双堆维护两半。

证明堆题时，要说明堆覆盖了所有可能的下一步候选，并且堆顶过期时不会被使用。本题的维护大小为 K 的小顶堆，堆顶就是当前第 K 大元素是选择顺序，数据不断到达，只关心当前第 K 大而非完整排序是候选集合来源。这道题的边界不能留到最后补丁式处理，实验时覆盖初始元素少于 K、重复值、K 为一、连续 add，记录堆顶是否过期、比较器是否让正确候选在堆顶、K 或窗口大小为边界值时是否稳定。

C++ 实现上，\code{std::priority_queue} 默认大顶堆，小顶堆需要比较器。堆中保存大对象时尽量保存下标或轻量记录；若不能删除任意元素，就用懒删除或过期检查。如果追问变成“这题是 Top K 的在线版本，和一次性排序的适用场景不同”，先判断原来的状态是否仍然覆盖新答案集合，再决定是微调边界、改 value 语义，还是换成另一类结构。

\subsubsection{LeetCode 857 雇佣 K 名工人的最低成本}

先每轮扫描所有候选找极值，确认答案选择的对象。本题的模型是工资比例由当前队伍中最高工资质量比决定，总成本等于比例乘质量和；暴力版慢在动态候选集合不断变化，却每次重新找最值。本题的优化点是按比例升序扫描，用大顶堆维护当前 K 个工人的最小质量和，因此堆中只能放足以比较下一步选择的轻量状态。

堆只保证堆顶，不保证内部全局有序。本题的按比例升序扫描，用大顶堆维护当前 K 个工人的最小质量和要转成堆顶可回答的问题；如果元素会过期，还要在真正使用堆顶前清理或跳过旧状态。放到“堆与动态极值”这条主线里看，关键原则是：堆把插入和弹出极值控制在对数复杂度。Top K 用固定大小堆维护边界，合并多路序列用堆保存每一路当前头部，数据流中位数用双堆维护两半。

证明堆题时，要说明堆覆盖了所有可能的下一步候选，并且堆顶过期时不会被使用。本题的按比例升序扫描，用大顶堆维护当前 K 个工人的最小质量和是选择顺序，工资比例由当前队伍中最高工资质量比决定，总成本等于比例乘质量和是候选集合来源。这道题的边界不能留到最后补丁式处理，实验时覆盖K 为一、比例相同、浮点精度、质量和更新，记录堆顶是否过期、比较器是否让正确候选在堆顶、K 或窗口大小为边界值时是否稳定。

C++ 实现上，\code{std::priority_queue} 默认大顶堆，小顶堆需要比较器。堆中保存大对象时尽量保存下标或轻量记录；若不能删除任意元素，就用懒删除或过期检查。如果追问变成“这是排序确定价格比例，堆维护可替换候选的组合题”，先判断原来的状态是否仍然覆盖新答案集合，再决定是微调边界、改 value 语义，还是换成另一类结构。

\subsubsection{LeetCode 871 最低加油次数}

先每轮扫描所有候选找极值，确认答案选择的对象。本题的模型是行驶过程中，经过的加油站形成可反悔的燃料候选；暴力版慢在动态候选集合不断变化，却每次重新找最值。本题的优化点是按位置扫描，油量不足到达下一点时，从已路过站点中取最大燃料补充，因此堆中只能放足以比较下一步选择的轻量状态。

堆只保证堆顶，不保证内部全局有序。本题的按位置扫描，油量不足到达下一点时，从已路过站点中取最大燃料补充要转成堆顶可回答的问题；如果元素会过期，还要在真正使用堆顶前清理或跳过旧状态。放到“堆与动态极值”这条主线里看，关键原则是：堆把插入和弹出极值控制在对数复杂度。Top K 用固定大小堆维护边界，合并多路序列用堆保存每一路当前头部，数据流中位数用双堆维护两半。

证明堆题时，要说明堆覆盖了所有可能的下一步候选，并且堆顶过期时不会被使用。本题的按位置扫描，油量不足到达下一点时，从已路过站点中取最大燃料补充是选择顺序，行驶过程中，经过的加油站形成可反悔的燃料候选是候选集合来源。这道题的边界不能留到最后补丁式处理，实验时覆盖起点油足够、无法到达第一个站、多个站同位置、目标作为终点事件，记录堆顶是否过期、比较器是否让正确候选在堆顶、K 或窗口大小为边界值时是否稳定。

C++ 实现上，\code{std::priority_queue} 默认大顶堆，小顶堆需要比较器。堆中保存大对象时尽量保存下标或轻量记录；若不能删除任意元素，就用懒删除或过期检查。如果追问变成“这是反悔贪心，堆里保存的是已经经过但尚未使用的资源”，先判断原来的状态是否仍然覆盖新答案集合，再决定是微调边界、改 value 语义，还是换成另一类结构。

\subsubsection{LeetCode 973 最接近原点的 K 个点}

先每轮扫描所有候选找极值，确认答案选择的对象。本题的模型是距离只需比较平方，避免开方误差和成本；暴力版慢在动态候选集合不断变化，却每次重新找最值。本题的优化点是大小为 K 的大顶堆保存当前最近 K 个点中最远者，因此堆中只能放足以比较下一步选择的轻量状态。

堆只保证堆顶，不保证内部全局有序。本题的大小为 K 的大顶堆保存当前最近 K 个点中最远者要转成堆顶可回答的问题；如果元素会过期，还要在真正使用堆顶前清理或跳过旧状态。放到“堆与动态极值”这条主线里看，关键原则是：堆把插入和弹出极值控制在对数复杂度。Top K 用固定大小堆维护边界，合并多路序列用堆保存每一路当前头部，数据流中位数用双堆维护两半。

证明堆题时，要说明堆覆盖了所有可能的下一步候选，并且堆顶过期时不会被使用。本题的大小为 K 的大顶堆保存当前最近 K 个点中最远者是选择顺序，距离只需比较平方，避免开方误差和成本是候选集合来源。这道题的边界不能留到最后补丁式处理，实验时覆盖K 为一、K 为全部、距离相同、坐标平方溢出，记录堆顶是否过期、比较器是否让正确候选在堆顶、K 或窗口大小为边界值时是否稳定。

C++ 实现上，\code{std::priority_queue} 默认大顶堆，小顶堆需要比较器。堆中保存大对象时尽量保存下标或轻量记录；若不能删除任意元素，就用懒删除或过期检查。如果追问变成“快速选择可平均线性，但堆更适合数据流”，先判断原来的状态是否仍然覆盖新答案集合，再决定是微调边界、改 value 语义，还是换成另一类结构。

\subsubsection{LeetCode 1046 最后一块石头的重量}

先每轮扫描所有候选找极值，确认答案选择的对象。本题的模型是每次都要取当前最重的两块石头碰撞；暴力版慢在动态候选集合不断变化，却每次重新找最值。本题的优化点是大顶堆保存所有重量，弹出两个最大值，若不同则把差值放回，因此堆中只能放足以比较下一步选择的轻量状态。

堆只保证堆顶，不保证内部全局有序。本题的大顶堆保存所有重量，弹出两个最大值，若不同则把差值放回要转成堆顶可回答的问题；如果元素会过期，还要在真正使用堆顶前清理或跳过旧状态。放到“堆与动态极值”这条主线里看，关键原则是：堆把插入和弹出极值控制在对数复杂度。Top K 用固定大小堆维护边界，合并多路序列用堆保存每一路当前头部，数据流中位数用双堆维护两半。

证明堆题时，要说明堆覆盖了所有可能的下一步候选，并且堆顶过期时不会被使用。本题的大顶堆保存所有重量，弹出两个最大值，若不同则把差值放回是选择顺序，每次都要取当前最重的两块石头碰撞是候选集合来源。这道题的边界不能留到最后补丁式处理，实验时覆盖没有石头、一块石头、两块相等、多次产生相同重量，记录堆顶是否过期、比较器是否让正确候选在堆顶、K 或窗口大小为边界值时是否稳定。

C++ 实现上，\code{std::priority_queue} 默认大顶堆，小顶堆需要比较器。堆中保存大对象时尽量保存下标或轻量记录；若不能删除任意元素，就用懒删除或过期检查。如果追问变成“最后一块石头 II 则是分组差值最小化，模型转为 0-1 背包”，先判断原来的状态是否仍然覆盖新答案集合，再决定是微调边界、改 value 语义，还是换成另一类结构。

\subsubsection{LeetCode 1353 最多可以参加的会议数目}

先每轮扫描所有候选找极值，确认答案选择的对象。本题的模型是每天只能参加一个正在进行的会议，应优先选择最早结束的会议；暴力版慢在动态候选集合不断变化，却每次重新找最值。本题的优化点是按开始日排序扫描日期，把当天可参加会议的结束日入小顶堆，弹出最早结束者，因此堆中只能放足以比较下一步选择的轻量状态。

堆只保证堆顶，不保证内部全局有序。本题的按开始日排序扫描日期，把当天可参加会议的结束日入小顶堆，弹出最早结束者要转成堆顶可回答的问题；如果元素会过期，还要在真正使用堆顶前清理或跳过旧状态。放到“堆与动态极值”这条主线里看，关键原则是：堆把插入和弹出极值控制在对数复杂度。Top K 用固定大小堆维护边界，合并多路序列用堆保存每一路当前头部，数据流中位数用双堆维护两半。

证明堆题时，要说明堆覆盖了所有可能的下一步候选，并且堆顶过期时不会被使用。本题的按开始日排序扫描日期，把当天可参加会议的结束日入小顶堆，弹出最早结束者是选择顺序，每天只能参加一个正在进行的会议，应优先选择最早结束的会议是候选集合来源。这道题的边界不能留到最后补丁式处理，实验时覆盖同一天多个会议、过期会议清理、空闲日期跳跃、结束日相同，记录堆顶是否过期、比较器是否让正确候选在堆顶、K 或窗口大小为边界值时是否稳定。

C++ 实现上，\code{std::priority_queue} 默认大顶堆，小顶堆需要比较器。堆中保存大对象时尽量保存下标或轻量记录；若不能删除任意元素，就用懒删除或过期检查。如果追问变成“这是扫描线和堆结合的调度题”，先判断原来的状态是否仍然覆盖新答案集合，再决定是微调边界、改 value 语义，还是换成另一类结构。

\subsubsection{LeetCode 1675 数组的最小偏移量}

先每轮扫描所有候选找极值，确认答案选择的对象。本题的模型是奇数可以先乘二，偶数可以不断除二，目标是缩小当前最大最小差；暴力版慢在动态候选集合不断变化，却每次重新找最值。本题的优化点是把所有数规范到可降低的偶数形态，用大顶堆维护当前最大值，同时记录当前最小值，因此堆中只能放足以比较下一步选择的轻量状态。

堆只保证堆顶，不保证内部全局有序。本题的把所有数规范到可降低的偶数形态，用大顶堆维护当前最大值，同时记录当前最小值要转成堆顶可回答的问题；如果元素会过期，还要在真正使用堆顶前清理或跳过旧状态。放到“堆与动态极值”这条主线里看，关键原则是：堆把插入和弹出极值控制在对数复杂度。Top K 用固定大小堆维护边界，合并多路序列用堆保存每一路当前头部，数据流中位数用双堆维护两半。

证明堆题时，要说明堆覆盖了所有可能的下一步候选，并且堆顶过期时不会被使用。本题的把所有数规范到可降低的偶数形态，用大顶堆维护当前最大值，同时记录当前最小值是选择顺序，奇数可以先乘二，偶数可以不断除二，目标是缩小当前最大最小差是候选集合来源。这道题的边界不能留到最后补丁式处理，实验时覆盖全奇数、全偶数、最大值变奇数停止、数值范围，记录堆顶是否过期、比较器是否让正确候选在堆顶、K 或窗口大小为边界值时是否稳定。

C++ 实现上，\code{std::priority_queue} 默认大顶堆，小顶堆需要比较器。堆中保存大对象时尽量保存下标或轻量记录；若不能删除任意元素，就用懒删除或过期检查。如果追问变成“这是堆维护动态极值和状态空间规范化的结合”，先判断原来的状态是否仍然覆盖新答案集合，再决定是微调边界、改 value 语义，还是换成另一类结构。

\subsubsection{LeetCode 1851 包含每个查询的最小区间}

先每轮扫描所有候选找极值，确认答案选择的对象。本题的模型是每个查询点需要找覆盖它的最短区间，区间和查询都可离线排序；暴力版慢在动态候选集合不断变化，却每次重新找最值。本题的优化点是按查询值升序扫描，把左端不大于查询的区间入堆，堆顶过期区间弹出，因此堆中只能放足以比较下一步选择的轻量状态。

堆只保证堆顶，不保证内部全局有序。本题的按查询值升序扫描，把左端不大于查询的区间入堆，堆顶过期区间弹出要转成堆顶可回答的问题；如果元素会过期，还要在真正使用堆顶前清理或跳过旧状态。放到“堆与动态极值”这条主线里看，关键原则是：堆把插入和弹出极值控制在对数复杂度。Top K 用固定大小堆维护边界，合并多路序列用堆保存每一路当前头部，数据流中位数用双堆维护两半。

证明堆题时，要说明堆覆盖了所有可能的下一步候选，并且堆顶过期时不会被使用。本题的按查询值升序扫描，把左端不大于查询的区间入堆，堆顶过期区间弹出是选择顺序，每个查询点需要找覆盖它的最短区间，区间和查询都可离线排序是候选集合来源。这道题的边界不能留到最后补丁式处理，实验时覆盖没有覆盖区间、区间长度相同、查询重复、右端过期，记录堆顶是否过期、比较器是否让正确候选在堆顶、K 或窗口大小为边界值时是否稳定。

C++ 实现上，\code{std::priority_queue} 默认大顶堆，小顶堆需要比较器。堆中保存大对象时尽量保存下标或轻量记录；若不能删除任意元素，就用懒删除或过期检查。如果追问变成“这是离线扫描线、堆和区间覆盖的组合题”，先判断原来的状态是否仍然覆盖新答案集合，再决定是微调边界、改 value 语义，还是换成另一类结构。

\subsection{链表与指针重连工作坊}

先用数组或多次扫描得到直观答案，再把操作改成节点重连。链表题的关键是保存后继、明确前驱和当前段边界，不能靠值交换掩盖指针关系。

\subsubsection{LeetCode 2 两数相加}

先把链表节点拷到数组里得到一个容易验证的版本，再回到节点重连。本题的模型是链表按低位到高位保存数字，本质是竖式加法而不是整数转换；暴力版能帮助确认节点顺序或节点身份。优化时关注维护两个游标和进位，任一链未结束或进位非零都继续生成节点，不要用值交换掩盖链表结构要求。

链表优化要把指针角色固定下来。本题用维护两个游标和进位，任一链未结束或进位非零都继续生成节点推进，至少要画出前驱、当前节点、后继节点和已处理链尾。每次改 next 前先保存后继，防止链断掉。放到“链表与指针重连”这条主线里看，关键原则是：优化来自哨兵节点、快慢指针、前驱指针和局部反转。链表题的核心不是值交换，而是保持断开和接回的顺序正确。

链表题的正确性靠结构不变量：已经处理好的链仍然连通，未处理部分仍然可达，二者之间只有明确连接点。本题的维护两个游标和进位，任一链未结束或进位非零都继续生成节点每执行一次，都不能破坏这三个条件。这道题的边界不能留到最后补丁式处理，实验时覆盖两链长度不同、最后进位、输入为零、不能用内置整数承载超长数字，每步画出前驱、当前、后继和下一段头。链表题不要只看输出值，还要检查节点是否丢失或成环。

C++ 实现上，平台管理输入节点时不要随意 \code{delete}。使用虚拟头节点统一头部边界；每次重连前保存后继。比较交点时比较节点地址，不比较节点值。如果追问变成“若链表高位在前，可以用栈反向处理或先反转链表”，先判断原来的状态是否仍然覆盖新答案集合，再决定是微调边界、改 value 语义，还是换成另一类结构。

\subsubsection{LeetCode 19 删除链表的倒数第 N 个结点}

先把链表节点拷到数组里得到一个容易验证的版本，再回到节点重连。本题的模型是倒数位置可以转成两个指针之间固定 n 步的距离；暴力版能帮助确认节点顺序或节点身份。优化时关注快指针先走 n 步，再让快慢指针同步移动，慢指针停在待删节点前驱，不要用值交换掩盖链表结构要求。

链表优化要把指针角色固定下来。本题用快指针先走 n 步，再让快慢指针同步移动，慢指针停在待删节点前驱推进，至少要画出前驱、当前节点、后继节点和已处理链尾。每次改 next 前先保存后继，防止链断掉。放到“链表与指针重连”这条主线里看，关键原则是：优化来自哨兵节点、快慢指针、前驱指针和局部反转。链表题的核心不是值交换，而是保持断开和接回的顺序正确。

链表题的正确性靠结构不变量：已经处理好的链仍然连通，未处理部分仍然可达，二者之间只有明确连接点。本题的快指针先走 n 步，再让快慢指针同步移动，慢指针停在待删节点前驱每执行一次，都不能破坏这三个条件。这道题的边界不能留到最后补丁式处理，实验时覆盖删除头节点、n 等于链表长度、只有一个节点、题目保证 n 合法，每步画出前驱、当前、后继和下一段头。链表题不要只看输出值，还要检查节点是否丢失或成环。

C++ 实现上，平台管理输入节点时不要随意 \code{delete}。使用虚拟头节点统一头部边界；每次重连前保存后继。比较交点时比较节点地址，不比较节点值。如果追问变成“若 n 不保证合法，需要在快指针预走阶段检测长度不足”，先判断原来的状态是否仍然覆盖新答案集合，再决定是微调边界、改 value 语义，还是换成另一类结构。

\subsubsection{LeetCode 21 合并两个有序链表}

先把链表节点拷到数组里得到一个容易验证的版本，再回到节点重连。本题的模型是两个链表已经各自有序，下一节点只能来自两个当前头中较小者；暴力版能帮助确认节点顺序或节点身份。优化时关注虚拟头统一结果链表，尾指针不断接上较小节点并推进来源链，不要用值交换掩盖链表结构要求。

链表优化要把指针角色固定下来。本题用虚拟头统一结果链表，尾指针不断接上较小节点并推进来源链推进，至少要画出前驱、当前节点、后继节点和已处理链尾。每次改 next 前先保存后继，防止链断掉。放到“链表与指针重连”这条主线里看，关键原则是：优化来自哨兵节点、快慢指针、前驱指针和局部反转。链表题的核心不是值交换，而是保持断开和接回的顺序正确。

链表题的正确性靠结构不变量：已经处理好的链仍然连通，未处理部分仍然可达，二者之间只有明确连接点。本题的虚拟头统一结果链表，尾指针不断接上较小节点并推进来源链每执行一次，都不能破坏这三个条件。这道题的边界不能留到最后补丁式处理，实验时覆盖某一条链为空、重复值、剩余整段接回、是否复用原节点，每步画出前驱、当前、后继和下一段头。链表题不要只看输出值，还要检查节点是否丢失或成环。

C++ 实现上，平台管理输入节点时不要随意 \code{delete}。使用虚拟头节点统一头部边界；每次重连前保存后继。比较交点时比较节点地址，不比较节点值。如果追问变成“合并 K 条链表可用分治或堆，核心仍是维护每条链的当前头”，先判断原来的状态是否仍然覆盖新答案集合，再决定是微调边界、改 value 语义，还是换成另一类结构。

\subsubsection{LeetCode 24 两两交换链表中的节点}

先把链表节点拷到数组里得到一个容易验证的版本，再回到节点重连。本题的模型是链表按固定长度为二的小组做局部重连；暴力版能帮助确认节点顺序或节点身份。优化时关注使用组前驱、第一节点、第二节点和下一组头，按顺序改指针并推进前驱，不要用值交换掩盖链表结构要求。

链表优化要把指针角色固定下来。本题用使用组前驱、第一节点、第二节点和下一组头，按顺序改指针并推进前驱推进，至少要画出前驱、当前节点、后继节点和已处理链尾。每次改 next 前先保存后继，防止链断掉。放到“链表与指针重连”这条主线里看，关键原则是：优化来自哨兵节点、快慢指针、前驱指针和局部反转。链表题的核心不是值交换，而是保持断开和接回的顺序正确。

链表题的正确性靠结构不变量：已经处理好的链仍然连通，未处理部分仍然可达，二者之间只有明确连接点。本题的使用组前驱、第一节点、第二节点和下一组头，按顺序改指针并推进前驱每执行一次，都不能破坏这三个条件。这道题的边界不能留到最后补丁式处理，实验时覆盖空链、单节点、奇数长度、交换后前驱更新到组尾，每步画出前驱、当前、后继和下一段头。链表题不要只看输出值，还要检查节点是否丢失或成环。

C++ 实现上，平台管理输入节点时不要随意 \code{delete}。使用虚拟头节点统一头部边界；每次重连前保存后继。比较交点时比较节点地址，不比较节点值。如果追问变成“K 个一组翻转是同一思想的推广，只是组内反转更长”，先判断原来的状态是否仍然覆盖新答案集合，再决定是微调边界、改 value 语义，还是换成另一类结构。

\subsubsection{LeetCode 25 K 个一组翻转链表}

先把链表节点拷到数组里得到一个容易验证的版本，再回到节点重连。本题的模型是链表被切成长度为 K 的连续组，只有完整组才反转；暴力版能帮助确认节点顺序或节点身份。优化时关注每轮先探测 K 个节点确认完整，再保存下一组头并做局部反转接回，不要用值交换掩盖链表结构要求。

链表优化要把指针角色固定下来。本题用每轮先探测 K 个节点确认完整，再保存下一组头并做局部反转接回推进，至少要画出前驱、当前节点、后继节点和已处理链尾。每次改 next 前先保存后继，防止链断掉。放到“链表与指针重连”这条主线里看，关键原则是：优化来自哨兵节点、快慢指针、前驱指针和局部反转。链表题的核心不是值交换，而是保持断开和接回的顺序正确。

链表题的正确性靠结构不变量：已经处理好的链仍然连通，未处理部分仍然可达，二者之间只有明确连接点。本题的每轮先探测 K 个节点确认完整，再保存下一组头并做局部反转接回每执行一次，都不能破坏这三个条件。这道题的边界不能留到最后补丁式处理，实验时覆盖不足 K 个不反转、K 为一、刚好整除、反转后头尾接错，每步画出前驱、当前、后继和下一段头。链表题不要只看输出值，还要检查节点是否丢失或成环。

C++ 实现上，平台管理输入节点时不要随意 \code{delete}。使用虚拟头节点统一头部边界；每次重连前保存后继。比较交点时比较节点地址，不比较节点值。如果追问变成“递归写法短但可能有栈风险，迭代写法更接近工程实现”，先判断原来的状态是否仍然覆盖新答案集合，再决定是微调边界、改 value 语义，还是换成另一类结构。

\subsubsection{LeetCode 61 旋转链表}

先把链表节点拷到数组里得到一个容易验证的版本，再回到节点重连。本题的模型是右旋 k 位等价于把链表尾部一段切到头部；暴力版能帮助确认节点顺序或节点身份。优化时关注先求长度并让 k 对长度取模，再找到新尾节点并重连成新头，不要用值交换掩盖链表结构要求。

链表优化要把指针角色固定下来。本题用先求长度并让 k 对长度取模，再找到新尾节点并重连成新头推进，至少要画出前驱、当前节点、后继节点和已处理链尾。每次改 next 前先保存后继，防止链断掉。放到“链表与指针重连”这条主线里看，关键原则是：优化来自哨兵节点、快慢指针、前驱指针和局部反转。链表题的核心不是值交换，而是保持断开和接回的顺序正确。

链表题的正确性靠结构不变量：已经处理好的链仍然连通，未处理部分仍然可达，二者之间只有明确连接点。本题的先求长度并让 k 对长度取模，再找到新尾节点并重连成新头每执行一次，都不能破坏这三个条件。这道题的边界不能留到最后补丁式处理，实验时覆盖空链、单节点、k 为零、k 大于长度、刚好整圈，每步画出前驱、当前、后继和下一段头。链表题不要只看输出值，还要检查节点是否丢失或成环。

C++ 实现上，平台管理输入节点时不要随意 \code{delete}。使用虚拟头节点统一头部边界；每次重连前保存后继。比较交点时比较节点地址，不比较节点值。如果追问变成“若本来把链表接成环，最后必须断开新尾的 next”，先判断原来的状态是否仍然覆盖新答案集合，再决定是微调边界、改 value 语义，还是换成另一类结构。

\subsubsection{LeetCode 82 删除排序链表中的重复元素 II}

先把链表节点拷到数组里得到一个容易验证的版本，再回到节点重连。本题的模型是有序链表中重复值连续出现，题目要求把所有重复值整组删除；暴力版能帮助确认节点顺序或节点身份。优化时关注虚拟头保存结果前驱，遇到一段重复值时跳过整段，否则前驱前进，不要用值交换掩盖链表结构要求。

链表优化要把指针角色固定下来。本题用虚拟头保存结果前驱，遇到一段重复值时跳过整段，否则前驱前进推进，至少要画出前驱、当前节点、后继节点和已处理链尾。每次改 next 前先保存后继，防止链断掉。放到“链表与指针重连”这条主线里看，关键原则是：优化来自哨兵节点、快慢指针、前驱指针和局部反转。链表题的核心不是值交换，而是保持断开和接回的顺序正确。

链表题的正确性靠结构不变量：已经处理好的链仍然连通，未处理部分仍然可达，二者之间只有明确连接点。本题的虚拟头保存结果前驱，遇到一段重复值时跳过整段，否则前驱前进每执行一次，都不能破坏这三个条件。这道题的边界不能留到最后补丁式处理，实验时覆盖头部重复、尾部重复、全部重复、没有重复，每步画出前驱、当前、后继和下一段头。链表题不要只看输出值，还要检查节点是否丢失或成环。

C++ 实现上，平台管理输入节点时不要随意 \code{delete}。使用虚拟头节点统一头部边界；每次重连前保存后继。比较交点时比较节点地址，不比较节点值。如果追问变成“若只保留一个重复值，就是删除重复元素的较简单版本”，先判断原来的状态是否仍然覆盖新答案集合，再决定是微调边界、改 value 语义，还是换成另一类结构。

\subsubsection{LeetCode 83 删除排序链表中的重复元素}

先把链表节点拷到数组里得到一个容易验证的版本，再回到节点重连。本题的模型是有序链表让重复节点相邻，只需保留每段相同值的第一个节点；暴力版能帮助确认节点顺序或节点身份。优化时关注当前节点和后继值相同就跳过后继，否则当前节点前进，不要用值交换掩盖链表结构要求。

链表优化要把指针角色固定下来。本题用当前节点和后继值相同就跳过后继，否则当前节点前进推进，至少要画出前驱、当前节点、后继节点和已处理链尾。每次改 next 前先保存后继，防止链断掉。放到“链表与指针重连”这条主线里看，关键原则是：优化来自哨兵节点、快慢指针、前驱指针和局部反转。链表题的核心不是值交换，而是保持断开和接回的顺序正确。

链表题的正确性靠结构不变量：已经处理好的链仍然连通，未处理部分仍然可达，二者之间只有明确连接点。本题的当前节点和后继值相同就跳过后继，否则当前节点前进每执行一次，都不能破坏这三个条件。这道题的边界不能留到最后补丁式处理，实验时覆盖空链、单节点、全重复、重复段在尾部，每步画出前驱、当前、后继和下一段头。链表题不要只看输出值，还要检查节点是否丢失或成环。

C++ 实现上，平台管理输入节点时不要随意 \code{delete}。使用虚拟头节点统一头部边界；每次重连前保存后继。比较交点时比较节点地址，不比较节点值。如果追问变成“和删除重复元素 II 的区别是是否整段删除”，先判断原来的状态是否仍然覆盖新答案集合，再决定是微调边界、改 value 语义，还是换成另一类结构。

\subsubsection{LeetCode 92 反转链表 II}

先把链表节点拷到数组里得到一个容易验证的版本，再回到节点重连。本题的模型是只反转链表中一段连续区间，区间外节点必须保持原连接；暴力版能帮助确认节点顺序或节点身份。优化时关注用虚拟头找到区间前驱，再用头插法把区间内部节点逐个移到前面，不要用值交换掩盖链表结构要求。

链表优化要把指针角色固定下来。本题用用虚拟头找到区间前驱，再用头插法把区间内部节点逐个移到前面推进，至少要画出前驱、当前节点、后继节点和已处理链尾。每次改 next 前先保存后继，防止链断掉。放到“链表与指针重连”这条主线里看，关键原则是：优化来自哨兵节点、快慢指针、前驱指针和局部反转。链表题的核心不是值交换，而是保持断开和接回的顺序正确。

链表题的正确性靠结构不变量：已经处理好的链仍然连通，未处理部分仍然可达，二者之间只有明确连接点。本题的用虚拟头找到区间前驱，再用头插法把区间内部节点逐个移到前面每执行一次，都不能破坏这三个条件。这道题的边界不能留到最后补丁式处理，实验时覆盖从头开始反转、只反转一个节点、反转到尾部、区间边界，每步画出前驱、当前、后继和下一段头。链表题不要只看输出值，还要检查节点是否丢失或成环。

C++ 实现上，平台管理输入节点时不要随意 \code{delete}。使用虚拟头节点统一头部边界；每次重连前保存后继。比较交点时比较节点地址，不比较节点值。如果追问变成“K 个一组反转是在多个固定长度区间上重复这一类局部重连”，先判断原来的状态是否仍然覆盖新答案集合，再决定是微调边界、改 value 语义，还是换成另一类结构。

\subsubsection{LeetCode 141 环形链表}

先把链表节点拷到数组里得到一个容易验证的版本，再回到节点重连。本题的模型是快慢指针在有环时必然相遇，无环时快指针先到空；暴力版能帮助确认节点顺序或节点身份。优化时关注不用哈希表也能常数空间检测环，不要用值交换掩盖链表结构要求。

链表优化要把指针角色固定下来。本题用不用哈希表也能常数空间检测环推进，至少要画出前驱、当前节点、后继节点和已处理链尾。每次改 next 前先保存后继，防止链断掉。放到“链表与指针重连”这条主线里看，关键原则是：优化来自哨兵节点、快慢指针、前驱指针和局部反转。链表题的核心不是值交换，而是保持断开和接回的顺序正确。

链表题的正确性靠结构不变量：已经处理好的链仍然连通，未处理部分仍然可达，二者之间只有明确连接点。本题的不用哈希表也能常数空间检测环每执行一次，都不能破坏这三个条件。这道题的边界不能留到最后补丁式处理，实验时覆盖空链、单节点自环、两节点环、尾部无环，每步画出前驱、当前、后继和下一段头。链表题不要只看输出值，还要检查节点是否丢失或成环。

C++ 实现上，平台管理输入节点时不要随意 \code{delete}。使用虚拟头节点统一头部边界；每次重连前保存后继。比较交点时比较节点地址，不比较节点值。如果追问变成“若要求环入口，用相遇后双指针同步走”，先判断原来的状态是否仍然覆盖新答案集合，再决定是微调边界、改 value 语义，还是换成另一类结构。

\subsubsection{LeetCode 142 环形链表 II}

先把链表节点拷到数组里得到一个容易验证的版本，再回到节点重连。本题的模型是相遇点到入口的距离关系来自快慢指针速度差；暴力版能帮助确认节点顺序或节点身份。优化时关注相遇后一个指针回头，两个指针每次走一步，再次相遇就是入口，不要用值交换掩盖链表结构要求。

链表优化要把指针角色固定下来。本题用相遇后一个指针回头，两个指针每次走一步，再次相遇就是入口推进，至少要画出前驱、当前节点、后继节点和已处理链尾。每次改 next 前先保存后继，防止链断掉。放到“链表与指针重连”这条主线里看，关键原则是：优化来自哨兵节点、快慢指针、前驱指针和局部反转。链表题的核心不是值交换，而是保持断开和接回的顺序正确。

链表题的正确性靠结构不变量：已经处理好的链仍然连通，未处理部分仍然可达，二者之间只有明确连接点。本题的相遇后一个指针回头，两个指针每次走一步，再次相遇就是入口每执行一次，都不能破坏这三个条件。这道题的边界不能留到最后补丁式处理，实验时覆盖入口在头节点、长尾短环、无环、单节点环，每步画出前驱、当前、后继和下一段头。链表题不要只看输出值，还要检查节点是否丢失或成环。

C++ 实现上，平台管理输入节点时不要随意 \code{delete}。使用虚拟头节点统一头部边界；每次重连前保存后继。比较交点时比较节点地址，不比较节点值。如果追问变成“也可用哈希集合，但空间更高”，先判断原来的状态是否仍然覆盖新答案集合，再决定是微调边界、改 value 语义，还是换成另一类结构。

\subsubsection{LeetCode 146 LRU 缓存}

先把链表节点拷到数组里得到一个容易验证的版本，再回到节点重连。本题的模型是哈希表负责按 key 定位，双向链表负责维护新旧顺序；暴力版能帮助确认节点顺序或节点身份。优化时关注访问和更新都把节点移动到头部，容量满时删除尾部最旧节点，不要用值交换掩盖链表结构要求。

链表优化要把指针角色固定下来。本题用访问和更新都把节点移动到头部，容量满时删除尾部最旧节点推进，至少要画出前驱、当前节点、后继节点和已处理链尾。每次改 next 前先保存后继，防止链断掉。放到“链表与指针重连”这条主线里看，关键原则是：优化来自哨兵节点、快慢指针、前驱指针和局部反转。链表题的核心不是值交换，而是保持断开和接回的顺序正确。

链表题的正确性靠结构不变量：已经处理好的链仍然连通，未处理部分仍然可达，二者之间只有明确连接点。本题的访问和更新都把节点移动到头部，容量满时删除尾部最旧节点每执行一次，都不能破坏这三个条件。这道题的边界不能留到最后补丁式处理，实验时覆盖容量为一、重复 put、get 不存在、更新已有值，每步画出前驱、当前、后继和下一段头。链表题不要只看输出值，还要检查节点是否丢失或成环。

C++ 实现上，平台管理输入节点时不要随意 \code{delete}。使用虚拟头节点统一头部边界；每次重连前保存后继。比较交点时比较节点地址，不比较节点值。如果追问变成“C++ 可用 \code{std::list} 加迭代器，也可手写双向链表理解所有权”，先判断原来的状态是否仍然覆盖新答案集合，再决定是微调边界、改 value 语义，还是换成另一类结构。

\subsubsection{LeetCode 148 排序链表}

先把链表节点拷到数组里得到一个容易验证的版本，再回到节点重连。本题的模型是链表不能随机访问，适合归并排序而不是快速排序数组 partition；暴力版能帮助确认节点顺序或节点身份。优化时关注自底向上迭代归并能避免递归栈，同时保持 O(n log n)，不要用值交换掩盖链表结构要求。

链表优化要把指针角色固定下来。本题用自底向上迭代归并能避免递归栈，同时保持 O(n log n)推进，至少要画出前驱、当前节点、后继节点和已处理链尾。每次改 next 前先保存后继，防止链断掉。放到“链表与指针重连”这条主线里看，关键原则是：优化来自哨兵节点、快慢指针、前驱指针和局部反转。链表题的核心不是值交换，而是保持断开和接回的顺序正确。

链表题的正确性靠结构不变量：已经处理好的链仍然连通，未处理部分仍然可达，二者之间只有明确连接点。本题的自底向上迭代归并能避免递归栈，同时保持 O(n log n)每执行一次，都不能破坏这三个条件。这道题的边界不能留到最后补丁式处理，实验时覆盖空链、单节点、重复值、奇数长度、切断子链，每步画出前驱、当前、后继和下一段头。链表题不要只看输出值，还要检查节点是否丢失或成环。

C++ 实现上，平台管理输入节点时不要随意 \code{delete}。使用虚拟头节点统一头部边界；每次重连前保存后继。比较交点时比较节点地址，不比较节点值。如果追问变成“若把值复制到数组排序，会改变节点身份语义”，先判断原来的状态是否仍然覆盖新答案集合，再决定是微调边界、改 value 语义，还是换成另一类结构。

\subsubsection{LeetCode 160 相交链表}

先把链表节点拷到数组里得到一个容易验证的版本，再回到节点重连。本题的模型是相交是节点地址相同，不是节点值相同；暴力版能帮助确认节点顺序或节点身份。优化时关注两个指针走完各自链表后切换到另一条链，走过总长度相同后相遇，不要用值交换掩盖链表结构要求。

链表优化要把指针角色固定下来。本题用两个指针走完各自链表后切换到另一条链，走过总长度相同后相遇推进，至少要画出前驱、当前节点、后继节点和已处理链尾。每次改 next 前先保存后继，防止链断掉。放到“链表与指针重连”这条主线里看，关键原则是：优化来自哨兵节点、快慢指针、前驱指针和局部反转。链表题的核心不是值交换，而是保持断开和接回的顺序正确。

链表题的正确性靠结构不变量：已经处理好的链仍然连通，未处理部分仍然可达，二者之间只有明确连接点。本题的两个指针走完各自链表后切换到另一条链，走过总长度相同后相遇每执行一次，都不能破坏这三个条件。这道题的边界不能留到最后补丁式处理，实验时覆盖不相交、头节点相交、尾节点相交、长度差很大，每步画出前驱、当前、后继和下一段头。链表题不要只看输出值，还要检查节点是否丢失或成环。

C++ 实现上，平台管理输入节点时不要随意 \code{delete}。使用虚拟头节点统一头部边界；每次重连前保存后继。比较交点时比较节点地址，不比较节点值。如果追问变成“哈希集合更直观但空间更高”，先判断原来的状态是否仍然覆盖新答案集合，再决定是微调边界、改 value 语义，还是换成另一类结构。

\subsubsection{LeetCode 206 反转链表}

先把链表节点拷到数组里得到一个容易验证的版本，再回到节点重连。本题的模型是每次把当前节点的 next 指向已经反转好的前缀头；暴力版能帮助确认节点顺序或节点身份。优化时关注先保存后继，再改指针，再推进 previous 和 current，不要用值交换掩盖链表结构要求。

链表优化要把指针角色固定下来。本题用先保存后继，再改指针，再推进 previous 和 current推进，至少要画出前驱、当前节点、后继节点和已处理链尾。每次改 next 前先保存后继，防止链断掉。放到“链表与指针重连”这条主线里看，关键原则是：优化来自哨兵节点、快慢指针、前驱指针和局部反转。链表题的核心不是值交换，而是保持断开和接回的顺序正确。

链表题的正确性靠结构不变量：已经处理好的链仍然连通，未处理部分仍然可达，二者之间只有明确连接点。本题的先保存后继，再改指针，再推进 previous 和 current每执行一次，都不能破坏这三个条件。这道题的边界不能留到最后补丁式处理，实验时覆盖空链、单节点、两个节点、不要丢失剩余链，每步画出前驱、当前、后继和下一段头。链表题不要只看输出值，还要检查节点是否丢失或成环。

C++ 实现上，平台管理输入节点时不要随意 \code{delete}。使用虚拟头节点统一头部边界；每次重连前保存后继。比较交点时比较节点地址，不比较节点值。如果追问变成“K 组反转是在这段局部反转外增加分组边界”，先判断原来的状态是否仍然覆盖新答案集合，再决定是微调边界、改 value 语义，还是换成另一类结构。

\subsubsection{LeetCode 234 回文链表}

先把链表节点拷到数组里得到一个容易验证的版本，再回到节点重连。本题的模型是链表回文要求前半段和后半段逆序后逐一相等；暴力版能帮助确认节点顺序或节点身份。优化时关注快慢指针找中点，反转后半段，再和前半段同步比较，不要用值交换掩盖链表结构要求。

链表优化要把指针角色固定下来。本题用快慢指针找中点，反转后半段，再和前半段同步比较推进，至少要画出前驱、当前节点、后继节点和已处理链尾。每次改 next 前先保存后继，防止链断掉。放到“链表与指针重连”这条主线里看，关键原则是：优化来自哨兵节点、快慢指针、前驱指针和局部反转。链表题的核心不是值交换，而是保持断开和接回的顺序正确。

链表题的正确性靠结构不变量：已经处理好的链仍然连通，未处理部分仍然可达，二者之间只有明确连接点。本题的快慢指针找中点，反转后半段，再和前半段同步比较每执行一次，都不能破坏这三个条件。这道题的边界不能留到最后补丁式处理，实验时覆盖空链、单节点、奇数长度、偶数长度、是否恢复链表，每步画出前驱、当前、后继和下一段头。链表题不要只看输出值，还要检查节点是否丢失或成环。

C++ 实现上，平台管理输入节点时不要随意 \code{delete}。使用虚拟头节点统一头部边界；每次重连前保存后继。比较交点时比较节点地址，不比较节点值。如果追问变成“如果不允许修改链表，可用栈保存前半段但空间变为线性”，先判断原来的状态是否仍然覆盖新答案集合，再决定是微调边界、改 value 语义，还是换成另一类结构。

\subsubsection{LeetCode 430 扁平化多级双向链表}

先把链表节点拷到数组里得到一个容易验证的版本，再回到节点重连。本题的模型是多级链表按深度优先顺序展开成一条双向链表；暴力版能帮助确认节点顺序或节点身份。优化时关注显式栈保存后续 next，遇到 child 先接 child，再在子链结束后恢复原 next，不要用值交换掩盖链表结构要求。

链表优化要把指针角色固定下来。本题用显式栈保存后续 next，遇到 child 先接 child，再在子链结束后恢复原 next推进，至少要画出前驱、当前节点、后继节点和已处理链尾。每次改 next 前先保存后继，防止链断掉。放到“链表与指针重连”这条主线里看，关键原则是：优化来自哨兵节点、快慢指针、前驱指针和局部反转。链表题的核心不是值交换，而是保持断开和接回的顺序正确。

链表题的正确性靠结构不变量：已经处理好的链仍然连通，未处理部分仍然可达，二者之间只有明确连接点。本题的显式栈保存后续 next，遇到 child 先接 child，再在子链结束后恢复原 next每执行一次，都不能破坏这三个条件。这道题的边界不能留到最后补丁式处理，实验时覆盖child 为空、next 和 child 同时存在、多层嵌套、prev 指针恢复，每步画出前驱、当前、后继和下一段头。链表题不要只看输出值，还要检查节点是否丢失或成环。

C++ 实现上，平台管理输入节点时不要随意 \code{delete}。使用虚拟头节点统一头部边界；每次重连前保存后继。比较交点时比较节点地址，不比较节点值。如果追问变成“它和二叉树前序展开类似，核心是保存未处理分支”，先判断原来的状态是否仍然覆盖新答案集合，再决定是微调边界、改 value 语义，还是换成另一类结构。

\subsection{二叉树与树形状态工作坊}

先想每个节点重复扫描子树会得到什么信息，再把这些信息压成节点向父节点返回的状态。树题要区分自顶向下约束、后序汇总和按层传播。

\subsubsection{LeetCode 94 二叉树中序遍历}

先想最直接的树遍历会在每个节点重复求什么。本题的模型是中序访问顺序是左子树、根、右子树，BST 中会得到有序序列；如果每个节点都重新扫描子树，慢点就会非常明显。本题要压缩的状态是用显式栈模拟一路向左，再回退访问根并转向右子树，也就是节点和父子方向之间真正传递的信息。

树题要先判定状态方向。本题的用显式栈模拟一路向左，再回退访问根并转向右子树可能是自顶向下的约束、后序向上的汇总，或队列中的层号。返回值或栈元素不能只有节点指针，还要带上题目需要的信息。放到“二叉树与树形状态”这条主线里看，关键原则是：树形 DP 把每个节点看成子问题根；BFS 把层次边界显式放进队列；BST 题利用中序有序或上下界约束。工程实现优先考虑显式栈避免深树调用栈风险。

树题常用对子树归纳。假设子节点状态已经正确，父节点只按用显式栈模拟一路向左，再回退访问根并转向右子树合并或传递信息。本题的中序访问顺序是左子树、根、右子树，BST 中会得到有序序列决定空节点、叶子和极端链式树如何处理。这道题的边界不能留到最后补丁式处理，实验时覆盖空树、单节点、链式左树、链式右树，尤其关注空节点、单节点、链式树和全负值。递归版本和显式栈版本可以互相对拍。

C++ 实现上，极端深树可能造成递归栈风险，工程写法可以用 \code{std::vector<TreeNode*>} 或队列显式保存状态。BST 边界最好用更宽整数或可选前驱值。如果追问变成“Morris 遍历可做到常数额外空间，但会临时改树结构”，先判断原来的状态是否仍然覆盖新答案集合，再决定是微调边界、改 value 语义，还是换成另一类结构。

\subsubsection{LeetCode 98 验证二叉搜索树}

先想最直接的树遍历会在每个节点重复求什么。本题的模型是BST 约束是全局上下界，不是只比较父子节点；如果每个节点都重新扫描子树，慢点就会非常明显。本题要压缩的状态是中序严格递增或显式传递上下界都可以；中序版本要保存前一个访问值，也就是节点和父子方向之间真正传递的信息。

树题要先判定状态方向。本题的中序严格递增或显式传递上下界都可以；中序版本要保存前一个访问值可能是自顶向下的约束、后序向上的汇总，或队列中的层号。返回值或栈元素不能只有节点指针，还要带上题目需要的信息。放到“二叉树与树形状态”这条主线里看，关键原则是：树形 DP 把每个节点看成子问题根；BFS 把层次边界显式放进队列；BST 题利用中序有序或上下界约束。工程实现优先考虑显式栈避免深树调用栈风险。

树题常用对子树归纳。假设子节点状态已经正确，父节点只按中序严格递增或显式传递上下界都可以；中序版本要保存前一个访问值合并或传递信息。本题的BST 约束是全局上下界，不是只比较父子节点决定空节点、叶子和极端链式树如何处理。这道题的边界不能留到最后补丁式处理，实验时覆盖重复值、int 最小最大值、局部合法但全局非法，尤其关注空节点、单节点、链式树和全负值。递归版本和显式栈版本可以互相对拍。

C++ 实现上，极端深树可能造成递归栈风险，工程写法可以用 \code{std::vector<TreeNode*>} 或队列显式保存状态。BST 边界最好用更宽整数或可选前驱值。如果追问变成“若允许重复值，需要明确重复放左还是放右”，先判断原来的状态是否仍然覆盖新答案集合，再决定是微调边界、改 value 语义，还是换成另一类结构。

\subsubsection{LeetCode 102 二叉树层序遍历}

先想最直接的树遍历会在每个节点重复求什么。本题的模型是队列在每轮开始时保存当前层所有节点；如果每个节点都重新扫描子树，慢点就会非常明显。本题要压缩的状态是记录 level size 固定层边界，避免下一层节点混入当前层，也就是节点和父子方向之间真正传递的信息。

树题要先判定状态方向。本题的记录 level size 固定层边界，避免下一层节点混入当前层可能是自顶向下的约束、后序向上的汇总，或队列中的层号。返回值或栈元素不能只有节点指针，还要带上题目需要的信息。放到“二叉树与树形状态”这条主线里看，关键原则是：树形 DP 把每个节点看成子问题根；BFS 把层次边界显式放进队列；BST 题利用中序有序或上下界约束。工程实现优先考虑显式栈避免深树调用栈风险。

树题常用对子树归纳。假设子节点状态已经正确，父节点只按记录 level size 固定层边界，避免下一层节点混入当前层合并或传递信息。本题的队列在每轮开始时保存当前层所有节点决定空节点、叶子和极端链式树如何处理。这道题的边界不能留到最后补丁式处理，实验时覆盖空树、只有根、最后一层不满、左右顺序，尤其关注空节点、单节点、链式树和全负值。递归版本和显式栈版本可以互相对拍。

C++ 实现上，极端深树可能造成递归栈风险，工程写法可以用 \code{std::vector<TreeNode*>} 或队列显式保存状态。BST 边界最好用更宽整数或可选前驱值。如果追问变成“锯齿形层序只是在每层输出顺序上增加方向控制”，先判断原来的状态是否仍然覆盖新答案集合，再决定是微调边界、改 value 语义，还是换成另一类结构。

\subsubsection{LeetCode 103 二叉树锯齿形层序遍历}

先想最直接的树遍历会在每个节点重复求什么。本题的模型是BFS 层边界不变，变化的是当前层写入结果的方向；如果每个节点都重新扫描子树，慢点就会非常明显。本题要压缩的状态是可以层内反转，也可以预分配数组按方向写入目标下标，也就是节点和父子方向之间真正传递的信息。

树题要先判定状态方向。本题的可以层内反转，也可以预分配数组按方向写入目标下标可能是自顶向下的约束、后序向上的汇总，或队列中的层号。返回值或栈元素不能只有节点指针，还要带上题目需要的信息。放到“二叉树与树形状态”这条主线里看，关键原则是：树形 DP 把每个节点看成子问题根；BFS 把层次边界显式放进队列；BST 题利用中序有序或上下界约束。工程实现优先考虑显式栈避免深树调用栈风险。

树题常用对子树归纳。假设子节点状态已经正确，父节点只按可以层内反转，也可以预分配数组按方向写入目标下标合并或传递信息。本题的BFS 层边界不变，变化的是当前层写入结果的方向决定空节点、叶子和极端链式树如何处理。这道题的边界不能留到最后补丁式处理，实验时覆盖单层、偶数层、奇数层、空树，尤其关注空节点、单节点、链式树和全负值。递归版本和显式栈版本可以互相对拍。

C++ 实现上，极端深树可能造成递归栈风险，工程写法可以用 \code{std::vector<TreeNode*>} 或队列显式保存状态。BST 边界最好用更宽整数或可选前驱值。如果追问变成“若要求从底向上输出，收集后反转层列表即可”，先判断原来的状态是否仍然覆盖新答案集合，再决定是微调边界、改 value 语义，还是换成另一类结构。

\subsubsection{LeetCode 104 二叉树最大深度}

先想最直接的树遍历会在每个节点重复求什么。本题的模型是深度可以看成层数，也可以作为栈中携带的状态；如果每个节点都重新扫描子树，慢点就会非常明显。本题要压缩的状态是显式栈保存节点和深度，避免极端链式树递归栈过深，也就是节点和父子方向之间真正传递的信息。

树题要先判定状态方向。本题的显式栈保存节点和深度，避免极端链式树递归栈过深可能是自顶向下的约束、后序向上的汇总，或队列中的层号。返回值或栈元素不能只有节点指针，还要带上题目需要的信息。放到“二叉树与树形状态”这条主线里看，关键原则是：树形 DP 把每个节点看成子问题根；BFS 把层次边界显式放进队列；BST 题利用中序有序或上下界约束。工程实现优先考虑显式栈避免深树调用栈风险。

树题常用对子树归纳。假设子节点状态已经正确，父节点只按显式栈保存节点和深度，避免极端链式树递归栈过深合并或传递信息。本题的深度可以看成层数，也可以作为栈中携带的状态决定空节点、叶子和极端链式树如何处理。这道题的边界不能留到最后补丁式处理，实验时覆盖空树返回零、根深度是一、链式树，尤其关注空节点、单节点、链式树和全负值。递归版本和显式栈版本可以互相对拍。

C++ 实现上，极端深树可能造成递归栈风险，工程写法可以用 \code{std::vector<TreeNode*>} 或队列显式保存状态。BST 边界最好用更宽整数或可选前驱值。如果追问变成“最小深度不能简单取左右最小，因为空子树不构成叶子路径”，先判断原来的状态是否仍然覆盖新答案集合，再决定是微调边界、改 value 语义，还是换成另一类结构。

\subsubsection{LeetCode 105 前序中序构造二叉树}

先想最直接的树遍历会在每个节点重复求什么。本题的模型是前序给根，中序把左右子树切开；如果每个节点都重新扫描子树，慢点就会非常明显。本题要压缩的状态是用哈希表记录中序下标，避免每次线性寻找根位置，也就是节点和父子方向之间真正传递的信息。

树题要先判定状态方向。本题的用哈希表记录中序下标，避免每次线性寻找根位置可能是自顶向下的约束、后序向上的汇总，或队列中的层号。返回值或栈元素不能只有节点指针，还要带上题目需要的信息。放到“二叉树与树形状态”这条主线里看，关键原则是：树形 DP 把每个节点看成子问题根；BFS 把层次边界显式放进队列；BST 题利用中序有序或上下界约束。工程实现优先考虑显式栈避免深树调用栈风险。

树题常用对子树归纳。假设子节点状态已经正确，父节点只按用哈希表记录中序下标，避免每次线性寻找根位置合并或传递信息。本题的前序给根，中序把左右子树切开决定空节点、叶子和极端链式树如何处理。这道题的边界不能留到最后补丁式处理，实验时覆盖节点值唯一、空子树、左右子树长度计算、输入不合法，尤其关注空节点、单节点、链式树和全负值。递归版本和显式栈版本可以互相对拍。

C++ 实现上，极端深树可能造成递归栈风险，工程写法可以用 \code{std::vector<TreeNode*>} 或队列显式保存状态。BST 边界最好用更宽整数或可选前驱值。如果追问变成“后序中序构造类似，只是根来自后序末尾”，先判断原来的状态是否仍然覆盖新答案集合，再决定是微调边界、改 value 语义，还是换成另一类结构。

\subsubsection{LeetCode 106 中序后序构造二叉树}

先想最直接的树遍历会在每个节点重复求什么。本题的模型是后序末尾是根，中序位置确定左右子树规模；如果每个节点都重新扫描子树，慢点就会非常明显。本题要压缩的状态是构造时要小心后序左右区间边界，避免切分错位，也就是节点和父子方向之间真正传递的信息。

树题要先判定状态方向。本题的构造时要小心后序左右区间边界，避免切分错位可能是自顶向下的约束、后序向上的汇总，或队列中的层号。返回值或栈元素不能只有节点指针，还要带上题目需要的信息。放到“二叉树与树形状态”这条主线里看，关键原则是：树形 DP 把每个节点看成子问题根；BFS 把层次边界显式放进队列；BST 题利用中序有序或上下界约束。工程实现优先考虑显式栈避免深树调用栈风险。

树题常用对子树归纳。假设子节点状态已经正确，父节点只按构造时要小心后序左右区间边界，避免切分错位合并或传递信息。本题的后序末尾是根，中序位置确定左右子树规模决定空节点、叶子和极端链式树如何处理。这道题的边界不能留到最后补丁式处理，实验时覆盖空树、单节点、全左链、全右链，尤其关注空节点、单节点、链式树和全负值。递归版本和显式栈版本可以互相对拍。

C++ 实现上，极端深树可能造成递归栈风险，工程写法可以用 \code{std::vector<TreeNode*>} 或队列显式保存状态。BST 边界最好用更宽整数或可选前驱值。如果追问变成“若给前序和后序，二叉树通常不能唯一确定，除非额外限制”，先判断原来的状态是否仍然覆盖新答案集合，再决定是微调边界、改 value 语义，还是换成另一类结构。

\subsubsection{LeetCode 108 有序数组转二叉搜索树}

先想最直接的树遍历会在每个节点重复求什么。本题的模型是有序数组中点作为根能让左右规模接近；如果每个节点都重新扫描子树，慢点就会非常明显。本题要压缩的状态是每次选择区间中点，左右子区间构造左右子树，也就是节点和父子方向之间真正传递的信息。

树题要先判定状态方向。本题的每次选择区间中点，左右子区间构造左右子树可能是自顶向下的约束、后序向上的汇总，或队列中的层号。返回值或栈元素不能只有节点指针，还要带上题目需要的信息。放到“二叉树与树形状态”这条主线里看，关键原则是：树形 DP 把每个节点看成子问题根；BFS 把层次边界显式放进队列；BST 题利用中序有序或上下界约束。工程实现优先考虑显式栈避免深树调用栈风险。

树题常用对子树归纳。假设子节点状态已经正确，父节点只按每次选择区间中点，左右子区间构造左右子树合并或传递信息。本题的有序数组中点作为根能让左右规模接近决定空节点、叶子和极端链式树如何处理。这道题的边界不能留到最后补丁式处理，实验时覆盖空区间、偶数长度选左中还是右中、高度平衡定义，尤其关注空节点、单节点、链式树和全负值。递归版本和显式栈版本可以互相对拍。

C++ 实现上，极端深树可能造成递归栈风险，工程写法可以用 \code{std::vector<TreeNode*>} 或队列显式保存状态。BST 边界最好用更宽整数或可选前驱值。如果追问变成“工程上递归可读，若严格避免递归可用队列保存区间和父指针”，先判断原来的状态是否仍然覆盖新答案集合，再决定是微调边界、改 value 语义，还是换成另一类结构。

\subsubsection{LeetCode 110 平衡二叉树}

先想最直接的树遍历会在每个节点重复求什么。本题的模型是每个节点需要知道子树高度以及是否已经失衡；如果每个节点都重新扫描子树，慢点就会非常明显。本题要压缩的状态是后序汇总时一旦子树失衡就向上返回失败标记，避免重复算高度，也就是节点和父子方向之间真正传递的信息。

树题要先判定状态方向。本题的后序汇总时一旦子树失衡就向上返回失败标记，避免重复算高度可能是自顶向下的约束、后序向上的汇总，或队列中的层号。返回值或栈元素不能只有节点指针，还要带上题目需要的信息。放到“二叉树与树形状态”这条主线里看，关键原则是：树形 DP 把每个节点看成子问题根；BFS 把层次边界显式放进队列；BST 题利用中序有序或上下界约束。工程实现优先考虑显式栈避免深树调用栈风险。

树题常用对子树归纳。假设子节点状态已经正确，父节点只按后序汇总时一旦子树失衡就向上返回失败标记，避免重复算高度合并或传递信息。本题的每个节点需要知道子树高度以及是否已经失衡决定空节点、叶子和极端链式树如何处理。这道题的边界不能留到最后补丁式处理，实验时覆盖空树、单节点、左右高度差正好一、链式树，尤其关注空节点、单节点、链式树和全负值。递归版本和显式栈版本可以互相对拍。

C++ 实现上，极端深树可能造成递归栈风险，工程写法可以用 \code{std::vector<TreeNode*>} 或队列显式保存状态。BST 边界最好用更宽整数或可选前驱值。如果追问变成“可以把返回值设计成 optional 高度，空表示不平衡”，先判断原来的状态是否仍然覆盖新答案集合，再决定是微调边界、改 value 语义，还是换成另一类结构。

\subsubsection{LeetCode 111 二叉树最小深度}

先想最直接的树遍历会在每个节点重复求什么。本题的模型是最小深度是到最近叶子的节点数，不是左右深度简单取最小；如果每个节点都重新扫描子树，慢点就会非常明显。本题要压缩的状态是BFS 第一次遇到叶子就是最小深度，因为按层扩展，也就是节点和父子方向之间真正传递的信息。

树题要先判定状态方向。本题的BFS 第一次遇到叶子就是最小深度，因为按层扩展可能是自顶向下的约束、后序向上的汇总，或队列中的层号。返回值或栈元素不能只有节点指针，还要带上题目需要的信息。放到“二叉树与树形状态”这条主线里看，关键原则是：树形 DP 把每个节点看成子问题根；BFS 把层次边界显式放进队列；BST 题利用中序有序或上下界约束。工程实现优先考虑显式栈避免深树调用栈风险。

树题常用对子树归纳。假设子节点状态已经正确，父节点只按BFS 第一次遇到叶子就是最小深度，因为按层扩展合并或传递信息。本题的最小深度是到最近叶子的节点数，不是左右深度简单取最小决定空节点、叶子和极端链式树如何处理。这道题的边界不能留到最后补丁式处理，实验时覆盖只有左子树、只有右子树、根为空、根是叶子，尤其关注空节点、单节点、链式树和全负值。递归版本和显式栈版本可以互相对拍。

C++ 实现上，极端深树可能造成递归栈风险，工程写法可以用 \code{std::vector<TreeNode*>} 或队列显式保存状态。BST 边界最好用更宽整数或可选前驱值。如果追问变成“DFS 写法必须特殊处理空子树不能作为最小路径”，先判断原来的状态是否仍然覆盖新答案集合，再决定是微调边界、改 value 语义，还是换成另一类结构。

\subsubsection{LeetCode 112 路径总和}

先想最直接的树遍历会在每个节点重复求什么。本题的模型是路径必须从根到叶子，状态是剩余目标值；如果每个节点都重新扫描子树，慢点就会非常明显。本题要压缩的状态是沿路径减去当前节点值，到叶子时检查剩余是否等于叶子值，也就是节点和父子方向之间真正传递的信息。

树题要先判定状态方向。本题的沿路径减去当前节点值，到叶子时检查剩余是否等于叶子值可能是自顶向下的约束、后序向上的汇总，或队列中的层号。返回值或栈元素不能只有节点指针，还要带上题目需要的信息。放到“二叉树与树形状态”这条主线里看，关键原则是：树形 DP 把每个节点看成子问题根；BFS 把层次边界显式放进队列；BST 题利用中序有序或上下界约束。工程实现优先考虑显式栈避免深树调用栈风险。

树题常用对子树归纳。假设子节点状态已经正确，父节点只按沿路径减去当前节点值，到叶子时检查剩余是否等于叶子值合并或传递信息。本题的路径必须从根到叶子，状态是剩余目标值决定空节点、叶子和极端链式树如何处理。这道题的边界不能留到最后补丁式处理，实验时覆盖负数节点、目标为零、非叶子提前达到目标、空树，尤其关注空节点、单节点、链式树和全负值。递归版本和显式栈版本可以互相对拍。

C++ 实现上，极端深树可能造成递归栈风险，工程写法可以用 \code{std::vector<TreeNode*>} 或队列显式保存状态。BST 边界最好用更宽整数或可选前驱值。如果追问变成“若要求列出所有路径，需要保存路径数组并在回退时弹出”，先判断原来的状态是否仍然覆盖新答案集合，再决定是微调边界、改 value 语义，还是换成另一类结构。

\subsubsection{LeetCode 113 路径总和 II}

先想最直接的树遍历会在每个节点重复求什么。本题的模型是相比判定题，状态还要保存当前路径；如果每个节点都重新扫描子树，慢点就会非常明显。本题要压缩的状态是到叶子且和满足时复制路径；回溯时恢复路径，也就是节点和父子方向之间真正传递的信息。

树题要先判定状态方向。本题的到叶子且和满足时复制路径；回溯时恢复路径可能是自顶向下的约束、后序向上的汇总，或队列中的层号。返回值或栈元素不能只有节点指针，还要带上题目需要的信息。放到“二叉树与树形状态”这条主线里看，关键原则是：树形 DP 把每个节点看成子问题根；BFS 把层次边界显式放进队列；BST 题利用中序有序或上下界约束。工程实现优先考虑显式栈避免深树调用栈风险。

树题常用对子树归纳。假设子节点状态已经正确，父节点只按到叶子且和满足时复制路径；回溯时恢复路径合并或传递信息。本题的相比判定题，状态还要保存当前路径决定空节点、叶子和极端链式树如何处理。这道题的边界不能留到最后补丁式处理，实验时覆盖多个答案、负数、空树、路径数组复制时机，尤其关注空节点、单节点、链式树和全负值。递归版本和显式栈版本可以互相对拍。

C++ 实现上，极端深树可能造成递归栈风险，工程写法可以用 \code{std::vector<TreeNode*>} 或队列显式保存状态。BST 边界最好用更宽整数或可选前驱值。如果追问变成“若路径不要求从根到叶，模型会变成前缀和加哈希”，先判断原来的状态是否仍然覆盖新答案集合，再决定是微调边界、改 value 语义，还是换成另一类结构。

\subsubsection{LeetCode 114 二叉树展开为链表}

先想最直接的树遍历会在每个节点重复求什么。本题的模型是目标是按前序顺序把树原地改成右指针链；如果每个节点都重新扫描子树，慢点就会非常明显。本题要压缩的状态是显式栈按右后左先入顺序处理，逐个接到前驱右侧，也就是节点和父子方向之间真正传递的信息。

树题要先判定状态方向。本题的显式栈按右后左先入顺序处理，逐个接到前驱右侧可能是自顶向下的约束、后序向上的汇总，或队列中的层号。返回值或栈元素不能只有节点指针，还要带上题目需要的信息。放到“二叉树与树形状态”这条主线里看，关键原则是：树形 DP 把每个节点看成子问题根；BFS 把层次边界显式放进队列；BST 题利用中序有序或上下界约束。工程实现优先考虑显式栈避免深树调用栈风险。

树题常用对子树归纳。假设子节点状态已经正确，父节点只按显式栈按右后左先入顺序处理，逐个接到前驱右侧合并或传递信息。本题的目标是按前序顺序把树原地改成右指针链决定空节点、叶子和极端链式树如何处理。这道题的边界不能留到最后补丁式处理，实验时覆盖空树、单节点、原有右子树不能丢、左指针置空，尤其关注空节点、单节点、链式树和全负值。递归版本和显式栈版本可以互相对拍。

C++ 实现上，极端深树可能造成递归栈风险，工程写法可以用 \code{std::vector<TreeNode*>} 或队列显式保存状态。BST 边界最好用更宽整数或可选前驱值。如果追问变成“也可用寻找左子树最右节点的原地方法”，先判断原来的状态是否仍然覆盖新答案集合，再决定是微调边界、改 value 语义，还是换成另一类结构。

\subsubsection{LeetCode 124 二叉树最大路径和}

先想最直接的树遍历会在每个节点重复求什么。本题的模型是路径可以从任意节点到任意节点，但不能分叉向父节点贡献两边；如果每个节点都重新扫描子树，慢点就会非常明显。本题要压缩的状态是每个节点向父亲贡献单边最大增益，全局答案可用左右两边加当前节点更新，也就是节点和父子方向之间真正传递的信息。

树题要先判定状态方向。本题的每个节点向父亲贡献单边最大增益，全局答案可用左右两边加当前节点更新可能是自顶向下的约束、后序向上的汇总，或队列中的层号。返回值或栈元素不能只有节点指针，还要带上题目需要的信息。放到“二叉树与树形状态”这条主线里看，关键原则是：树形 DP 把每个节点看成子问题根；BFS 把层次边界显式放进队列；BST 题利用中序有序或上下界约束。工程实现优先考虑显式栈避免深树调用栈风险。

树题常用对子树归纳。假设子节点状态已经正确，父节点只按每个节点向父亲贡献单边最大增益，全局答案可用左右两边加当前节点更新合并或传递信息。本题的路径可以从任意节点到任意节点，但不能分叉向父节点贡献两边决定空节点、叶子和极端链式树如何处理。这道题的边界不能留到最后补丁式处理，实验时覆盖全负数、单节点、左右增益为负时取零、答案初始值，尤其关注空节点、单节点、链式树和全负值。递归版本和显式栈版本可以互相对拍。

C++ 实现上，极端深树可能造成递归栈风险，工程写法可以用 \code{std::vector<TreeNode*>} 或队列显式保存状态。BST 边界最好用更宽整数或可选前驱值。如果追问变成“若路径必须从根到叶，状态和答案位置完全不同”，先判断原来的状态是否仍然覆盖新答案集合，再决定是微调边界、改 value 语义，还是换成另一类结构。

\subsubsection{LeetCode 199 二叉树右视图}

先想最直接的树遍历会在每个节点重复求什么。本题的模型是每层从右侧能看到的就是该层最后访问或最先访问的右侧节点；如果每个节点都重新扫描子树，慢点就会非常明显。本题要压缩的状态是BFS 固定层边界，取每层最后一个；或 DFS 按右优先首次到达深度，也就是节点和父子方向之间真正传递的信息。

树题要先判定状态方向。本题的BFS 固定层边界，取每层最后一个；或 DFS 按右优先首次到达深度可能是自顶向下的约束、后序向上的汇总，或队列中的层号。返回值或栈元素不能只有节点指针，还要带上题目需要的信息。放到“二叉树与树形状态”这条主线里看，关键原则是：树形 DP 把每个节点看成子问题根；BFS 把层次边界显式放进队列；BST 题利用中序有序或上下界约束。工程实现优先考虑显式栈避免深树调用栈风险。

树题常用对子树归纳。假设子节点状态已经正确，父节点只按BFS 固定层边界，取每层最后一个；或 DFS 按右优先首次到达深度合并或传递信息。本题的每层从右侧能看到的就是该层最后访问或最先访问的右侧节点决定空节点、叶子和极端链式树如何处理。这道题的边界不能留到最后补丁式处理，实验时覆盖空树、左链、右链、左右交错，尤其关注空节点、单节点、链式树和全负值。递归版本和显式栈版本可以互相对拍。

C++ 实现上，极端深树可能造成递归栈风险，工程写法可以用 \code{std::vector<TreeNode*>} 或队列显式保存状态。BST 边界最好用更宽整数或可选前驱值。如果追问变成“若求左视图，只改变层内取值方向”，先判断原来的状态是否仍然覆盖新答案集合，再决定是微调边界、改 value 语义，还是换成另一类结构。

\subsubsection{LeetCode 208 实现 Trie}

先想最直接的树遍历会在每个节点重复求什么。本题的模型是前缀树把字符串公共前缀共享成路径；如果每个节点都重新扫描子树，慢点就会非常明显。本题要压缩的状态是插入和查询都逐字符沿边移动，不存在的边按需要创建，也就是节点和父子方向之间真正传递的信息。

树题要先判定状态方向。本题的插入和查询都逐字符沿边移动，不存在的边按需要创建可能是自顶向下的约束、后序向上的汇总，或队列中的层号。返回值或栈元素不能只有节点指针，还要带上题目需要的信息。放到“二叉树与树形状态”这条主线里看，关键原则是：树形 DP 把每个节点看成子问题根；BFS 把层次边界显式放进队列；BST 题利用中序有序或上下界约束。工程实现优先考虑显式栈避免深树调用栈风险。

树题常用对子树归纳。假设子节点状态已经正确，父节点只按插入和查询都逐字符沿边移动，不存在的边按需要创建合并或传递信息。本题的前缀树把字符串公共前缀共享成路径决定空节点、叶子和极端链式树如何处理。这道题的边界不能留到最后补丁式处理，实验时覆盖空字符串约定、单词结束标记、前缀存在不等于单词存在，尤其关注空节点、单节点、链式树和全负值。递归版本和显式栈版本可以互相对拍。

C++ 实现上，极端深树可能造成递归栈风险，工程写法可以用 \code{std::vector<TreeNode*>} 或队列显式保存状态。BST 边界最好用更宽整数或可选前驱值。如果追问变成“若字符集很大，用哈希子节点；小写字母用数组更快”，先判断原来的状态是否仍然覆盖新答案集合，再决定是微调边界、改 value 语义，还是换成另一类结构。

\subsubsection{LeetCode 211 添加与搜索单词}

先想最直接的树遍历会在每个节点重复求什么。本题的模型是点号通配符让查询可能分叉，但前缀树仍能剪掉不匹配前缀；如果每个节点都重新扫描子树，慢点就会非常明显。本题要压缩的状态是普通字符沿唯一边走，点号枚举当前节点所有孩子，也就是节点和父子方向之间真正传递的信息。

树题要先判定状态方向。本题的普通字符沿唯一边走，点号枚举当前节点所有孩子可能是自顶向下的约束、后序向上的汇总，或队列中的层号。返回值或栈元素不能只有节点指针，还要带上题目需要的信息。放到“二叉树与树形状态”这条主线里看，关键原则是：树形 DP 把每个节点看成子问题根；BFS 把层次边界显式放进队列；BST 题利用中序有序或上下界约束。工程实现优先考虑显式栈避免深树调用栈风险。

树题常用对子树归纳。假设子节点状态已经正确，父节点只按普通字符沿唯一边走，点号枚举当前节点所有孩子合并或传递信息。本题的点号通配符让查询可能分叉，但前缀树仍能剪掉不匹配前缀决定空节点、叶子和极端链式树如何处理。这道题的边界不能留到最后补丁式处理，实验时覆盖连续通配符、空结果、单词结束标记、字符集大小，尤其关注空节点、单节点、链式树和全负值。递归版本和显式栈版本可以互相对拍。

C++ 实现上，极端深树可能造成递归栈风险，工程写法可以用 \code{std::vector<TreeNode*>} 或队列显式保存状态。BST 边界最好用更宽整数或可选前驱值。如果追问变成“大量通配符会接近遍历整棵 Trie，需要规模约束”，先判断原来的状态是否仍然覆盖新答案集合，再决定是微调边界、改 value 语义，还是换成另一类结构。

\subsubsection{LeetCode 226 翻转二叉树}

先想最直接的树遍历会在每个节点重复求什么。本题的模型是每个节点交换左右孩子即可，遍历顺序不影响最终结构；如果每个节点都重新扫描子树，慢点就会非常明显。本题要压缩的状态是显式队列或栈逐个节点交换，避免递归深度，也就是节点和父子方向之间真正传递的信息。

树题要先判定状态方向。本题的显式队列或栈逐个节点交换，避免递归深度可能是自顶向下的约束、后序向上的汇总，或队列中的层号。返回值或栈元素不能只有节点指针，还要带上题目需要的信息。放到“二叉树与树形状态”这条主线里看，关键原则是：树形 DP 把每个节点看成子问题根；BFS 把层次边界显式放进队列；BST 题利用中序有序或上下界约束。工程实现优先考虑显式栈避免深树调用栈风险。

树题常用对子树归纳。假设子节点状态已经正确，父节点只按显式队列或栈逐个节点交换，避免递归深度合并或传递信息。本题的每个节点交换左右孩子即可，遍历顺序不影响最终结构决定空节点、叶子和极端链式树如何处理。这道题的边界不能留到最后补丁式处理，实验时覆盖空树、单节点、只有一侧子树、重复值，尤其关注空节点、单节点、链式树和全负值。递归版本和显式栈版本可以互相对拍。

C++ 实现上，极端深树可能造成递归栈风险，工程写法可以用 \code{std::vector<TreeNode*>} 或队列显式保存状态。BST 边界最好用更宽整数或可选前驱值。如果追问变成“这题简单但适合训练树节点指针修改”，先判断原来的状态是否仍然覆盖新答案集合，再决定是微调边界、改 value 语义，还是换成另一类结构。

\subsubsection{LeetCode 230 二叉搜索树中第 K 小}

先想最直接的树遍历会在每个节点重复求什么。本题的模型是BST 中序遍历按升序访问节点；如果每个节点都重新扫描子树，慢点就会非常明显。本题要压缩的状态是显式栈中序，访问第 k 个节点时返回，也就是节点和父子方向之间真正传递的信息。

树题要先判定状态方向。本题的显式栈中序，访问第 k 个节点时返回可能是自顶向下的约束、后序向上的汇总，或队列中的层号。返回值或栈元素不能只有节点指针，还要带上题目需要的信息。放到“二叉树与树形状态”这条主线里看，关键原则是：树形 DP 把每个节点看成子问题根；BFS 把层次边界显式放进队列；BST 题利用中序有序或上下界约束。工程实现优先考虑显式栈避免深树调用栈风险。

树题常用对子树归纳。假设子节点状态已经正确，父节点只按显式栈中序，访问第 k 个节点时返回合并或传递信息。本题的BST 中序遍历按升序访问节点决定空节点、叶子和极端链式树如何处理。这道题的边界不能留到最后补丁式处理，实验时覆盖k 为一、k 为节点数、树不平衡、是否有重复值，尤其关注空节点、单节点、链式树和全负值。递归版本和显式栈版本可以互相对拍。

C++ 实现上，极端深树可能造成递归栈风险，工程写法可以用 \code{std::vector<TreeNode*>} 或队列显式保存状态。BST 边界最好用更宽整数或可选前驱值。如果追问变成“若频繁查询，可在节点维护子树大小成为顺序统计树”，先判断原来的状态是否仍然覆盖新答案集合，再决定是微调边界、改 value 语义，还是换成另一类结构。

\subsubsection{LeetCode 235 二叉搜索树最近公共祖先}

先想最直接的树遍历会在每个节点重复求什么。本题的模型是BST 的值域关系能决定两个目标在当前节点的哪一侧；如果每个节点都重新扫描子树，慢点就会非常明显。本题要压缩的状态是两个值都小去左，都大去右，否则当前节点就是分叉点，也就是节点和父子方向之间真正传递的信息。

树题要先判定状态方向。本题的两个值都小去左，都大去右，否则当前节点就是分叉点可能是自顶向下的约束、后序向上的汇总，或队列中的层号。返回值或栈元素不能只有节点指针，还要带上题目需要的信息。放到“二叉树与树形状态”这条主线里看，关键原则是：树形 DP 把每个节点看成子问题根；BFS 把层次边界显式放进队列；BST 题利用中序有序或上下界约束。工程实现优先考虑显式栈避免深树调用栈风险。

树题常用对子树归纳。假设子节点状态已经正确，父节点只按两个值都小去左，都大去右，否则当前节点就是分叉点合并或传递信息。本题的BST 的值域关系能决定两个目标在当前节点的哪一侧决定空节点、叶子和极端链式树如何处理。这道题的边界不能留到最后补丁式处理，实验时覆盖一个节点是另一个祖先、目标顺序、重复值约定，尤其关注空节点、单节点、链式树和全负值。递归版本和显式栈版本可以互相对拍。

C++ 实现上，极端深树可能造成递归栈风险，工程写法可以用 \code{std::vector<TreeNode*>} 或队列显式保存状态。BST 边界最好用更宽整数或可选前驱值。如果追问变成“普通二叉树不能用值域剪枝，需要父指针或后序状态”，先判断原来的状态是否仍然覆盖新答案集合，再决定是微调边界、改 value 语义，还是换成另一类结构。

\subsubsection{LeetCode 236 二叉树最近公共祖先}

先想最直接的树遍历会在每个节点重复求什么。本题的模型是普通树没有有序性，只能依靠祖先关系；如果每个节点都重新扫描子树，慢点就会非常明显。本题要压缩的状态是父指针集合或后序返回目标命中状态都可行，也就是节点和父子方向之间真正传递的信息。

树题要先判定状态方向。本题的父指针集合或后序返回目标命中状态都可行可能是自顶向下的约束、后序向上的汇总，或队列中的层号。返回值或栈元素不能只有节点指针，还要带上题目需要的信息。放到“二叉树与树形状态”这条主线里看，关键原则是：树形 DP 把每个节点看成子问题根；BFS 把层次边界显式放进队列；BST 题利用中序有序或上下界约束。工程实现优先考虑显式栈避免深树调用栈风险。

树题常用对子树归纳。假设子节点状态已经正确，父节点只按父指针集合或后序返回目标命中状态都可行合并或传递信息。本题的普通树没有有序性，只能依靠祖先关系决定空节点、叶子和极端链式树如何处理。这道题的边界不能留到最后补丁式处理，实验时覆盖目标不存在、p 等于 q、一个是另一个祖先、比较节点地址，尤其关注空节点、单节点、链式树和全负值。递归版本和显式栈版本可以互相对拍。

C++ 实现上，极端深树可能造成递归栈风险，工程写法可以用 \code{std::vector<TreeNode*>} 或队列显式保存状态。BST 边界最好用更宽整数或可选前驱值。如果追问变成“多次 LCA 查询可预处理倍增祖先表”，先判断原来的状态是否仍然覆盖新答案集合，再决定是微调边界、改 value 语义，还是换成另一类结构。

\subsubsection{LeetCode 337 打家劫舍 III}

先想最直接的树遍历会在每个节点重复求什么。本题的模型是树上相邻节点不能同时选，每个节点需要偷和不偷两个状态；如果每个节点都重新扫描子树，慢点就会非常明显。本题要压缩的状态是后序汇总：偷当前则孩子不能偷，不偷当前则孩子可偷可不偷，也就是节点和父子方向之间真正传递的信息。

树题要先判定状态方向。本题的后序汇总：偷当前则孩子不能偷，不偷当前则孩子可偷可不偷可能是自顶向下的约束、后序向上的汇总，或队列中的层号。返回值或栈元素不能只有节点指针，还要带上题目需要的信息。放到“二叉树与树形状态”这条主线里看，关键原则是：树形 DP 把每个节点看成子问题根；BFS 把层次边界显式放进队列；BST 题利用中序有序或上下界约束。工程实现优先考虑显式栈避免深树调用栈风险。

树题常用对子树归纳。假设子节点状态已经正确，父节点只按后序汇总：偷当前则孩子不能偷，不偷当前则孩子可偷可不偷合并或传递信息。本题的树上相邻节点不能同时选，每个节点需要偷和不偷两个状态决定空节点、叶子和极端链式树如何处理。这道题的边界不能留到最后补丁式处理，实验时覆盖空树、负值是否可能、链式树、递归深度，尤其关注空节点、单节点、链式树和全负值。递归版本和显式栈版本可以互相对拍。

C++ 实现上，极端深树可能造成递归栈风险，工程写法可以用 \code{std::vector<TreeNode*>} 或队列显式保存状态。BST 边界最好用更宽整数或可选前驱值。如果追问变成“可用显式栈后序保存节点状态，避免调用栈”，先判断原来的状态是否仍然覆盖新答案集合，再决定是微调边界、改 value 语义，还是换成另一类结构。

\subsubsection{LeetCode 450 删除二叉搜索树中的节点}

先想最直接的树遍历会在每个节点重复求什么。本题的模型是删除节点后仍要保持 BST 中序有序；如果每个节点都重新扫描子树，慢点就会非常明显。本题要压缩的状态是若有两个孩子，用后继或前驱替换当前值，再删除那个后继节点，也就是节点和父子方向之间真正传递的信息。

树题要先判定状态方向。本题的若有两个孩子，用后继或前驱替换当前值，再删除那个后继节点可能是自顶向下的约束、后序向上的汇总，或队列中的层号。返回值或栈元素不能只有节点指针，还要带上题目需要的信息。放到“二叉树与树形状态”这条主线里看，关键原则是：树形 DP 把每个节点看成子问题根；BFS 把层次边界显式放进队列；BST 题利用中序有序或上下界约束。工程实现优先考虑显式栈避免深树调用栈风险。

树题常用对子树归纳。假设子节点状态已经正确，父节点只按若有两个孩子，用后继或前驱替换当前值，再删除那个后继节点合并或传递信息。本题的删除节点后仍要保持 BST 中序有序决定空节点、叶子和极端链式树如何处理。这道题的边界不能留到最后补丁式处理，实验时覆盖删除叶子、一个孩子、两个孩子、删除根节点，尤其关注空节点、单节点、链式树和全负值。递归版本和显式栈版本可以互相对拍。

C++ 实现上，极端深树可能造成递归栈风险，工程写法可以用 \code{std::vector<TreeNode*>} 或队列显式保存状态。BST 边界最好用更宽整数或可选前驱值。如果追问变成“工程实现要清楚是移动值还是重连节点”，先判断原来的状态是否仍然覆盖新答案集合，再决定是微调边界、改 value 语义，还是换成另一类结构。

\subsubsection{LeetCode 543 二叉树直径}

先想最直接的树遍历会在每个节点重复求什么。本题的模型是直径经过某节点时等于左高加右高，但向父亲只能贡献单边高度；如果每个节点都重新扫描子树，慢点就会非常明显。本题要压缩的状态是后序遍历同时更新全局直径并返回高度，也就是节点和父子方向之间真正传递的信息。

树题要先判定状态方向。本题的后序遍历同时更新全局直径并返回高度可能是自顶向下的约束、后序向上的汇总，或队列中的层号。返回值或栈元素不能只有节点指针，还要带上题目需要的信息。放到“二叉树与树形状态”这条主线里看，关键原则是：树形 DP 把每个节点看成子问题根；BFS 把层次边界显式放进队列；BST 题利用中序有序或上下界约束。工程实现优先考虑显式栈避免深树调用栈风险。

树题常用对子树归纳。假设子节点状态已经正确，父节点只按后序遍历同时更新全局直径并返回高度合并或传递信息。本题的直径经过某节点时等于左高加右高，但向父亲只能贡献单边高度决定空节点、叶子和极端链式树如何处理。这道题的边界不能留到最后补丁式处理，实验时覆盖空树、单节点、链式树、边数还是节点数定义，尤其关注空节点、单节点、链式树和全负值。递归版本和显式栈版本可以互相对拍。

C++ 实现上，极端深树可能造成递归栈风险，工程写法可以用 \code{std::vector<TreeNode*>} 或队列显式保存状态。BST 边界最好用更宽整数或可选前驱值。如果追问变成“最大路径和是同型题，但贡献值和全局更新有负数处理”，先判断原来的状态是否仍然覆盖新答案集合，再决定是微调边界、改 value 语义，还是换成另一类结构。

\subsection{图建模、BFS 与 DFS工作坊}

先定义节点、边和状态，再决定 BFS、DFS、拓扑或并查集。图题失败常常不是搜索模板错，而是状态维度不完整或邻居生成过慢。

\subsubsection{LeetCode 127 单词接龙}

先画图，不先选搜索模板。本题的模型是每个单词是节点，一次修改一个字符形成边；朴素做法通常从多个起点反复搜索，或者把状态维度压得太少。优化首先要围绕通配模式桶加速邻居查询，BFS 第一次到达终点就是最短转换长度明确节点、边和访问标记。

图状态由节点和附加资源共同组成。本题的通配模式桶加速邻居查询，BFS 第一次到达终点就是最短转换长度要决定访问标记何时设置、状态是否只按位置去重、邻居如何生成。建模错了，BFS 或 DFS 再标准也会错。放到“图建模、BFS 与 DFS”这条主线里看，关键原则是：无权最短步数用 BFS，连通块和状态检测用 DFS 或显式栈，有向无环依赖用拓扑排序。多源问题把所有源同时入队，状态带资源的问题不能只按位置去重。

图搜索证明要回到可达性或层次。一次搜索必须覆盖从起点按每个单词是节点，一次修改一个字符形成边可达的全部状态，同时不能越过非法边；通配模式桶加速邻居查询，BFS 第一次到达终点就是最短转换长度保证重复状态不会反复入队或入栈。这道题的边界不能留到最后补丁式处理，实验时覆盖终点不在词表、起点和终点相同、桶重复扫描、单词长度一致，并打印入队时的完整状态。若同一位置但资源不同的状态被合并，很多图题会出现隐蔽错误。

C++ 实现上，连续编号图用 \code{std::vector<std::vector<int>>}，字符串节点先编号或使用哈希表。BFS 通常入队时标记访问，方向数组用固定常量集中维护。如果追问变成“双向 BFS 可以进一步降低搜索空间”，先判断原来的状态是否仍然覆盖新答案集合，再决定是微调边界、改 value 语义，还是换成另一类结构。

\subsubsection{LeetCode 130 被围绕的区域}

先画图，不先选搜索模板。本题的模型是只有不连通到边界的 O 才会被翻转；朴素做法通常从多个起点反复搜索，或者把状态维度压得太少。优化首先要围绕从边界 O 出发标记安全区域，再翻转剩余 O明确节点、边和访问标记。

图状态由节点和附加资源共同组成。本题的从边界 O 出发标记安全区域，再翻转剩余 O要决定访问标记何时设置、状态是否只按位置去重、邻居如何生成。建模错了，BFS 或 DFS 再标准也会错。放到“图建模、BFS 与 DFS”这条主线里看，关键原则是：无权最短步数用 BFS，连通块和状态检测用 DFS 或显式栈，有向无环依赖用拓扑排序。多源问题把所有源同时入队，状态带资源的问题不能只按位置去重。

图搜索证明要回到可达性或层次。一次搜索必须覆盖从起点按只有不连通到边界的 O 才会被翻转可达的全部状态，同时不能越过非法边；从边界 O 出发标记安全区域，再翻转剩余 O保证重复状态不会反复入队或入栈。这道题的边界不能留到最后补丁式处理，实验时覆盖全边界、无 O、单行单列、内部岛屿，并打印入队时的完整状态。若同一位置但资源不同的状态被合并，很多图题会出现隐蔽错误。

C++ 实现上，连续编号图用 \code{std::vector<std::vector<int>>}，字符串节点先编号或使用哈希表。BFS 通常入队时标记访问，方向数组用固定常量集中维护。如果追问变成“这是反向思维：找不能被翻转的区域比直接判断每块区域更简单”，先判断原来的状态是否仍然覆盖新答案集合，再决定是微调边界、改 value 语义，还是换成另一类结构。

\subsubsection{LeetCode 133 克隆图}

先画图，不先选搜索模板。本题的模型是每个原节点需要唯一对应一个新节点，边按原图连接；朴素做法通常从多个起点反复搜索，或者把状态维度压得太少。优化首先要围绕哈希表保存原节点到克隆节点的映射，BFS 或 DFS 遍历图明确节点、边和访问标记。

图状态由节点和附加资源共同组成。本题的哈希表保存原节点到克隆节点的映射，BFS 或 DFS 遍历图要决定访问标记何时设置、状态是否只按位置去重、邻居如何生成。建模错了，BFS 或 DFS 再标准也会错。放到“图建模、BFS 与 DFS”这条主线里看，关键原则是：无权最短步数用 BFS，连通块和状态检测用 DFS 或显式栈，有向无环依赖用拓扑排序。多源问题把所有源同时入队，状态带资源的问题不能只按位置去重。

图搜索证明要回到可达性或层次。一次搜索必须覆盖从起点按每个原节点需要唯一对应一个新节点，边按原图连接可达的全部状态，同时不能越过非法边；哈希表保存原节点到克隆节点的映射，BFS 或 DFS 遍历图保证重复状态不会反复入队或入栈。这道题的边界不能留到最后补丁式处理，实验时覆盖空图、自环、重边、环结构、不能按节点值唯一假设，并打印入队时的完整状态。若同一位置但资源不同的状态被合并，很多图题会出现隐蔽错误。

C++ 实现上，连续编号图用 \code{std::vector<std::vector<int>>}，字符串节点先编号或使用哈希表。BFS 通常入队时标记访问，方向数组用固定常量集中维护。如果追问变成“若图节点值不唯一，必须用地址作为键”，先判断原来的状态是否仍然覆盖新答案集合，再决定是微调边界、改 value 语义，还是换成另一类结构。

\subsubsection{LeetCode 200 岛屿数量}

先画图，不先选搜索模板。本题的模型是陆地格子是节点，四方向相邻是边，岛屿是连通块；朴素做法通常从多个起点反复搜索，或者把状态维度压得太少。优化首先要围绕扫描遇到未访问陆地就启动搜索并标记整块明确节点、边和访问标记。

图状态由节点和附加资源共同组成。本题的扫描遇到未访问陆地就启动搜索并标记整块要决定访问标记何时设置、状态是否只按位置去重、邻居如何生成。建模错了，BFS 或 DFS 再标准也会错。放到“图建模、BFS 与 DFS”这条主线里看，关键原则是：无权最短步数用 BFS，连通块和状态检测用 DFS 或显式栈，有向无环依赖用拓扑排序。多源问题把所有源同时入队，状态带资源的问题不能只按位置去重。

图搜索证明要回到可达性或层次。一次搜索必须覆盖从起点按陆地格子是节点，四方向相邻是边，岛屿是连通块可达的全部状态，同时不能越过非法边；扫描遇到未访问陆地就启动搜索并标记整块保证重复状态不会反复入队或入栈。这道题的边界不能留到最后补丁式处理，实验时覆盖空网格、全水、全陆、斜对角不连通、是否允许修改输入，并打印入队时的完整状态。若同一位置但资源不同的状态被合并，很多图题会出现隐蔽错误。

C++ 实现上，连续编号图用 \code{std::vector<std::vector<int>>}，字符串节点先编号或使用哈希表。BFS 通常入队时标记访问，方向数组用固定常量集中维护。如果追问变成“也可用并查集合并相邻陆地，适合动态岛屿扩展”，先判断原来的状态是否仍然覆盖新答案集合，再决定是微调边界、改 value 语义，还是换成另一类结构。

\subsubsection{LeetCode 207 课程表}

先画图，不先选搜索模板。本题的模型是课程是节点，先修关系是有向边，问题是判断是否有环；朴素做法通常从多个起点反复搜索，或者把状态维度压得太少。优化首先要围绕入度拓扑排序不断学习当前无依赖课程明确节点、边和访问标记。

图状态由节点和附加资源共同组成。本题的入度拓扑排序不断学习当前无依赖课程要决定访问标记何时设置、状态是否只按位置去重、邻居如何生成。建模错了，BFS 或 DFS 再标准也会错。放到“图建模、BFS 与 DFS”这条主线里看，关键原则是：无权最短步数用 BFS，连通块和状态检测用 DFS 或显式栈，有向无环依赖用拓扑排序。多源问题把所有源同时入队，状态带资源的问题不能只按位置去重。

图搜索证明要回到可达性或层次。一次搜索必须覆盖从起点按课程是节点，先修关系是有向边，问题是判断是否有环可达的全部状态，同时不能越过非法边；入度拓扑排序不断学习当前无依赖课程保证重复状态不会反复入队或入栈。这道题的边界不能留到最后补丁式处理，实验时覆盖边方向、孤立课程、环、自依赖、重复边，并打印入队时的完整状态。若同一位置但资源不同的状态被合并，很多图题会出现隐蔽错误。

C++ 实现上，连续编号图用 \code{std::vector<std::vector<int>>}，字符串节点先编号或使用哈希表。BFS 通常入队时标记访问，方向数组用固定常量集中维护。如果追问变成“DFS 三色标记也能判环，但要明确访问状态”，先判断原来的状态是否仍然覆盖新答案集合，再决定是微调边界、改 value 语义，还是换成另一类结构。

\subsubsection{LeetCode 210 课程表 II}

先画图，不先选搜索模板。本题的模型是在判环基础上还要输出一个可行拓扑序；朴素做法通常从多个起点反复搜索，或者把状态维度压得太少。优化首先要围绕每弹出一个入度为零课程，就把它加入顺序并删除出边明确节点、边和访问标记。

图状态由节点和附加资源共同组成。本题的每弹出一个入度为零课程，就把它加入顺序并删除出边要决定访问标记何时设置、状态是否只按位置去重、邻居如何生成。建模错了，BFS 或 DFS 再标准也会错。放到“图建模、BFS 与 DFS”这条主线里看，关键原则是：无权最短步数用 BFS，连通块和状态检测用 DFS 或显式栈，有向无环依赖用拓扑排序。多源问题把所有源同时入队，状态带资源的问题不能只按位置去重。

图搜索证明要回到可达性或层次。一次搜索必须覆盖从起点按在判环基础上还要输出一个可行拓扑序可达的全部状态，同时不能越过非法边；每弹出一个入度为零课程，就把它加入顺序并删除出边保证重复状态不会反复入队或入栈。这道题的边界不能留到最后补丁式处理，实验时覆盖有环返回空数组、多个可行顺序、边方向，并打印入队时的完整状态。若同一位置但资源不同的状态被合并，很多图题会出现隐蔽错误。

C++ 实现上，连续编号图用 \code{std::vector<std::vector<int>>}，字符串节点先编号或使用哈希表。BFS 通常入队时标记访问，方向数组用固定常量集中维护。如果追问变成“若需要字典序最小拓扑序，把队列换成小顶堆”，先判断原来的状态是否仍然覆盖新答案集合，再决定是微调边界、改 value 语义，还是换成另一类结构。

\subsubsection{LeetCode 417 太平洋大西洋水流问题}

先画图，不先选搜索模板。本题的模型是从每个格子向海搜索很慢，反向从海边沿非递减高度搜索；朴素做法通常从多个起点反复搜索，或者把状态维度压得太少。优化首先要围绕分别标记能到太平洋和大西洋的格子，交集就是答案明确节点、边和访问标记。

图状态由节点和附加资源共同组成。本题的分别标记能到太平洋和大西洋的格子，交集就是答案要决定访问标记何时设置、状态是否只按位置去重、邻居如何生成。建模错了，BFS 或 DFS 再标准也会错。放到“图建模、BFS 与 DFS”这条主线里看，关键原则是：无权最短步数用 BFS，连通块和状态检测用 DFS 或显式栈，有向无环依赖用拓扑排序。多源问题把所有源同时入队，状态带资源的问题不能只按位置去重。

图搜索证明要回到可达性或层次。一次搜索必须覆盖从起点按从每个格子向海搜索很慢，反向从海边沿非递减高度搜索可达的全部状态，同时不能越过非法边；分别标记能到太平洋和大西洋的格子，交集就是答案保证重复状态不会反复入队或入栈。这道题的边界不能留到最后补丁式处理，实验时覆盖单行单列、平台高度、边界重复入队、方向条件反向，并打印入队时的完整状态。若同一位置但资源不同的状态被合并，很多图题会出现隐蔽错误。

C++ 实现上，连续编号图用 \code{std::vector<std::vector<int>>}，字符串节点先编号或使用哈希表。BFS 通常入队时标记访问，方向数组用固定常量集中维护。如果追问变成“反向搜索是很多可达性题的重要技巧”，先判断原来的状态是否仍然覆盖新答案集合，再决定是微调边界、改 value 语义，还是换成另一类结构。

\subsubsection{LeetCode 695 岛屿的最大面积}

先画图，不先选搜索模板。本题的模型是面积就是一个连通块中陆地格子的数量；朴素做法通常从多个起点反复搜索，或者把状态维度压得太少。优化首先要围绕每发现新陆地就搜索完整连通块并计数，取最大值明确节点、边和访问标记。

图状态由节点和附加资源共同组成。本题的每发现新陆地就搜索完整连通块并计数，取最大值要决定访问标记何时设置、状态是否只按位置去重、邻居如何生成。建模错了，BFS 或 DFS 再标准也会错。放到“图建模、BFS 与 DFS”这条主线里看，关键原则是：无权最短步数用 BFS，连通块和状态检测用 DFS 或显式栈，有向无环依赖用拓扑排序。多源问题把所有源同时入队，状态带资源的问题不能只按位置去重。

图搜索证明要回到可达性或层次。一次搜索必须覆盖从起点按面积就是一个连通块中陆地格子的数量可达的全部状态，同时不能越过非法边；每发现新陆地就搜索完整连通块并计数，取最大值保证重复状态不会反复入队或入栈。这道题的边界不能留到最后补丁式处理，实验时覆盖全水、全陆、细长岛、是否允许修改输入，并打印入队时的完整状态。若同一位置但资源不同的状态被合并，很多图题会出现隐蔽错误。

C++ 实现上，连续编号图用 \code{std::vector<std::vector<int>>}，字符串节点先编号或使用哈希表。BFS 通常入队时标记访问，方向数组用固定常量集中维护。如果追问变成“与岛屿数量相比，只是连通块聚合值不同”，先判断原来的状态是否仍然覆盖新答案集合，再决定是微调边界、改 value 语义，还是换成另一类结构。

\subsubsection{LeetCode 752 打开转盘锁}

先画图，不先选搜索模板。本题的模型是四位密码是状态，每次转动一位形成边；朴素做法通常从多个起点反复搜索，或者把状态维度压得太少。优化首先要围绕BFS 从起点扩展，死亡状态不能入队，访问状态避免重复明确节点、边和访问标记。

图状态由节点和附加资源共同组成。本题的BFS 从起点扩展，死亡状态不能入队，访问状态避免重复要决定访问标记何时设置、状态是否只按位置去重、邻居如何生成。建模错了，BFS 或 DFS 再标准也会错。放到“图建模、BFS 与 DFS”这条主线里看，关键原则是：无权最短步数用 BFS，连通块和状态检测用 DFS 或显式栈，有向无环依赖用拓扑排序。多源问题把所有源同时入队，状态带资源的问题不能只按位置去重。

图搜索证明要回到可达性或层次。一次搜索必须覆盖从起点按四位密码是状态，每次转动一位形成边可达的全部状态，同时不能越过非法边；BFS 从起点扩展，死亡状态不能入队，访问状态避免重复保证重复状态不会反复入队或入栈。这道题的边界不能留到最后补丁式处理，实验时覆盖起点死亡、目标就是起点、数字环绕、死亡集合很大，并打印入队时的完整状态。若同一位置但资源不同的状态被合并，很多图题会出现隐蔽错误。

C++ 实现上，连续编号图用 \code{std::vector<std::vector<int>>}，字符串节点先编号或使用哈希表。BFS 通常入队时标记访问，方向数组用固定常量集中维护。如果追问变成“双向 BFS 能明显减少状态展开”，先判断原来的状态是否仍然覆盖新答案集合，再决定是微调边界、改 value 语义，还是换成另一类结构。

\subsubsection{LeetCode 864 获取所有钥匙的最短路径}

先画图，不先选搜索模板。本题的模型是位置相同但钥匙集合不同，未来能开的门不同；朴素做法通常从多个起点反复搜索，或者把状态维度压得太少。优化首先要围绕BFS 状态包含行、列和钥匙掩码，访问标记也必须包含掩码明确节点、边和访问标记。

图状态由节点和附加资源共同组成。本题的BFS 状态包含行、列和钥匙掩码，访问标记也必须包含掩码要决定访问标记何时设置、状态是否只按位置去重、邻居如何生成。建模错了，BFS 或 DFS 再标准也会错。放到“图建模、BFS 与 DFS”这条主线里看，关键原则是：无权最短步数用 BFS，连通块和状态检测用 DFS 或显式栈，有向无环依赖用拓扑排序。多源问题把所有源同时入队，状态带资源的问题不能只按位置去重。

图搜索证明要回到可达性或层次。一次搜索必须覆盖从起点按位置相同但钥匙集合不同，未来能开的门不同可达的全部状态，同时不能越过非法边；BFS 状态包含行、列和钥匙掩码，访问标记也必须包含掩码保证重复状态不会反复入队或入栈。这道题的边界不能留到最后补丁式处理，实验时覆盖起点旁边有门、钥匙顺序、不可达、钥匙数量决定掩码大小，并打印入队时的完整状态。若同一位置但资源不同的状态被合并，很多图题会出现隐蔽错误。

C++ 实现上，连续编号图用 \code{std::vector<std::vector<int>>}，字符串节点先编号或使用哈希表。BFS 通常入队时标记访问，方向数组用固定常量集中维护。如果追问变成“这是状态图维度不能只按位置去重的典型题”，先判断原来的状态是否仍然覆盖新答案集合，再决定是微调边界、改 value 语义，还是换成另一类结构。

\subsubsection{LeetCode 994 腐烂的橘子}

先画图，不先选搜索模板。本题的模型是所有初始腐烂橘子同时作为 BFS 源点；朴素做法通常从多个起点反复搜索，或者把状态维度压得太少。优化首先要围绕按层扩散，分钟数等于最后一个新鲜橘子被首次访问的层数明确节点、边和访问标记。

图状态由节点和附加资源共同组成。本题的按层扩散，分钟数等于最后一个新鲜橘子被首次访问的层数要决定访问标记何时设置、状态是否只按位置去重、邻居如何生成。建模错了，BFS 或 DFS 再标准也会错。放到“图建模、BFS 与 DFS”这条主线里看，关键原则是：无权最短步数用 BFS，连通块和状态检测用 DFS 或显式栈，有向无环依赖用拓扑排序。多源问题把所有源同时入队，状态带资源的问题不能只按位置去重。

图搜索证明要回到可达性或层次。一次搜索必须覆盖从起点按所有初始腐烂橘子同时作为 BFS 源点可达的全部状态，同时不能越过非法边；按层扩散，分钟数等于最后一个新鲜橘子被首次访问的层数保证重复状态不会反复入队或入栈。这道题的边界不能留到最后补丁式处理，实验时覆盖没有新鲜橘子、没有腐烂源、隔离新鲜橘子、空格阻断，并打印入队时的完整状态。若同一位置但资源不同的状态被合并，很多图题会出现隐蔽错误。

C++ 实现上，连续编号图用 \code{std::vector<std::vector<int>>}，字符串节点先编号或使用哈希表。BFS 通常入队时标记访问，方向数组用固定常量集中维护。如果追问变成“多源 BFS 也适用于最近零距离、地图扩散等题”，先判断原来的状态是否仍然覆盖新答案集合，再决定是微调边界、改 value 语义，还是换成另一类结构。

\subsubsection{LeetCode 1091 二进制矩阵中的最短路径}

先画图，不先选搜索模板。本题的模型是八方向网格无权边，答案是从左上到右下的最少格子数；朴素做法通常从多个起点反复搜索，或者把状态维度压得太少。优化首先要围绕BFS 入队时标记访问，层数或距离随状态保存明确节点、边和访问标记。

图状态由节点和附加资源共同组成。本题的BFS 入队时标记访问，层数或距离随状态保存要决定访问标记何时设置、状态是否只按位置去重、邻居如何生成。建模错了，BFS 或 DFS 再标准也会错。放到“图建模、BFS 与 DFS”这条主线里看，关键原则是：无权最短步数用 BFS，连通块和状态检测用 DFS 或显式栈，有向无环依赖用拓扑排序。多源问题把所有源同时入队，状态带资源的问题不能只按位置去重。

图搜索证明要回到可达性或层次。一次搜索必须覆盖从起点按八方向网格无权边，答案是从左上到右下的最少格子数可达的全部状态，同时不能越过非法边；BFS 入队时标记访问，层数或距离随状态保存保证重复状态不会反复入队或入栈。这道题的边界不能留到最后补丁式处理，实验时覆盖起点或终点阻塞、单格、八方向越界、无路可达，并打印入队时的完整状态。若同一位置但资源不同的状态被合并，很多图题会出现隐蔽错误。

C++ 实现上，连续编号图用 \code{std::vector<std::vector<int>>}，字符串节点先编号或使用哈希表。BFS 通常入队时标记访问，方向数组用固定常量集中维护。如果追问变成“边权相同才能用 BFS，若每格有代价就要最短路”，先判断原来的状态是否仍然覆盖新答案集合，再决定是微调边界、改 value 语义，还是换成另一类结构。

\subsubsection{LeetCode 1293 网格中的最短路径}

先画图，不先选搜索模板。本题的模型是走到同一格时，剩余可消除障碍数量不同会影响未来可达性；朴素做法通常从多个起点反复搜索，或者把状态维度压得太少。优化首先要围绕BFS 状态包含位置和剩余消除次数，visited 记录该维度或记录到达该格的最大剩余次数明确节点、边和访问标记。

图状态由节点和附加资源共同组成。本题的BFS 状态包含位置和剩余消除次数，visited 记录该维度或记录到达该格的最大剩余次数要决定访问标记何时设置、状态是否只按位置去重、邻居如何生成。建模错了，BFS 或 DFS 再标准也会错。放到“图建模、BFS 与 DFS”这条主线里看，关键原则是：无权最短步数用 BFS，连通块和状态检测用 DFS 或显式栈，有向无环依赖用拓扑排序。多源问题把所有源同时入队，状态带资源的问题不能只按位置去重。

图搜索证明要回到可达性或层次。一次搜索必须覆盖从起点按走到同一格时，剩余可消除障碍数量不同会影响未来可达性可达的全部状态，同时不能越过非法边；BFS 状态包含位置和剩余消除次数，visited 记录该维度或记录到达该格的最大剩余次数保证重复状态不会反复入队或入栈。这道题的边界不能留到最后补丁式处理，实验时覆盖起点终点、k 很大、障碍密集、二维 visited 错误剪枝，并打印入队时的完整状态。若同一位置但资源不同的状态被合并，很多图题会出现隐蔽错误。

C++ 实现上，连续编号图用 \code{std::vector<std::vector<int>>}，字符串节点先编号或使用哈希表。BFS 通常入队时标记访问，方向数组用固定常量集中维护。如果追问变成“它和钥匙题一样说明状态维度不能丢”，先判断原来的状态是否仍然覆盖新答案集合，再决定是微调边界、改 value 语义，还是换成另一类结构。

\subsection{最短路与带权图工作坊}

先确认目标是步数最少还是代价最小，再检查边权是否非负。非负权用 Dijkstra，等权用 BFS，限制边数用分层松弛或 Bellman-Ford 思想。

\subsubsection{LeetCode 505 迷宫 II}

先区分步数和代价。本题的模型是小球一旦选择方向就会滚到墙前停止，边权是滚动经过的格子数；若暴力枚举路径，路径数量会爆炸。优化点邻居生成不是走一步，而是沿四个方向模拟滚动到停点，再用 Dijkstra 按累计距离松弛说明下一步扩展必须按某种代价顺序组织，而不是普通队列随便推进。

最短路状态至少包含当前节点和当前代价。本题的邻居生成不是走一步，而是沿四个方向模拟滚动到停点，再用 Dijkstra 按累计距离松弛决定优先队列按什么排序，以及旧状态如何被识别。非负边条件、代价合成方式和不可达处理都要写清楚。放到“最短路与带权图”这条主线里看，关键原则是：非负权图用 Dijkstra，小顶堆保存当前距离最小的候选；边权为一用 BFS；有负边要考虑 Bellman-Ford 或重新建模。路径代价若是最大边权，也可用 Dijkstra 或二分答案。

最短路证明依赖扩展顺序。若本题使用 Dijkstra，就要说明非负代价让弹出的未过期状态已经最优；若用二分可达性，就要说明阈值越大可用边只会增加。这道题的边界不能留到最后补丁式处理，实验时覆盖起点就是终点、目标无法停下、绕远路更短、同一停点多次被松弛，打印每次弹出的代价和是否过期。若普通 BFS 与 Dijkstra 结果不同，要能解释边权条件造成的差异。

C++ 实现上，Dijkstra 用邻接表保存 \code{(to, weight)}，距离数组用 \code{long long}。由于 \code{priority_queue} 没有 decrease-key，压入新状态后弹出时跳过旧状态。如果追问变成“若只问能否到达，可用 BFS 或 DFS；若问最少距离，就必须处理带权边”，先判断原来的状态是否仍然覆盖新答案集合，再决定是微调边界、改 value 语义，还是换成另一类结构。

\subsubsection{LeetCode 743 网络延迟时间}

先区分步数和代价。本题的模型是有向正权图从起点到所有节点的最短路最大值；若暴力枚举路径，路径数量会爆炸。优化点Dijkstra 求每个节点最短距离，最后取最大，若有无穷则不可达说明下一步扩展必须按某种代价顺序组织，而不是普通队列随便推进。

最短路状态至少包含当前节点和当前代价。本题的Dijkstra 求每个节点最短距离，最后取最大，若有无穷则不可达决定优先队列按什么排序，以及旧状态如何被识别。非负边条件、代价合成方式和不可达处理都要写清楚。放到“最短路与带权图”这条主线里看，关键原则是：非负权图用 Dijkstra，小顶堆保存当前距离最小的候选；边权为一用 BFS；有负边要考虑 Bellman-Ford 或重新建模。路径代价若是最大边权，也可用 Dijkstra 或二分答案。

最短路证明依赖扩展顺序。若本题使用 Dijkstra，就要说明非负代价让弹出的未过期状态已经最优；若用二分可达性，就要说明阈值越大可用边只会增加。这道题的边界不能留到最后补丁式处理，实验时覆盖节点编号从一开始、不可达、重边、过期堆元素，打印每次弹出的代价和是否过期。若普通 BFS 与 Dijkstra 结果不同，要能解释边权条件造成的差异。

C++ 实现上，Dijkstra 用邻接表保存 \code{(to, weight)}，距离数组用 \code{long long}。由于 \code{priority_queue} 没有 decrease-key，压入新状态后弹出时跳过旧状态。如果追问变成“边权相同才可用普通 BFS”，先判断原来的状态是否仍然覆盖新答案集合，再决定是微调边界、改 value 语义，还是换成另一类结构。

\subsubsection{LeetCode 778 水位上升的泳池中游泳}

先区分步数和代价。本题的模型是到达某格子的代价是路径上最大高度，而不是高度和；若暴力枚举路径，路径数量会爆炸。优化点用 Dijkstra 按当前路径最大高度扩展，松弛时取 max 当前代价和邻格高度说明下一步扩展必须按某种代价顺序组织，而不是普通队列随便推进。

最短路状态至少包含当前节点和当前代价。本题的用 Dijkstra 按当前路径最大高度扩展，松弛时取 max 当前代价和邻格高度决定优先队列按什么排序，以及旧状态如何被识别。非负边条件、代价合成方式和不可达处理都要写清楚。放到“最短路与带权图”这条主线里看，关键原则是：非负权图用 Dijkstra，小顶堆保存当前距离最小的候选；边权为一用 BFS；有负边要考虑 Bellman-Ford 或重新建模。路径代价若是最大边权，也可用 Dijkstra 或二分答案。

最短路证明依赖扩展顺序。若本题使用 Dijkstra，就要说明非负代价让弹出的未过期状态已经最优；若用二分可达性，就要说明阈值越大可用边只会增加。这道题的边界不能留到最后补丁式处理，实验时覆盖单格、起点高度很高、需要绕路、多个相同高度，打印每次弹出的代价和是否过期。若普通 BFS 与 Dijkstra 结果不同，要能解释边权条件造成的差异。

C++ 实现上，Dijkstra 用邻接表保存 \code{(to, weight)}，距离数组用 \code{long long}。由于 \code{priority_queue} 没有 decrease-key，压入新状态后弹出时跳过旧状态。如果追问变成“也可以按水位二分可达，或按高度排序并查集合并”，先判断原来的状态是否仍然覆盖新答案集合，再决定是微调边界、改 value 语义，还是换成另一类结构。

\subsubsection{LeetCode 787 K 站中转内最便宜的航班}

先区分步数和代价。本题的模型是路径代价受中转次数限制，不能只按费用最小弹出节点后永久确定；若暴力枚举路径，路径数量会爆炸。优化点用按边数分层的 Bellman-Ford 或带层数状态的优先队列，状态包含城市和已用边数说明下一步扩展必须按某种代价顺序组织，而不是普通队列随便推进。

最短路状态至少包含当前节点和当前代价。本题的用按边数分层的 Bellman-Ford 或带层数状态的优先队列，状态包含城市和已用边数决定优先队列按什么排序，以及旧状态如何被识别。非负边条件、代价合成方式和不可达处理都要写清楚。放到“最短路与带权图”这条主线里看，关键原则是：非负权图用 Dijkstra，小顶堆保存当前距离最小的候选；边权为一用 BFS；有负边要考虑 Bellman-Ford 或重新建模。路径代价若是最大边权，也可用 Dijkstra 或二分答案。

最短路证明依赖扩展顺序。若本题使用 Dijkstra，就要说明非负代价让弹出的未过期状态已经最优；若用二分可达性，就要说明阈值越大可用边只会增加。这道题的边界不能留到最后补丁式处理，实验时覆盖起点终点相同、不可达、K 为零、便宜路径中转过多，打印每次弹出的代价和是否过期。若普通 BFS 与 Dijkstra 结果不同，要能解释边权条件造成的差异。

C++ 实现上，Dijkstra 用邻接表保存 \code{(to, weight)}，距离数组用 \code{long long}。由于 \code{priority_queue} 没有 decrease-key，压入新状态后弹出时跳过旧状态。如果追问变成“若去掉中转限制，普通 Dijkstra 才能直接按费用组织扩展”，先判断原来的状态是否仍然覆盖新答案集合，再决定是微调边界、改 value 语义，还是换成另一类结构。

\subsubsection{LeetCode 1102 得分最高的路径}

先区分步数和代价。本题的模型是路径得分是沿途最小格子值，要最大化这个最小值；若暴力枚举路径，路径数量会爆炸。优化点用最大堆按当前路径得分扩展，进入邻格时取当前得分和邻格值的较小者说明下一步扩展必须按某种代价顺序组织，而不是普通队列随便推进。

最短路状态至少包含当前节点和当前代价。本题的用最大堆按当前路径得分扩展，进入邻格时取当前得分和邻格值的较小者决定优先队列按什么排序，以及旧状态如何被识别。非负边条件、代价合成方式和不可达处理都要写清楚。放到“最短路与带权图”这条主线里看，关键原则是：非负权图用 Dijkstra，小顶堆保存当前距离最小的候选；边权为一用 BFS；有负边要考虑 Bellman-Ford 或重新建模。路径代价若是最大边权，也可用 Dijkstra 或二分答案。

最短路证明依赖扩展顺序。若本题使用 Dijkstra，就要说明非负代价让弹出的未过期状态已经最优；若用二分可达性，就要说明阈值越大可用边只会增加。这道题的边界不能留到最后补丁式处理，实验时覆盖单格、起点或终点很小、必须绕路、多个相同得分路径，打印每次弹出的代价和是否过期。若普通 BFS 与 Dijkstra 结果不同，要能解释边权条件造成的差异。

C++ 实现上，Dijkstra 用邻接表保存 \code{(to, weight)}，距离数组用 \code{long long}。由于 \code{priority_queue} 没有 decrease-key，压入新状态后弹出时跳过旧状态。如果追问变成“也可把阈值二分后检查连通性，或按值从高到低并查集合并”，先判断原来的状态是否仍然覆盖新答案集合，再决定是微调边界、改 value 语义，还是换成另一类结构。

\subsubsection{LeetCode 1334 阈值距离内邻居最少的城市}

先区分步数和代价。本题的模型是每个城市都需要知道到其他城市的最短距离是否不超过阈值；若暴力枚举路径，路径数量会爆炸。优化点节点数较小时可用 Floyd，边稀疏时也可从每个点跑 Dijkstra说明下一步扩展必须按某种代价顺序组织，而不是普通队列随便推进。

最短路状态至少包含当前节点和当前代价。本题的节点数较小时可用 Floyd，边稀疏时也可从每个点跑 Dijkstra决定优先队列按什么排序，以及旧状态如何被识别。非负边条件、代价合成方式和不可达处理都要写清楚。放到“最短路与带权图”这条主线里看，关键原则是：非负权图用 Dijkstra，小顶堆保存当前距离最小的候选；边权为一用 BFS；有负边要考虑 Bellman-Ford 或重新建模。路径代价若是最大边权，也可用 Dijkstra 或二分答案。

最短路证明依赖扩展顺序。若本题使用 Dijkstra，就要说明非负代价让弹出的未过期状态已经最优；若用二分可达性，就要说明阈值越大可用边只会增加。这道题的边界不能留到最后补丁式处理，实验时覆盖距离等于阈值、并列时选编号最大、不可达、无向边，打印每次弹出的代价和是否过期。若普通 BFS 与 Dijkstra 结果不同，要能解释边权条件造成的差异。

C++ 实现上，Dijkstra 用邻接表保存 \code{(to, weight)}，距离数组用 \code{long long}。由于 \code{priority_queue} 没有 decrease-key，压入新状态后弹出时跳过旧状态。如果追问变成“这是多源最短路选择问题，算法取决于 n、边数和查询密度”，先判断原来的状态是否仍然覆盖新答案集合，再决定是微调边界、改 value 语义，还是换成另一类结构。

\subsubsection{LeetCode 1368 使网格图至少有一条有效路径的最小代价}

先区分步数和代价。本题的模型是沿着格子箭头走代价为零，改方向走代价为一；若暴力枚举路径，路径数量会爆炸。优化点边权只有零和一，用 0-1 BFS：零代价邻居入队首，一代价邻居入队尾说明下一步扩展必须按某种代价顺序组织，而不是普通队列随便推进。

最短路状态至少包含当前节点和当前代价。本题的边权只有零和一，用 0-1 BFS：零代价邻居入队首，一代价邻居入队尾决定优先队列按什么排序，以及旧状态如何被识别。非负边条件、代价合成方式和不可达处理都要写清楚。放到“最短路与带权图”这条主线里看，关键原则是：非负权图用 Dijkstra，小顶堆保存当前距离最小的候选；边权为一用 BFS；有负边要考虑 Bellman-Ford 或重新建模。路径代价若是最大边权，也可用 Dijkstra 或二分答案。

最短路证明依赖扩展顺序。若本题使用 Dijkstra，就要说明非负代价让弹出的未过期状态已经最优；若用二分可达性，就要说明阈值越大可用边只会增加。这道题的边界不能留到最后补丁式处理，实验时覆盖单格、箭头指向越界、多个零代价路径、过期状态，打印每次弹出的代价和是否过期。若普通 BFS 与 Dijkstra 结果不同，要能解释边权条件造成的差异。

C++ 实现上，Dijkstra 用邻接表保存 \code{(to, weight)}，距离数组用 \code{long long}。由于 \code{priority_queue} 没有 decrease-key，压入新状态后弹出时跳过旧状态。如果追问变成“普通 BFS 按步数扩展，不等于总代价最小”，先判断原来的状态是否仍然覆盖新答案集合，再决定是微调边界、改 value 语义，还是换成另一类结构。

\subsubsection{LeetCode 1514 概率最大的路径}

先区分步数和代价。本题的模型是路径概率是边概率乘积，目标是最大乘积路径；若暴力枚举路径，路径数量会爆炸。优化点用大顶堆每次扩展当前概率最高节点，乘以边概率得到候选概率说明下一步扩展必须按某种代价顺序组织，而不是普通队列随便推进。

最短路状态至少包含当前节点和当前代价。本题的用大顶堆每次扩展当前概率最高节点，乘以边概率得到候选概率决定优先队列按什么排序，以及旧状态如何被识别。非负边条件、代价合成方式和不可达处理都要写清楚。放到“最短路与带权图”这条主线里看，关键原则是：非负权图用 Dijkstra，小顶堆保存当前距离最小的候选；边权为一用 BFS；有负边要考虑 Bellman-Ford 或重新建模。路径代价若是最大边权，也可用 Dijkstra 或二分答案。

最短路证明依赖扩展顺序。若本题使用 Dijkstra，就要说明非负代价让弹出的未过期状态已经最优；若用二分可达性，就要说明阈值越大可用边只会增加。这道题的边界不能留到最后补丁式处理，实验时覆盖不可达、概率为零、起点等于终点、浮点比较，打印每次弹出的代价和是否过期。若普通 BFS 与 Dijkstra 结果不同，要能解释边权条件造成的差异。

C++ 实现上，Dijkstra 用邻接表保存 \code{(to, weight)}，距离数组用 \code{long long}。由于 \code{priority_queue} 没有 decrease-key，压入新状态后弹出时跳过旧状态。如果追问变成“取负对数后可转成最短路，但直接最大堆更直观”，先判断原来的状态是否仍然覆盖新答案集合，再决定是微调边界、改 value 语义，还是换成另一类结构。

\subsubsection{LeetCode 1631 最小体力消耗路径}

先区分步数和代价。本题的模型是路径代价是沿途最大高度差，不是边权和；若暴力枚举路径，路径数量会爆炸。优化点Dijkstra 的松弛改成取当前代价和新边代价的最大值说明下一步扩展必须按某种代价顺序组织，而不是普通队列随便推进。

最短路状态至少包含当前节点和当前代价。本题的Dijkstra 的松弛改成取当前代价和新边代价的最大值决定优先队列按什么排序，以及旧状态如何被识别。非负边条件、代价合成方式和不可达处理都要写清楚。放到“最短路与带权图”这条主线里看，关键原则是：非负权图用 Dijkstra，小顶堆保存当前距离最小的候选；边权为一用 BFS；有负边要考虑 Bellman-Ford 或重新建模。路径代价若是最大边权，也可用 Dijkstra 或二分答案。

最短路证明依赖扩展顺序。若本题使用 Dijkstra，就要说明非负代价让弹出的未过期状态已经最优；若用二分可达性，就要说明阈值越大可用边只会增加。这道题的边界不能留到最后补丁式处理，实验时覆盖单格、平坦网格、必须绕路、多个相同代价，打印每次弹出的代价和是否过期。若普通 BFS 与 Dijkstra 结果不同，要能解释边权条件造成的差异。

C++ 实现上，Dijkstra 用邻接表保存 \code{(to, weight)}，距离数组用 \code{long long}。由于 \code{priority_queue} 没有 decrease-key，压入新状态后弹出时跳过旧状态。如果追问变成“也可二分体力上限后 BFS 检查可达”，先判断原来的状态是否仍然覆盖新答案集合，再决定是微调边界、改 value 语义，还是换成另一类结构。

\subsubsection{LeetCode 1928 规定时间内到达终点的最小花费}

先区分步数和代价。本题的模型是路径同时有时间约束和费用目标，单一距离无法描述状态；若暴力枚举路径，路径数量会爆炸。优化点状态包含节点和已用时间，DP 或 Dijkstra 在时间维度内松弛最小费用说明下一步扩展必须按某种代价顺序组织，而不是普通队列随便推进。

最短路状态至少包含当前节点和当前代价。本题的状态包含节点和已用时间，DP 或 Dijkstra 在时间维度内松弛最小费用决定优先队列按什么排序，以及旧状态如何被识别。非负边条件、代价合成方式和不可达处理都要写清楚。放到“最短路与带权图”这条主线里看，关键原则是：非负权图用 Dijkstra，小顶堆保存当前距离最小的候选；边权为一用 BFS；有负边要考虑 Bellman-Ford 或重新建模。路径代价若是最大边权，也可用 Dijkstra 或二分答案。

最短路证明依赖扩展顺序。若本题使用 Dijkstra，就要说明非负代价让弹出的未过期状态已经最优；若用二分可达性，就要说明阈值越大可用边只会增加。这道题的边界不能留到最后补丁式处理，实验时覆盖时间刚好到达、费用高但时间短、不可达、状态数量由 maxTime 决定，打印每次弹出的代价和是否过期。若普通 BFS 与 Dijkstra 结果不同，要能解释边权条件造成的差异。

C++ 实现上，Dijkstra 用邻接表保存 \code{(to, weight)}，距离数组用 \code{long long}。由于 \code{priority_queue} 没有 decrease-key，压入新状态后弹出时跳过旧状态。如果追问变成“这是资源约束最短路，不能只保存每个节点一个最小费用”，先判断原来的状态是否仍然覆盖新答案集合，再决定是微调边界、改 value 语义，还是换成另一类结构。

\subsubsection{LeetCode 1976 到达目的地的方案数}

先区分步数和代价。本题的模型是需要在最短距离之外统计有多少条最短路径；若暴力枚举路径，路径数量会爆炸。优化点Dijkstra 松弛时若得到更短距离就重置方案数，若等于最短距离就累加方案数说明下一步扩展必须按某种代价顺序组织，而不是普通队列随便推进。

最短路状态至少包含当前节点和当前代价。本题的Dijkstra 松弛时若得到更短距离就重置方案数，若等于最短距离就累加方案数决定优先队列按什么排序，以及旧状态如何被识别。非负边条件、代价合成方式和不可达处理都要写清楚。放到“最短路与带权图”这条主线里看，关键原则是：非负权图用 Dijkstra，小顶堆保存当前距离最小的候选；边权为一用 BFS；有负边要考虑 Bellman-Ford 或重新建模。路径代价若是最大边权，也可用 Dijkstra 或二分答案。

最短路证明依赖扩展顺序。若本题使用 Dijkstra，就要说明非负代价让弹出的未过期状态已经最优；若用二分可达性，就要说明阈值越大可用边只会增加。这道题的边界不能留到最后补丁式处理，实验时覆盖多条等长路径、取模、旧堆状态、不可达，打印每次弹出的代价和是否过期。若普通 BFS 与 Dijkstra 结果不同，要能解释边权条件造成的差异。

C++ 实现上，Dijkstra 用邻接表保存 \code{(to, weight)}，距离数组用 \code{long long}。由于 \code{priority_queue} 没有 decrease-key，压入新状态后弹出时跳过旧状态。如果追问变成“这是最短路和计数 DP 的结合，更新顺序必须和距离松弛一致”，先判断原来的状态是否仍然覆盖新答案集合，再决定是微调边界、改 value 语义，还是换成另一类结构。

\subsubsection{LeetCode 2045 到达目的地的第二短时间}

先区分步数和代价。本题的模型是每个节点需要保留最短和次短到达步数，再考虑红绿灯等待时间；若暴力枚举路径，路径数量会爆炸。优化点BFS 式扩展边数或时间，只有候选时间不同且能进入前两名时才入队说明下一步扩展必须按某种代价顺序组织，而不是普通队列随便推进。

最短路状态至少包含当前节点和当前代价。本题的BFS 式扩展边数或时间，只有候选时间不同且能进入前两名时才入队决定优先队列按什么排序，以及旧状态如何被识别。非负边条件、代价合成方式和不可达处理都要写清楚。放到“最短路与带权图”这条主线里看，关键原则是：非负权图用 Dijkstra，小顶堆保存当前距离最小的候选；边权为一用 BFS；有负边要考虑 Bellman-Ford 或重新建模。路径代价若是最大边权，也可用 Dijkstra 或二分答案。

最短路证明依赖扩展顺序。若本题使用 Dijkstra，就要说明非负代价让弹出的未过期状态已经最优；若用二分可达性，就要说明阈值越大可用边只会增加。这道题的边界不能留到最后补丁式处理，实验时覆盖第二短不能等于最短、红灯等待、无向图往返、到达终点后继续找次短，打印每次弹出的代价和是否过期。若普通 BFS 与 Dijkstra 结果不同，要能解释边权条件造成的差异。

C++ 实现上，Dijkstra 用邻接表保存 \code{(to, weight)}，距离数组用 \code{long long}。由于 \code{priority_queue} 没有 decrease-key，压入新状态后弹出时跳过旧状态。如果追问变成“它说明最短路状态有时不是一个最优值，而是前两个不同值”，先判断原来的状态是否仍然覆盖新答案集合，再决定是微调边界、改 value 语义，还是换成另一类结构。

\subsection{并查集与动态连通性工作坊}

先用搜索回答连通性，再发现查询和合并交替出现。并查集只维护集合代表，如果题目还要路径、距离、比例或奇偶性，就必须扩展状态。

\subsubsection{LeetCode 305 岛屿数量 II}

先用搜索回答每次连通性查询，观察合并和查询如何交替出现。本题的模型是陆地逐步加入，每次加入只可能影响相邻陆地所在集合；慢点在每次都重新遍历已有连接。并查集只适合新增陆地先让岛屿数加一，再和四邻已存在陆地合并，合并成功则岛屿数减一这类集合代表问题，若题目要路径本身就不能硬套。

并查集状态是代表元和集合附加信息。本题的新增陆地先让岛屿数加一，再和四邻已存在陆地合并，合并成功则岛屿数减一只要求合并集合和查询同组关系；如果还要大小、边数、比例或奇偶性，就要在合并根时同步维护额外数组。放到“并查集与动态连通性”这条主线里看，关键原则是：路径压缩把查找路径上的节点直接挂到根，按大小或按秩合并控制树高。邮箱、等式、边、网格都可以映射成元素集合。

并查集正确性来自等价关系。合并两个代表后，两个集合中任意元素的代表都会一致；不合并时，集合关系保持不变。本题的新增陆地先让岛屿数加一，再和四邻已存在陆地合并，合并成功则岛屿数减一必须只依赖这种代表关系。这道题的边界不能留到最后补丁式处理，实验时覆盖重复加入同一格、边界位置、四邻属于同一岛、空操作列表，打印每个元素的代表元和集合大小。重复合并、已经连通的边和字符串编号最容易暴露实现问题。

C++ 实现上，并查集类中成员访问保持 \code{this->}，\code{find} 可以写成迭代路径压缩避免深递归。若维护集合大小、边数或权值差，合并根时同步更新。如果追问变成“这是动态连通性，比每次重新 BFS 更适合并查集”，先判断原来的状态是否仍然覆盖新答案集合，再决定是微调边界、改 value 语义，还是换成另一类结构。

\subsubsection{LeetCode 399 除法求值}

先用搜索回答每次连通性查询，观察合并和查询如何交替出现。本题的模型是变量之间的除法关系可以看成带权连通分量；慢点在每次都重新遍历已有连接。并查集只适合并查集维护节点到根的乘积权值，合并时根据等式调整根之间权重这类集合代表问题，若题目要路径本身就不能硬套。

并查集状态是代表元和集合附加信息。本题的并查集维护节点到根的乘积权值，合并时根据等式调整根之间权重只要求合并集合和查询同组关系；如果还要大小、边数、比例或奇偶性，就要在合并根时同步维护额外数组。放到“并查集与动态连通性”这条主线里看，关键原则是：路径压缩把查找路径上的节点直接挂到根，按大小或按秩合并控制树高。邮箱、等式、边、网格都可以映射成元素集合。

并查集正确性来自等价关系。合并两个代表后，两个集合中任意元素的代表都会一致；不合并时，集合关系保持不变。本题的并查集维护节点到根的乘积权值，合并时根据等式调整根之间权重必须只依赖这种代表关系。这道题的边界不能留到最后补丁式处理，实验时覆盖查询未知变量、自除、矛盾数据、路径压缩时权值同步更新，打印每个元素的代表元和集合大小。重复合并、已经连通的边和字符串编号最容易暴露实现问题。

C++ 实现上，并查集类中成员访问保持 \code{this->}，\code{find} 可以写成迭代路径压缩避免深递归。若维护集合大小、边数或权值差，合并根时同步更新。如果追问变成“也可建图 DFS，但多次查询时带权并查集更适合”，先判断原来的状态是否仍然覆盖新答案集合，再决定是微调边界、改 value 语义，还是换成另一类结构。

\subsubsection{LeetCode 547 省份数量}

先用搜索回答每次连通性查询，观察合并和查询如何交替出现。本题的模型是城市连接关系形成无向图，省份就是连通分量；慢点在每次都重新遍历已有连接。并查集只适合并查集合并所有直接连接的城市，最后统计根数量这类集合代表问题，若题目要路径本身就不能硬套。

并查集状态是代表元和集合附加信息。本题的并查集合并所有直接连接的城市，最后统计根数量只要求合并集合和查询同组关系；如果还要大小、边数、比例或奇偶性，就要在合并根时同步维护额外数组。放到“并查集与动态连通性”这条主线里看，关键原则是：路径压缩把查找路径上的节点直接挂到根，按大小或按秩合并控制树高。邮箱、等式、边、网格都可以映射成元素集合。

并查集正确性来自等价关系。合并两个代表后，两个集合中任意元素的代表都会一致；不合并时，集合关系保持不变。本题的并查集合并所有直接连接的城市，最后统计根数量必须只依赖这种代表关系。这道题的边界不能留到最后补丁式处理，实验时覆盖邻接矩阵对称、自连接、孤立城市、只扫描上三角，打印每个元素的代表元和集合大小。重复合并、已经连通的边和字符串编号最容易暴露实现问题。

C++ 实现上，并查集类中成员访问保持 \code{this->}，\code{find} 可以写成迭代路径压缩避免深递归。若维护集合大小、边数或权值差，合并根时同步更新。如果追问变成“BFS 也可做；并查集更突出动态合并”，先判断原来的状态是否仍然覆盖新答案集合，再决定是微调边界、改 value 语义，还是换成另一类结构。

\subsubsection{LeetCode 684 冗余连接}

先用搜索回答每次连通性查询，观察合并和查询如何交替出现。本题的模型是树加一条边会形成唯一环；慢点在每次都重新遍历已有连接。并查集只适合按顺序加边，若一条边两端已连通，则它就是冗余边这类集合代表问题，若题目要路径本身就不能硬套。

并查集状态是代表元和集合附加信息。本题的按顺序加边，若一条边两端已连通，则它就是冗余边只要求合并集合和查询同组关系；如果还要大小、边数、比例或奇偶性，就要在合并根时同步维护额外数组。放到“并查集与动态连通性”这条主线里看，关键原则是：路径压缩把查找路径上的节点直接挂到根，按大小或按秩合并控制树高。邮箱、等式、边、网格都可以映射成元素集合。

并查集正确性来自等价关系。合并两个代表后，两个集合中任意元素的代表都会一致；不合并时，集合关系保持不变。本题的按顺序加边，若一条边两端已连通，则它就是冗余边必须只依赖这种代表关系。这道题的边界不能留到最后补丁式处理，实验时覆盖节点编号、最后一条成环边、重复边、并查集初始化大小，打印每个元素的代表元和集合大小。重复合并、已经连通的边和字符串编号最容易暴露实现问题。

C++ 实现上，并查集类中成员访问保持 \code{this->}，\code{find} 可以写成迭代路径压缩避免深递归。若维护集合大小、边数或权值差，合并根时同步更新。如果追问变成“有向冗余连接需要同时处理入度为二和环，复杂度更高”，先判断原来的状态是否仍然覆盖新答案集合，再决定是微调边界、改 value 语义，还是换成另一类结构。

\subsubsection{LeetCode 685 冗余连接 II}

先用搜索回答每次连通性查询，观察合并和查询如何交替出现。本题的模型是有向图中额外边可能导致某节点入度为二，也可能导致环；慢点在每次都重新遍历已有连接。并查集只适合先记录入度为二的两条候选边，再用并查集排除其中一条测试是否成树这类集合代表问题，若题目要路径本身就不能硬套。

并查集状态是代表元和集合附加信息。本题的先记录入度为二的两条候选边，再用并查集排除其中一条测试是否成树只要求合并集合和查询同组关系；如果还要大小、边数、比例或奇偶性，就要在合并根时同步维护额外数组。放到“并查集与动态连通性”这条主线里看，关键原则是：路径压缩把查找路径上的节点直接挂到根，按大小或按秩合并控制树高。邮箱、等式、边、网格都可以映射成元素集合。

并查集正确性来自等价关系。合并两个代表后，两个集合中任意元素的代表都会一致；不合并时，集合关系保持不变。本题的先记录入度为二的两条候选边，再用并查集排除其中一条测试是否成树必须只依赖这种代表关系。这道题的边界不能留到最后补丁式处理，实验时覆盖只有环、只有入度二、两者同时存在、返回最后出现的边，打印每个元素的代表元和集合大小。重复合并、已经连通的边和字符串编号最容易暴露实现问题。

C++ 实现上，并查集类中成员访问保持 \code{this->}，\code{find} 可以写成迭代路径压缩避免深递归。若维护集合大小、边数或权值差，合并根时同步更新。如果追问变成“它比无向冗余连接多了有向树入度约束”，先判断原来的状态是否仍然覆盖新答案集合，再决定是微调边界、改 value 语义，还是换成另一类结构。

\subsubsection{LeetCode 721 账户合并}

先用搜索回答每次连通性查询，观察合并和查询如何交替出现。本题的模型是共享邮箱说明账户属于同一人，邮箱之间形成连通分量；慢点在每次都重新遍历已有连接。并查集只适合给邮箱编号并查集合并，同根邮箱收集后排序输出这类集合代表问题，若题目要路径本身就不能硬套。

并查集状态是代表元和集合附加信息。本题的给邮箱编号并查集合并，同根邮箱收集后排序输出只要求合并集合和查询同组关系；如果还要大小、边数、比例或奇偶性，就要在合并根时同步维护额外数组。放到“并查集与动态连通性”这条主线里看，关键原则是：路径压缩把查找路径上的节点直接挂到根，按大小或按秩合并控制树高。邮箱、等式、边、网格都可以映射成元素集合。

并查集正确性来自等价关系。合并两个代表后，两个集合中任意元素的代表都会一致；不合并时，集合关系保持不变。本题的给邮箱编号并查集合并，同根邮箱收集后排序输出必须只依赖这种代表关系。这道题的边界不能留到最后补丁式处理，实验时覆盖同名不同人、一个账户多个邮箱、重复邮箱、输出排序，打印每个元素的代表元和集合大小。重复合并、已经连通的边和字符串编号最容易暴露实现问题。

C++ 实现上，并查集类中成员访问保持 \code{this->}，\code{find} 可以写成迭代路径压缩避免深递归。若维护集合大小、边数或权值差，合并根时同步更新。如果追问变成“姓名不能作为合并依据，只能作为输出字段”，先判断原来的状态是否仍然覆盖新答案集合，再决定是微调边界、改 value 语义，还是换成另一类结构。

\subsubsection{LeetCode 947 移除最多的同行或同列石头}

先用搜索回答每次连通性查询，观察合并和查询如何交替出现。本题的模型是同一行或同一列的石头属于同一连通分量，每个分量至少留一块；慢点在每次都重新遍历已有连接。并查集只适合把行和列编号后合并，答案是石头数减去连通分量数这类集合代表问题，若题目要路径本身就不能硬套。

并查集状态是代表元和集合附加信息。本题的把行和列编号后合并，答案是石头数减去连通分量数只要求合并集合和查询同组关系；如果还要大小、边数、比例或奇偶性，就要在合并根时同步维护额外数组。放到“并查集与动态连通性”这条主线里看，关键原则是：路径压缩把查找路径上的节点直接挂到根，按大小或按秩合并控制树高。邮箱、等式、边、网格都可以映射成元素集合。

并查集正确性来自等价关系。合并两个代表后，两个集合中任意元素的代表都会一致；不合并时，集合关系保持不变。本题的把行和列编号后合并，答案是石头数减去连通分量数必须只依赖这种代表关系。这道题的边界不能留到最后补丁式处理，实验时覆盖单块石头、行列编号冲突、多个分量、重复坐标，打印每个元素的代表元和集合大小。重复合并、已经连通的边和字符串编号最容易暴露实现问题。

C++ 实现上，并查集类中成员访问保持 \code{this->}，\code{find} 可以写成迭代路径压缩避免深递归。若维护集合大小、边数或权值差，合并根时同步更新。如果追问变成“也可把石头当节点建图，但行列并查集更省边”，先判断原来的状态是否仍然覆盖新答案集合，再决定是微调边界、改 value 语义，还是换成另一类结构。

\subsubsection{LeetCode 959 由斜杠划分区域}

先用搜索回答每次连通性查询，观察合并和查询如何交替出现。本题的模型是每个网格可拆成四个小三角，斜杠决定内部哪些三角相连；慢点在每次都重新遍历已有连接。并查集只适合给每格四个区域编号，先按字符合并内部，再合并相邻格边界区域这类集合代表问题，若题目要路径本身就不能硬套。

并查集状态是代表元和集合附加信息。本题的给每格四个区域编号，先按字符合并内部，再合并相邻格边界区域只要求合并集合和查询同组关系；如果还要大小、边数、比例或奇偶性，就要在合并根时同步维护额外数组。放到“并查集与动态连通性”这条主线里看，关键原则是：路径压缩把查找路径上的节点直接挂到根，按大小或按秩合并控制树高。邮箱、等式、边、网格都可以映射成元素集合。

并查集正确性来自等价关系。合并两个代表后，两个集合中任意元素的代表都会一致；不合并时，集合关系保持不变。本题的给每格四个区域编号，先按字符合并内部，再合并相邻格边界区域必须只依赖这种代表关系。这道题的边界不能留到最后补丁式处理，实验时覆盖空格、正斜杠、反斜杠、边界合并、区域计数，打印每个元素的代表元和集合大小。重复合并、已经连通的边和字符串编号最容易暴露实现问题。

C++ 实现上，并查集类中成员访问保持 \code{this->}，\code{find} 可以写成迭代路径压缩避免深递归。若维护集合大小、边数或权值差，合并根时同步更新。如果追问变成“这是把几何连通性离散成并查集元素的典型题”，先判断原来的状态是否仍然覆盖新答案集合，再决定是微调边界、改 value 语义，还是换成另一类结构。

\subsubsection{LeetCode 990 等式方程的可满足性}

先用搜索回答每次连通性查询，观察合并和查询如何交替出现。本题的模型是相等关系形成集合，不等关系要求两端不在同一集合；慢点在每次都重新遍历已有连接。并查集只适合先合并所有等式，再检查每个不等式是否冲突这类集合代表问题，若题目要路径本身就不能硬套。

并查集状态是代表元和集合附加信息。本题的先合并所有等式，再检查每个不等式是否冲突只要求合并集合和查询同组关系；如果还要大小、边数、比例或奇偶性，就要在合并根时同步维护额外数组。放到“并查集与动态连通性”这条主线里看，关键原则是：路径压缩把查找路径上的节点直接挂到根，按大小或按秩合并控制树高。邮箱、等式、边、网格都可以映射成元素集合。

并查集正确性来自等价关系。合并两个代表后，两个集合中任意元素的代表都会一致；不合并时，集合关系保持不变。本题的先合并所有等式，再检查每个不等式是否冲突必须只依赖这种代表关系。这道题的边界不能留到最后补丁式处理，实验时覆盖自等式、自不等式、传递相等、多组变量，打印每个元素的代表元和集合大小。重复合并、已经连通的边和字符串编号最容易暴露实现问题。

C++ 实现上，并查集类中成员访问保持 \code{this->}，\code{find} 可以写成迭代路径压缩避免深递归。若维护集合大小、边数或权值差，合并根时同步更新。如果追问变成“先处理不等式会漏掉后续等式带来的传递冲突”，先判断原来的状态是否仍然覆盖新答案集合，再决定是微调边界、改 value 语义，还是换成另一类结构。

\subsubsection{LeetCode 1061 按字典序排列最小的等效字符串}

先用搜索回答每次连通性查询，观察合并和查询如何交替出现。本题的模型是等价字符形成集合，每个集合应选择字典序最小代表；慢点在每次都重新遍历已有连接。并查集只适合合并两个字符集合时让较小根成为代表，扫描目标串时替换为根这类集合代表问题，若题目要路径本身就不能硬套。

并查集状态是代表元和集合附加信息。本题的合并两个字符集合时让较小根成为代表，扫描目标串时替换为根只要求合并集合和查询同组关系；如果还要大小、边数、比例或奇偶性，就要在合并根时同步维护额外数组。放到“并查集与动态连通性”这条主线里看，关键原则是：路径压缩把查找路径上的节点直接挂到根，按大小或按秩合并控制树高。邮箱、等式、边、网格都可以映射成元素集合。

并查集正确性来自等价关系。合并两个代表后，两个集合中任意元素的代表都会一致；不合并时，集合关系保持不变。本题的合并两个字符集合时让较小根成为代表，扫描目标串时替换为根必须只依赖这种代表关系。这道题的边界不能留到最后补丁式处理，实验时覆盖传递等价、重复等价关系、自等价、未出现字符，打印每个元素的代表元和集合大小。重复合并、已经连通的边和字符串编号最容易暴露实现问题。

C++ 实现上，并查集类中成员访问保持 \code{this->}，\code{find} 可以写成迭代路径压缩避免深递归。若维护集合大小、边数或权值差，合并根时同步更新。如果追问变成“如果集合代表还要保存其他权值，合并规则需要同步维护”，先判断原来的状态是否仍然覆盖新答案集合，再决定是微调边界、改 value 语义，还是换成另一类结构。

\subsubsection{LeetCode 1202 交换字符串中的元素}

先用搜索回答每次连通性查询，观察合并和查询如何交替出现。本题的模型是可交换下标形成连通分量，同一分量内字符可以任意重排；慢点在每次都重新遍历已有连接。并查集只适合并查集合并所有可交换下标，对每个分量收集字符并按字典序分配回最小下标这类集合代表问题，若题目要路径本身就不能硬套。

并查集状态是代表元和集合附加信息。本题的并查集合并所有可交换下标，对每个分量收集字符并按字典序分配回最小下标只要求合并集合和查询同组关系；如果还要大小、边数、比例或奇偶性，就要在合并根时同步维护额外数组。放到“并查集与动态连通性”这条主线里看，关键原则是：路径压缩把查找路径上的节点直接挂到根，按大小或按秩合并控制树高。邮箱、等式、边、网格都可以映射成元素集合。

并查集正确性来自等价关系。合并两个代表后，两个集合中任意元素的代表都会一致；不合并时，集合关系保持不变。本题的并查集合并所有可交换下标，对每个分量收集字符并按字典序分配回最小下标必须只依赖这种代表关系。这道题的边界不能留到最后补丁式处理，实验时覆盖没有交换、多个分量、重复字符、分量下标排序，打印每个元素的代表元和集合大小。重复合并、已经连通的边和字符串编号最容易暴露实现问题。

C++ 实现上，并查集类中成员访问保持 \code{this->}，\code{find} 可以写成迭代路径压缩避免深递归。若维护集合大小、边数或权值差，合并根时同步更新。如果追问变成“它展示并查集不仅能回答连通性，还能把分量内资源重新组织”，先判断原来的状态是否仍然覆盖新答案集合，再决定是微调边界、改 value 语义，还是换成另一类结构。

\subsubsection{LeetCode 1319 连通网络的操作次数}

先用搜索回答每次连通性查询，观察合并和查询如何交替出现。本题的模型是多余网线可以挪用来连接不同连通分量；慢点在每次都重新遍历已有连接。并查集只适合先判断边数至少为 n 减一，再用并查集统计连通分量，答案是分量数减一这类集合代表问题，若题目要路径本身就不能硬套。

并查集状态是代表元和集合附加信息。本题的先判断边数至少为 n 减一，再用并查集统计连通分量，答案是分量数减一只要求合并集合和查询同组关系；如果还要大小、边数、比例或奇偶性，就要在合并根时同步维护额外数组。放到“并查集与动态连通性”这条主线里看，关键原则是：路径压缩把查找路径上的节点直接挂到根，按大小或按秩合并控制树高。邮箱、等式、边、网格都可以映射成元素集合。

并查集正确性来自等价关系。合并两个代表后，两个集合中任意元素的代表都会一致；不合并时，集合关系保持不变。本题的先判断边数至少为 n 减一，再用并查集统计连通分量，答案是分量数减一必须只依赖这种代表关系。这道题的边界不能留到最后补丁式处理，实验时覆盖边数不足、已经连通、重复边、多分量，打印每个元素的代表元和集合大小。重复合并、已经连通的边和字符串编号最容易暴露实现问题。

C++ 实现上，并查集类中成员访问保持 \code{this->}，\code{find} 可以写成迭代路径压缩避免深递归。若维护集合大小、边数或权值差，合并根时同步更新。如果追问变成“合并失败的边代表冗余资源，但只要总边数足够即可”，先判断原来的状态是否仍然覆盖新答案集合，再决定是微调边界、改 value 语义，还是换成另一类结构。

\subsubsection{LeetCode 1579 保证图可完全遍历}

先用搜索回答每次连通性查询，观察合并和查询如何交替出现。本题的模型是Alice 和 Bob 各自需要图连通，公共边应优先服务两人；慢点在每次都重新遍历已有连接。并查集只适合先处理类型三公共边同步合并两个并查集，再分别处理各自边，最后检查两图连通这类集合代表问题，若题目要路径本身就不能硬套。

并查集状态是代表元和集合附加信息。本题的先处理类型三公共边同步合并两个并查集，再分别处理各自边，最后检查两图连通只要求合并集合和查询同组关系；如果还要大小、边数、比例或奇偶性，就要在合并根时同步维护额外数组。放到“并查集与动态连通性”这条主线里看，关键原则是：路径压缩把查找路径上的节点直接挂到根，按大小或按秩合并控制树高。邮箱、等式、边、网格都可以映射成元素集合。

并查集正确性来自等价关系。合并两个代表后，两个集合中任意元素的代表都会一致；不合并时，集合关系保持不变。本题的先处理类型三公共边同步合并两个并查集，再分别处理各自边，最后检查两图连通必须只依赖这种代表关系。这道题的边界不能留到最后补丁式处理，实验时覆盖公共边冗余、某一方不连通、边数很多、节点从一开始编号，打印每个元素的代表元和集合大小。重复合并、已经连通的边和字符串编号最容易暴露实现问题。

C++ 实现上，并查集类中成员访问保持 \code{this->}，\code{find} 可以写成迭代路径压缩避免深递归。若维护集合大小、边数或权值差，合并根时同步更新。如果追问变成“这是双并查集和资源优先级的组合”，先判断原来的状态是否仍然覆盖新答案集合，再决定是微调边界、改 value 语义，还是换成另一类结构。

\subsection{回溯、枚举与剪枝工作坊}

先写搜索树：路径是什么、选择列表是什么、什么时候结束。然后用排序、剩余容量、同层去重或合法性预处理剪掉不可能分支。

\subsubsection{LeetCode 17 电话号码的字母组合}

先画前三层搜索树。本题的模型是每个数字对应若干字母，答案是所有按位选择形成的字符串；暴力搜索要明确路径、选择列表和终止条件。优化点是暴力就是完整笛卡尔积，优化重点是路径长度和下标同步推进，避免在每层重新拼接大量中间字符串，也就是哪些分支在展开前已经不可能产生合法答案。

回溯状态由路径和当前位置共同定义。本题的暴力就是完整笛卡尔积，优化重点是路径长度和下标同步推进，避免在每层重新拼接大量中间字符串要落成剪枝条件或同层去重条件。剪枝只能删掉不可能成功的分支，不能删掉只是暂时看起来不优的分支。放到“回溯、枚举与剪枝”这条主线里看，关键原则是：剪枝来自容量、剩余长度、排序后去重、合法性预检查和提前终止。组合题关注起点推进，排列题关注元素是否使用，分割题关注当前片段是否合法。

回溯证明分两半：搜索树覆盖所有合法答案，同层去重或剪枝不会删除合法答案。本题的每个数字对应若干字母，答案是所有按位选择形成的字符串给出选择来源，暴力就是完整笛卡尔积，优化重点是路径长度和下标同步推进，避免在每层重新拼接大量中间字符串给出提前停止的理由。这道题的边界不能留到最后补丁式处理，实验时覆盖空输入、包含 7 或 9、输入中不应出现 0 和 1、输出规模指数增长，统计搜索节点数量和剪枝次数。重复元素题要打印同一层跳过哪些值，确认没有误删合法路径。

C++ 实现上，路径容器修改要成对出现：加入、搜索、撤销。重复元素题先排序，再用同层去重条件；输出规模大时，复杂度要把答案数量算进去。如果追问变成“若字符映射来自键盘以外的字典，只需要替换候选表，搜索骨架不变”，先判断原来的状态是否仍然覆盖新答案集合，再决定是微调边界、改 value 语义，还是换成另一类结构。

\subsubsection{LeetCode 22 括号生成}

先画前三层搜索树。本题的模型是合法括号串的前缀中右括号数量不能超过左括号数量；暴力搜索要明确路径、选择列表和终止条件。优化点是状态保存已用左括号和右括号数量，只有满足约束的选择才继续递归，也就是哪些分支在展开前已经不可能产生合法答案。

回溯状态由路径和当前位置共同定义。本题的状态保存已用左括号和右括号数量，只有满足约束的选择才继续递归要落成剪枝条件或同层去重条件。剪枝只能删掉不可能成功的分支，不能删掉只是暂时看起来不优的分支。放到“回溯、枚举与剪枝”这条主线里看，关键原则是：剪枝来自容量、剩余长度、排序后去重、合法性预检查和提前终止。组合题关注起点推进，排列题关注元素是否使用，分割题关注当前片段是否合法。

回溯证明分两半：搜索树覆盖所有合法答案，同层去重或剪枝不会删除合法答案。本题的合法括号串的前缀中右括号数量不能超过左括号数量给出选择来源，状态保存已用左括号和右括号数量，只有满足约束的选择才继续递归给出提前停止的理由。这道题的边界不能留到最后补丁式处理，实验时覆盖n 为零或一、右括号提前超过左括号、输出规模是卡特兰数，统计搜索节点数量和剪枝次数。重复元素题要打印同一层跳过哪些值，确认没有误删合法路径。

C++ 实现上，路径容器修改要成对出现：加入、搜索、撤销。重复元素题先排序，再用同层去重条件；输出规模大时，复杂度要把答案数量算进去。如果追问变成“若有多种括号类型，状态还需要栈保存类型匹配关系”，先判断原来的状态是否仍然覆盖新答案集合，再决定是微调边界、改 value 语义，还是换成另一类结构。

\subsubsection{LeetCode 39 组合总和}

先画前三层搜索树。本题的模型是候选数字可重复使用，答案不关心排列顺序；暴力搜索要明确路径、选择列表和终止条件。优化点是排序后按起点枚举，递归时仍传当前下标表示可以继续使用同一数字，也就是哪些分支在展开前已经不可能产生合法答案。

回溯状态由路径和当前位置共同定义。本题的排序后按起点枚举，递归时仍传当前下标表示可以继续使用同一数字要落成剪枝条件或同层去重条件。剪枝只能删掉不可能成功的分支，不能删掉只是暂时看起来不优的分支。放到“回溯、枚举与剪枝”这条主线里看，关键原则是：剪枝来自容量、剩余长度、排序后去重、合法性预检查和提前终止。组合题关注起点推进，排列题关注元素是否使用，分割题关注当前片段是否合法。

回溯证明分两半：搜索树覆盖所有合法答案，同层去重或剪枝不会删除合法答案。本题的候选数字可重复使用，答案不关心排列顺序给出选择来源，排序后按起点枚举，递归时仍传当前下标表示可以继续使用同一数字给出提前停止的理由。这道题的边界不能留到最后补丁式处理，实验时覆盖目标小于最小数、刚好等于、候选有重复时题意如何处理、输出规模，统计搜索节点数量和剪枝次数。重复元素题要打印同一层跳过哪些值，确认没有误删合法路径。

C++ 实现上，路径容器修改要成对出现：加入、搜索、撤销。重复元素题先排序，再用同层去重条件；输出规模大时，复杂度要把答案数量算进去。如果追问变成“组合总和 II 每个数只能用一次，起点推进和去重条件都不同”，先判断原来的状态是否仍然覆盖新答案集合，再决定是微调边界、改 value 语义，还是换成另一类结构。

\subsubsection{LeetCode 40 组合总和 II}

先画前三层搜索树。本题的模型是每个候选只能使用一次，输入中可能存在重复值；暴力搜索要明确路径、选择列表和终止条件。优化点是排序后递归下一层从 i 加一开始，同层遇到相同值要跳过，也就是哪些分支在展开前已经不可能产生合法答案。

回溯状态由路径和当前位置共同定义。本题的排序后递归下一层从 i 加一开始，同层遇到相同值要跳过要落成剪枝条件或同层去重条件。剪枝只能删掉不可能成功的分支，不能删掉只是暂时看起来不优的分支。放到“回溯、枚举与剪枝”这条主线里看，关键原则是：剪枝来自容量、剩余长度、排序后去重、合法性预检查和提前终止。组合题关注起点推进，排列题关注元素是否使用，分割题关注当前片段是否合法。

回溯证明分两半：搜索树覆盖所有合法答案，同层去重或剪枝不会删除合法答案。本题的每个候选只能使用一次，输入中可能存在重复值给出选择来源，排序后递归下一层从 i 加一开始，同层遇到相同值要跳过给出提前停止的理由。这道题的边界不能留到最后补丁式处理，实验时覆盖全重复、目标为零、同层去重写成跨层去重、答案顺序不要求，统计搜索节点数量和剪枝次数。重复元素题要打印同一层跳过哪些值，确认没有误删合法路径。

C++ 实现上，路径容器修改要成对出现：加入、搜索、撤销。重复元素题先排序，再用同层去重条件；输出规模大时，复杂度要把答案数量算进去。如果追问变成“子集 II 使用同样的同层去重，只是每个路径节点都可能成为答案”，先判断原来的状态是否仍然覆盖新答案集合，再决定是微调边界、改 value 语义，还是换成另一类结构。

\subsubsection{LeetCode 46 全排列}

先画前三层搜索树。本题的模型是排列关心元素顺序，每个位置从所有未使用元素中选择一个；暴力搜索要明确路径、选择列表和终止条件。优化点是路径长度达到 n 时产生答案，used 数组保证同一元素不重复使用，也就是哪些分支在展开前已经不可能产生合法答案。

回溯状态由路径和当前位置共同定义。本题的路径长度达到 n 时产生答案，used 数组保证同一元素不重复使用要落成剪枝条件或同层去重条件。剪枝只能删掉不可能成功的分支，不能删掉只是暂时看起来不优的分支。放到“回溯、枚举与剪枝”这条主线里看，关键原则是：剪枝来自容量、剩余长度、排序后去重、合法性预检查和提前终止。组合题关注起点推进，排列题关注元素是否使用，分割题关注当前片段是否合法。

回溯证明分两半：搜索树覆盖所有合法答案，同层去重或剪枝不会删除合法答案。本题的排列关心元素顺序，每个位置从所有未使用元素中选择一个给出选择来源，路径长度达到 n 时产生答案，used 数组保证同一元素不重复使用给出提前停止的理由。这道题的边界不能留到最后补丁式处理，实验时覆盖空数组、单元素、输出规模为 n 阶乘、元素值唯一假设，统计搜索节点数量和剪枝次数。重复元素题要打印同一层跳过哪些值，确认没有误删合法路径。

C++ 实现上，路径容器修改要成对出现：加入、搜索、撤销。重复元素题先排序，再用同层去重条件；输出规模大时，复杂度要把答案数量算进去。如果追问变成“若有重复元素，需要排序并在同层跳过等价选择”，先判断原来的状态是否仍然覆盖新答案集合，再决定是微调边界、改 value 语义，还是换成另一类结构。

\subsubsection{LeetCode 47 全排列 II}

先画前三层搜索树。本题的模型是重复元素会让不同下标选择生成相同排列；暴力搜索要明确路径、选择列表和终止条件。优化点是排序后使用 used 数组，同层遇到相同值且前一个相同值未使用时跳过，也就是哪些分支在展开前已经不可能产生合法答案。

回溯状态由路径和当前位置共同定义。本题的排序后使用 used 数组，同层遇到相同值且前一个相同值未使用时跳过要落成剪枝条件或同层去重条件。剪枝只能删掉不可能成功的分支，不能删掉只是暂时看起来不优的分支。放到“回溯、枚举与剪枝”这条主线里看，关键原则是：剪枝来自容量、剩余长度、排序后去重、合法性预检查和提前终止。组合题关注起点推进，排列题关注元素是否使用，分割题关注当前片段是否合法。

回溯证明分两半：搜索树覆盖所有合法答案，同层去重或剪枝不会删除合法答案。本题的重复元素会让不同下标选择生成相同排列给出选择来源，排序后使用 used 数组，同层遇到相同值且前一个相同值未使用时跳过给出提前停止的理由。这道题的边界不能留到最后补丁式处理，实验时覆盖全重复、部分重复、同层去重条件写反、输出顺序不要求，统计搜索节点数量和剪枝次数。重复元素题要打印同一层跳过哪些值，确认没有误删合法路径。

C++ 实现上，路径容器修改要成对出现：加入、搜索、撤销。重复元素题先排序，再用同层去重条件；输出规模大时，复杂度要把答案数量算进去。如果追问变成“子集 II 和组合总和 II 也是同层去重，但起点推进方式不同”，先判断原来的状态是否仍然覆盖新答案集合，再决定是微调边界、改 value 语义，还是换成另一类结构。

\subsubsection{LeetCode 51 N 皇后}

先画前三层搜索树。本题的模型是每一行放一个皇后，列和两条对角线不能冲突；暴力搜索要明确路径、选择列表和终止条件。优化点是逐行搜索，用列集合、主对角线集合和副对角线集合快速判断合法性，也就是哪些分支在展开前已经不可能产生合法答案。

回溯状态由路径和当前位置共同定义。本题的逐行搜索，用列集合、主对角线集合和副对角线集合快速判断合法性要落成剪枝条件或同层去重条件。剪枝只能删掉不可能成功的分支，不能删掉只是暂时看起来不优的分支。放到“回溯、枚举与剪枝”这条主线里看，关键原则是：剪枝来自容量、剩余长度、排序后去重、合法性预检查和提前终止。组合题关注起点推进，排列题关注元素是否使用，分割题关注当前片段是否合法。

回溯证明分两半：搜索树覆盖所有合法答案，同层去重或剪枝不会删除合法答案。本题的每一行放一个皇后，列和两条对角线不能冲突给出选择来源，逐行搜索，用列集合、主对角线集合和副对角线集合快速判断合法性给出提前停止的理由。这道题的边界不能留到最后补丁式处理，实验时覆盖n 为一、n 为二或三无解、对角线编号、回退时状态恢复，统计搜索节点数量和剪枝次数。重复元素题要打印同一层跳过哪些值，确认没有误删合法路径。

C++ 实现上，路径容器修改要成对出现：加入、搜索、撤销。重复元素题先排序，再用同层去重条件；输出规模大时，复杂度要把答案数量算进去。如果追问变成“可用位掩码压缩三个约束集合，提高常数性能”，先判断原来的状态是否仍然覆盖新答案集合，再决定是微调边界、改 value 语义，还是换成另一类结构。

\subsubsection{LeetCode 78 子集}

先画前三层搜索树。本题的模型是每个元素都有选或不选两种选择，答案是所有路径节点；暴力搜索要明确路径、选择列表和终止条件。优化点是回溯按起点向后枚举，避免不同顺序生成同一集合，也就是哪些分支在展开前已经不可能产生合法答案。

回溯状态由路径和当前位置共同定义。本题的回溯按起点向后枚举，避免不同顺序生成同一集合要落成剪枝条件或同层去重条件。剪枝只能删掉不可能成功的分支，不能删掉只是暂时看起来不优的分支。放到“回溯、枚举与剪枝”这条主线里看，关键原则是：剪枝来自容量、剩余长度、排序后去重、合法性预检查和提前终止。组合题关注起点推进，排列题关注元素是否使用，分割题关注当前片段是否合法。

回溯证明分两半：搜索树覆盖所有合法答案，同层去重或剪枝不会删除合法答案。本题的每个元素都有选或不选两种选择，答案是所有路径节点给出选择来源，回溯按起点向后枚举，避免不同顺序生成同一集合给出提前停止的理由。这道题的边界不能留到最后补丁式处理，实验时覆盖空数组、单元素、输出数量为二的 n 次方、顺序不要求，统计搜索节点数量和剪枝次数。重复元素题要打印同一层跳过哪些值，确认没有误删合法路径。

C++ 实现上，路径容器修改要成对出现：加入、搜索、撤销。重复元素题先排序，再用同层去重条件；输出规模大时，复杂度要把答案数量算进去。如果追问变成“也可用位掩码枚举所有子集，适合 n 较小的场景”，先判断原来的状态是否仍然覆盖新答案集合，再决定是微调边界、改 value 语义，还是换成另一类结构。

\subsubsection{LeetCode 79 单词搜索}

先画前三层搜索树。本题的模型是网格路径不能重复使用同一格，状态包含位置、匹配下标和访问标记；暴力搜索要明确路径、选择列表和终止条件。优化点是剪枝包括首字符不匹配、剩余长度不足、字符频次不足和越界访问，也就是哪些分支在展开前已经不可能产生合法答案。

回溯状态由路径和当前位置共同定义。本题的剪枝包括首字符不匹配、剩余长度不足、字符频次不足和越界访问要落成剪枝条件或同层去重条件。剪枝只能删掉不可能成功的分支，不能删掉只是暂时看起来不优的分支。放到“回溯、枚举与剪枝”这条主线里看，关键原则是：剪枝来自容量、剩余长度、排序后去重、合法性预检查和提前终止。组合题关注起点推进，排列题关注元素是否使用，分割题关注当前片段是否合法。

回溯证明分两半：搜索树覆盖所有合法答案，同层去重或剪枝不会删除合法答案。本题的网格路径不能重复使用同一格，状态包含位置、匹配下标和访问标记给出选择来源，剪枝包括首字符不匹配、剩余长度不足、字符频次不足和越界访问给出提前停止的理由。这道题的边界不能留到最后补丁式处理，实验时覆盖单字符单词、路径需要回退、重复字母、单元格不能复用，统计搜索节点数量和剪枝次数。重复元素题要打印同一层跳过哪些值，确认没有误删合法路径。

C++ 实现上，路径容器修改要成对出现：加入、搜索、撤销。重复元素题先排序，再用同层去重条件；输出规模大时，复杂度要把答案数量算进去。如果追问变成“若要查很多单词，应使用 Trie 合并搜索前缀”，先判断原来的状态是否仍然覆盖新答案集合，再决定是微调边界、改 value 语义，还是换成另一类结构。

\subsubsection{LeetCode 90 子集 II}

先画前三层搜索树。本题的模型是重复元素会让不同下标路径生成相同集合；暴力搜索要明确路径、选择列表和终止条件。优化点是排序后在同一层跳过相同值，保留不同层使用重复值的可能，也就是哪些分支在展开前已经不可能产生合法答案。

回溯状态由路径和当前位置共同定义。本题的排序后在同一层跳过相同值，保留不同层使用重复值的可能要落成剪枝条件或同层去重条件。剪枝只能删掉不可能成功的分支，不能删掉只是暂时看起来不优的分支。放到“回溯、枚举与剪枝”这条主线里看，关键原则是：剪枝来自容量、剩余长度、排序后去重、合法性预检查和提前终止。组合题关注起点推进，排列题关注元素是否使用，分割题关注当前片段是否合法。

回溯证明分两半：搜索树覆盖所有合法答案，同层去重或剪枝不会删除合法答案。本题的重复元素会让不同下标路径生成相同集合给出选择来源，排序后在同一层跳过相同值，保留不同层使用重复值的可能给出提前停止的理由。这道题的边界不能留到最后补丁式处理，实验时覆盖全重复、空集、重复值相邻、去重条件写成同层而不是全局，统计搜索节点数量和剪枝次数。重复元素题要打印同一层跳过哪些值，确认没有误删合法路径。

C++ 实现上，路径容器修改要成对出现：加入、搜索、撤销。重复元素题先排序，再用同层去重条件；输出规模大时，复杂度要把答案数量算进去。如果追问变成“组合总和 II 使用同样的同层去重思想”，先判断原来的状态是否仍然覆盖新答案集合，再决定是微调边界、改 value 语义，还是换成另一类结构。

\subsubsection{LeetCode 93 复原 IP 地址}

先画前三层搜索树。本题的模型是字符串要切成四段，每段必须是 0 到 255 且无非法前导零；暴力搜索要明确路径、选择列表和终止条件。优化点是按段数和剩余长度剪枝，每次尝试长度一到三的片段，也就是哪些分支在展开前已经不可能产生合法答案。

回溯状态由路径和当前位置共同定义。本题的按段数和剩余长度剪枝，每次尝试长度一到三的片段要落成剪枝条件或同层去重条件。剪枝只能删掉不可能成功的分支，不能删掉只是暂时看起来不优的分支。放到“回溯、枚举与剪枝”这条主线里看，关键原则是：剪枝来自容量、剩余长度、排序后去重、合法性预检查和提前终止。组合题关注起点推进，排列题关注元素是否使用，分割题关注当前片段是否合法。

回溯证明分两半：搜索树覆盖所有合法答案，同层去重或剪枝不会删除合法答案。本题的字符串要切成四段，每段必须是 0 到 255 且无非法前导零给出选择来源，按段数和剩余长度剪枝，每次尝试长度一到三的片段给出提前停止的理由。这道题的边界不能留到最后补丁式处理，实验时覆盖前导零、数值超过 255、剩余字符太多或太少、刚好四段，统计搜索节点数量和剪枝次数。重复元素题要打印同一层跳过哪些值，确认没有误删合法路径。

C++ 实现上，路径容器修改要成对出现：加入、搜索、撤销。重复元素题先排序，再用同层去重条件；输出规模大时，复杂度要把答案数量算进去。如果追问变成“分割回文串也是前缀分割，但合法性检查不同”，先判断原来的状态是否仍然覆盖新答案集合，再决定是微调边界、改 value 语义，还是换成另一类结构。

\subsubsection{LeetCode 131 分割回文串}

先画前三层搜索树。本题的模型是每次选择一个从当前位置开始的回文前缀；暴力搜索要明确路径、选择列表和终止条件。优化点是预处理回文 DP 表，让回溯合法性检查为常数，也就是哪些分支在展开前已经不可能产生合法答案。

回溯状态由路径和当前位置共同定义。本题的预处理回文 DP 表，让回溯合法性检查为常数要落成剪枝条件或同层去重条件。剪枝只能删掉不可能成功的分支，不能删掉只是暂时看起来不优的分支。放到“回溯、枚举与剪枝”这条主线里看，关键原则是：剪枝来自容量、剩余长度、排序后去重、合法性预检查和提前终止。组合题关注起点推进，排列题关注元素是否使用，分割题关注当前片段是否合法。

回溯证明分两半：搜索树覆盖所有合法答案，同层去重或剪枝不会删除合法答案。本题的每次选择一个从当前位置开始的回文前缀给出选择来源，预处理回文 DP 表，让回溯合法性检查为常数给出提前停止的理由。这道题的边界不能留到最后补丁式处理，实验时覆盖空串、单字符、全相同字符、重复分割路径，统计搜索节点数量和剪枝次数。重复元素题要打印同一层跳过哪些值，确认没有误删合法路径。

C++ 实现上，路径容器修改要成对出现：加入、搜索、撤销。重复元素题先排序，再用同层去重条件；输出规模大时，复杂度要把答案数量算进去。如果追问变成“若只求最少切割次数，可改成 DP 最短切分”，先判断原来的状态是否仍然覆盖新答案集合，再决定是微调边界、改 value 语义，还是换成另一类结构。

\subsubsection{LeetCode 212 单词搜索 II}

先画前三层搜索树。本题的模型是多个单词在同一棋盘上搜索，公共前缀可以共享；暴力搜索要明确路径、选择列表和终止条件。优化点是先建 Trie，再从每个格子回溯；路径到达单词结束节点就收集答案，也就是哪些分支在展开前已经不可能产生合法答案。

回溯状态由路径和当前位置共同定义。本题的先建 Trie，再从每个格子回溯；路径到达单词结束节点就收集答案要落成剪枝条件或同层去重条件。剪枝只能删掉不可能成功的分支，不能删掉只是暂时看起来不优的分支。放到“回溯、枚举与剪枝”这条主线里看，关键原则是：剪枝来自容量、剩余长度、排序后去重、合法性预检查和提前终止。组合题关注起点推进，排列题关注元素是否使用，分割题关注当前片段是否合法。

回溯证明分两半：搜索树覆盖所有合法答案，同层去重或剪枝不会删除合法答案。本题的多个单词在同一棋盘上搜索，公共前缀可以共享给出选择来源，先建 Trie，再从每个格子回溯；路径到达单词结束节点就收集答案给出提前停止的理由。这道题的边界不能留到最后补丁式处理，实验时覆盖重复单词、同一格不能复用、前缀是完整单词、剪枝删除空子树，统计搜索节点数量和剪枝次数。重复元素题要打印同一层跳过哪些值，确认没有误删合法路径。

C++ 实现上，路径容器修改要成对出现：加入、搜索、撤销。重复元素题先排序，再用同层去重条件；输出规模大时，复杂度要把答案数量算进去。如果追问变成“这是单词搜索从单目标扩展到多目标后的结构升级”，先判断原来的状态是否仍然覆盖新答案集合，再决定是微调边界、改 value 语义，还是换成另一类结构。

\subsubsection{LeetCode 216 组合总和 III}

先画前三层搜索树。本题的模型是从一到九中选 k 个数和为 n，每个数最多用一次；暴力搜索要明确路径、选择列表和终止条件。优化点是按递增起点枚举，利用剩余个数和剩余和剪枝，也就是哪些分支在展开前已经不可能产生合法答案。

回溯状态由路径和当前位置共同定义。本题的按递增起点枚举，利用剩余个数和剩余和剪枝要落成剪枝条件或同层去重条件。剪枝只能删掉不可能成功的分支，不能删掉只是暂时看起来不优的分支。放到“回溯、枚举与剪枝”这条主线里看，关键原则是：剪枝来自容量、剩余长度、排序后去重、合法性预检查和提前终止。组合题关注起点推进，排列题关注元素是否使用，分割题关注当前片段是否合法。

回溯证明分两半：搜索树覆盖所有合法答案，同层去重或剪枝不会删除合法答案。本题的从一到九中选 k 个数和为 n，每个数最多用一次给出选择来源，按递增起点枚举，利用剩余个数和剩余和剪枝给出提前停止的理由。这道题的边界不能留到最后补丁式处理，实验时覆盖n 太小、n 太大、k 为零、路径长度已经超过 k，统计搜索节点数量和剪枝次数。重复元素题要打印同一层跳过哪些值，确认没有误删合法路径。

C++ 实现上，路径容器修改要成对出现：加入、搜索、撤销。重复元素题先排序，再用同层去重条件；输出规模大时，复杂度要把答案数量算进去。如果追问变成“可预估剩余最小和最大和进一步剪枝”，先判断原来的状态是否仍然覆盖新答案集合，再决定是微调边界、改 value 语义，还是换成另一类结构。

\subsubsection{LeetCode 698 划分为 k 个相等的子集}

先画前三层搜索树。本题的模型是每个元素要分配到 k 个桶中，所有桶目标和相同；暴力搜索要明确路径、选择列表和终止条件。优化点是先判断总和可整除并排序降序，再用桶剩余容量回溯放置元素，也就是哪些分支在展开前已经不可能产生合法答案。

回溯状态由路径和当前位置共同定义。本题的先判断总和可整除并排序降序，再用桶剩余容量回溯放置元素要落成剪枝条件或同层去重条件。剪枝只能删掉不可能成功的分支，不能删掉只是暂时看起来不优的分支。放到“回溯、枚举与剪枝”这条主线里看，关键原则是：剪枝来自容量、剩余长度、排序后去重、合法性预检查和提前终止。组合题关注起点推进，排列题关注元素是否使用，分割题关注当前片段是否合法。

回溯证明分两半：搜索树覆盖所有合法答案，同层去重或剪枝不会删除合法答案。本题的每个元素要分配到 k 个桶中，所有桶目标和相同给出选择来源，先判断总和可整除并排序降序，再用桶剩余容量回溯放置元素给出提前停止的理由。这道题的边界不能留到最后补丁式处理，实验时覆盖存在元素大于目标、重复元素、空桶对称、k 为一，统计搜索节点数量和剪枝次数。重复元素题要打印同一层跳过哪些值，确认没有误删合法路径。

C++ 实现上，路径容器修改要成对出现：加入、搜索、撤销。重复元素题先排序，再用同层去重条件；输出规模大时，复杂度要把答案数量算进去。如果追问变成“状态压缩 DP 也能做，适合用掩码表示已选元素集合”，先判断原来的状态是否仍然覆盖新答案集合，再决定是微调边界、改 value 语义，还是换成另一类结构。

\subsection{动态规划与状态定义工作坊}

先写递归搜索并找重复状态，再给状态命名。状态必须包含未来决策所需信息，遍历顺序必须保证依赖状态已经算完。

\subsubsection{LeetCode 53 最大子数组和}

先写递归或枚举，把重复子问题暴露出来。本题的模型是状态是必须以当前位置结尾的最大连续和；慢点不是递归形式，而是相同状态被反复求。本题需要命名的状态来自若前一个后缀为负，接上它只会拖累当前元素，因此可以重新开始，状态定义必须让未来决策不再依赖完整历史。

DP 状态必须足够描述未来。本题的若前一个后缀为负，接上它只会拖累当前元素，因此可以重新开始要写成数组下标、容量、前缀长度、持有状态或区间边界。初始化对应空状态，遍历顺序对应依赖方向。放到“动态规划与状态定义”这条主线里看，关键原则是：状态定义要包含足够信息但不能过度膨胀。线性 DP 看最后一步选择，二维 DP 看两个前缀或网格位置，背包看物品阶段和容量，区间 DP 看最后合并点。

DP 证明通常看最后一步。任意最优解或合法方案的最后一步必在转移枚举中；每个转移来源又已经按遍历顺序算完。本题的若前一个后缀为负，接上它只会拖累当前元素，因此可以重新开始要支撑这个依赖关系。这道题的边界不能留到最后补丁式处理，实验时覆盖全负数组、单元素、和溢出、答案不能初始化为零，先打印完整 DP 表，再比较空间压缩版本。若压缩版不同，优先检查遍历顺序和旧值保存。

C++ 实现上，DP 表初始化要对应状态语义，不可达哨兵不能溢出。空间压缩前先写完整表；压缩后要说清读到的是上一轮、本轮左侧还是左上角旧值。如果追问变成“若要求返回区间下标，同步记录重新开始位置和最优左右端”，先判断原来的状态是否仍然覆盖新答案集合，再决定是微调边界、改 value 语义，还是换成另一类结构。

\subsubsection{LeetCode 62 不同路径}

先写递归或枚举，把重复子问题暴露出来。本题的模型是网格状态由上方和左方两种最后一步转移而来；慢点不是递归形式，而是相同状态被反复求。本题需要命名的状态来自二维表可压缩为一维，因为当前行只依赖上一行同列和当前行左侧，状态定义必须让未来决策不再依赖完整历史。

DP 状态必须足够描述未来。本题的二维表可压缩为一维，因为当前行只依赖上一行同列和当前行左侧要写成数组下标、容量、前缀长度、持有状态或区间边界。初始化对应空状态，遍历顺序对应依赖方向。放到“动态规划与状态定义”这条主线里看，关键原则是：状态定义要包含足够信息但不能过度膨胀。线性 DP 看最后一步选择，二维 DP 看两个前缀或网格位置，背包看物品阶段和容量，区间 DP 看最后合并点。

DP 证明通常看最后一步。任意最优解或合法方案的最后一步必在转移枚举中；每个转移来源又已经按遍历顺序算完。本题的二维表可压缩为一维，因为当前行只依赖上一行同列和当前行左侧要支撑这个依赖关系。这道题的边界不能留到最后补丁式处理，实验时覆盖一行、一列、较大网格结果溢出、障碍物版本边界，先打印完整 DP 表，再比较空间压缩版本。若压缩版不同，优先检查遍历顺序和旧值保存。

C++ 实现上，DP 表初始化要对应状态语义，不可达哨兵不能溢出。空间压缩前先写完整表；压缩后要说清读到的是上一轮、本轮左侧还是左上角旧值。如果追问变成“带障碍物时遇到障碍状态清零，第一行第一列要按阻断处理”，先判断原来的状态是否仍然覆盖新答案集合，再决定是微调边界、改 value 语义，还是换成另一类结构。

\subsubsection{LeetCode 63 不同路径 II}

先写递归或枚举，把重复子问题暴露出来。本题的模型是障碍物让某些格子的路径数变为零；慢点不是递归形式，而是相同状态被反复求。本题需要命名的状态来自滚动数组中遇到障碍直接把当前位置清零，否则累加左侧，状态定义必须让未来决策不再依赖完整历史。

DP 状态必须足够描述未来。本题的滚动数组中遇到障碍直接把当前位置清零，否则累加左侧要写成数组下标、容量、前缀长度、持有状态或区间边界。初始化对应空状态，遍历顺序对应依赖方向。放到“动态规划与状态定义”这条主线里看，关键原则是：状态定义要包含足够信息但不能过度膨胀。线性 DP 看最后一步选择，二维 DP 看两个前缀或网格位置，背包看物品阶段和容量，区间 DP 看最后合并点。

DP 证明通常看最后一步。任意最优解或合法方案的最后一步必在转移枚举中；每个转移来源又已经按遍历顺序算完。本题的滚动数组中遇到障碍直接把当前位置清零，否则累加左侧要支撑这个依赖关系。这道题的边界不能留到最后补丁式处理，实验时覆盖起点或终点是障碍、整行被截断、单行障碍，先打印完整 DP 表，再比较空间压缩版本。若压缩版不同，优先检查遍历顺序和旧值保存。

C++ 实现上，DP 表初始化要对应状态语义，不可达哨兵不能溢出。空间压缩前先写完整表；压缩后要说清读到的是上一轮、本轮左侧还是左上角旧值。如果追问变成“若允许向上或向左走，图会有环，普通填表不再成立”，先判断原来的状态是否仍然覆盖新答案集合，再决定是微调边界、改 value 语义，还是换成另一类结构。

\subsubsection{LeetCode 64 最小路径和}

先写递归或枚举，把重复子问题暴露出来。本题的模型是状态是到达当前格子的最小代价，最后一步来自上方或左方；慢点不是递归形式，而是相同状态被反复求。本题需要命名的状态来自先初始化第一行和第一列，再处理内部，主转移保持简单，状态定义必须让未来决策不再依赖完整历史。

DP 状态必须足够描述未来。本题的先初始化第一行和第一列，再处理内部，主转移保持简单要写成数组下标、容量、前缀长度、持有状态或区间边界。初始化对应空状态，遍历顺序对应依赖方向。放到“动态规划与状态定义”这条主线里看，关键原则是：状态定义要包含足够信息但不能过度膨胀。线性 DP 看最后一步选择，二维 DP 看两个前缀或网格位置，背包看物品阶段和容量，区间 DP 看最后合并点。

DP 证明通常看最后一步。任意最优解或合法方案的最后一步必在转移枚举中；每个转移来源又已经按遍历顺序算完。本题的先初始化第一行和第一列，再处理内部，主转移保持简单要支撑这个依赖关系。这道题的边界不能留到最后补丁式处理，实验时覆盖单行、单列、权值为零、路径和溢出，先打印完整 DP 表，再比较空间压缩版本。若压缩版不同，优先检查遍历顺序和旧值保存。

C++ 实现上，DP 表初始化要对应状态语义，不可达哨兵不能溢出。空间压缩前先写完整表；压缩后要说清读到的是上一轮、本轮左侧还是左上角旧值。如果追问变成“若允许四方向移动，变成最短路而不是普通网格 DP”，先判断原来的状态是否仍然覆盖新答案集合，再决定是微调边界、改 value 语义，还是换成另一类结构。

\subsubsection{LeetCode 70 爬楼梯}

先写递归或枚举，把重复子问题暴露出来。本题的模型是最后一步来自前一阶或前两阶，是最小的重叠子问题模型；慢点不是递归形式，而是相同状态被反复求。本题需要命名的状态来自滚动变量保存最近两个状态即可，不需要完整数组，状态定义必须让未来决策不再依赖完整历史。

DP 状态必须足够描述未来。本题的滚动变量保存最近两个状态即可，不需要完整数组要写成数组下标、容量、前缀长度、持有状态或区间边界。初始化对应空状态，遍历顺序对应依赖方向。放到“动态规划与状态定义”这条主线里看，关键原则是：状态定义要包含足够信息但不能过度膨胀。线性 DP 看最后一步选择，二维 DP 看两个前缀或网格位置，背包看物品阶段和容量，区间 DP 看最后合并点。

DP 证明通常看最后一步。任意最优解或合法方案的最后一步必在转移枚举中；每个转移来源又已经按遍历顺序算完。本题的滚动变量保存最近两个状态即可，不需要完整数组要支撑这个依赖关系。这道题的边界不能留到最后补丁式处理，实验时覆盖n 为一、n 为二、题目是否允许零阶、结果范围，先打印完整 DP 表，再比较空间压缩版本。若压缩版不同，优先检查遍历顺序和旧值保存。

C++ 实现上，DP 表初始化要对应状态语义，不可达哨兵不能溢出。空间压缩前先写完整表；压缩后要说清读到的是上一轮、本轮左侧还是左上角旧值。如果追问变成“若每次可走一到 k 阶，转移变成固定宽度窗口求和”，先判断原来的状态是否仍然覆盖新答案集合，再决定是微调边界、改 value 语义，还是换成另一类结构。

\subsubsection{LeetCode 72 编辑距离}

先写递归或枚举，把重复子问题暴露出来。本题的模型是二维状态表示两个前缀之间的最少编辑操作；慢点不是递归形式，而是相同状态被反复求。本题需要命名的状态来自末尾字符相同继承左上；不同则枚举插入、删除、替换，状态定义必须让未来决策不再依赖完整历史。

DP 状态必须足够描述未来。本题的末尾字符相同继承左上；不同则枚举插入、删除、替换要写成数组下标、容量、前缀长度、持有状态或区间边界。初始化对应空状态，遍历顺序对应依赖方向。放到“动态规划与状态定义”这条主线里看，关键原则是：状态定义要包含足够信息但不能过度膨胀。线性 DP 看最后一步选择，二维 DP 看两个前缀或网格位置，背包看物品阶段和容量，区间 DP 看最后合并点。

DP 证明通常看最后一步。任意最优解或合法方案的最后一步必在转移枚举中；每个转移来源又已经按遍历顺序算完。本题的末尾字符相同继承左上；不同则枚举插入、删除、替换要支撑这个依赖关系。这道题的边界不能留到最后补丁式处理，实验时覆盖空串、完全相同、完全不同、操作代价不同，先打印完整 DP 表，再比较空间压缩版本。若压缩版不同，优先检查遍历顺序和旧值保存。

C++ 实现上，DP 表初始化要对应状态语义，不可达哨兵不能溢出。空间压缩前先写完整表；压缩后要说清读到的是上一轮、本轮左侧还是左上角旧值。如果追问变成“若只允许插入删除，问题退化到和最长公共子序列相关”，先判断原来的状态是否仍然覆盖新答案集合，再决定是微调边界、改 value 语义，还是换成另一类结构。

\subsubsection{LeetCode 121 买卖股票的最佳时机}

先写递归或枚举，把重复子问题暴露出来。本题的模型是卖出日固定时，最佳买入日是此前最低价格；慢点不是递归形式，而是相同状态被反复求。本题需要命名的状态来自扫描维护历史最低价和最大利润，一次交易无需复杂状态机，状态定义必须让未来决策不再依赖完整历史。

DP 状态必须足够描述未来。本题的扫描维护历史最低价和最大利润，一次交易无需复杂状态机要写成数组下标、容量、前缀长度、持有状态或区间边界。初始化对应空状态，遍历顺序对应依赖方向。放到“动态规划与状态定义”这条主线里看，关键原则是：状态定义要包含足够信息但不能过度膨胀。线性 DP 看最后一步选择，二维 DP 看两个前缀或网格位置，背包看物品阶段和容量，区间 DP 看最后合并点。

DP 证明通常看最后一步。任意最优解或合法方案的最后一步必在转移枚举中；每个转移来源又已经按遍历顺序算完。本题的扫描维护历史最低价和最大利润，一次交易无需复杂状态机要支撑这个依赖关系。这道题的边界不能留到最后补丁式处理，实验时覆盖单日价格、一直下跌、一直上涨、价格相同，先打印完整 DP 表，再比较空间压缩版本。若压缩版不同，优先检查遍历顺序和旧值保存。

C++ 实现上，DP 表初始化要对应状态语义，不可达哨兵不能溢出。空间压缩前先写完整表；压缩后要说清读到的是上一轮、本轮左侧还是左上角旧值。如果追问变成“多次交易、冷冻期和手续费版本都要扩展成状态机”，先判断原来的状态是否仍然覆盖新答案集合，再决定是微调边界、改 value 语义，还是换成另一类结构。

\subsubsection{LeetCode 123 买卖股票 III}

先写递归或枚举，把重复子问题暴露出来。本题的模型是最多两次交易，状态需要包含交易次数和持有状态；慢点不是递归形式，而是相同状态被反复求。本题需要命名的状态来自维护 buy1、sell1、buy2、sell2 四个状态，按天更新，状态定义必须让未来决策不再依赖完整历史。

DP 状态必须足够描述未来。本题的维护 buy1、sell1、buy2、sell2 四个状态，按天更新要写成数组下标、容量、前缀长度、持有状态或区间边界。初始化对应空状态，遍历顺序对应依赖方向。放到“动态规划与状态定义”这条主线里看，关键原则是：状态定义要包含足够信息但不能过度膨胀。线性 DP 看最后一步选择，二维 DP 看两个前缀或网格位置，背包看物品阶段和容量，区间 DP 看最后合并点。

DP 证明通常看最后一步。任意最优解或合法方案的最后一步必在转移枚举中；每个转移来源又已经按遍历顺序算完。本题的维护 buy1、sell1、buy2、sell2 四个状态，按天更新要支撑这个依赖关系。这道题的边界不能留到最后补丁式处理，实验时覆盖价格为空、只够一次交易、两次交易间不能重叠，先打印完整 DP 表，再比较空间压缩版本。若压缩版不同，优先检查遍历顺序和旧值保存。

C++ 实现上，DP 表初始化要对应状态语义，不可达哨兵不能溢出。空间压缩前先写完整表；压缩后要说清读到的是上一轮、本轮左侧还是左上角旧值。如果追问变成“最多 k 次交易是它的泛化，k 很大时退化为无限次交易”，先判断原来的状态是否仍然覆盖新答案集合，再决定是微调边界、改 value 语义，还是换成另一类结构。

\subsubsection{LeetCode 139 单词拆分}

先写递归或枚举，把重复子问题暴露出来。本题的模型是状态表示前缀能否由字典单词拼出；慢点不是递归形式，而是相同状态被反复求。本题需要命名的状态来自枚举最后一个单词的起点，左侧前缀可达且子串在字典中则当前可达，状态定义必须让未来决策不再依赖完整历史。

DP 状态必须足够描述未来。本题的枚举最后一个单词的起点，左侧前缀可达且子串在字典中则当前可达要写成数组下标、容量、前缀长度、持有状态或区间边界。初始化对应空状态，遍历顺序对应依赖方向。放到“动态规划与状态定义”这条主线里看，关键原则是：状态定义要包含足够信息但不能过度膨胀。线性 DP 看最后一步选择，二维 DP 看两个前缀或网格位置，背包看物品阶段和容量，区间 DP 看最后合并点。

DP 证明通常看最后一步。任意最优解或合法方案的最后一步必在转移枚举中；每个转移来源又已经按遍历顺序算完。本题的枚举最后一个单词的起点，左侧前缀可达且子串在字典中则当前可达要支撑这个依赖关系。这道题的边界不能留到最后补丁式处理，实验时覆盖空串、重复单词、长单词、substr 复制成本，先打印完整 DP 表，再比较空间压缩版本。若压缩版不同，优先检查遍历顺序和旧值保存。

C++ 实现上，DP 表初始化要对应状态语义，不可达哨兵不能溢出。空间压缩前先写完整表；压缩后要说清读到的是上一轮、本轮左侧还是左上角旧值。如果追问变成“若要输出所有句子，需要 DFS 加记忆化或前驱列表”，先判断原来的状态是否仍然覆盖新答案集合，再决定是微调边界、改 value 语义，还是换成另一类结构。

\subsubsection{LeetCode 152 乘积最大子数组}

先写递归或枚举，把重复子问题暴露出来。本题的模型是负数会让最大乘积和最小乘积互换；慢点不是递归形式，而是相同状态被反复求。本题需要命名的状态来自同时维护以当前位置结尾的最大和最小乘积，状态定义必须让未来决策不再依赖完整历史。

DP 状态必须足够描述未来。本题的同时维护以当前位置结尾的最大和最小乘积要写成数组下标、容量、前缀长度、持有状态或区间边界。初始化对应空状态，遍历顺序对应依赖方向。放到“动态规划与状态定义”这条主线里看，关键原则是：状态定义要包含足够信息但不能过度膨胀。线性 DP 看最后一步选择，二维 DP 看两个前缀或网格位置，背包看物品阶段和容量，区间 DP 看最后合并点。

DP 证明通常看最后一步。任意最优解或合法方案的最后一步必在转移枚举中；每个转移来源又已经按遍历顺序算完。本题的同时维护以当前位置结尾的最大和最小乘积要支撑这个依赖关系。这道题的边界不能留到最后补丁式处理，实验时覆盖零、负数个数奇偶、单元素、乘积溢出，先打印完整 DP 表，再比较空间压缩版本。若压缩版不同，优先检查遍历顺序和旧值保存。

C++ 实现上，DP 表初始化要对应状态语义，不可达哨兵不能溢出。空间压缩前先写完整表；压缩后要说清读到的是上一轮、本轮左侧还是左上角旧值。如果追问变成“如果要求返回区间，同步记录状态来源”，先判断原来的状态是否仍然覆盖新答案集合，再决定是微调边界、改 value 语义，还是换成另一类结构。

\subsubsection{LeetCode 198 打家劫舍}

先写递归或枚举，把重复子问题暴露出来。本题的模型是每间房只有偷和不偷两种结局，偷当前就不能偷前一间；慢点不是递归形式，而是相同状态被反复求。本题需要命名的状态来自滚动保存前一状态和前两状态，状态定义必须让未来决策不再依赖完整历史。

DP 状态必须足够描述未来。本题的滚动保存前一状态和前两状态要写成数组下标、容量、前缀长度、持有状态或区间边界。初始化对应空状态，遍历顺序对应依赖方向。放到“动态规划与状态定义”这条主线里看，关键原则是：状态定义要包含足够信息但不能过度膨胀。线性 DP 看最后一步选择，二维 DP 看两个前缀或网格位置，背包看物品阶段和容量，区间 DP 看最后合并点。

DP 证明通常看最后一步。任意最优解或合法方案的最后一步必在转移枚举中；每个转移来源又已经按遍历顺序算完。本题的滚动保存前一状态和前两状态要支撑这个依赖关系。这道题的边界不能留到最后补丁式处理，实验时覆盖空数组、单间、全零、金额很大，先打印完整 DP 表，再比较空间压缩版本。若压缩版不同，优先检查遍历顺序和旧值保存。

C++ 实现上，DP 表初始化要对应状态语义，不可达哨兵不能溢出。空间压缩前先写完整表；压缩后要说清读到的是上一轮、本轮左侧还是左上角旧值。如果追问变成“环形房屋要拆成不含首和不含尾两个线性问题”，先判断原来的状态是否仍然覆盖新答案集合，再决定是微调边界、改 value 语义，还是换成另一类结构。

\subsubsection{LeetCode 221 最大正方形}

先写递归或枚举，把重复子问题暴露出来。本题的模型是以当前位置为右下角的最大正方形由上、左、左上三个状态限制；慢点不是递归形式，而是相同状态被反复求。本题需要命名的状态来自当前位置为一时取三者最小加一，否则为零，状态定义必须让未来决策不再依赖完整历史。

DP 状态必须足够描述未来。本题的当前位置为一时取三者最小加一，否则为零要写成数组下标、容量、前缀长度、持有状态或区间边界。初始化对应空状态，遍历顺序对应依赖方向。放到“动态规划与状态定义”这条主线里看，关键原则是：状态定义要包含足够信息但不能过度膨胀。线性 DP 看最后一步选择，二维 DP 看两个前缀或网格位置，背包看物品阶段和容量，区间 DP 看最后合并点。

DP 证明通常看最后一步。任意最优解或合法方案的最后一步必在转移枚举中；每个转移来源又已经按遍历顺序算完。本题的当前位置为一时取三者最小加一，否则为零要支撑这个依赖关系。这道题的边界不能留到最后补丁式处理，实验时覆盖全零、全一、单行单列、字符一和数字一不要混淆，先打印完整 DP 表，再比较空间压缩版本。若压缩版不同，优先检查遍历顺序和旧值保存。

C++ 实现上，DP 表初始化要对应状态语义，不可达哨兵不能溢出。空间压缩前先写完整表；压缩后要说清读到的是上一轮、本轮左侧还是左上角旧值。如果追问变成“最大矩形则需要按行构造柱状图再用单调栈”，先判断原来的状态是否仍然覆盖新答案集合，再决定是微调边界、改 value 语义，还是换成另一类结构。

\subsubsection{LeetCode 300 最长递增子序列}

先写递归或枚举，把重复子问题暴露出来。本题的模型是二次 DP 固定结尾，二分优化维护每个长度的最小结尾；慢点不是递归形式，而是相同状态被反复求。本题需要命名的状态来自tails 不是实际答案序列，而是未来扩展潜力表，状态定义必须让未来决策不再依赖完整历史。

DP 状态必须足够描述未来。本题的tails 不是实际答案序列，而是未来扩展潜力表要写成数组下标、容量、前缀长度、持有状态或区间边界。初始化对应空状态，遍历顺序对应依赖方向。放到“动态规划与状态定义”这条主线里看，关键原则是：状态定义要包含足够信息但不能过度膨胀。线性 DP 看最后一步选择，二维 DP 看两个前缀或网格位置，背包看物品阶段和容量，区间 DP 看最后合并点。

DP 证明通常看最后一步。任意最优解或合法方案的最后一步必在转移枚举中；每个转移来源又已经按遍历顺序算完。本题的tails 不是实际答案序列，而是未来扩展潜力表要支撑这个依赖关系。这道题的边界不能留到最后补丁式处理，实验时覆盖重复元素、严格递增和非递减区别、空数组，先打印完整 DP 表，再比较空间压缩版本。若压缩版不同，优先检查遍历顺序和旧值保存。

C++ 实现上，DP 表初始化要对应状态语义，不可达哨兵不能溢出。空间压缩前先写完整表；压缩后要说清读到的是上一轮、本轮左侧还是左上角旧值。如果追问变成“若要还原路径，需要额外前驱数组，不能只看 tails”，先判断原来的状态是否仍然覆盖新答案集合，再决定是微调边界、改 value 语义，还是换成另一类结构。

\subsubsection{LeetCode 309 最佳买卖股票含冷冻期}

先写递归或枚举，把重复子问题暴露出来。本题的模型是冷冻期让买入来源受昨天卖出状态限制；慢点不是递归形式，而是相同状态被反复求。本题需要命名的状态来自状态机维护持有、刚卖出、休息三种日终状态，状态定义必须让未来决策不再依赖完整历史。

DP 状态必须足够描述未来。本题的状态机维护持有、刚卖出、休息三种日终状态要写成数组下标、容量、前缀长度、持有状态或区间边界。初始化对应空状态，遍历顺序对应依赖方向。放到“动态规划与状态定义”这条主线里看，关键原则是：状态定义要包含足够信息但不能过度膨胀。线性 DP 看最后一步选择，二维 DP 看两个前缀或网格位置，背包看物品阶段和容量，区间 DP 看最后合并点。

DP 证明通常看最后一步。任意最优解或合法方案的最后一步必在转移枚举中；每个转移来源又已经按遍历顺序算完。本题的状态机维护持有、刚卖出、休息三种日终状态要支撑这个依赖关系。这道题的边界不能留到最后补丁式处理，实验时覆盖空价格、一天、连续上涨、卖出后不能次日买入，先打印完整 DP 表，再比较空间压缩版本。若压缩版不同，优先检查遍历顺序和旧值保存。

C++ 实现上，DP 表初始化要对应状态语义，不可达哨兵不能溢出。空间压缩前先写完整表；压缩后要说清读到的是上一轮、本轮左侧还是左上角旧值。如果追问变成“手续费版本可以在买入或卖出转移中扣费”，先判断原来的状态是否仍然覆盖新答案集合，再决定是微调边界、改 value 语义，还是换成另一类结构。

\subsubsection{LeetCode 714 买卖股票含手续费}

先写递归或枚举，把重复子问题暴露出来。本题的模型是手续费改变交易收益，仍可用持有和不持有状态；慢点不是递归形式，而是相同状态被反复求。本题需要命名的状态来自卖出时扣费或买入时扣费都可以，但不要重复扣，状态定义必须让未来决策不再依赖完整历史。

DP 状态必须足够描述未来。本题的卖出时扣费或买入时扣费都可以，但不要重复扣要写成数组下标、容量、前缀长度、持有状态或区间边界。初始化对应空状态，遍历顺序对应依赖方向。放到“动态规划与状态定义”这条主线里看，关键原则是：状态定义要包含足够信息但不能过度膨胀。线性 DP 看最后一步选择，二维 DP 看两个前缀或网格位置，背包看物品阶段和容量，区间 DP 看最后合并点。

DP 证明通常看最后一步。任意最优解或合法方案的最后一步必在转移枚举中；每个转移来源又已经按遍历顺序算完。本题的卖出时扣费或买入时扣费都可以，但不要重复扣要支撑这个依赖关系。这道题的边界不能留到最后补丁式处理，实验时覆盖手续费为零、价格震荡、小利润不足手续费，先打印完整 DP 表，再比较空间压缩版本。若压缩版不同，优先检查遍历顺序和旧值保存。

C++ 实现上，DP 表初始化要对应状态语义，不可达哨兵不能溢出。空间压缩前先写完整表；压缩后要说清读到的是上一轮、本轮左侧还是左上角旧值。如果追问变成“这是状态机比局部贪心更稳的例子”，先判断原来的状态是否仍然覆盖新答案集合，再决定是微调边界、改 value 语义，还是换成另一类结构。

\subsubsection{LeetCode 1143 最长公共子序列}

先写递归或枚举，把重复子问题暴露出来。本题的模型是状态表示两个前缀的最长公共子序列长度；慢点不是递归形式，而是相同状态被反复求。本题需要命名的状态来自末尾相同取左上加一，不同取上方和左方最大，状态定义必须让未来决策不再依赖完整历史。

DP 状态必须足够描述未来。本题的末尾相同取左上加一，不同取上方和左方最大要写成数组下标、容量、前缀长度、持有状态或区间边界。初始化对应空状态，遍历顺序对应依赖方向。放到“动态规划与状态定义”这条主线里看，关键原则是：状态定义要包含足够信息但不能过度膨胀。线性 DP 看最后一步选择，二维 DP 看两个前缀或网格位置，背包看物品阶段和容量，区间 DP 看最后合并点。

DP 证明通常看最后一步。任意最优解或合法方案的最后一步必在转移枚举中；每个转移来源又已经按遍历顺序算完。本题的末尾相同取左上加一，不同取上方和左方最大要支撑这个依赖关系。这道题的边界不能留到最后补丁式处理，实验时覆盖空串、完全相同、无公共字符、子序列不是子串，先打印完整 DP 表，再比较空间压缩版本。若压缩版不同，优先检查遍历顺序和旧值保存。

C++ 实现上，DP 表初始化要对应状态语义，不可达哨兵不能溢出。空间压缩前先写完整表；压缩后要说清读到的是上一轮、本轮左侧还是左上角旧值。如果追问变成“若要求还原序列，需要从表右下角反向追踪选择”，先判断原来的状态是否仍然覆盖新答案集合，再决定是微调边界、改 value 语义，还是换成另一类结构。

\subsubsection{LeetCode 1235 规划兼职工作}

先写递归或枚举，把重复子问题暴露出来。本题的模型是每份工作有开始、结束和收益，选择不冲突工作使收益最大；慢点不是递归形式，而是相同状态被反复求。本题需要命名的状态来自按结束时间排序，dp[i] 在不选当前和选当前加上兼容前驱之间取最大，前驱用二分找，状态定义必须让未来决策不再依赖完整历史。

DP 状态必须足够描述未来。本题的按结束时间排序，dp[i] 在不选当前和选当前加上兼容前驱之间取最大，前驱用二分找要写成数组下标、容量、前缀长度、持有状态或区间边界。初始化对应空状态，遍历顺序对应依赖方向。放到“动态规划与状态定义”这条主线里看，关键原则是：状态定义要包含足够信息但不能过度膨胀。线性 DP 看最后一步选择，二维 DP 看两个前缀或网格位置，背包看物品阶段和容量，区间 DP 看最后合并点。

DP 证明通常看最后一步。任意最优解或合法方案的最后一步必在转移枚举中；每个转移来源又已经按遍历顺序算完。本题的按结束时间排序，dp[i] 在不选当前和选当前加上兼容前驱之间取最大，前驱用二分找要支撑这个依赖关系。这道题的边界不能留到最后补丁式处理，实验时覆盖结束等于开始是否兼容、收益相同、工作排序、空列表，先打印完整 DP 表，再比较空间压缩版本。若压缩版不同，优先检查遍历顺序和旧值保存。

C++ 实现上，DP 表初始化要对应状态语义，不可达哨兵不能溢出。空间压缩前先写完整表；压缩后要说清读到的是上一轮、本轮左侧还是左上角旧值。如果追问变成“这是排序、二分和 DP 的组合，不是按收益贪心”，先判断原来的状态是否仍然覆盖新答案集合，再决定是微调边界、改 value 语义，还是换成另一类结构。

\subsection{背包模型工作坊}

先枚举物品选择，再把剩余容量作为状态。0-1 背包倒序容量，完全背包正序容量，组合和排列的循环顺序不能混用。

\subsubsection{LeetCode 279 完全平方数}

先枚举每个物品选或不选、选几次。本题的模型是完全平方数可以看成无限使用的物品，目标是凑出 n 的最少数量；暴力慢在相同阶段和容量反复出现。优化点完全背包最少物品数，或把数字看成图节点做 BFS决定容量循环方向、状态值含义和是否允许当前物品在同一轮被再次使用。

背包状态围绕阶段和容量。本题的完全背包最少物品数，或把数字看成图节点做 BFS决定倒序还是正序、求最值还是计数、组合还是排列。循环顺序不是写法偏好，而是题意的一部分。放到“背包模型”这条主线里看，关键原则是：0-1 背包每个物品只能用一次，一维压缩时容量倒序；完全背包物品可无限使用，容量正序。计数组合和计数排列的循环顺序不同。

背包证明围绕当前物品使用次数。0-1 只有选和不选，完全背包允许同一物品继续转移。本题的完全背包最少物品数，或把数字看成图节点做 BFS要与循环方向一致，否则会把题意改掉。这道题的边界不能留到最后补丁式处理，实验时覆盖n 为零或一、完全平方数本身、较大 n，并故意把容量方向写反构造最小反例。这样能确认循环方向来自题意，不是背模板。

C++ 实现上，一维背包数组的容量方向要和题意一致。方案数用 \code{long long} 或题目要求的取模；容量来自数值大小时，复杂度是伪多项式。如果追问变成“数学四平方定理可给更强解法，但面试 DP 更通用”，先判断原来的状态是否仍然覆盖新答案集合，再决定是微调边界、改 value 语义，还是换成另一类结构。

\subsubsection{LeetCode 322 零钱兑换}

先枚举每个物品选或不选、选几次。本题的模型是金额是容量，硬币可无限使用，目标是最少硬币数；暴力慢在相同阶段和容量反复出现。优化点dp[value] 从 value 减 coin 的已知状态转移决定容量循环方向、状态值含义和是否允许当前物品在同一轮被再次使用。

背包状态围绕阶段和容量。本题的dp[value] 从 value 减 coin 的已知状态转移决定倒序还是正序、求最值还是计数、组合还是排列。循环顺序不是写法偏好，而是题意的一部分。放到“背包模型”这条主线里看，关键原则是：0-1 背包每个物品只能用一次，一维压缩时容量倒序；完全背包物品可无限使用，容量正序。计数组合和计数排列的循环顺序不同。

背包证明围绕当前物品使用次数。0-1 只有选和不选，完全背包允许同一物品继续转移。本题的dp[value] 从 value 减 coin 的已知状态转移要与循环方向一致，否则会把题意改掉。这道题的边界不能留到最后补丁式处理，实验时覆盖无法凑出、amount 为零、硬币面额大于目标、哨兵溢出，并故意把容量方向写反构造最小反例。这样能确认循环方向来自题意，不是背模板。

C++ 实现上，一维背包数组的容量方向要和题意一致。方案数用 \code{long long} 或题目要求的取模；容量来自数值大小时，复杂度是伪多项式。如果追问变成“零钱兑换 II 问组合数，循环语义不同”，先判断原来的状态是否仍然覆盖新答案集合，再决定是微调边界、改 value 语义，还是换成另一类结构。

\subsubsection{LeetCode 377 组合总和 IV}

先枚举每个物品选或不选、选几次。本题的模型是题目统计排列数，选择顺序不同算不同方案；暴力慢在相同阶段和容量反复出现。优化点外层遍历容量，内层遍历数字，让每个容量的方案按最后一个数累加决定容量循环方向、状态值含义和是否允许当前物品在同一轮被再次使用。

背包状态围绕阶段和容量。本题的外层遍历容量，内层遍历数字，让每个容量的方案按最后一个数累加决定倒序还是正序、求最值还是计数、组合还是排列。循环顺序不是写法偏好，而是题意的一部分。放到“背包模型”这条主线里看，关键原则是：0-1 背包每个物品只能用一次，一维压缩时容量倒序；完全背包物品可无限使用，容量正序。计数组合和计数排列的循环顺序不同。

背包证明围绕当前物品使用次数。0-1 只有选和不选，完全背包允许同一物品继续转移。本题的外层遍历容量，内层遍历数字，让每个容量的方案按最后一个数累加要与循环方向一致，否则会把题意改掉。这道题的边界不能留到最后补丁式处理，实验时覆盖target 为零、数字大于目标、方案数溢出、循环顺序，并故意把容量方向写反构造最小反例。这样能确认循环方向来自题意，不是背模板。

C++ 实现上，一维背包数组的容量方向要和题意一致。方案数用 \code{long long} 或题目要求的取模；容量来自数值大小时，复杂度是伪多项式。如果追问变成“若顺序不同不算新方案，就变成完全背包组合数，循环顺序要交换”，先判断原来的状态是否仍然覆盖新答案集合，再决定是微调边界、改 value 语义，还是换成另一类结构。

\subsubsection{LeetCode 416 分割等和子集}

先枚举每个物品选或不选、选几次。本题的模型是总和一半是背包容量，每个数只能使用一次；暴力慢在相同阶段和容量反复出现。优化点0-1 背包可达性，一维容量必须倒序决定容量循环方向、状态值含义和是否允许当前物品在同一轮被再次使用。

背包状态围绕阶段和容量。本题的0-1 背包可达性，一维容量必须倒序决定倒序还是正序、求最值还是计数、组合还是排列。循环顺序不是写法偏好，而是题意的一部分。放到“背包模型”这条主线里看，关键原则是：0-1 背包每个物品只能用一次，一维压缩时容量倒序；完全背包物品可无限使用，容量正序。计数组合和计数排列的循环顺序不同。

背包证明围绕当前物品使用次数。0-1 只有选和不选，完全背包允许同一物品继续转移。本题的0-1 背包可达性，一维容量必须倒序要与循环方向一致，否则会把题意改掉。这道题的边界不能留到最后补丁式处理，实验时覆盖总和为奇数、目标为零、数值较大、复杂度依赖 target，并故意把容量方向写反构造最小反例。这样能确认循环方向来自题意，不是背模板。

C++ 实现上，一维背包数组的容量方向要和题意一致。方案数用 \code{long long} 或题目要求的取模；容量来自数值大小时，复杂度是伪多项式。如果追问变成“负数会破坏普通容量模型，需要重新建模”，先判断原来的状态是否仍然覆盖新答案集合，再决定是微调边界、改 value 语义，还是换成另一类结构。

\subsubsection{LeetCode 474 一和零}

先枚举每个物品选或不选、选几次。本题的模型是每个字符串消耗零和一两种容量，每个字符串最多选一次；暴力慢在相同阶段和容量反复出现。优化点二维 0-1 背包，两个容量都要倒序遍历决定容量循环方向、状态值含义和是否允许当前物品在同一轮被再次使用。

背包状态围绕阶段和容量。本题的二维 0-1 背包，两个容量都要倒序遍历决定倒序还是正序、求最值还是计数、组合还是排列。循环顺序不是写法偏好，而是题意的一部分。放到“背包模型”这条主线里看，关键原则是：0-1 背包每个物品只能用一次，一维压缩时容量倒序；完全背包物品可无限使用，容量正序。计数组合和计数排列的循环顺序不同。

背包证明围绕当前物品使用次数。0-1 只有选和不选，完全背包允许同一物品继续转移。本题的二维 0-1 背包，两个容量都要倒序遍历要与循环方向一致，否则会把题意改掉。这道题的边界不能留到最后补丁式处理，实验时覆盖字符串全零或全一、容量为零、倒序方向、计数预处理，并故意把容量方向写反构造最小反例。这样能确认循环方向来自题意，不是背模板。

C++ 实现上，一维背包数组的容量方向要和题意一致。方案数用 \code{long long} 或题目要求的取模；容量来自数值大小时，复杂度是伪多项式。如果追问变成“这是背包从一维容量扩展到二维容量的代表”，先判断原来的状态是否仍然覆盖新答案集合，再决定是微调边界、改 value 语义，还是换成另一类结构。

\subsubsection{LeetCode 494 目标和}

先枚举每个物品选或不选、选几次。本题的模型是正负号选择可代数转换为选若干数凑出特定正和；暴力慢在相同阶段和容量反复出现。优化点若 S 加 target 为奇数或目标绝对值过大，直接无解决定容量循环方向、状态值含义和是否允许当前物品在同一轮被再次使用。

背包状态围绕阶段和容量。本题的若 S 加 target 为奇数或目标绝对值过大，直接无解决定倒序还是正序、求最值还是计数、组合还是排列。循环顺序不是写法偏好，而是题意的一部分。放到“背包模型”这条主线里看，关键原则是：0-1 背包每个物品只能用一次，一维压缩时容量倒序；完全背包物品可无限使用，容量正序。计数组合和计数排列的循环顺序不同。

背包证明围绕当前物品使用次数。0-1 只有选和不选，完全背包允许同一物品继续转移。本题的若 S 加 target 为奇数或目标绝对值过大，直接无解要与循环方向一致，否则会把题意改掉。这道题的边界不能留到最后补丁式处理，实验时覆盖零元素会让方案数翻倍、目标为负、容量为零，并故意把容量方向写反构造最小反例。这样能确认循环方向来自题意，不是背模板。

C++ 实现上，一维背包数组的容量方向要和题意一致。方案数用 \code{long long} 或题目要求的取模；容量来自数值大小时，复杂度是伪多项式。如果追问变成“也可用哈希 DP 保存所有可能和，但背包转换更高效”，先判断原来的状态是否仍然覆盖新答案集合，再决定是微调边界、改 value 语义，还是换成另一类结构。

\subsubsection{LeetCode 518 零钱兑换 II}

先枚举每个物品选或不选、选几次。本题的模型是硬币无限使用，问组合数而不是最少数量；暴力慢在相同阶段和容量反复出现。优化点外层遍历硬币、内层金额正序，保证组合按固定硬币顺序统计决定容量循环方向、状态值含义和是否允许当前物品在同一轮被再次使用。

背包状态围绕阶段和容量。本题的外层遍历硬币、内层金额正序，保证组合按固定硬币顺序统计决定倒序还是正序、求最值还是计数、组合还是排列。循环顺序不是写法偏好，而是题意的一部分。放到“背包模型”这条主线里看，关键原则是：0-1 背包每个物品只能用一次，一维压缩时容量倒序；完全背包物品可无限使用，容量正序。计数组合和计数排列的循环顺序不同。

背包证明围绕当前物品使用次数。0-1 只有选和不选，完全背包允许同一物品继续转移。本题的外层遍历硬币、内层金额正序，保证组合按固定硬币顺序统计要与循环方向一致，否则会把题意改掉。这道题的边界不能留到最后补丁式处理，实验时覆盖amount 为零、无硬币、组合数较大、循环顺序反例，并故意把容量方向写反构造最小反例。这样能确认循环方向来自题意，不是背模板。

C++ 实现上，一维背包数组的容量方向要和题意一致。方案数用 \code{long long} 或题目要求的取模；容量来自数值大小时，复杂度是伪多项式。如果追问变成“若外层金额内层硬币，会统计排列数”，先判断原来的状态是否仍然覆盖新答案集合，再决定是微调边界、改 value 语义，还是换成另一类结构。

\subsubsection{LeetCode 879 盈利计划}

先枚举每个物品选或不选、选几次。本题的模型是每个项目消耗人数并贡献利润，利润达到下限后可截断；暴力慢在相同阶段和容量反复出现。优化点二维 DP 用人数和已达到利润表示方案数，利润维度超过目标就压到目标决定容量循环方向、状态值含义和是否允许当前物品在同一轮被再次使用。

背包状态围绕阶段和容量。本题的二维 DP 用人数和已达到利润表示方案数，利润维度超过目标就压到目标决定倒序还是正序、求最值还是计数、组合还是排列。循环顺序不是写法偏好，而是题意的一部分。放到“背包模型”这条主线里看，关键原则是：0-1 背包每个物品只能用一次，一维压缩时容量倒序；完全背包物品可无限使用，容量正序。计数组合和计数排列的循环顺序不同。

背包证明围绕当前物品使用次数。0-1 只有选和不选，完全背包允许同一物品继续转移。本题的二维 DP 用人数和已达到利润表示方案数，利润维度超过目标就压到目标要与循环方向一致，否则会把题意改掉。这道题的边界不能留到最后补丁式处理，实验时覆盖minProfit 为零、人数不足、方案数取模、项目只能选一次，并故意把容量方向写反构造最小反例。这样能确认循环方向来自题意，不是背模板。

C++ 实现上，一维背包数组的容量方向要和题意一致。方案数用 \code{long long} 或题目要求的取模；容量来自数值大小时，复杂度是伪多项式。如果追问变成“这是 0-1 背包的计数扩展，目标维度需要饱和处理”，先判断原来的状态是否仍然覆盖新答案集合，再决定是微调边界、改 value 语义，还是换成另一类结构。

\subsubsection{LeetCode 956 最高的广告牌}

先枚举每个物品选或不选、选几次。本题的模型是两组钢筋高度相等时得到广告牌，高度差是关键状态；暴力慢在相同阶段和容量反复出现。优化点DP 保存某个高度差下较矮一侧的最大高度，新增钢筋可放高侧、低侧或不用决定容量循环方向、状态值含义和是否允许当前物品在同一轮被再次使用。

背包状态围绕阶段和容量。本题的DP 保存某个高度差下较矮一侧的最大高度，新增钢筋可放高侧、低侧或不用决定倒序还是正序、求最值还是计数、组合还是排列。循环顺序不是写法偏好，而是题意的一部分。放到“背包模型”这条主线里看，关键原则是：0-1 背包每个物品只能用一次，一维压缩时容量倒序；完全背包物品可无限使用，容量正序。计数组合和计数排列的循环顺序不同。

背包证明围绕当前物品使用次数。0-1 只有选和不选，完全背包允许同一物品继续转移。本题的DP 保存某个高度差下较矮一侧的最大高度，新增钢筋可放高侧、低侧或不用要与循环方向一致，否则会把题意改掉。这道题的边界不能留到最后补丁式处理，实验时覆盖空数组、无法平衡、重复长度、差值更新顺序，并故意把容量方向写反构造最小反例。这样能确认循环方向来自题意，不是背模板。

C++ 实现上，一维背包数组的容量方向要和题意一致。方案数用 \code{long long} 或题目要求的取模；容量来自数值大小时，复杂度是伪多项式。如果追问变成“这是背包从容量转成差值状态的代表”，先判断原来的状态是否仍然覆盖新答案集合，再决定是微调边界、改 value 语义，还是换成另一类结构。

\subsubsection{LeetCode 1049 最后一块石头的重量 II}

先枚举每个物品选或不选、选几次。本题的模型是把石头分成两堆，最终差值等于两堆重量差；暴力慢在相同阶段和容量反复出现。优化点0-1 背包找不超过总和一半的最大可达重量决定容量循环方向、状态值含义和是否允许当前物品在同一轮被再次使用。

背包状态围绕阶段和容量。本题的0-1 背包找不超过总和一半的最大可达重量决定倒序还是正序、求最值还是计数、组合还是排列。循环顺序不是写法偏好，而是题意的一部分。放到“背包模型”这条主线里看，关键原则是：0-1 背包每个物品只能用一次，一维压缩时容量倒序；完全背包物品可无限使用，容量正序。计数组合和计数排列的循环顺序不同。

背包证明围绕当前物品使用次数。0-1 只有选和不选，完全背包允许同一物品继续转移。本题的0-1 背包找不超过总和一半的最大可达重量要与循环方向一致，否则会把题意改掉。这道题的边界不能留到最后补丁式处理，实验时覆盖单块石头、总和为奇偶、多个相同重量、容量方向，并故意把容量方向写反构造最小反例。这样能确认循环方向来自题意，不是背模板。

C++ 实现上，一维背包数组的容量方向要和题意一致。方案数用 \code{long long} 或题目要求的取模；容量来自数值大小时，复杂度是伪多项式。如果追问变成“它和分割等和子集同源，只是目标从是否可达变成最小差”，先判断原来的状态是否仍然覆盖新答案集合，再决定是微调边界、改 value 语义，还是换成另一类结构。

\subsubsection{LeetCode 1155 掷骰子等于目标和的方法数}

先枚举每个物品选或不选、选几次。本题的模型是多个骰子每个只能取一到 faces 的点数，问目标和方案数；暴力慢在相同阶段和容量反复出现。优化点按骰子数量分阶段，容量为当前点数和，枚举当前骰子的点数转移决定容量循环方向、状态值含义和是否允许当前物品在同一轮被再次使用。

背包状态围绕阶段和容量。本题的按骰子数量分阶段，容量为当前点数和，枚举当前骰子的点数转移决定倒序还是正序、求最值还是计数、组合还是排列。循环顺序不是写法偏好，而是题意的一部分。放到“背包模型”这条主线里看，关键原则是：0-1 背包每个物品只能用一次，一维压缩时容量倒序；完全背包物品可无限使用，容量正序。计数组合和计数排列的循环顺序不同。

背包证明围绕当前物品使用次数。0-1 只有选和不选，完全背包允许同一物品继续转移。本题的按骰子数量分阶段，容量为当前点数和，枚举当前骰子的点数转移要与循环方向一致，否则会把题意改掉。这道题的边界不能留到最后补丁式处理，实验时覆盖target 太小或太大、取模、骰子数为一、滚动数组方向，并故意把容量方向写反构造最小反例。这样能确认循环方向来自题意，不是背模板。

C++ 实现上，一维背包数组的容量方向要和题意一致。方案数用 \code{long long} 或题目要求的取模；容量来自数值大小时，复杂度是伪多项式。如果追问变成“这是有界选择计数，和完全背包不能混”，先判断原来的状态是否仍然覆盖新答案集合，再决定是微调边界、改 value 语义，还是换成另一类结构。

\subsubsection{LeetCode 1449 数位成本和为目标值的最大数字}

先枚举每个物品选或不选、选几次。本题的模型是每个数字可重复选择，容量是成本，目标是组成字典序和长度最大的数字；暴力慢在相同阶段和容量反复出现。优化点完全背包先最大化位数，再按数字从大到小贪心还原答案决定容量循环方向、状态值含义和是否允许当前物品在同一轮被再次使用。

背包状态围绕阶段和容量。本题的完全背包先最大化位数，再按数字从大到小贪心还原答案决定倒序还是正序、求最值还是计数、组合还是排列。循环顺序不是写法偏好，而是题意的一部分。放到“背包模型”这条主线里看，关键原则是：0-1 背包每个物品只能用一次，一维压缩时容量倒序；完全背包物品可无限使用，容量正序。计数组合和计数排列的循环顺序不同。

背包证明围绕当前物品使用次数。0-1 只有选和不选，完全背包允许同一物品继续转移。本题的完全背包先最大化位数，再按数字从大到小贪心还原答案要与循环方向一致，否则会把题意改掉。这道题的边界不能留到最后补丁式处理，实验时覆盖无法组成目标、多个数字同成本、target 为零、字符串比较，并故意把容量方向写反构造最小反例。这样能确认循环方向来自题意，不是背模板。

C++ 实现上，一维背包数组的容量方向要和题意一致。方案数用 \code{long long} 或题目要求的取模；容量来自数值大小时，复杂度是伪多项式。如果追问变成“这题说明背包状态值不一定是数值收益，也可以是长度和字典序”，先判断原来的状态是否仍然覆盖新答案集合，再决定是微调边界、改 value 语义，还是换成另一类结构。

\subsection{贪心与交换证明工作坊}

先把选择空间完整枚举，再寻找可以交换到最优解前面的局部选择。贪心必须能证明当前选择不会让后续资源变差，而不是只看排序后代码短。

\subsubsection{LeetCode 45 跳跃游戏 II}

先把选择顺序完整枚举，哪怕这个版本只能处理很小输入。本题的模型是每一次跳跃可以看成 BFS 的一层，当前层内所有位置共享同一次跳数；真正要证明的是扫描当前层时维护下一层最远边界，到达当前层末尾才增加跳数为什么不会伤害后续选择。没有这一步证明，排序或局部选择都只是猜测。

贪心状态要表达当前方案为什么不落后。本题围绕扫描当前层时维护下一层最远边界，到达当前层末尾才增加跳数保留局部最优边界；排序只是把可比较的候选排到一起，证明仍要落在交换或领先关系上。放到“贪心与交换证明”这条主线里看，关键原则是：贪心需要找到可交换的局部最优：最早结束、最小右端、当前最远覆盖、当前最大收益或最小代价。排序只是手段，不是证明。

贪心证明要么用交换，要么用领先性。围绕扫描当前层时维护下一层最远边界，到达当前层末尾才增加跳数说明把任意最优解的第一步换成贪心选择不会变差，或说明当前方案在同等步数下始终不落后。这道题的边界不能留到最后补丁式处理，实验时除了数组长度一、刚好到达、存在多个可选跳点、不要在每个位置都加跳数，还要故意替换成一个看似合理但错误的排序规则，构造反例。能解释错误贪心为何失败，才算理解正确贪心。

C++ 实现上，比较器必须满足严格弱序，相同关键字的处理要服务于证明。涉及区间端点、收益和代价时，排序键要写成业务含义，不要只凭直觉写小于号。如果追问变成“若问能否到达，只需维护全局最远可达位置”，先判断原来的状态是否仍然覆盖新答案集合，再决定是微调边界、改 value 语义，还是换成另一类结构。

\subsubsection{LeetCode 55 跳跃游戏}

先把选择顺序完整枚举，哪怕这个版本只能处理很小输入。本题的模型是扫描过程中只关心当前能到达的最远位置；真正要证明的是如果下标超过最远可达，任何之前路径都不能到达它；否则用它扩展边界为什么不会伤害后续选择。没有这一步证明，排序或局部选择都只是猜测。

贪心状态要表达当前方案为什么不落后。本题围绕如果下标超过最远可达，任何之前路径都不能到达它；否则用它扩展边界保留局部最优边界；排序只是把可比较的候选排到一起，证明仍要落在交换或领先关系上。放到“贪心与交换证明”这条主线里看，关键原则是：贪心需要找到可交换的局部最优：最早结束、最小右端、当前最远覆盖、当前最大收益或最小代价。排序只是手段，不是证明。

贪心证明要么用交换，要么用领先性。围绕如果下标超过最远可达，任何之前路径都不能到达它；否则用它扩展边界说明把任意最优解的第一步换成贪心选择不会变差，或说明当前方案在同等步数下始终不落后。这道题的边界不能留到最后补丁式处理，实验时除了第一个元素为零、数组长度一、恰好到达最后、远超最后，还要故意替换成一个看似合理但错误的排序规则，构造反例。能解释错误贪心为何失败，才算理解正确贪心。

C++ 实现上，比较器必须满足严格弱序，相同关键字的处理要服务于证明。涉及区间端点、收益和代价时，排序键要写成业务含义，不要只凭直觉写小于号。如果追问变成“求最少跳数时把可达区间看成 BFS 层”，先判断原来的状态是否仍然覆盖新答案集合，再决定是微调边界、改 value 语义，还是换成另一类结构。

\subsubsection{LeetCode 122 买卖股票 II}

先把选择顺序完整枚举，哪怕这个版本只能处理很小输入。本题的模型是允许多次交易且不能同时持有，所有正收益差都可以收集；真正要证明的是把长上涨段拆成每天买卖收益相同，因此累加正差值最优为什么不会伤害后续选择。没有这一步证明，排序或局部选择都只是猜测。

贪心状态要表达当前方案为什么不落后。本题围绕把长上涨段拆成每天买卖收益相同，因此累加正差值最优保留局部最优边界；排序只是把可比较的候选排到一起，证明仍要落在交换或领先关系上。放到“贪心与交换证明”这条主线里看，关键原则是：贪心需要找到可交换的局部最优：最早结束、最小右端、当前最远覆盖、当前最大收益或最小代价。排序只是手段，不是证明。

贪心证明要么用交换，要么用领先性。围绕把长上涨段拆成每天买卖收益相同，因此累加正差值最优说明把任意最优解的第一步换成贪心选择不会变差，或说明当前方案在同等步数下始终不落后。这道题的边界不能留到最后补丁式处理，实验时除了下降段、平台价格、只有一天、手续费不存在，还要故意替换成一个看似合理但错误的排序规则，构造反例。能解释错误贪心为何失败，才算理解正确贪心。

C++ 实现上，比较器必须满足严格弱序，相同关键字的处理要服务于证明。涉及区间端点、收益和代价时，排序键要写成业务含义，不要只凭直觉写小于号。如果追问变成“有手续费时不能简单累加正差，需要状态机或调整贪心阈值”，先判断原来的状态是否仍然覆盖新答案集合，再决定是微调边界、改 value 语义，还是换成另一类结构。

\subsubsection{LeetCode 134 加油站}

先把选择顺序完整枚举，哪怕这个版本只能处理很小输入。本题的模型是绕行一圈是否可行取决于总油量差和局部剩余油量；真正要证明的是总和为负必不可行；扫描时若当前剩余为负，起点只能放到下一站为什么不会伤害后续选择。没有这一步证明，排序或局部选择都只是猜测。

贪心状态要表达当前方案为什么不落后。本题围绕总和为负必不可行；扫描时若当前剩余为负，起点只能放到下一站保留局部最优边界；排序只是把可比较的候选排到一起，证明仍要落在交换或领先关系上。放到“贪心与交换证明”这条主线里看，关键原则是：贪心需要找到可交换的局部最优：最早结束、最小右端、当前最远覆盖、当前最大收益或最小代价。排序只是手段，不是证明。

贪心证明要么用交换，要么用领先性。围绕总和为负必不可行；扫描时若当前剩余为负，起点只能放到下一站说明把任意最优解的第一步换成贪心选择不会变差，或说明当前方案在同等步数下始终不落后。这道题的边界不能留到最后补丁式处理，实验时除了总和刚好为零、多个可行起点、某站油量为零、消耗大于补给，还要故意替换成一个看似合理但错误的排序规则，构造反例。能解释错误贪心为何失败，才算理解正确贪心。

C++ 实现上，比较器必须满足严格弱序，相同关键字的处理要服务于证明。涉及区间端点、收益和代价时，排序键要写成业务含义，不要只凭直觉写小于号。如果追问变成“它的领先性证明是前一段任意站点都不能作为起点”，先判断原来的状态是否仍然覆盖新答案集合，再决定是微调边界、改 value 语义，还是换成另一类结构。

\subsubsection{LeetCode 135 分发糖果}

先把选择顺序完整枚举，哪怕这个版本只能处理很小输入。本题的模型是相邻评分约束需要同时满足左邻和右邻两个方向；真正要证明的是左到右满足递增关系，右到左满足反向递增关系，最后取两侧要求最大值为什么不会伤害后续选择。没有这一步证明，排序或局部选择都只是猜测。

贪心状态要表达当前方案为什么不落后。本题围绕左到右满足递增关系，右到左满足反向递增关系，最后取两侧要求最大值保留局部最优边界；排序只是把可比较的候选排到一起，证明仍要落在交换或领先关系上。放到“贪心与交换证明”这条主线里看，关键原则是：贪心需要找到可交换的局部最优：最早结束、最小右端、当前最远覆盖、当前最大收益或最小代价。排序只是手段，不是证明。

贪心证明要么用交换，要么用领先性。围绕左到右满足递增关系，右到左满足反向递增关系，最后取两侧要求最大值说明把任意最优解的第一步换成贪心选择不会变差，或说明当前方案在同等步数下始终不落后。这道题的边界不能留到最后补丁式处理，实验时除了评分相等、严格递增、严格递减、山峰和谷底，还要故意替换成一个看似合理但错误的排序规则，构造反例。能解释错误贪心为何失败，才算理解正确贪心。

C++ 实现上，比较器必须满足严格弱序，相同关键字的处理要服务于证明。涉及区间端点、收益和代价时，排序键要写成业务含义，不要只凭直觉写小于号。如果追问变成“这是局部约束两遍扫描，不能只用一个方向贪心”，先判断原来的状态是否仍然覆盖新答案集合，再决定是微调边界、改 value 语义，还是换成另一类结构。

\subsubsection{LeetCode 169 多数元素}

先把选择顺序完整枚举，哪怕这个版本只能处理很小输入。本题的模型是超过一半的元素可以和其他元素两两抵消后仍然剩下；真正要证明的是Boyer-Moore 投票维护候选和计数为什么不会伤害后续选择。没有这一步证明，排序或局部选择都只是猜测。

贪心状态要表达当前方案为什么不落后。本题围绕Boyer-Moore 投票维护候选和计数保留局部最优边界；排序只是把可比较的候选排到一起，证明仍要落在交换或领先关系上。放到“贪心与交换证明”这条主线里看，关键原则是：贪心需要找到可交换的局部最优：最早结束、最小右端、当前最远覆盖、当前最大收益或最小代价。排序只是手段，不是证明。

贪心证明要么用交换，要么用领先性。围绕Boyer-Moore 投票维护候选和计数说明把任意最优解的第一步换成贪心选择不会变差，或说明当前方案在同等步数下始终不落后。这道题的边界不能留到最后补丁式处理，实验时除了题目是否保证存在多数、计数归零、全相同，还要故意替换成一个看似合理但错误的排序规则，构造反例。能解释错误贪心为何失败，才算理解正确贪心。

C++ 实现上，比较器必须满足严格弱序，相同关键字的处理要服务于证明。涉及区间端点、收益和代价时，排序键要写成业务含义，不要只凭直觉写小于号。如果追问变成“若找超过三分之一的元素，需要维护两个候选”，先判断原来的状态是否仍然覆盖新答案集合，再决定是微调边界、改 value 语义，还是换成另一类结构。

\subsubsection{LeetCode 316 去除重复字母}

先把选择顺序完整枚举，哪怕这个版本只能处理很小输入。本题的模型是每个字符必须保留一次，同时结果字典序尽量小；真正要证明的是单调栈维护当前结果，若栈顶更大且后面还会出现，就可以弹出再等后面放回为什么不会伤害后续选择。没有这一步证明，排序或局部选择都只是猜测。

贪心状态要表达当前方案为什么不落后。本题围绕单调栈维护当前结果，若栈顶更大且后面还会出现，就可以弹出再等后面放回保留局部最优边界；排序只是把可比较的候选排到一起，证明仍要落在交换或领先关系上。放到“贪心与交换证明”这条主线里看，关键原则是：贪心需要找到可交换的局部最优：最早结束、最小右端、当前最远覆盖、当前最大收益或最小代价。排序只是手段，不是证明。

贪心证明要么用交换，要么用领先性。围绕单调栈维护当前结果，若栈顶更大且后面还会出现，就可以弹出再等后面放回说明把任意最优解的第一步换成贪心选择不会变差，或说明当前方案在同等步数下始终不落后。这道题的边界不能留到最后补丁式处理，实验时除了字符只出现一次、重复字符、弹出后 visited 恢复、字典序比较，还要故意替换成一个看似合理但错误的排序规则，构造反例。能解释错误贪心为何失败，才算理解正确贪心。

C++ 实现上，比较器必须满足严格弱序，相同关键字的处理要服务于证明。涉及区间端点、收益和代价时，排序键要写成业务含义，不要只凭直觉写小于号。如果追问变成“402 移掉 K 位数字也是单调栈删除，但删除预算和必须保留集合不同”，先判断原来的状态是否仍然覆盖新答案集合，再决定是微调边界、改 value 语义，还是换成另一类结构。

\subsubsection{LeetCode 406 根据身高重建队列}

先把选择顺序完整枚举，哪怕这个版本只能处理很小输入。本题的模型是高个子的位置先确定后，不会被矮个子影响其前方高个数量；真正要证明的是按身高降序、k 升序排序，再按 k 插入当前结果位置为什么不会伤害后续选择。没有这一步证明，排序或局部选择都只是猜测。

贪心状态要表达当前方案为什么不落后。本题围绕按身高降序、k 升序排序，再按 k 插入当前结果位置保留局部最优边界；排序只是把可比较的候选排到一起，证明仍要落在交换或领先关系上。放到“贪心与交换证明”这条主线里看，关键原则是：贪心需要找到可交换的局部最优：最早结束、最小右端、当前最远覆盖、当前最大收益或最小代价。排序只是手段，不是证明。

贪心证明要么用交换，要么用领先性。围绕按身高降序、k 升序排序，再按 k 插入当前结果位置说明把任意最优解的第一步换成贪心选择不会变差，或说明当前方案在同等步数下始终不落后。这道题的边界不能留到最后补丁式处理，实验时除了相同身高、k 为零、插入到末尾、结果稳定性，还要故意替换成一个看似合理但错误的排序规则，构造反例。能解释错误贪心为何失败，才算理解正确贪心。

C++ 实现上，比较器必须满足严格弱序，相同关键字的处理要服务于证明。涉及区间端点、收益和代价时，排序键要写成业务含义，不要只凭直觉写小于号。如果追问变成“这是排序后插入的构造型贪心，证明依赖高个子先固定”，先判断原来的状态是否仍然覆盖新答案集合，再决定是微调边界、改 value 语义，还是换成另一类结构。

\subsubsection{LeetCode 435 无重叠区间}

先把选择顺序完整枚举，哪怕这个版本只能处理很小输入。本题的模型是保留最多不重叠区间等价于删除最少区间；真正要证明的是按终点升序选择，结束越早给后续留下空间越多为什么不会伤害后续选择。没有这一步证明，排序或局部选择都只是猜测。

贪心状态要表达当前方案为什么不落后。本题围绕按终点升序选择，结束越早给后续留下空间越多保留局部最优边界；排序只是把可比较的候选排到一起，证明仍要落在交换或领先关系上。放到“贪心与交换证明”这条主线里看，关键原则是：贪心需要找到可交换的局部最优：最早结束、最小右端、当前最远覆盖、当前最大收益或最小代价。排序只是手段，不是证明。

贪心证明要么用交换，要么用领先性。围绕按终点升序选择，结束越早给后续留下空间越多说明把任意最优解的第一步换成贪心选择不会变差，或说明当前方案在同等步数下始终不落后。这道题的边界不能留到最后补丁式处理，实验时除了端点相接是否重叠、完全覆盖、相同终点、空列表，还要故意替换成一个看似合理但错误的排序规则，构造反例。能解释错误贪心为何失败，才算理解正确贪心。

C++ 实现上，比较器必须满足严格弱序，相同关键字的处理要服务于证明。涉及区间端点、收益和代价时，排序键要写成业务含义，不要只凭直觉写小于号。如果追问变成“用交换论证说明最早结束的选择不劣于其他第一选择”，先判断原来的状态是否仍然覆盖新答案集合，再决定是微调边界、改 value 语义，还是换成另一类结构。

\subsubsection{LeetCode 452 用最少数量的箭引爆气球}

先把选择顺序完整枚举，哪怕这个版本只能处理很小输入。本题的模型是一支箭的位置可以覆盖所有包含该位置的区间；真正要证明的是按右端点排序，当前箭放在最早结束气球的右端为什么不会伤害后续选择。没有这一步证明，排序或局部选择都只是猜测。

贪心状态要表达当前方案为什么不落后。本题围绕按右端点排序，当前箭放在最早结束气球的右端保留局部最优边界；排序只是把可比较的候选排到一起，证明仍要落在交换或领先关系上。放到“贪心与交换证明”这条主线里看，关键原则是：贪心需要找到可交换的局部最优：最早结束、最小右端、当前最远覆盖、当前最大收益或最小代价。排序只是手段，不是证明。

贪心证明要么用交换，要么用领先性。围绕按右端点排序，当前箭放在最早结束气球的右端说明把任意最优解的第一步换成贪心选择不会变差，或说明当前方案在同等步数下始终不落后。这道题的边界不能留到最后补丁式处理，实验时除了端点相同、负坐标、区间包含、闭区间边界，还要故意替换成一个看似合理但错误的排序规则，构造反例。能解释错误贪心为何失败，才算理解正确贪心。

C++ 实现上，比较器必须满足严格弱序，相同关键字的处理要服务于证明。涉及区间端点、收益和代价时，排序键要写成业务含义，不要只凭直觉写小于号。如果追问变成“它和无重叠区间是同一类最早结束贪心”，先判断原来的状态是否仍然覆盖新答案集合，再决定是微调边界、改 value 语义，还是换成另一类结构。

\subsubsection{LeetCode 630 课程表 III}

先把选择顺序完整枚举，哪怕这个版本只能处理很小输入。本题的模型是每门课有持续时间和截止日，已选课程总时长不能超过当前截止日；真正要证明的是按截止日排序，先选当前课；若超时，就用大顶堆丢掉已选中耗时最长的课为什么不会伤害后续选择。没有这一步证明，排序或局部选择都只是猜测。

贪心状态要表达当前方案为什么不落后。本题围绕按截止日排序，先选当前课；若超时，就用大顶堆丢掉已选中耗时最长的课保留局部最优边界；排序只是把可比较的候选排到一起，证明仍要落在交换或领先关系上。放到“贪心与交换证明”这条主线里看，关键原则是：贪心需要找到可交换的局部最优：最早结束、最小右端、当前最远覆盖、当前最大收益或最小代价。排序只是手段，不是证明。

贪心证明要么用交换，要么用领先性。围绕按截止日排序，先选当前课；若超时，就用大顶堆丢掉已选中耗时最长的课说明把任意最优解的第一步换成贪心选择不会变差，或说明当前方案在同等步数下始终不落后。这道题的边界不能留到最后补丁式处理，实验时除了课程无法单独完成、截止日相同、替换已选课程、空列表，还要故意替换成一个看似合理但错误的排序规则，构造反例。能解释错误贪心为何失败，才算理解正确贪心。

C++ 实现上，比较器必须满足严格弱序，相同关键字的处理要服务于证明。涉及区间端点、收益和代价时，排序键要写成业务含义，不要只凭直觉写小于号。如果追问变成“这是反悔贪心：接受后在约束被破坏时删最差候选”，先判断原来的状态是否仍然覆盖新答案集合，再决定是微调边界、改 value 语义，还是换成另一类结构。

\subsubsection{LeetCode 646 最长数对链}

先把选择顺序完整枚举，哪怕这个版本只能处理很小输入。本题的模型是每个数对像一个区间，选择链要求前一个右端小于后一个左端；真正要证明的是按右端升序选择，当前右端越小越容易接后续数对为什么不会伤害后续选择。没有这一步证明，排序或局部选择都只是猜测。

贪心状态要表达当前方案为什么不落后。本题围绕按右端升序选择，当前右端越小越容易接后续数对保留局部最优边界；排序只是把可比较的候选排到一起，证明仍要落在交换或领先关系上。放到“贪心与交换证明”这条主线里看，关键原则是：贪心需要找到可交换的局部最优：最早结束、最小右端、当前最远覆盖、当前最大收益或最小代价。排序只是手段，不是证明。

贪心证明要么用交换，要么用领先性。围绕按右端升序选择，当前右端越小越容易接后续数对说明把任意最优解的第一步换成贪心选择不会变差，或说明当前方案在同等步数下始终不落后。这道题的边界不能留到最后补丁式处理，实验时除了端点相等不允许连接、完全覆盖、负数端点、空数组，还要故意替换成一个看似合理但错误的排序规则，构造反例。能解释错误贪心为何失败，才算理解正确贪心。

C++ 实现上，比较器必须满足严格弱序，相同关键字的处理要服务于证明。涉及区间端点、收益和代价时，排序键要写成业务含义，不要只凭直觉写小于号。如果追问变成“也可以 DP，但贪心用最早结束证明更简洁”，先判断原来的状态是否仍然覆盖新答案集合，再决定是微调边界、改 value 语义，还是换成另一类结构。

\subsubsection{LeetCode 763 划分字母区间}

先把选择顺序完整枚举，哪怕这个版本只能处理很小输入。本题的模型是每个片段必须覆盖其中所有字符的最后出现位置；真正要证明的是扫描维护当前片段最远右端，到达右端即可切分为什么不会伤害后续选择。没有这一步证明，排序或局部选择都只是猜测。

贪心状态要表达当前方案为什么不落后。本题围绕扫描维护当前片段最远右端，到达右端即可切分保留局部最优边界；排序只是把可比较的候选排到一起，证明仍要落在交换或领先关系上。放到“贪心与交换证明”这条主线里看，关键原则是：贪心需要找到可交换的局部最优：最早结束、最小右端、当前最远覆盖、当前最大收益或最小代价。排序只是手段，不是证明。

贪心证明要么用交换，要么用领先性。围绕扫描维护当前片段最远右端，到达右端即可切分说明把任意最优解的第一步换成贪心选择不会变差，或说明当前方案在同等步数下始终不落后。这道题的边界不能留到最后补丁式处理，实验时除了单字符、所有字符相同、字符多次跨段，还要故意替换成一个看似合理但错误的排序规则，构造反例。能解释错误贪心为何失败，才算理解正确贪心。

C++ 实现上，比较器必须满足严格弱序，相同关键字的处理要服务于证明。涉及区间端点、收益和代价时，排序键要写成业务含义，不要只凭直觉写小于号。如果追问变成“这是最早合法切分贪心，能最大化片段数量”，先判断原来的状态是否仍然覆盖新答案集合，再决定是微调边界、改 value 语义，还是换成另一类结构。

\subsection{区间、排序与合并工作坊}

先两两比较或完整重排，再利用排序后相交区间连续出现的性质。动态区间题还要关注前驱、后继、端点闭开和覆盖关系。

\subsubsection{LeetCode 56 合并区间}

先两两比较区间，确认相交、覆盖或冲突的定义。本题的模型是排序后可能与当前区间相交的只会是结果尾部；暴力慢在每个区间都回头看太多无关区间。优化点按起点排序，扫描时维护结果最后区间的右端点通常依赖排序后相邻关系，必须先定义端点闭开。

区间结构的状态是当前右端、结果尾区间或动态有序表的相邻关系。本题的按起点排序，扫描时维护结果最后区间的右端点只允许丢掉已经不可能再冲突的区间；端点相等是否合并要先定义。放到“区间、排序与合并”这条主线里看，关键原则是：合并按起点排序，调度按终点排序，覆盖题常用起点升序且终点降序。扫描时维护当前最远右端、结果尾区间或已选区间终点。

区间题证明依赖排序后的邻接关系。若当前区间无法影响结果尾部或相邻区间，它就不会再影响更远的区间。本题的按起点排序，扫描时维护结果最后区间的右端点要说明哪些区间已经安全归档。这道题的边界不能留到最后补丁式处理，实验时覆盖端点相接、完全覆盖、无重叠、空列表、输入未排序，画出排序后的区间序列和当前结果尾部。动态版本要专门测前驱、后继和端点相接。

C++ 实现上，区间可用 \code{std::vector<std::array<int, 2>>} 或结构体表达端点，排序比较器要处理相同起点。动态区间需要 \code{std::map} 的前驱后继，而不是哈希表。如果追问变成“若区间是半开区间，端点相等不一定合并”，先判断原来的状态是否仍然覆盖新答案集合，再决定是微调边界、改 value 语义，还是换成另一类结构。

\subsubsection{LeetCode 57 插入区间}

先两两比较区间，确认相交、覆盖或冲突的定义。本题的模型是新区间只会影响与它相交的一段连续区间；暴力慢在每个区间都回头看太多无关区间。优化点先追加左侧完全不相交区间，再合并重叠段，最后追加右侧通常依赖排序后相邻关系，必须先定义端点闭开。

区间结构的状态是当前右端、结果尾区间或动态有序表的相邻关系。本题的先追加左侧完全不相交区间，再合并重叠段，最后追加右侧只允许丢掉已经不可能再冲突的区间；端点相等是否合并要先定义。放到“区间、排序与合并”这条主线里看，关键原则是：合并按起点排序，调度按终点排序，覆盖题常用起点升序且终点降序。扫描时维护当前最远右端、结果尾区间或已选区间终点。

区间题证明依赖排序后的邻接关系。若当前区间无法影响结果尾部或相邻区间，它就不会再影响更远的区间。本题的先追加左侧完全不相交区间，再合并重叠段，最后追加右侧要说明哪些区间已经安全归档。这道题的边界不能留到最后补丁式处理，实验时覆盖新区间在最前、最后、覆盖全部、刚好相邻、原列表为空，画出排序后的区间序列和当前结果尾部。动态版本要专门测前驱、后继和端点相接。

C++ 实现上，区间可用 \code{std::vector<std::array<int, 2>>} 或结构体表达端点，排序比较器要处理相同起点。动态区间需要 \code{std::map} 的前驱后继，而不是哈希表。如果追问变成“若输入不是有序且不重叠，必须先排序或换模型”，先判断原来的状态是否仍然覆盖新答案集合，再决定是微调边界、改 value 语义，还是换成另一类结构。

\subsubsection{LeetCode 252 会议室}

先两两比较区间，确认相交、覆盖或冲突的定义。本题的模型是每个会议占用一个时间区间，任意重叠都表示不能全部参加；暴力慢在每个区间都回头看太多无关区间。优化点按开始时间排序，只需检查相邻会议是否冲突通常依赖排序后相邻关系，必须先定义端点闭开。

区间结构的状态是当前右端、结果尾区间或动态有序表的相邻关系。本题的按开始时间排序，只需检查相邻会议是否冲突只允许丢掉已经不可能再冲突的区间；端点相等是否合并要先定义。放到“区间、排序与合并”这条主线里看，关键原则是：合并按起点排序，调度按终点排序，覆盖题常用起点升序且终点降序。扫描时维护当前最远右端、结果尾区间或已选区间终点。

区间题证明依赖排序后的邻接关系。若当前区间无法影响结果尾部或相邻区间，它就不会再影响更远的区间。本题的按开始时间排序，只需检查相邻会议是否冲突要说明哪些区间已经安全归档。这道题的边界不能留到最后补丁式处理，实验时覆盖端点相等是否冲突、空列表、单会议、完全重叠，画出排序后的区间序列和当前结果尾部。动态版本要专门测前驱、后继和端点相接。

C++ 实现上，区间可用 \code{std::vector<std::array<int, 2>>} 或结构体表达端点，排序比较器要处理相同起点。动态区间需要 \code{std::map} 的前驱后继，而不是哈希表。如果追问变成“若问最少会议室，需要维护同时进行的会议数量”，先判断原来的状态是否仍然覆盖新答案集合，再决定是微调边界、改 value 语义，还是换成另一类结构。

\subsubsection{LeetCode 253 会议室 II}

先两两比较区间，确认相交、覆盖或冲突的定义。本题的模型是同时进行的会议数量决定最少房间数；暴力慢在每个区间都回头看太多无关区间。优化点按开始时间扫描，用小顶堆维护当前会议的最早结束时间通常依赖排序后相邻关系，必须先定义端点闭开。

区间结构的状态是当前右端、结果尾区间或动态有序表的相邻关系。本题的按开始时间扫描，用小顶堆维护当前会议的最早结束时间只允许丢掉已经不可能再冲突的区间；端点相等是否合并要先定义。放到“区间、排序与合并”这条主线里看，关键原则是：合并按起点排序，调度按终点排序，覆盖题常用起点升序且终点降序。扫描时维护当前最远右端、结果尾区间或已选区间终点。

区间题证明依赖排序后的邻接关系。若当前区间无法影响结果尾部或相邻区间，它就不会再影响更远的区间。本题的按开始时间扫描，用小顶堆维护当前会议的最早结束时间要说明哪些区间已经安全归档。这道题的边界不能留到最后补丁式处理，实验时覆盖会议端点相接、多个同一时间开始、空列表、堆顶释放条件，画出排序后的区间序列和当前结果尾部。动态版本要专门测前驱、后继和端点相接。

C++ 实现上，区间可用 \code{std::vector<std::array<int, 2>>} 或结构体表达端点，排序比较器要处理相同起点。动态区间需要 \code{std::map} 的前驱后继，而不是哈希表。如果追问变成“也可以把开始和结束分别排序，用双指针统计并发数”，先判断原来的状态是否仍然覆盖新答案集合，再决定是微调边界、改 value 语义，还是换成另一类结构。

\subsubsection{LeetCode 352 将数据流变为多个不相交区间}

先两两比较区间，确认相交、覆盖或冲突的定义。本题的模型是数字逐个到达，需要动态维护已覆盖的有序不相交区间；暴力慢在每个区间都回头看太多无关区间。优化点用有序表找到新值附近的前驱和后继，判断是否连接、合并或新建区间通常依赖排序后相邻关系，必须先定义端点闭开。

区间结构的状态是当前右端、结果尾区间或动态有序表的相邻关系。本题的用有序表找到新值附近的前驱和后继，判断是否连接、合并或新建区间只允许丢掉已经不可能再冲突的区间；端点相等是否合并要先定义。放到“区间、排序与合并”这条主线里看，关键原则是：合并按起点排序，调度按终点排序，覆盖题常用起点升序且终点降序。扫描时维护当前最远右端、结果尾区间或已选区间终点。

区间题证明依赖排序后的邻接关系。若当前区间无法影响结果尾部或相邻区间，它就不会再影响更远的区间。本题的用有序表找到新值附近的前驱和后继，判断是否连接、合并或新建区间要说明哪些区间已经安全归档。这道题的边界不能留到最后补丁式处理，实验时覆盖重复插入、连接左右两个区间、插入最小最大值、相邻端点，画出排序后的区间序列和当前结果尾部。动态版本要专门测前驱、后继和端点相接。

C++ 实现上，区间可用 \code{std::vector<std::array<int, 2>>} 或结构体表达端点，排序比较器要处理相同起点。动态区间需要 \code{std::map} 的前驱后继，而不是哈希表。如果追问变成“这类动态区间不能用普通数组扫描长期维护”，先判断原来的状态是否仍然覆盖新答案集合，再决定是微调边界、改 value 语义，还是换成另一类结构。

\subsubsection{LeetCode 436 寻找右区间}

先两两比较区间，确认相交、覆盖或冲突的定义。本题的模型是每个区间需要寻找起点不小于当前终点的最小起点区间；暴力慢在每个区间都回头看太多无关区间。优化点把所有起点和原下标排序，对每个终点做 lower bound通常依赖排序后相邻关系，必须先定义端点闭开。

区间结构的状态是当前右端、结果尾区间或动态有序表的相邻关系。本题的把所有起点和原下标排序，对每个终点做 lower bound只允许丢掉已经不可能再冲突的区间；端点相等是否合并要先定义。放到“区间、排序与合并”这条主线里看，关键原则是：合并按起点排序，调度按终点排序，覆盖题常用起点升序且终点降序。扫描时维护当前最远右端、结果尾区间或已选区间终点。

区间题证明依赖排序后的邻接关系。若当前区间无法影响结果尾部或相邻区间，它就不会再影响更远的区间。本题的把所有起点和原下标排序，对每个终点做 lower bound要说明哪些区间已经安全归档。这道题的边界不能留到最后补丁式处理，实验时覆盖不存在右区间、起点相同、负端点、单区间，画出排序后的区间序列和当前结果尾部。动态版本要专门测前驱、后继和端点相接。

C++ 实现上，区间可用 \code{std::vector<std::array<int, 2>>} 或结构体表达端点，排序比较器要处理相同起点。动态区间需要 \code{std::map} 的前驱后继，而不是哈希表。如果追问变成“如果区间动态插入，则需要有序表而不是一次排序数组”，先判断原来的状态是否仍然覆盖新答案集合，再决定是微调边界、改 value 语义，还是换成另一类结构。

\subsubsection{LeetCode 729 我的日程安排表 I}

先两两比较区间，确认相交、覆盖或冲突的定义。本题的模型是每次插入新区间时，需要判断它是否和已有区间重叠；暴力慢在每个区间都回头看太多无关区间。优化点有序表按起点存区间，只检查新区间的前驱和后继通常依赖排序后相邻关系，必须先定义端点闭开。

区间结构的状态是当前右端、结果尾区间或动态有序表的相邻关系。本题的有序表按起点存区间，只检查新区间的前驱和后继只允许丢掉已经不可能再冲突的区间；端点相等是否合并要先定义。放到“区间、排序与合并”这条主线里看，关键原则是：合并按起点排序，调度按终点排序，覆盖题常用起点升序且终点降序。扫描时维护当前最远右端、结果尾区间或已选区间终点。

区间题证明依赖排序后的邻接关系。若当前区间无法影响结果尾部或相邻区间，它就不会再影响更远的区间。本题的有序表按起点存区间，只检查新区间的前驱和后继要说明哪些区间已经安全归档。这道题的边界不能留到最后补丁式处理，实验时覆盖端点相接、插入最前最后、完全覆盖已有区间、动态多次插入，画出排序后的区间序列和当前结果尾部。动态版本要专门测前驱、后继和端点相接。

C++ 实现上，区间可用 \code{std::vector<std::array<int, 2>>} 或结构体表达端点，排序比较器要处理相同起点。动态区间需要 \code{std::map} 的前驱后继，而不是哈希表。如果追问变成“哈希表无法回答前驱后继，动态区间要用有序结构”，先判断原来的状态是否仍然覆盖新答案集合，再决定是微调边界、改 value 语义，还是换成另一类结构。

\subsubsection{LeetCode 731 我的日程安排表 II}

先两两比较区间，确认相交、覆盖或冲突的定义。本题的模型是允许双重预订但不允许三重预订，插入会改变区间重叠层数；暴力慢在每个区间都回头看太多无关区间。优化点保存已有预订和已经双重预订的区间；新预订若碰到双重区间则失败，否则记录新产生的重叠通常依赖排序后相邻关系，必须先定义端点闭开。

区间结构的状态是当前右端、结果尾区间或动态有序表的相邻关系。本题的保存已有预订和已经双重预订的区间；新预订若碰到双重区间则失败，否则记录新产生的重叠只允许丢掉已经不可能再冲突的区间；端点相等是否合并要先定义。放到“区间、排序与合并”这条主线里看，关键原则是：合并按起点排序，调度按终点排序，覆盖题常用起点升序且终点降序。扫描时维护当前最远右端、结果尾区间或已选区间终点。

区间题证明依赖排序后的邻接关系。若当前区间无法影响结果尾部或相邻区间，它就不会再影响更远的区间。本题的保存已有预订和已经双重预订的区间；新预订若碰到双重区间则失败，否则记录新产生的重叠要说明哪些区间已经安全归档。这道题的边界不能留到最后补丁式处理，实验时覆盖端点相接、多个局部重叠、完全覆盖、回滚失败插入，画出排序后的区间序列和当前结果尾部。动态版本要专门测前驱、后继和端点相接。

C++ 实现上，区间可用 \code{std::vector<std::array<int, 2>>} 或结构体表达端点，排序比较器要处理相同起点。动态区间需要 \code{std::map} 的前驱后继，而不是哈希表。如果追问变成“扫描线差分也能做，但要处理失败时恢复状态”，先判断原来的状态是否仍然覆盖新答案集合，再决定是微调边界、改 value 语义，还是换成另一类结构。

\subsubsection{LeetCode 986 区间列表的交集}

先两两比较区间，确认相交、覆盖或冲突的定义。本题的模型是两个有序且不重叠的区间列表，交集只可能来自当前两个区间；暴力慢在每个区间都回头看太多无关区间。优化点每次比较两个当前区间的重叠部分，然后推进右端较小的区间通常依赖排序后相邻关系，必须先定义端点闭开。

区间结构的状态是当前右端、结果尾区间或动态有序表的相邻关系。本题的每次比较两个当前区间的重叠部分，然后推进右端较小的区间只允许丢掉已经不可能再冲突的区间；端点相等是否合并要先定义。放到“区间、排序与合并”这条主线里看，关键原则是：合并按起点排序，调度按终点排序，覆盖题常用起点升序且终点降序。扫描时维护当前最远右端、结果尾区间或已选区间终点。

区间题证明依赖排序后的邻接关系。若当前区间无法影响结果尾部或相邻区间，它就不会再影响更远的区间。本题的每次比较两个当前区间的重叠部分，然后推进右端较小的区间要说明哪些区间已经安全归档。这道题的边界不能留到最后补丁式处理，实验时覆盖空列表、端点相接、一个区间覆盖多个区间、无交集，画出排序后的区间序列和当前结果尾部。动态版本要专门测前驱、后继和端点相接。

C++ 实现上，区间可用 \code{std::vector<std::array<int, 2>>} 或结构体表达端点，排序比较器要处理相同起点。动态区间需要 \code{std::map} 的前驱后继，而不是哈希表。如果追问变成“这是区间双指针，不需要把两个列表合并排序”，先判断原来的状态是否仍然覆盖新答案集合，再决定是微调边界、改 value 语义，还是换成另一类结构。

\subsubsection{LeetCode 1094 拼车}

先两两比较区间，确认相交、覆盖或冲突的定义。本题的模型是乘客上下车事件构成沿路线位置变化的载客量；暴力慢在每个区间都回头看太多无关区间。优化点用差分数组或扫描线记录每个位置人数变化，累计人数不能超过容量通常依赖排序后相邻关系，必须先定义端点闭开。

区间结构的状态是当前右端、结果尾区间或动态有序表的相邻关系。本题的用差分数组或扫描线记录每个位置人数变化，累计人数不能超过容量只允许丢掉已经不可能再冲突的区间；端点相等是否合并要先定义。放到“区间、排序与合并”这条主线里看，关键原则是：合并按起点排序，调度按终点排序，覆盖题常用起点升序且终点降序。扫描时维护当前最远右端、结果尾区间或已选区间终点。

区间题证明依赖排序后的邻接关系。若当前区间无法影响结果尾部或相邻区间，它就不会再影响更远的区间。本题的用差分数组或扫描线记录每个位置人数变化，累计人数不能超过容量要说明哪些区间已经安全归档。这道题的边界不能留到最后补丁式处理，实验时覆盖同站上下车、容量刚好、站点范围、输入未排序，画出排序后的区间序列和当前结果尾部。动态版本要专门测前驱、后继和端点相接。

C++ 实现上，区间可用 \code{std::vector<std::array<int, 2>>} 或结构体表达端点，排序比较器要处理相同起点。动态区间需要 \code{std::map} 的前驱后继，而不是哈希表。如果追问变成“航班预订统计也是差分思想，只是只需要最终每段累加结果”，先判断原来的状态是否仍然覆盖新答案集合，再决定是微调边界、改 value 语义，还是换成另一类结构。

\subsubsection{LeetCode 1109 航班预订统计}

先两两比较区间，确认相交、覆盖或冲突的定义。本题的模型是每条预订给一个连续航班区间增加座位数；暴力慢在每个区间都回头看太多无关区间。优化点差分数组在区间起点加人数，终点后一位减人数，最后前缀还原每航班总数通常依赖排序后相邻关系，必须先定义端点闭开。

区间结构的状态是当前右端、结果尾区间或动态有序表的相邻关系。本题的差分数组在区间起点加人数，终点后一位减人数，最后前缀还原每航班总数只允许丢掉已经不可能再冲突的区间；端点相等是否合并要先定义。放到“区间、排序与合并”这条主线里看，关键原则是：合并按起点排序，调度按终点排序，覆盖题常用起点升序且终点降序。扫描时维护当前最远右端、结果尾区间或已选区间终点。

区间题证明依赖排序后的邻接关系。若当前区间无法影响结果尾部或相邻区间，它就不会再影响更远的区间。本题的差分数组在区间起点加人数，终点后一位减人数，最后前缀还原每航班总数要说明哪些区间已经安全归档。这道题的边界不能留到最后补丁式处理，实验时覆盖区间到最后一个航班、单点区间、多条重叠预订、下标从一开始，画出排序后的区间序列和当前结果尾部。动态版本要专门测前驱、后继和端点相接。

C++ 实现上，区间可用 \code{std::vector<std::array<int, 2>>} 或结构体表达端点，排序比较器要处理相同起点。动态区间需要 \code{std::map} 的前驱后继，而不是哈希表。如果追问变成“这是静态区间加法，不需要线段树”，先判断原来的状态是否仍然覆盖新答案集合，再决定是微调边界、改 value 语义，还是换成另一类结构。

\subsubsection{LeetCode 1229 安排会议日程}

先两两比较区间，确认相交、覆盖或冲突的定义。本题的模型是两个人的可用时间段都有序，目标是找最早长度足够的交集；暴力慢在每个区间都回头看太多无关区间。优化点双指针比较当前两个区间交集，若长度不足就推进结束更早的一侧通常依赖排序后相邻关系，必须先定义端点闭开。

区间结构的状态是当前右端、结果尾区间或动态有序表的相邻关系。本题的双指针比较当前两个区间交集，若长度不足就推进结束更早的一侧只允许丢掉已经不可能再冲突的区间；端点相等是否合并要先定义。放到“区间、排序与合并”这条主线里看，关键原则是：合并按起点排序，调度按终点排序，覆盖题常用起点升序且终点降序。扫描时维护当前最远右端、结果尾区间或已选区间终点。

区间题证明依赖排序后的邻接关系。若当前区间无法影响结果尾部或相邻区间，它就不会再影响更远的区间。本题的双指针比较当前两个区间交集，若长度不足就推进结束更早的一侧要说明哪些区间已经安全归档。这道题的边界不能留到最后补丁式处理，实验时覆盖无可行时间、刚好等于 duration、区间端点相接、输入需要先排序，画出排序后的区间序列和当前结果尾部。动态版本要专门测前驱、后继和端点相接。

C++ 实现上，区间可用 \code{std::vector<std::array<int, 2>>} 或结构体表达端点，排序比较器要处理相同起点。动态区间需要 \code{std::map} 的前驱后继，而不是哈希表。如果追问变成“它和区间列表交集相同，但额外带最早可行和长度约束”，先判断原来的状态是否仍然覆盖新答案集合，再决定是微调边界、改 value 语义，还是换成另一类结构。

\subsubsection{LeetCode 1288 删除被覆盖区间}

先两两比较区间，确认相交、覆盖或冲突的定义。本题的模型是被覆盖要求另一区间起点不晚且终点不早；暴力慢在每个区间都回头看太多无关区间。优化点按起点升序、终点降序排序，扫描维护最大右端通常依赖排序后相邻关系，必须先定义端点闭开。

区间结构的状态是当前右端、结果尾区间或动态有序表的相邻关系。本题的按起点升序、终点降序排序，扫描维护最大右端只允许丢掉已经不可能再冲突的区间；端点相等是否合并要先定义。放到“区间、排序与合并”这条主线里看，关键原则是：合并按起点排序，调度按终点排序，覆盖题常用起点升序且终点降序。扫描时维护当前最远右端、结果尾区间或已选区间终点。

区间题证明依赖排序后的邻接关系。若当前区间无法影响结果尾部或相邻区间，它就不会再影响更远的区间。本题的按起点升序、终点降序排序，扫描维护最大右端要说明哪些区间已经安全归档。这道题的边界不能留到最后补丁式处理，实验时覆盖同起点区间、完全覆盖、无覆盖、端点相等，画出排序后的区间序列和当前结果尾部。动态版本要专门测前驱、后继和端点相接。

C++ 实现上，区间可用 \code{std::vector<std::array<int, 2>>} 或结构体表达端点，排序比较器要处理相同起点。动态区间需要 \code{std::map} 的前驱后继，而不是哈希表。如果追问变成“终点降序是关键，否则同起点小区间会先出现造成误判”，先判断原来的状态是否仍然覆盖新答案集合，再决定是微调边界、改 value 语义，还是换成另一类结构。

\subsection{位运算与状态压缩工作坊}

先用计数或哈希保证正确，再寻找位运算的代数抵消、按位统计或函数图结构。位题的难点是证明被抵消的信息确实不影响答案。

\subsubsection{LeetCode 29 两数相除}

先用计数、集合或普通 DP 写出清楚版本。本题的模型是除法可以看成不断减去除数的倍增值；位运算只是把状态压进二进制表示。优化点把除数按二进制倍增，优先减去不超过被除数的最大倍数并累加商必须说明每一位或每次异或到底保留了什么信息。

位状态要给每一位固定含义。本题的把除数按二进制倍增，优先减去不超过被除数的最大倍数并累加商如果依赖异或抵消，就要说明成对元素为什么完全消失；如果依赖掩码，就要说明同一位置不同掩码为什么不能合并。放到“位运算与状态压缩”这条主线里看，关键原则是：异或解决成对抵消，按位计数处理出现三次，掩码表示访问集合或可用集合，低位提取用于枚举候选。

位运算证明要逐位说明。异或、按位计数或掩码转移都不能只靠直觉；本题的把除数按二进制倍增，优先减去不超过被除数的最大倍数并累加商必须保证被压掉的信息不会影响除法可以看成不断减去除数的倍增值要求的答案。这道题的边界不能留到最后补丁式处理，实验时覆盖除数为一、被除数为 int 最小值、结果溢出、符号处理，把数字写成二进制逐位检查。负数、最高位、空掩码和全集掩码都要单独测。

C++ 实现上，移位时优先使用无符号或足够宽的整数，位数超过 31 时考虑 \code{long long} 或 \code{std::bitset}。位运算表达式加括号，避免优先级误读。如果追问变成“它和快速幂类似，都利用二进制拆分减少重复加减”，先判断原来的状态是否仍然覆盖新答案集合，再决定是微调边界、改 value 语义，还是换成另一类结构。

\subsubsection{LeetCode 136 只出现一次的数字}

先用计数、集合或普通 DP 写出清楚版本。本题的模型是成对元素可以用异或抵消，剩下单个元素；位运算只是把状态压进二进制表示。优化点异或满足交换结合，扫描顺序不影响答案必须说明每一位或每次异或到底保留了什么信息。

位状态要给每一位固定含义。本题的异或满足交换结合，扫描顺序不影响答案如果依赖异或抵消，就要说明成对元素为什么完全消失；如果依赖掩码，就要说明同一位置不同掩码为什么不能合并。放到“位运算与状态压缩”这条主线里看，关键原则是：异或解决成对抵消，按位计数处理出现三次，掩码表示访问集合或可用集合，低位提取用于枚举候选。

位运算证明要逐位说明。异或、按位计数或掩码转移都不能只靠直觉；本题的异或满足交换结合，扫描顺序不影响答案必须保证被压掉的信息不会影响成对元素可以用异或抵消，剩下单个元素要求的答案。这道题的边界不能留到最后补丁式处理，实验时覆盖零、负数、唯一值在开头或结尾、所有其他值恰好两次，把数字写成二进制逐位检查。负数、最高位、空掩码和全集掩码都要单独测。

C++ 实现上，移位时优先使用无符号或足够宽的整数，位数超过 31 时考虑 \code{long long} 或 \code{std::bitset}。位运算表达式加括号，避免优先级误读。如果追问变成“若其他值出现三次，要改为按位计数取模”，先判断原来的状态是否仍然覆盖新答案集合，再决定是微调边界、改 value 语义，还是换成另一类结构。

\subsubsection{LeetCode 137 只出现一次的数字 II}

先用计数、集合或普通 DP 写出清楚版本。本题的模型是除一个数字外，其余数字都出现三次，单纯异或不能抵消；位运算只是把状态压进二进制表示。优化点逐位统计一的数量，对三取模后还原目标数字必须说明每一位或每次异或到底保留了什么信息。

位状态要给每一位固定含义。本题的逐位统计一的数量，对三取模后还原目标数字如果依赖异或抵消，就要说明成对元素为什么完全消失；如果依赖掩码，就要说明同一位置不同掩码为什么不能合并。放到“位运算与状态压缩”这条主线里看，关键原则是：异或解决成对抵消，按位计数处理出现三次，掩码表示访问集合或可用集合，低位提取用于枚举候选。

位运算证明要逐位说明。异或、按位计数或掩码转移都不能只靠直觉；本题的逐位统计一的数量，对三取模后还原目标数字必须保证被压掉的信息不会影响除一个数字外，其余数字都出现三次，单纯异或不能抵消要求的答案。这道题的边界不能留到最后补丁式处理，实验时覆盖负数、符号位、目标为零、所有其他数恰好三次，把数字写成二进制逐位检查。负数、最高位、空掩码和全集掩码都要单独测。

C++ 实现上，移位时优先使用无符号或足够宽的整数，位数超过 31 时考虑 \code{long long} 或 \code{std::bitset}。位运算表达式加括号，避免优先级误读。如果追问变成“若其他数出现 k 次，思路仍是按位计数取模”，先判断原来的状态是否仍然覆盖新答案集合，再决定是微调边界、改 value 语义，还是换成另一类结构。

\subsubsection{LeetCode 190 颠倒二进制位}

先用计数、集合或普通 DP 写出清楚版本。本题的模型是每一位的新位置由原位置镜像决定；位运算只是把状态压进二进制表示。优化点循环三十二次，把结果左移并接上当前最低位，再右移输入必须说明每一位或每次异或到底保留了什么信息。

位状态要给每一位固定含义。本题的循环三十二次，把结果左移并接上当前最低位，再右移输入如果依赖异或抵消，就要说明成对元素为什么完全消失；如果依赖掩码，就要说明同一位置不同掩码为什么不能合并。放到“位运算与状态压缩”这条主线里看，关键原则是：异或解决成对抵消，按位计数处理出现三次，掩码表示访问集合或可用集合，低位提取用于枚举候选。

位运算证明要逐位说明。异或、按位计数或掩码转移都不能只靠直觉；本题的循环三十二次，把结果左移并接上当前最低位，再右移输入必须保证被压掉的信息不会影响每一位的新位置由原位置镜像决定要求的答案。这道题的边界不能留到最后补丁式处理，实验时覆盖最高位、最低位、全零、全一、无符号语义，把数字写成二进制逐位检查。负数、最高位、空掩码和全集掩码都要单独测。

C++ 实现上，移位时优先使用无符号或足够宽的整数，位数超过 31 时考虑 \code{long long} 或 \code{std::bitset}。位运算表达式加括号，避免优先级误读。如果追问变成“若多次调用，可以按字节预处理反转表”，先判断原来的状态是否仍然覆盖新答案集合，再决定是微调边界、改 value 语义，还是换成另一类结构。

\subsubsection{LeetCode 191 位 1 的个数}

先用计数、集合或普通 DP 写出清楚版本。本题的模型是每次清掉最低位的一可以让循环次数等于一的个数；位运算只是把状态压进二进制表示。优化点使用 n \& (n - 1) 删除最低位一，比逐位检查更直接必须说明每一位或每次异或到底保留了什么信息。

位状态要给每一位固定含义。本题的使用 n \& (n - 1) 删除最低位一，比逐位检查更直接如果依赖异或抵消，就要说明成对元素为什么完全消失；如果依赖掩码，就要说明同一位置不同掩码为什么不能合并。放到“位运算与状态压缩”这条主线里看，关键原则是：异或解决成对抵消，按位计数处理出现三次，掩码表示访问集合或可用集合，低位提取用于枚举候选。

位运算证明要逐位说明。异或、按位计数或掩码转移都不能只靠直觉；本题的使用 n \& (n - 1) 删除最低位一，比逐位检查更直接必须保证被压掉的信息不会影响每次清掉最低位的一可以让循环次数等于一的个数要求的答案。这道题的边界不能留到最后补丁式处理，实验时覆盖零、只有最高位、一的个数很多、无符号输入，把数字写成二进制逐位检查。负数、最高位、空掩码和全集掩码都要单独测。

C++ 实现上，移位时优先使用无符号或足够宽的整数，位数超过 31 时考虑 \code{long long} 或 \code{std::bitset}。位运算表达式加括号，避免优先级误读。如果追问变成“比特位计数可以把这个过程 DP 化到一整段数”，先判断原来的状态是否仍然覆盖新答案集合，再决定是微调边界、改 value 语义，还是换成另一类结构。

\subsubsection{LeetCode 201 数字范围按位与}

先用计数、集合或普通 DP 写出清楚版本。本题的模型是区间内所有数按位与后，只保留左右端公共二进制前缀；位运算只是把状态压进二进制表示。优化点不断右移 left 和 right，直到二者相等，再把公共前缀移回原位必须说明每一位或每次异或到底保留了什么信息。

位状态要给每一位固定含义。本题的不断右移 left 和 right，直到二者相等，再把公共前缀移回原位如果依赖异或抵消，就要说明成对元素为什么完全消失；如果依赖掩码，就要说明同一位置不同掩码为什么不能合并。放到“位运算与状态压缩”这条主线里看，关键原则是：异或解决成对抵消，按位计数处理出现三次，掩码表示访问集合或可用集合，低位提取用于枚举候选。

位运算证明要逐位说明。异或、按位计数或掩码转移都不能只靠直觉；本题的不断右移 left 和 right，直到二者相等，再把公共前缀移回原位必须保证被压掉的信息不会影响区间内所有数按位与后，只保留左右端公共二进制前缀要求的答案。这道题的边界不能留到最后补丁式处理，实验时覆盖left 等于 right、区间跨过二的幂边界、结果为零，把数字写成二进制逐位检查。负数、最高位、空掩码和全集掩码都要单独测。

C++ 实现上，移位时优先使用无符号或足够宽的整数，位数超过 31 时考虑 \code{long long} 或 \code{std::bitset}。位运算表达式加括号，避免优先级误读。如果追问变成“也可以不断清除 right 的最低位一，直到 right 不大于 left”，先判断原来的状态是否仍然覆盖新答案集合，再决定是微调边界、改 value 语义，还是换成另一类结构。

\subsubsection{LeetCode 231 2 的幂}

先用计数、集合或普通 DP 写出清楚版本。本题的模型是二的幂在二进制中只有一个一；位运算只是把状态压进二进制表示。优化点正数 n 满足 n 与 n 减一按位与为零必须说明每一位或每次异或到底保留了什么信息。

位状态要给每一位固定含义。本题的正数 n 满足 n 与 n 减一按位与为零如果依赖异或抵消，就要说明成对元素为什么完全消失；如果依赖掩码，就要说明同一位置不同掩码为什么不能合并。放到“位运算与状态压缩”这条主线里看，关键原则是：异或解决成对抵消，按位计数处理出现三次，掩码表示访问集合或可用集合，低位提取用于枚举候选。

位运算证明要逐位说明。异或、按位计数或掩码转移都不能只靠直觉；本题的正数 n 满足 n 与 n 减一按位与为零必须保证被压掉的信息不会影响二的幂在二进制中只有一个一要求的答案。这道题的边界不能留到最后补丁式处理，实验时覆盖n 为零、负数、int 最小值、最高位，把数字写成二进制逐位检查。负数、最高位、空掩码和全集掩码都要单独测。

C++ 实现上，移位时优先使用无符号或足够宽的整数，位数超过 31 时考虑 \code{long long} 或 \code{std::bitset}。位运算表达式加括号，避免优先级误读。如果追问变成“判断四的幂还要额外限制一所在的位置为偶数位”，先判断原来的状态是否仍然覆盖新答案集合，再决定是微调边界、改 value 语义，还是换成另一类结构。

\subsubsection{LeetCode 260 只出现一次的数字 III}

先用计数、集合或普通 DP 写出清楚版本。本题的模型是两个只出现一次的数异或后至少有一位不同；位运算只是把状态压进二进制表示。优化点用最低不同位把数组分成两组，每组内部再异或抵消必须说明每一位或每次异或到底保留了什么信息。

位状态要给每一位固定含义。本题的用最低不同位把数组分成两组，每组内部再异或抵消如果依赖异或抵消，就要说明成对元素为什么完全消失；如果依赖掩码，就要说明同一位置不同掩码为什么不能合并。放到“位运算与状态压缩”这条主线里看，关键原则是：异或解决成对抵消，按位计数处理出现三次，掩码表示访问集合或可用集合，低位提取用于枚举候选。

位运算证明要逐位说明。异或、按位计数或掩码转移都不能只靠直觉；本题的用最低不同位把数组分成两组，每组内部再异或抵消必须保证被压掉的信息不会影响两个只出现一次的数异或后至少有一位不同要求的答案。这道题的边界不能留到最后补丁式处理，实验时覆盖两个答案一正一负、最低位分组、其他数恰好两次，把数字写成二进制逐位检查。负数、最高位、空掩码和全集掩码都要单独测。

C++ 实现上，移位时优先使用无符号或足够宽的整数，位数超过 31 时考虑 \code{long long} 或 \code{std::bitset}。位运算表达式加括号，避免优先级误读。如果追问变成“它是在异或抵消基础上增加分组位”，先判断原来的状态是否仍然覆盖新答案集合，再决定是微调边界、改 value 语义，还是换成另一类结构。

\subsubsection{LeetCode 268 丢失的数字}

先用计数、集合或普通 DP 写出清楚版本。本题的模型是下标零到 n 与数组值相比，缺失值会在异或抵消后剩下；位运算只是把状态压进二进制表示。优化点把所有下标和所有值一起异或，重复出现的数抵消必须说明每一位或每次异或到底保留了什么信息。

位状态要给每一位固定含义。本题的把所有下标和所有值一起异或，重复出现的数抵消如果依赖异或抵消，就要说明成对元素为什么完全消失；如果依赖掩码，就要说明同一位置不同掩码为什么不能合并。放到“位运算与状态压缩”这条主线里看，关键原则是：异或解决成对抵消，按位计数处理出现三次，掩码表示访问集合或可用集合，低位提取用于枚举候选。

位运算证明要逐位说明。异或、按位计数或掩码转移都不能只靠直觉；本题的把所有下标和所有值一起异或，重复出现的数抵消必须保证被压掉的信息不会影响下标零到 n 与数组值相比，缺失值会在异或抵消后剩下要求的答案。这道题的边界不能留到最后补丁式处理，实验时覆盖缺失零、缺失 n、数组长度一、也可用求和但要防溢出，把数字写成二进制逐位检查。负数、最高位、空掩码和全集掩码都要单独测。

C++ 实现上，移位时优先使用无符号或足够宽的整数，位数超过 31 时考虑 \code{long long} 或 \code{std::bitset}。位运算表达式加括号，避免优先级误读。如果追问变成“异或法适合展示值域完整且只有一个缺口的结构”，先判断原来的状态是否仍然覆盖新答案集合，再决定是微调边界、改 value 语义，还是换成另一类结构。

\subsubsection{LeetCode 287 寻找重复数}

先用计数、集合或普通 DP 写出清楚版本。本题的模型是数组值域使下标到值形成函数图，重复数是环入口；位运算只是把状态压进二进制表示。优化点快慢指针不修改数组且常数空间；值域二分用抽屉原理必须说明每一位或每次异或到底保留了什么信息。

位状态要给每一位固定含义。本题的快慢指针不修改数组且常数空间；值域二分用抽屉原理如果依赖异或抵消，就要说明成对元素为什么完全消失；如果依赖掩码，就要说明同一位置不同掩码为什么不能合并。放到“位运算与状态压缩”这条主线里看，关键原则是：异或解决成对抵消，按位计数处理出现三次，掩码表示访问集合或可用集合，低位提取用于枚举候选。

位运算证明要逐位说明。异或、按位计数或掩码转移都不能只靠直觉；本题的快慢指针不修改数组且常数空间；值域二分用抽屉原理必须保证被压掉的信息不会影响数组值域使下标到值形成函数图，重复数是环入口要求的答案。这道题的边界不能留到最后补丁式处理，实验时覆盖重复数出现多次、不能排序、不能改数组、长度为 n 加一，把数字写成二进制逐位检查。负数、最高位、空掩码和全集掩码都要单独测。

C++ 实现上，移位时优先使用无符号或足够宽的整数，位数超过 31 时考虑 \code{long long} 或 \code{std::bitset}。位运算表达式加括号，避免优先级误读。如果追问变成“快慢指针要求值都能作为合法下标”，先判断原来的状态是否仍然覆盖新答案集合，再决定是微调边界、改 value 语义，还是换成另一类结构。

\subsubsection{LeetCode 318 最大单词长度乘积}

先用计数、集合或普通 DP 写出清楚版本。本题的模型是两个单词没有公共字母时才可相乘，字母集合可压成位掩码；位运算只是把状态压进二进制表示。优化点每个单词生成 26 位掩码，两个掩码按位与为零表示无交集必须说明每一位或每次异或到底保留了什么信息。

位状态要给每一位固定含义。本题的每个单词生成 26 位掩码，两个掩码按位与为零表示无交集如果依赖异或抵消，就要说明成对元素为什么完全消失；如果依赖掩码，就要说明同一位置不同掩码为什么不能合并。放到“位运算与状态压缩”这条主线里看，关键原则是：异或解决成对抵消，按位计数处理出现三次，掩码表示访问集合或可用集合，低位提取用于枚举候选。

位运算证明要逐位说明。异或、按位计数或掩码转移都不能只靠直觉；本题的每个单词生成 26 位掩码，两个掩码按位与为零表示无交集必须保证被压掉的信息不会影响两个单词没有公共字母时才可相乘，字母集合可压成位掩码要求的答案。这道题的边界不能留到最后补丁式处理，实验时覆盖重复字母、空字符串、多个单词同掩码、乘积范围，把数字写成二进制逐位检查。负数、最高位、空掩码和全集掩码都要单独测。

C++ 实现上，移位时优先使用无符号或足够宽的整数，位数超过 31 时考虑 \code{long long} 或 \code{std::bitset}。位运算表达式加括号，避免优先级误读。如果追问变成“若字符集大于整数位数，需要 bitset 或哈希集合”，先判断原来的状态是否仍然覆盖新答案集合，再决定是微调边界、改 value 语义，还是换成另一类结构。

\subsubsection{LeetCode 338 比特位计数}

先用计数、集合或普通 DP 写出清楚版本。本题的模型是每个数的一的个数可由去掉最低位一后的较小数转移；位运算只是把状态压进二进制表示。优化点dp[x] = dp[x \& (x - 1)] + 1，或 dp[x] = dp[x >> 1] + (x \& 1)必须说明每一位或每次异或到底保留了什么信息。

位状态要给每一位固定含义。本题的dp[x] = dp[x \& (x - 1)] + 1，或 dp[x] = dp[x >> 1] + (x \& 1)如果依赖异或抵消，就要说明成对元素为什么完全消失；如果依赖掩码，就要说明同一位置不同掩码为什么不能合并。放到“位运算与状态压缩”这条主线里看，关键原则是：异或解决成对抵消，按位计数处理出现三次，掩码表示访问集合或可用集合，低位提取用于枚举候选。

位运算证明要逐位说明。异或、按位计数或掩码转移都不能只靠直觉；本题的dp[x] = dp[x \& (x - 1)] + 1，或 dp[x] = dp[x >> 1] + (x \& 1)必须保证被压掉的信息不会影响每个数的一的个数可由去掉最低位一后的较小数转移要求的答案。这道题的边界不能留到最后补丁式处理，实验时覆盖n 为零、二的幂、连续区间、表达式优先级，把数字写成二进制逐位检查。负数、最高位、空掩码和全集掩码都要单独测。

C++ 实现上，移位时优先使用无符号或足够宽的整数，位数超过 31 时考虑 \code{long long} 或 \code{std::bitset}。位运算表达式加括号，避免优先级误读。如果追问变成“它展示位运算和 DP 的结合，而不是逐个数循环数位”，先判断原来的状态是否仍然覆盖新答案集合，再决定是微调边界、改 value 语义，还是换成另一类结构。

\subsubsection{LeetCode 371 两整数之和}

先用计数、集合或普通 DP 写出清楚版本。本题的模型是不用加号时，加法可以拆成无进位和与进位；位运算只是把状态压进二进制表示。优化点异或得到无进位和，与运算后左移得到进位，迭代直到进位为零必须说明每一位或每次异或到底保留了什么信息。

位状态要给每一位固定含义。本题的异或得到无进位和，与运算后左移得到进位，迭代直到进位为零如果依赖异或抵消，就要说明成对元素为什么完全消失；如果依赖掩码，就要说明同一位置不同掩码为什么不能合并。放到“位运算与状态压缩”这条主线里看，关键原则是：异或解决成对抵消，按位计数处理出现三次，掩码表示访问集合或可用集合，低位提取用于枚举候选。

位运算证明要逐位说明。异或、按位计数或掩码转移都不能只靠直觉；本题的异或得到无进位和，与运算后左移得到进位，迭代直到进位为零必须保证被压掉的信息不会影响不用加号时，加法可以拆成无进位和与进位要求的答案。这道题的边界不能留到最后补丁式处理，实验时覆盖正负数、进位传播、有符号溢出、需要使用无符号中间值，把数字写成二进制逐位检查。负数、最高位、空掩码和全集掩码都要单独测。

C++ 实现上，移位时优先使用无符号或足够宽的整数，位数超过 31 时考虑 \code{long long} 或 \code{std::bitset}。位运算表达式加括号，避免优先级误读。如果追问变成“这是位级算术模拟，和异或抵消类题目的语义不同”，先判断原来的状态是否仍然覆盖新答案集合，再决定是微调边界、改 value 语义，还是换成另一类结构。

\subsubsection{LeetCode 393 UTF-8 编码验证}

先用计数、集合或普通 DP 写出清楚版本。本题的模型是每个字节的高位模式决定一个字符需要的后续字节数量；位运算只是把状态压进二进制表示。优化点用位掩码判断首字节类型，并维护还需要多少个 continuation 字节必须说明每一位或每次异或到底保留了什么信息。

位状态要给每一位固定含义。本题的用位掩码判断首字节类型，并维护还需要多少个 continuation 字节如果依赖异或抵消，就要说明成对元素为什么完全消失；如果依赖掩码，就要说明同一位置不同掩码为什么不能合并。放到“位运算与状态压缩”这条主线里看，关键原则是：异或解决成对抵消，按位计数处理出现三次，掩码表示访问集合或可用集合，低位提取用于枚举候选。

位运算证明要逐位说明。异或、按位计数或掩码转移都不能只靠直觉；本题的用位掩码判断首字节类型，并维护还需要多少个 continuation 字节必须保证被压掉的信息不会影响每个字节的高位模式决定一个字符需要的后续字节数量要求的答案。这道题的边界不能留到最后补丁式处理，实验时覆盖非法首字节、后续字节不足、后续字节格式错误、单字节字符，把数字写成二进制逐位检查。负数、最高位、空掩码和全集掩码都要单独测。

C++ 实现上，移位时优先使用无符号或足够宽的整数，位数超过 31 时考虑 \code{long long} 或 \code{std::bitset}。位运算表达式加括号，避免优先级误读。如果追问变成“这是位模式解析题，重点是协议规则和状态机”，先判断原来的状态是否仍然覆盖新答案集合，再决定是微调边界、改 value 语义，还是换成另一类结构。

\subsubsection{LeetCode 421 数组中两个数的最大异或值}

先用计数、集合或普通 DP 写出清楚版本。本题的模型是最大异或值由高位到低位尽量取一决定；位运算只是把状态压进二进制表示。优化点逐位尝试把当前前缀答案设为一，用哈希集合验证是否存在两个前缀配对必须说明每一位或每次异或到底保留了什么信息。

位状态要给每一位固定含义。本题的逐位尝试把当前前缀答案设为一，用哈希集合验证是否存在两个前缀配对如果依赖异或抵消，就要说明成对元素为什么完全消失；如果依赖掩码，就要说明同一位置不同掩码为什么不能合并。放到“位运算与状态压缩”这条主线里看，关键原则是：异或解决成对抵消，按位计数处理出现三次，掩码表示访问集合或可用集合，低位提取用于枚举候选。

位运算证明要逐位说明。异或、按位计数或掩码转移都不能只靠直觉；本题的逐位尝试把当前前缀答案设为一，用哈希集合验证是否存在两个前缀配对必须保证被压掉的信息不会影响最大异或值由高位到低位尽量取一决定要求的答案。这道题的边界不能留到最后补丁式处理，实验时覆盖零、重复数、最高位、负数语义，把数字写成二进制逐位检查。负数、最高位、空掩码和全集掩码都要单独测。

C++ 实现上，移位时优先使用无符号或足够宽的整数，位数超过 31 时考虑 \code{long long} 或 \code{std::bitset}。位运算表达式加括号，避免优先级误读。如果追问变成“也可以建二进制 Trie，每个数尽量走相反位”，先判断原来的状态是否仍然覆盖新答案集合，再决定是微调边界、改 value 语义，还是换成另一类结构。

\begin{principlebox}{工作坊使用方式}
不要一次性把本节当题解背完。建议每次选一个题型，先遮住优化落点，只看题面模型和暴力基线，自己写出慢在哪里；再看结构观察和正确性；最后用边界实验做对拍。能完成这个过程，才说明题目进入了自己的知识体系。
\end{principlebox}
